\chapt{उत्तरकाण्डः}

\sect{प्रथमः सर्गः}

\twolineshloka
{प्राप्तराज्यस्य रामस्य राक्षसानां वधे कृते}
{आजग्मुरृषयः सर्वे राघवं प्रतिनन्दितुम्} %||7-1-1||

\twolineshloka
{कौशिकोऽथ यवक्रीतो रैभ्यश्च्यवन एव च}
{कण्वो मेधातिथेः पुत्रः पूर्वस्यां दिशि ये श्रिताः} %||7-1-2||

\twolineshloka
{स्वस्त्यात्रेयश्च भगवान्नमुचिः प्रमुचुस्तथा}
{आजग्मुस्ते सहागस्त्या ये श्रिता दक्षिणां दिशम्} %||7-1-3||

\twolineshloka
{पृषद्गुः कवषो धौम्यो रौद्रेयश्च महानृषिः}
{तेऽप्याजग्मुः सशिष्या वै ये श्रिताः पश्चिमां दिशम्} %||7-1-4||

\twolineshloka
{वसिष्ठः कश्यपोऽथात्रिर्विश्वामित्रोऽथ गौतमः}
{जमदग्निर्भरद्वाजस्तेऽपि सप्तमहर्षयः} %||7-1-5||

\twolineshloka
{सम्प्राप्यैते महात्मानो राघवस्य निवेशनम्}
{विष्ठिताः प्रतिहारार्थं हुताशनसमप्रभाः} %||7-1-6||

\twolineshloka
{प्रतिहारस्ततस्तूर्णमगस्त्यवचनादथ}
{समीपं राघवस्याशु प्रविवेश महात्मनः} %||7-1-7||

\twolineshloka
{स रामं दृश्य सहसा पूर्णचन्द्रसमद्युतिम्}
{अगस्त्यं कथयामास सम्प्रातमृषिभिः सह} %||7-1-8||

\twolineshloka
{श्रुत्वा प्राप्तान्मुनींस्तांस्तु बालसूर्यसमप्रभान्}
{तदोवाच नृपो द्वाःस्थं प्रवेशय यथासुखम्} %||7-1-9||

\twolineshloka
{दृष्ट्वा प्राप्तान्मुनींस्तांस्तु प्रत्युत्थाय कृताञ्जलिः}
{रामोऽभिवाद्य प्रयत आसनान्यादिदेश ह} %||7-1-10||

\twolineshloka
{तेषु काञ्चनचित्रेषु स्वास्तीर्णेषु सुखेषु च}
{यथार्हमुपविष्टास्ते आसनेष्वृषिपुङ्गवाः} %||7-1-11||

\twolineshloka
{रामेण कुशलं पृष्टाः सशिष्याः सपुरोगमाः}
{महर्षयो वेदविदो रामं वचनमब्रुवन्} %||7-1-12||

\twolineshloka
{कुशलं नो महाबाहो सर्वत्र रघुनन्दन}
{त्वां तु दिष्ट्या कुशलिनं पश्यामो हतशात्रवम्} %||7-1-13||

\twolineshloka
{न हि भारः स ते राम रावणो राक्षसेश्वरः}
{सधनुस्त्वं हि लोकांस्त्रीन्विजयेथा न संशयः} %||7-1-14||

\twolineshloka
{दिष्ट्या त्वया हतो राम रावणः पुत्रपौत्रवान्}
{दिष्ट्या विजयिनं त्वाद्य पश्यामः सह भार्यया} %||7-1-15||

\twolineshloka
{दिष्ट्या प्रहस्तो विकटो विरूपाक्षो महोदरः}
{अकम्पनश्च दुर्धर्षो निहतास्ते निशाचराः} %||7-1-16||

\twolineshloka
{यस्य प्रमाणाद्विपुलं प्रमाणं नेह विद्यते}
{दिष्ट्या ते समरे राम कुम्भकर्णो निपातितः} %||7-1-17||

\twolineshloka
{दिष्ट्या त्वं राक्षसेन्द्रेण द्वन्द्वयुद्धमुपागतः}
{देवतानामवध्येन विजयं प्राप्तवानसि} %||7-1-18||

\twolineshloka
{सङ्ख्ये तस्य न किञ्चित्तु रावणस्य पराभवः}
{द्वन्द्वयुद्धमनुप्राप्तो दिष्ट्या ते रावणिर्हतः} %||7-1-19||

\twolineshloka
{दिष्ट्या तस्य महाबाहो कालस्येवाभिधावतः}
{मुक्तः सुररिपोर्वीर प्राप्तश्च विजयस्त्वया} %||7-1-20||

\twolineshloka
{विस्मयस्त्वेष नः सौम्य संश्रुत्येन्द्रजितं हतम्}
{अवध्यः सर्वभूतानां महामायाधरो युधि} %||7-1-21||

\twolineshloka
{दत्त्वा पुण्यामिमां वीर सौम्यामभयदक्षिणाम्}
{दिष्ट्या वर्धसि काकुत्स्थ जयेनामित्रकर्शन} %||7-1-22||

\twolineshloka
{श्रुत्वा तु वचनं तेषामृषीणां भावितात्मनाम्}
{विस्मयं परमं गत्वा रामः प्राञ्जलिरब्रवीत्} %||7-1-23||

\twolineshloka
{भवन्तः कुम्भकर्णं च रावणं च निशाचरम्}
{अतिक्रम्य महावीर्यौ किं प्रशंसथ रावणिम्} %||7-1-24||

\twolineshloka
{महोदरं प्रहस्तं च विरूपाक्षं च राक्षसम्}
{अतिक्रम्य महावीर्यान्किं प्रशंसथ रावणिम्} %||7-1-25||

\twolineshloka
{कीदृशो वै प्रभावोऽस्य किं बलं कः पराक्रमः}
{केन वा कारणेनैष रावणादतिरिच्यते} %||7-1-26||

\threelineshloka
{शक्यं यदि मया श्रोतुं न खल्वाज्ञापयामि वः}
{यदि गुह्यं न चेद्वक्तुं श्रोतुमिच्छामि कथ्यताम्}
{कथं शक्रो जितस्तेन कथं लब्धवरश्च सः} %||7-2-27||


{॥इत्यार्षे श्रीमद्रामायणे वाल्मीकीये आदिकाव्ये उत्तरकाण्डे प्रथमः सर्गः॥}

\sect{द्वितीयः सर्गः}

\twolineshloka
{तस्य तद्वचनं श्रुत्वा राघवस्य महात्मनः}
{कुम्भयोनिर्महातेजा वाक्यमेतदुवाच ह} %||7-2-1||

\twolineshloka
{शृणु राजन्यथावृत्तं यस्य तेजोबलं महत्}
{जघान च रिपून्युद्धे यथावध्यश्च शत्रुभिः} %||7-2-2||

\twolineshloka
{अहं ते रावणस्येदं कुलं जन्म च राघव}
{वरप्रदानं च तथा तस्मै दत्तं ब्रवीमि ते} %||7-2-3||

\twolineshloka
{पुरा कृतयुगे राम प्रजापतिसुतः प्रभुः}
{पुलस्त्यो नाम ब्रह्मर्षिः साक्षादिव पितामहः} %||7-2-4||

\twolineshloka
{नानुकीर्त्या गुणास्तस्य धर्मतः शीलतस्तथा}
{प्रजापतेः पुत्र इति वक्तुं शक्यं हि नामतः} %||7-2-5||

\twolineshloka
{स तु धर्मप्रसङ्गेन मेरोः पार्श्वे महागिरेः}
{तृणबिन्द्वाश्रमं गत्वा न्यवसन्मुनिपुङ्गवः} %||7-2-6||

\twolineshloka
{तपस्तेपे स धर्मात्मा स्वाध्यायनियतेन्द्रियः}
{गत्वाश्रमपदं तस्य विघ्नं कुर्वन्ति कन्यकाः} %||7-2-7||

\twolineshloka
{देवपन्नगकन्याश्च राजर्षितनयाश्च याः}
{क्रीडन्त्योऽप्सरसश्चैव तं देशमुपपेदिरे} %||7-2-8||

\twolineshloka
{सर्वर्तुषूपभोग्यत्वाद्रम्यत्वात्काननस्य च}
{नित्यशस्तास्तु तं देशं गत्वा क्रीडन्ति कन्यकाः} %||7-2-9||

\twolineshloka
{अथ रुष्टो महातेजा व्याजहार महामुनिः}
{या मे दर्शनमागच्छेत्सा गर्भं धारयिष्यति} %||7-2-10||

\twolineshloka
{तास्तु सर्वाः प्रतिगताः श्रुत्वा वाक्यं महात्मनः}
{ब्रह्मशापभयाद्भीतास्तं देशं नोपचक्रमुः} %||7-2-11||

\twolineshloka
{तृणबिन्दोस्तु राजर्षेस्तनया न शृणोति तत्}
{गत्वाश्रमपदं तस्य विचचार सुनिर्भया} %||7-2-12||

\twolineshloka
{तस्मिन्नेव तु काले स प्राजापत्यो महानृषिः}
{स्वाध्यायमकरोत्तत्र तपसा द्योतितप्रभः} %||7-2-13||

\twolineshloka
{सा तु वेदध्वनिं श्रुत्वा दृष्ट्वा चैव तपोधनम्}
{अभवत्पाण्डुदेहा सा सुव्यञ्जितशरीरजा} %||7-2-14||

\twolineshloka
{दृष्ट्वा परमसंविग्ना सा तु तद्रूपमात्मनः}
{इदं मे किं न्विति ज्ञात्वा पितुर्गत्वाग्रतः स्थिता} %||7-2-15||

\twolineshloka
{तां तु दृष्ट्वा तथा भूतां तृणबिन्दुरथाब्रवीत्}
{किं त्वमेतत्त्वसदृशं धारयस्यात्मनो वपुः} %||7-2-16||

\twolineshloka
{सा तु कृत्वाञ्जलिं दीना कन्योवाच तपोधनम्}
{न जाने कारणं तात येन मे रूपमीदृशम्} %||7-2-17||

\twolineshloka
{किं तु पूर्वं गतास्म्येका महर्षेर्भावितात्मनः}
{पुलस्त्यस्याश्रमं दिव्यमन्वेष्टुं स्वसखीजनम्} %||7-2-18||

\twolineshloka
{न च पश्याम्यहं तत्र काञ्चिदप्यागतां सखीम्}
{रूपस्य तु विपर्यासं दृष्ट्वा चाहमिहागता} %||7-2-19||

\twolineshloka
{तृणबिन्दुस्तु राजर्षिस्तपसा द्योतितप्रभः}
{ध्यानं विवेश तच्चापि अपश्यदृषिकर्मजम्} %||7-2-20||

\twolineshloka
{स तु विज्ञाय तं शापं महर्षेर्भावितात्मनः}
{गृहीत्वा तनयां गत्वा पुलस्त्यमिदमब्रवीत्} %||7-2-21||

\twolineshloka
{भगवंस्तनयां मे त्वं गुणैः स्वैरेव भूषिताम्}
{भिक्षां प्रतिगृहाणेमां महर्षे स्वयमुद्यताम्} %||7-2-22||

\twolineshloka
{तपश्चरणयुक्तस्य श्राम्यमाणेन्द्रियस्य ते}
{शुश्रूषातत्परा नित्यं भविष्यति न संशयः} %||7-2-23||

\twolineshloka
{तं ब्रुवाणं तु तद्वाक्यं राजर्षिं धार्मिकं तदा}
{जिघृक्षुरब्रवीत्कन्यां बाढमित्येव स द्विजः} %||7-2-24||

\threelineshloka
{दत्त्वा तु स गतो राजा स्वमाश्रमपदं तदा}
{सापि तत्रावसत्कन्या तोषयन्ती पतिं गुणैः}
{प्रीतः स तु महातेजा वाक्यमेतदुवाच ह} %||7-2-25||

\threelineshloka
{परितुष्टोऽस्मि भद्रं ते गुणानां सम्पदा भृशम्}
{तस्मात्ते विरमाम्यद्य पुत्रमात्मसमं गुणैः}
{उभयोर्वंशकर्तारं पौलस्त्य इति विश्रुतम्} %||7-2-26||

\twolineshloka
{यस्मात्तु विश्रुतो वेदस्त्वयेहाभ्यस्यतो मम}
{तस्मात्स विश्रवा नाम भविष्यति न संशयः} %||7-2-27||

\twolineshloka
{एवमुक्ता तु सा कन्या प्रहृष्टेनान्तरात्मना}
{अचिरेणैव कालेन सूता विश्रवसं सुतम्} %||7-2-28||

\twolineshloka
{स तु लोकत्रये ख्यातः शौचधर्मसमन्वितः}
{पितेव तपसा युक्तो विश्रवा मुनिपुङ्गवः} %||7-3-29||


{॥इत्यार्षे श्रीमद्रामायणे वाल्मीकीये आदिकाव्ये उत्तरकाण्डे द्वितीयः सर्गः॥}

\sect{तृतीयः सर्गः}

\twolineshloka
{अथ पुत्रः पुलस्त्यस्य विश्रवा मुनिपुङ्गवः}
{अचिरेणैव कालेन पितेव तपसि स्थितः} %||7-3-1||

\twolineshloka
{सत्यवाञ्शीलवान्दक्षः स्वाध्यायनिरतः शुचिः}
{सर्वभोगेष्वसंसक्तो नित्यं धर्मपरायणः} %||7-3-2||

\twolineshloka
{ज्ञात्वा तस्य तु तद्वृत्तं भरद्वाजो महानृषिः}
{ददौ विश्रवसे भार्यां स्वां सुतां देववर्णिनीम्} %||7-3-3||

\twolineshloka
{प्रतिगृह्य तु धर्मेण भरद्वाजसुतां तदा}
{मुदा परमया युक्तो विश्रवा मुनिपुङ्गवः} %||7-3-4||

\twolineshloka
{स तस्यां वीर्यसम्पन्नमपत्यं परमाद्भुतम्}
{जनयामास धर्मात्मा सर्वैर्ब्रह्मगुणैर्युतम्} %||7-3-5||

\twolineshloka
{तस्मिञ्जाते तु संहृष्टः स बभूव पितामहः}
{नाम चास्याकरोत्प्रीतः सार्धं देवर्षिभिस्तदा} %||7-3-6||

\twolineshloka
{यस्माद्विश्रवसोऽपत्यं सादृश्याद्विश्रवा इव}
{तस्माद्वैश्रवणो नाम भविष्यत्येष विश्रुतः} %||7-3-7||

\twolineshloka
{स तु वैश्रवणस्तत्र तपोवनगतस्तदा}
{अवर्धत महातेजा हुताहुतिरिवानलः} %||7-3-8||

\twolineshloka
{तस्याश्रमपदस्थस्य बुद्धिर्जज्ञे महात्मनः}
{चरिष्ये नियतो धर्मं धर्मो हि परमा गतिः} %||7-3-9||

\twolineshloka
{स तु वर्षसहस्राणि तपस्तप्त्वा महावने}
{पूर्णे वर्षसहस्रे तु तं तं विधिमवर्तत} %||7-3-10||

\twolineshloka
{जलाशी मारुताहारो निराहारस्तथैव च}
{एवं वर्षसहस्राणि जग्मुस्तान्येव वर्षवत्} %||7-3-11||

\twolineshloka
{अथ प्रीतो महातेजाः सेन्द्रैः सुरगणैः सह}
{गत्वा तस्याश्रमपदं ब्रह्मेदं वाक्यमब्रवीत्} %||7-3-12||

\twolineshloka
{परितुष्टोऽस्मि ते वत्स कर्मणानेन सुव्रत}
{वरं वृणीष्व भद्रं ते वरार्हस्त्वं हि मे मतः} %||7-3-13||

\twolineshloka
{अथाब्रवीद्वैश्रवणः पितामहमुपस्थितम्}
{भगवँल्लोकपालत्वमिच्छेयं वित्तरक्षणम्} %||7-3-14||

\twolineshloka
{ततोऽब्रवीद्वैश्रवणं परितुष्टेन चेतसा}
{ब्रह्मा सुरगणैः सार्धं बाढमित्येव हृष्टवत्} %||7-3-15||

\twolineshloka
{अहं हि लोकपालानां चतुर्थं स्रष्टुमुद्यतः}
{यमेन्द्रवरुणानां हि पदं यत्तव चेप्सितम्} %||7-3-16||

\twolineshloka
{तत्कृतं गच्छ धर्मज्ञ धनेशत्वमवाप्नुहि}
{यमेन्द्रवरुणानां हि चतुर्थोऽद्य भविष्यसि} %||7-3-17||

\twolineshloka
{एतच्च पुष्पकं नाम विमानं सूर्यसंनिभम्}
{प्रतिगृह्णीष्व यानार्थं त्रिदशैः समतां व्रज} %||7-3-18||

\twolineshloka
{स्वस्ति तेऽस्तु गमिष्यामः सर्व एव यथागतम्}
{कृतकृत्या वयं तात दत्त्वा तव महावरम्} %||7-3-19||

\twolineshloka
{गतेषु ब्रह्मपूर्वेषु देवेष्वथ नभस्तलम्}
{धनेशः पितरं प्राह विनयात्प्रणतो वचः} %||7-3-20||

\twolineshloka
{भगवँल्लब्धवानस्मि वरं कमलयोनितः}
{निवासं न तु मे देवो विदधे स प्रजापतिः} %||7-3-21||

\twolineshloka
{तत्पश्य भगवन्कञ्चिद्देशं वासाय नः प्रभो}
{न च पीडा भवेद्यत्र प्राणिनो यस्य कस्यचित्} %||7-3-22||

\twolineshloka
{एवमुक्तस्तु पुत्रेण विश्रवा मुनिपुङ्गवः}
{वचनं प्राह धर्मज्ञ श्रूयतामिति धर्मवित्} %||7-3-23||

\twolineshloka
{लङ्का नाम पुरी रम्या निर्मिता विश्वकर्मणा}
{राक्षसानां निवासार्थं यथेन्द्रस्यामरावती} %||7-3-24||

\threelineshloka
{रमणीया पुरी सा हि रुक्मवैदूर्यतोरणा}
{राक्षसैः सा परित्यक्ता पुरा विष्णुभयार्दितैः}
{शून्या रक्षोगणैः सर्वै रसातलतलं गतैः} %||7-3-25||

\twolineshloka
{स त्वं तत्र निवासाय रोचयस्व मतिं स्वकाम्}
{निर्दोषस्तत्र ते वासो न च बाधास्ति कस्यचित्} %||7-3-26||

\twolineshloka
{एतच्छ्रुत्वा तु धर्मात्मा धर्मिष्ठं वचनं पितुः}
{निवेशयामास तदा लङ्कां पर्वतमूर्धनि} %||7-3-27||

\twolineshloka
{नैरृतानां सहस्रैस्तु हृष्टैः प्रमुदितैः सदा}
{अचिरेणैककालेन सम्पूर्णा तस्य शासनात्} %||7-3-28||

\twolineshloka
{अथ तत्रावसत्प्रीतो धर्मात्मा नैरृताधिपः}
{समुद्रपरिधानायां लङ्कायां विश्रवात्मजः} %||7-3-29||

\twolineshloka
{काले काले विनीतात्मा पुष्पकेण धनेश्वरः}
{अभ्यगच्छत्सुसंहृष्टः पितरं मातरं च सः} %||7-3-30||

\fourlineindentedshloka
{स देवगन्धर्वगणैरभिष्टुत}
{स्तथैव सिद्धैः सह चारणैरपि}
{गभस्तिभिः सूर्य इवौजसा वृतः}
{ पितुः समीपं प्रययौ श्रिया वृतः} %||7-4-31||


{॥इत्यार्षे श्रीमद्रामायणे वाल्मीकीये आदिकाव्ये उत्तरकाण्डे तृतीयः सर्गः॥}

\sect{चतुर्थः सर्गः}

\twolineshloka
{श्रुत्वागस्त्येरितं वाक्यं रामो विस्मयमागतः}
{पूर्वमासीत्तु लङ्कायां रक्षसामिति सम्भवः} %||7-4-1||

\twolineshloka
{ततः शिरः कम्पयित्वा त्रेताग्निसमविग्रहम्}
{अगस्त्यं तं मुहुर्दृष्ट्वा स्मयमानोऽभ्यभाषत} %||7-4-2||

\twolineshloka
{भगवन्पूर्वमप्येषा लङ्कासीत्पिशिताशिनाम्}
{इतीदं भवतः श्रुत्वा विस्मयो जनितो मम} %||7-4-3||

\twolineshloka
{पुलस्त्यवंशादुद्भूता राक्षसा इति नः श्रुतम्}
{इदानीमन्यतश्चापि सम्भवः कीर्तितस्त्वया} %||7-4-4||

\twolineshloka
{रावणात्कुम्भकर्णाच्च प्रहस्ताद्विकटादपि}
{रावणस्य च पुत्रेभ्यः किं नु ते बलवत्तराः} %||7-4-5||

\twolineshloka
{क एषां पूर्वको ब्रह्मन्किंनामा किन्तपोबलः}
{अपराधं च कं प्राप्य विष्णुना द्राविताः पुरा} %||7-4-6||

\twolineshloka
{एतद्विस्तरतः सर्वं कथयस्व ममानघ}
{कौतूहलं कृतं मह्यं नुद भानुर्यथा तमः} %||7-4-7||

\twolineshloka
{राघवस्य तु तच्छ्रुत्वा संस्कारालङ्कृतं वचः}
{ईषद्विस्मयमानस्तमगस्त्यः प्राह राघवम्} %||7-4-8||

\twolineshloka
{प्रजापतिः पुरा सृष्ट्वा अपः सलिलसम्भवः}
{तासां गोपायने सत्त्वानसृजत्पद्मसम्भवः} %||7-4-9||

\twolineshloka
{ते सत्त्वाः सत्त्वकर्तारं विनीतवदुपस्थिताः}
{किं कुर्म इति भाषन्तः क्षुत्पिपासाभयार्दिताः} %||7-4-10||

\twolineshloka
{प्रजापतिस्तु तान्याह सत्त्वानि प्रहसन्निव}
{आभाष्य वाचा यत्नेन रक्षध्वमिति मानदः} %||7-4-11||

\twolineshloka
{रक्षाम इति तत्रान्यैर्यक्षामेति तथापरैः}
{भुङ्क्षिताभुङ्क्षितैरुक्तस्ततस्तानाह भूतकृत्} %||7-4-12||

\twolineshloka
{रक्षाम इति यैरुक्तं राक्षसास्ते भवन्तु वः}
{यक्षाम इति यैरुक्तं ते वै यक्षा भवन्तु वः} %||7-4-13||

\twolineshloka
{तत्र हेतिः प्रहेतिश्च भ्रातरौ राक्षसर्षभौ}
{मधुकैटभसङ्काशौ बभूवतुररिन्दमौ} %||7-4-14||

\twolineshloka
{प्रहेतिर्धार्मिकस्तत्र न दारान्सोऽभिकाङ्क्षति}
{हेतिर्दारक्रियार्थं तु यत्नं परमथाकरोत्} %||7-4-15||

\twolineshloka
{स कालभगिनीं कन्यां भयां नाम भयावहाम्}
{उदावहदमेयात्मा स्वयमेव महामतिः} %||7-4-16||

\twolineshloka
{स तस्यां जनयामास हेती राक्षसपुङ्गवः}
{पुत्रं पुत्रवतां श्रेष्ठो विद्युत्केश इति श्रुतम्} %||7-4-17||

\twolineshloka
{विद्युत्केशो हेतिपुत्रः प्रदीप्ताग्निसमप्रभः}
{व्यवर्धत महातेजास्तोयमध्य इवाम्बुजम्} %||7-4-18||

\twolineshloka
{स यदा यौवनं भद्रमनुप्राप्तो निशाचरः}
{ततो दारक्रियां तस्य कर्तुं व्यवसितः पिता} %||7-4-19||

\twolineshloka
{सन्ध्यादुहितरं सोऽथ सन्ध्यातुल्यां प्रभावतः}
{वरयामास पुत्रार्थं हेती राक्षसपुङ्गवः} %||7-4-20||

\twolineshloka
{अवश्यमेव दातव्या परस्मै सेति सन्ध्यया}
{चिन्तयित्वा सुता दत्ता विद्युत्केशाय राघव} %||7-4-21||

\twolineshloka
{सन्ध्यायास्तनयां लब्ध्वा विद्युत्केशो निशाचरः}
{रमते स तया सार्धं पौलोम्या मघवानिव} %||7-4-22||

\twolineshloka
{केनचित्त्वथ कालेन राम सालकटङ्कटा}
{विद्युत्केशाद्गर्भमाप घनराजिरिवार्णवात्} %||7-4-23||

\twolineshloka
{ततः सा राक्षसी गर्भं घनगर्भसमप्रभम्}
{प्रसूता मन्दरं गत्वा गङ्गा गर्भमिवाग्निजम्} %||7-4-24||

\twolineshloka
{तमुत्सृज्य तु सा गर्भं विद्युत्केशाद्रतार्थिनी}
{रेमे सा पतिना सार्धं विस्मृत्य सुतमात्मजम्} %||7-4-25||

\twolineshloka
{तयोत्सृष्टः स तु शिशुः शरदर्कसमद्युतिः}
{पाणिमास्ये समाधाय रुरोद घनराडिव} %||7-4-26||

\twolineshloka
{अथोपरिष्टाद्गच्छन्वै वृषभस्थो हरः प्रभुः}
{अपश्यदुमया सार्धं रुदन्तं राक्षसात्मजम्} %||7-4-27||

\twolineshloka
{कारुण्यभावात्पार्वत्या भवस्त्रिपुरहा ततः}
{तं राक्षसात्मजं चक्रे मातुरेव वयः समम्} %||7-4-28||

\twolineshloka
{अमरं चैव तं कृत्वा महादेवोऽक्षयोऽव्ययः}
{पुरमाकाशगं प्रादात्पार्वत्याः प्रियकाम्यया} %||7-4-29||

\threelineshloka
{उमयापि वरो दत्तो राक्षसीनां नृपात्मज}
{सद्योपलब्धिर्गर्भस्य प्रसूतिः सद्य एव च}
{सद्य एव वयःप्राप्तिर्मातुरेव वयः समम्} %||7-4-30||

\fourlineindentedshloka
{ततः सुकेशो वरदानगर्वितः}
{ श्रियं प्रभोः प्राप्य हरस्य पार्श्वतः}
{चचार सर्वत्र महामतिः खगः}
{ खगं पुरं प्राप्य पुरन्दरो यथा} %||7-5-31||


{॥इत्यार्षे श्रीमद्रामायणे वाल्मीकीये आदिकाव्ये उत्तरकाण्डे चतुर्थः सर्गः॥}

\sect{पञ्चमः सर्गः}

\twolineshloka
{सुकेशं धार्मिकं दृष्ट्वा वरलब्धं च राक्षसम्}
{ग्रामणीर्नाम गन्धर्वो विश्वावसुसमप्रभः} %||7-5-1||

\twolineshloka
{तस्य देववती नाम द्वितीया श्रीरिवात्मजा}
{तां सुकेशाय धर्मेण ददौ दक्षः श्रियं यथा} %||7-5-2||

\twolineshloka
{वरदानकृतैश्वर्यं सा तं प्राप्य पतिं प्रियम्}
{आसीद्देववती तुष्टा धनं प्राप्येव निर्धनः} %||7-5-3||

\twolineshloka
{स तया सह संयुक्तो रराज रजनीचरः}
{अञ्जनादभिनिष्क्रान्तः करेण्वेव महागजः} %||7-5-4||

\threelineshloka
{देववत्यां सुकेशस्तु जनयामास राघव}
{त्रींस्त्रिनेत्रसमान्पुत्रान्राक्षसान्राक्षसाधिपः}
{माल्यवन्तं सुमालिं च मालिं च बलिनां वरम्} %||7-5-5||

\twolineshloka
{त्रयो लोका इवाव्यग्राः स्थितास्त्रय इवाग्नयः}
{त्रयो मन्त्रा इवात्युग्रास्त्रयो घोरा इवामयाः} %||7-5-6||

\twolineshloka
{त्रयः सुकेशस्य सुतास्त्रेताग्निसमवर्चसः}
{विवृद्धिमगमंस्तत्र व्याधयोपेक्षिता इव} %||7-5-7||

\twolineshloka
{वरप्राप्तिं पितुस्ते तु ज्ञात्वैश्वर्यं ततो महत्}
{तपस्तप्तुं गता मेरुं भ्रातरः कृतनिश्चयाः} %||7-5-8||

\twolineshloka
{प्रगृह्य नियमान्घोरान्राक्षसा नृपसत्तम}
{विचेरुस्ते तपो घोरं सर्वभूतभयावहम्} %||7-5-9||

\twolineshloka
{सत्यार्जवदमोपेतैस्तपोभिर्भुवि दुष्करैः}
{सन्तापयन्तस्त्रीँल्लोकान्सदेवासुरमानुषान्} %||7-5-10||

\twolineshloka
{ततो विभुश्चतुर्वक्त्रो विमानवरमास्थितः}
{सुकेशपुत्रानामन्त्र्य वरदोऽस्मीत्यभाषत} %||7-5-11||

\twolineshloka
{ब्रह्माणं वरदं ज्ञात्वा सेन्द्रैर्देवगणैर्वृतम्}
{ऊचुः प्राञ्जलयः सर्वे वेपमाना इव द्रुमाः} %||7-5-12||

\threelineshloka
{तपसाराधितो देव यदि नो दिशसे वरम्}
{अजेयाः शत्रुहन्तारस्तथैव चिरजीविनः}
{प्रभविष्णवो भवामेति परस्परमनुव्रताः} %||7-5-13||

\twolineshloka
{एवं भविष्यतीत्युक्त्वा सुकेशतनयान्प्रभुः}
{प्रययौ ब्रह्मलोकाय ब्रह्मा ब्राह्मणवत्सलः} %||7-5-14||

\twolineshloka
{वरं लब्ध्वा ततः सर्वे राम रात्रिञ्चरास्तदा}
{सुरासुरान्प्रबाधन्ते वरदानात्सुनिर्भयाः} %||7-5-15||

\twolineshloka
{तैर्वध्यमानास्त्रिदशाः सर्षिसङ्घाः सचारणाः}
{त्रातारं नाधिगच्छन्ति निरयस्था यथा नराः} %||7-5-16||

\twolineshloka
{अथ ते विश्वकर्माणं शिल्पिनां वरमव्ययम्}
{ऊचुः समेत्य संहृष्टा राक्षसा रघुसत्तम} %||7-5-17||

\twolineshloka
{गृहकर्ता भवानेव देवानां हृदयेप्सितम्}
{अस्माकमपि तावत्त्वं गृहं कुरु महामते} %||7-5-18||

\twolineshloka
{हिमवन्तं समाश्रित्य मेरुं मन्दरमेव वा}
{महेश्वरगृहप्रख्यं गृहं नः क्रियतां महत्} %||7-5-19||

\twolineshloka
{विश्वकर्मा ततस्तेषां राक्षसानां महाभुजः}
{निवासं कथयामास शक्रस्येवामरावतीम्} %||7-5-20||

\threelineshloka
{दक्षिणस्योदधेस्तीरे त्रिकूटो नाम पर्वतः}
{शिखरे तस्य शैलस्य मध्यमेऽम्बुदसंनिभे}
{शकुनैरपि दुष्प्रापे टङ्कच्छिन्नचतुर्दिशि} %||7-5-21||

\twolineshloka
{त्रिंशद्योजनविस्तीर्णा स्वर्णप्राकारतोरणा}
{मया लङ्केति नगरी शक्राज्ञप्तेन निर्मिता} %||7-5-22||

\twolineshloka
{तस्यां वसत दुर्धर्षाः पुर्यां राक्षससत्तमाः}
{अमरावतीं समासाद्य सेन्द्रा इव दिवौकसः} %||7-5-23||

\twolineshloka
{लङ्कादुर्गं समासाद्य राक्षसैर्बहुभिर्वृताः}
{भविष्यथ दुराधर्षाः शत्रूणां शत्रुसूदनाः} %||7-5-24||

\twolineshloka
{विश्वकर्मवचः श्रुत्वा ततस्ते राम राक्षसाः}
{सहस्रानुचरा गत्वा लङ्कां तामवसन्पुरीम्} %||7-5-25||

\twolineshloka
{दृढप्राकारपरिखां हैमैर्गृहशतैर्वृताम्}
{लङ्कामवाप्य ते हृष्टा विहरन्ति निशाचराः} %||7-5-26||

\twolineshloka
{नर्मदा नाम गन्धर्वी नानाधर्मसमेधिता}
{तस्याः कन्यात्रयं ह्यासीद्धीश्रीकीर्तिसमद्युति} %||7-5-27||

\twolineshloka
{ज्येष्ठक्रमेण सा तेषां राक्षसानामराक्षसी}
{कन्यास्ताः प्रददौ हृष्टा पूर्णचन्द्रनिभाननाः} %||7-5-28||

\twolineshloka
{त्रयाणां राक्षसेन्द्राणां तिस्रो गन्धर्वकन्यकाः}
{मात्रा दत्ता महाभागा नक्षत्रे भगदैवते} %||7-5-29||

\twolineshloka
{कृतदारास्तु ते राम सुकेशतनयाः प्रभो}
{भार्याभिः सह चिक्रीडुरप्सरोभिरिवामराः} %||7-5-30||

\twolineshloka
{तत्र माल्यवतो भार्या सुन्दरी नाम सुन्दरी}
{स तस्यां जनयामास यदपत्यं निबोध तत्} %||7-5-31||

\threelineshloka
{वज्रमुष्टिर्विरूपाक्षो दुर्मुखश्चैव राक्षसः}
{सुप्तघ्नो यज्ञकोपश्च मत्तोन्मत्तौ तथैव च}
{अनला चाभवत्कन्या सुन्दर्यां राम सुन्दरी} %||7-5-32||

\twolineshloka
{सुमालिनोऽपि भार्यासीत्पूर्णचन्द्रनिभानना}
{नाम्ना केतुमती नाम प्राणेभ्योऽपि गरीयसी} %||7-5-33||

\twolineshloka
{सुमाली जनयामास यदपत्यं निशाचरः}
{केतुमत्यां महाराज तन्निबोधानुपूर्वशः} %||7-5-34||

\twolineshloka
{प्रहस्तोऽकम्पनैश्चैव विकटः कालकार्मुकः}
{धूम्राक्शश्चाथ दण्डश्च सुपार्श्वश्च महाबलः} %||7-5-35||

\threelineshloka
{संह्रादिः प्रघसश्चैव भासकर्णश्च राक्षसः}
{राका पुष्पोत्कटा चैव कैकसी च शुचिस्मिता}
{कुम्भीनसी च इत्येते सुमालेः प्रसवाः स्मृताः} %||7-5-36||

\twolineshloka
{मालेस्तु वसुदा नाम गन्धर्वी रूपशालिनी}
{भार्यासीत्पद्मपत्राक्षी स्वक्षी यक्षीवरोपमा} %||7-5-37||

\twolineshloka
{सुमालेरनुजस्तस्यां जनयामास यत्प्रभो}
{अपत्यं कथ्यमानं तन्मया त्वं शृणु राघव} %||7-5-38||

\twolineshloka
{अनलश्चानिलश्चैव हरः सम्पातिरेव च}
{एते विभीषणामात्या मालेयास्ते निशाचराः} %||7-5-39||

\fourlineindentedshloka
{ततस्तु ते राक्षसपुङ्गवास्त्रयो}
{ निशाचरैः पुत्रशतैश्च संवृताः}
{सुरान्सहेन्द्रानृषिनागदानवा}
{न्बबाधिरे ते बलवीर्यदर्पिताः} %||7-5-40||

\fourlineindentedshloka
{जगद्भ्रमन्तोऽनिलवद्दुरासदा}
{ रणे च मृत्युप्रतिमाः समाहिताः}
{वरप्रदानादभिगर्विता भृशम्}
{ क्रतुक्रियाणां प्रशमङ्कराः सदा} %||7-6-41||


{॥इत्यार्षे श्रीमद्रामायणे वाल्मीकीये आदिकाव्ये उत्तरकाण्डे पञ्चमः सर्गः॥}

\sect{षष्ठः सर्गः}

\twolineshloka
{तैर्वध्यमाना देवाश्च ऋषयश्च तपोधनाः}
{भयार्ताः शरणं जग्मुर्देवदेवं महेश्वरम्} %||7-6-1||

\twolineshloka
{ते समेत्य तु कामारिं त्रिपुरारिं त्रिलोचनम्}
{ऊचुः प्राञ्जलयो देवा भयगद्गदभाषिणः} %||7-6-2||

\twolineshloka
{सुकेशपुत्रैर्भगवन्पितामहवरोद्धतैः}
{प्रजाध्यक्ष प्रजाः सर्वा बाध्यन्ते रिपुबाधन} %||7-6-3||

\twolineshloka
{शरण्यान्यशरण्यानि आश्रमाणि कृतानि नः}
{स्वर्गाच्च च्यावितः शक्रः स्वर्गे क्रीडन्ति शक्रवत्} %||7-6-4||

\twolineshloka
{अहं विष्णुरहं रुद्रो ब्रह्माहं देवराडहम्}
{अहं यमोऽहं वरुणश्चन्द्रोऽहं रविरप्यहम्} %||7-6-5||

\twolineshloka
{इति ते राक्षसा देव वरदानेन दर्पिताः}
{बाधन्ते समरोद्धर्षा ये च तेषां पुरःसराः} %||7-6-6||

\twolineshloka
{तन्नो देवभयार्तानामभयं दातुमर्हसि}
{अशिवं वपुरास्थाय जहि दैवतकण्टकान्} %||7-6-7||

\twolineshloka
{इत्युक्तस्तु सुरैः सर्वैः कपर्दी नीललोहितः}
{सुकेशं प्रति सापेक्ष आह देवगणान्प्रभुः} %||7-6-8||

\twolineshloka
{नाहं तान्निहनिष्यामि अवध्या मम तेऽसुराः}
{किं तु मन्त्रं प्रदास्यामि यो वै तान्निहनिष्यति} %||7-6-9||

\twolineshloka
{एवमेव समुद्योगं पुरस्कृत्य सुरर्षभाः}
{गच्छन्तु शरणं विष्णुं हनिष्यति स तान्प्रभुः} %||7-6-10||

\twolineshloka
{ततस्ते जयशब्देन प्रतिनन्द्य महेश्वरम्}
{विष्णोः समीपमाजग्मुर्निशाचरभयार्दिताः} %||7-6-11||

\twolineshloka
{शङ्खचक्रधरं देवं प्रणम्य बहुमान्य च}
{ऊचुः सम्भ्रान्तवद्वाक्यं सुकेशतनयार्दिताः} %||7-6-12||

\twolineshloka
{सुकेशतनयैर्देवत्रिभिस्त्रेताग्निसंनिभैः}
{आक्रम्य वरदानेन स्थानान्यपहृतानि नः} %||7-6-13||

\twolineshloka
{लङ्का नाम पुरी दुर्गा त्रिकूटशिखरे स्थिता}
{तत्र स्थिताः प्रबाधन्ते सर्वान्नः क्षणदाचराः} %||7-6-14||

\twolineshloka
{स त्वमस्मत्प्रियार्थं तु जहि तान्मधुसूदन}
{चक्रकृत्तास्यकमलान्निवेदय यमाय वै} %||7-6-15||

\twolineshloka
{भयेष्वभयदोऽस्माकं नान्योऽस्ति भवता समः}
{नुद त्वं नो भयं देव नीहारमिव भास्करः} %||7-6-16||

\twolineshloka
{इत्येवं दैवतैरुक्तो देवदेवो जनार्दनः}
{अभयं भयदोऽरीणां दत्त्वा देवानुवाच ह} %||7-6-17||

\twolineshloka
{सुकेशं राक्षसं जाने ईशान वरदर्पितम्}
{तांश्चास्य तनयाञ्जाने येषां ज्येष्ठः स माल्यवान्} %||7-6-18||

\twolineshloka
{तानहं समतिक्रान्तमर्यादान्राक्षसाधमान्}
{सूदयिष्यामि सङ्ग्रामे सुरा भवत विज्वराः} %||7-6-19||

\twolineshloka
{इत्युक्तास्ते सुराः सर्वे विष्णुना प्रभविष्णुना}
{यथा वासं ययुर्हृष्टाः प्रशमन्तो जनार्दनम्} %||7-6-20||

\twolineshloka
{विबुधानां समुद्योगं माल्यवान्स निशाचरः}
{श्रुत्वा तौ भ्रातरौ वीराविदं वचनमब्रवीत्} %||7-6-21||

\twolineshloka
{अमरा ऋषयश्चैव संहत्य किल शङ्करम्}
{अस्मद्वधं परीप्सन्त इदमूचुस्त्रिलोचनम्} %||7-6-22||

\twolineshloka
{सुकेशतनया देव वरदानबलोद्धताः}
{बाधन्तेऽस्मान्समुद्युक्ता घोररूपाः पदे पदे} %||7-6-23||

\twolineshloka
{राक्षसैरभिभूताः स्म न शक्ताः स्म उमापते}
{स्वेषु वेश्मसु संस्थातुं भयात्तेषां दुरात्मनाम्} %||7-6-24||

\twolineshloka
{तदस्माकं हितार्थे त्वं जहि तांस्तांस्त्रिलोचन}
{राक्षसान्हुङ्कृतेनैव दह प्रदहतां वर} %||7-6-25||

\twolineshloka
{इत्येवं त्रिदशैरुक्तो निशम्यान्धकसूदनः}
{शिरः करं च धुन्वान इदं वचनमब्रवीत्} %||7-6-26||

\twolineshloka
{अवध्या मम ते देवाः सुकेशतनया रणे}
{मन्त्रं तु वः प्रदास्यामि यो वै तान्निहनिष्यति} %||7-6-27||

\twolineshloka
{यः स चक्रगदापाणिः पीतवासा जनार्दनः}
{हनिष्यति स तान्युद्धे शरणं तं प्रपद्यथ} %||7-6-28||

\twolineshloka
{हरान्नावाप्य ते कामं कामारिमभिवाद्य च}
{नारायणालयं प्राप्तास्तस्मै सर्वं न्यवेदयन्} %||7-6-29||

\twolineshloka
{ततो नारायणेनोक्ता देवा इन्द्रपुरोगमाः}
{सुरारीन्सूदयिष्यामि सुरा भवत विज्वराः} %||7-6-30||

\twolineshloka
{देवानां भयभीतानां हरिणा राक्षसर्षभौ}
{प्रतिज्ञातो वधोऽस्माकं तच्चिन्तयथ यत्क्षमम्} %||7-6-31||

\twolineshloka
{हिरण्यकशिपोर्मृत्युरन्येषां च सुरद्विषाम्}
{दुःखं नारायणं जेतुं यो नो हन्तुमभीप्सति} %||7-6-32||

\twolineshloka
{ततः सुमाली माली च श्रुत्वा माल्यवतो वचः}
{ऊचतुर्भ्रातरं ज्येष्ठं भगांशाविव वासवम्} %||7-6-33||

\twolineshloka
{स्वधीतं दत्तमिष्टं च ऐश्वर्यं परिपालितम्}
{आयुर्निरामयं प्राप्तं स्वधर्मः स्थापितश्च नः} %||7-6-34||

\twolineshloka
{देवसागरमक्षोभ्यं शस्त्रौघैः प्रविगाह्य च}
{जिता देवा रणे नित्यं न नो मृत्युकृतं भयम्} %||7-6-35||

\twolineshloka
{नारायणश्च रुद्रश्च शक्रश्चापि यमस्तथा}
{अस्माकं प्रमुखे स्थातुं सर्व एव हि बिभ्यति} %||7-6-36||

\twolineshloka
{विष्णोर्दोषश्च नास्त्यत्र कारणं राक्षसेश्वर}
{देवानामेव दोषेण विष्णोः प्रचलितं मनः} %||7-6-37||

\twolineshloka
{तस्मादद्य समुद्युक्ताः सर्वसैन्यसमावृताः}
{देवानेव जिघांसामो येभ्यो दोषः समुत्थितः} %||7-6-38||

\threelineshloka
{इति माली सुमाली च माल्यवानग्रजः प्रभुः}
{उद्योगं घोषयित्वाथ राक्षसाः सर्व एव ते}
{युद्धाय निर्ययुः क्रुद्धा जम्भवृत्रबला इव} %||7-6-39||

\twolineshloka
{स्यन्दनैर्वारणेन्द्रैश्च हयैश्च गिरिसंनिभैः}
{खरैर्गोभिरथोष्ट्रैश्च शिंशुमारैर्भुजं गमैः} %||7-6-40||

\twolineshloka
{मकरैः कच्छपैर्मीनैर्विहङ्गैर्गरुडोपमैः}
{सिंहैर्व्याघ्रैर्वराहैश्च सृमरैश्चमरैरपि} %||7-6-41||

\twolineshloka
{त्यक्त्वा लङ्कां ततः सर्वे राक्षसा बलगर्विताः}
{प्रयाता देवलोकाय योद्धुं दैवतशत्रवः} %||7-6-42||

\twolineshloka
{लङ्काविपर्ययं दृष्ट्वा यानि लङ्कालयान्यथ}
{भूतानि भयदर्शीनि विमनस्कानि सर्वशः} %||7-6-43||

\twolineshloka
{भौमास्तथान्तरिक्षाश्च कालाज्ञप्ता भयावहाः}
{उत्पाता राक्षसेन्द्राणामभावायोत्थिता द्रुतम्} %||7-6-44||

\twolineshloka
{अस्थीनि मेघा वर्षन्ति उष्णं शोणितमेव च}
{वेलां समुद्रोऽप्युत्क्रान्तश्चलन्ते चाचलोत्तमाः} %||7-6-45||

\twolineshloka
{अट्टहासान्विमुञ्चन्तो घननादसमस्वनान्}
{भूताः परिपतन्ति स्म नृत्यमानाः सहस्रशः} %||7-6-46||

\twolineshloka
{गृध्रचक्रं महच्चापि ज्वलनोद्गारिभिर्मुखैः}
{राक्षसानामुपरि वै भ्रमते कालचक्रवत्} %||7-6-47||

\twolineshloka
{तानचिन्त्यमहोत्पातान्राक्षसा बलगर्विताः}
{यन्त्येव न निवर्तन्ते मृत्युपाशावपाशिताः} %||7-6-48||

\twolineshloka
{माल्यवांश्च सुमाली च माली च रजनीचराः}
{आसन्पुरःसरास्तेषां क्रतूनामिव पावकाः} %||7-6-49||

\twolineshloka
{माल्यवन्तं तु ते सर्वे माल्यवन्तमिवाचलम्}
{निशाचरा आश्रयन्ते धातारमिव देहिनः} %||7-6-50||

\twolineshloka
{तद्बलं राक्षसेन्द्राणां महाभ्रघननादितम्}
{जयेप्सया देवलोकं ययौ माली वशे स्थितम्} %||7-6-51||

\twolineshloka
{राक्षसानां समुद्योगं तं तु नारायणः प्रभुः}
{देवदूतादुपश्रुत्य दध्रे युद्धे ततो मनः} %||7-6-52||

\fourlineindentedshloka
{स देवसिद्धर्षिमहोरगैश्च}
{ गन्धर्वमुख्याप्सरसोपगीतः}
{समाससादामरशत्रुसैन्यम्}
{ चक्रासिसीरप्रवरादिधारी} %||7-6-53||

\fourlineindentedshloka
{सुपर्णपक्षानिलनुन्नपक्षम्}
{ भ्रमत्पताकं प्रविकीर्णशस्त्रम्}
{चचाल तद्राक्षसराजसैन्यम्}
{ चलोपलो नील इवाचलेन्द्रः} %||7-6-54||

\fourlineindentedshloka
{तथ शितैः शोणितमांसरूषितैर्-}
{ युगान्तवैश्वानरतुल्यविग्रहैः}
{निशाचराः सम्परिवार्य माधवम्}
{ वरायुधैर्निर्बिभिदुः सहस्रशः} %||7-7-55||


{॥इत्यार्षे श्रीमद्रामायणे वाल्मीकीये आदिकाव्ये उत्तरकाण्डे षष्ठः सर्गः॥}

\sect{सप्तमः सर्गः}

\twolineshloka
{नारायणगिरिं ते तु गर्जन्तो राक्षसाम्बुदाः}
{अवर्षन्निषुवर्षेण वर्षेणाद्रिमिवाम्बुदाः} %||7-7-1||

\twolineshloka
{श्यामावदातस्तैर्विष्णुर्नीलैर्नक्तञ्चरोत्तमैः}
{वृतोऽञ्जनगिरीवासीद्वर्षमाणैः पयोधरैः} %||7-7-2||

\twolineshloka
{शलभा इव केदारं मशका इव पर्वतम्}
{यथामृतघटं जीवा मकरा इव चार्णवम्} %||7-7-3||

\twolineshloka
{तथा रक्षोधनुर्मुक्ता वज्रानिलमनोजवाः}
{हरिं विशन्ति स्म शरा लोकास्तमिव पर्यये} %||7-7-4||

\twolineshloka
{स्यन्दनैः स्यन्दनगता गजैश्च गजधूर्गताः}
{अश्वारोहाः सदश्वैश्च पादाताश्चाम्बरे चराः} %||7-7-5||

\twolineshloka
{राक्षसेन्द्रा गिरिनिभाः शरशक्त्यृष्टितोमरैः}
{निरुच्छ्वासं हरिं चक्रुः प्राणायाम इव द्विजम्} %||7-7-6||

\twolineshloka
{निशाचरैस्तुद्यमानो मीनैरिव महातिमिः}
{शार्ङ्गमायम्य गात्राणि राक्षसानां महाहवे} %||7-7-7||

\twolineshloka
{शरैः पूर्णायतोत्सृष्टैर्वज्रवक्त्रैर्मनोजवैः}
{चिच्छेद तिलशो विष्णुः शतशोऽथ सहस्रशः} %||7-7-8||

\twolineshloka
{विद्राव्य शरवर्षं तं वर्षं वायुरिवोत्थितम्}
{पाञ्चजन्यं महाशङ्खं प्रदध्मौ पुरुषोत्तमः} %||7-7-9||

\twolineshloka
{सोऽम्बुजो हरिणा ध्मातः सर्वप्राणेन शङ्खराट्}
{ररास भीमनिह्रादो युगान्ते जलदो यथा} %||7-7-10||

\twolineshloka
{शङ्खराजरवः सोऽथ त्रासयामास राक्षसान्}
{मृगराज इवारण्ये समदानिव कुञ्जरान्} %||7-7-11||

\twolineshloka
{न शेकुरश्वाः संस्थातुं विमदाः कुञ्जराभवन्}
{स्यन्दनेभ्यश्च्युता योधाः शङ्खरावितदुर्बलाः} %||7-7-12||

\twolineshloka
{शार्ङ्गचापविनिर्मुक्ता वज्रतुल्याननाः शराः}
{विदार्य तानि रक्षांसि सुपुङ्खा विविशुः क्षितिम्} %||7-7-13||

\twolineshloka
{भिद्यमानाः शरैश्चान्ये नारायणधनुश्च्युतैः}
{निपेतू राक्षसा भीमाः शैला वज्रहता इव} %||7-7-14||

\twolineshloka
{व्रणैर्व्रणकरारीणामधोक्षजशरोद्भवैः}
{असृक्क्षरन्ति धाराभिः स्वर्णधारामिवाचलाः} %||7-7-15||

\twolineshloka
{शङ्खराजरवश्चापि शार्ङ्गचापरवस्तथा}
{राक्षसानां रवांश्चापि ग्रसते वैष्णवो रवः} %||7-7-16||

\twolineshloka
{सूर्यादिव करा घोरा ऊर्मयः सागरादिव}
{पर्वतादिव नागेन्द्रा वार्योघा इव चाम्बुदात्} %||7-7-17||

\twolineshloka
{तथा बाणा विनिर्मुक्ताः शार्ङ्गान्नरायणेरिताः}
{निर्धावन्तीषवस्तूर्णं शतशोऽथ सहस्रशः} %||7-7-18||

\twolineshloka
{शरभेण यथा सिंहाः सिंहेन द्विरदा यथा}
{द्विरदेन यथा व्याघ्रा व्याघ्रेण द्वीपिनो यथा} %||7-7-19||

\twolineshloka
{द्वीपिना च यथा श्वानः शुना मार्जारका यथा}
{मार्जारेण यथा सर्पाः सर्पेण च यथाखवः} %||7-7-20||

\twolineshloka
{तथा ते राक्षसा युद्धे विष्णुना प्रभविष्णुना}
{द्रवन्ति द्राविताश्चैव शायिताश्च महीतले} %||7-7-21||

\twolineshloka
{राक्षसानां सहस्राणि निहत्य मधुसूदनः}
{वारिजं नादयामास तोयदं सुरराडिव} %||7-7-22||

\twolineshloka
{नारायणशरग्रस्तं शङ्खनादसुविह्वलम्}
{ययौ लङ्कामभिमुखं प्रभग्नं राक्षसं बलम्} %||7-7-23||

\twolineshloka
{प्रभग्ने राक्षसबले नारायणशराहते}
{सुमाली शरवर्षेण आववार रणे हरिम्} %||7-7-24||

\twolineshloka
{उत्क्षिप्य हेमाभरणं करं करमिव द्विपः}
{ररास राक्षसो हर्षात्सतडित्तोयदो यथा} %||7-7-25||

\twolineshloka
{सुमालेर्नर्दतस्तस्य शिरो ज्वलितकुण्डलम्}
{चिच्छेद यन्तुरश्वाश्च भ्रान्तास्तस्य तु रक्षसः} %||7-7-26||

\twolineshloka
{तैरश्वैर्भ्राम्यते भ्रान्तैः सुमाली राक्षसेश्वरः}
{इन्द्रियाश्वैर्यथा भ्रान्तैर्धृतिहीनो यथा नरः} %||7-7-27||

\threelineshloka
{माली चाभ्यद्रवद्युद्धे प्रगृह्य सशरं धनुः}
{मालेर्धनुश्च्युता बाणाः कार्तस्वरविभूषिताः}
{विविशुर्हरिमासाद्य क्रौञ्चं पत्ररथा इव} %||7-7-28||

\twolineshloka
{अर्द्यमानः शरैः सोऽथ मालिमुक्तैः सहस्रशः}
{चुक्षुभे न रणे विष्णुर्जितेन्द्रिय इवाधिभिः} %||7-7-29||

\twolineshloka
{अथ मौर्वी स्वनं कृत्वा भगवान्भूतभावनः}
{मालिनं प्रति बाणौघान्ससर्जासिगदाधरः} %||7-7-30||

\twolineshloka
{ते मालिदेहमासाद्य वज्रविद्युत्प्रभाः शराः}
{पिबन्ति रुधिरं तस्य नागा इव पुरामृतम्} %||7-7-31||

\twolineshloka
{मालिनं विमुखं कृत्वा मालिमौलिं हरिर्बलात्}
{रथं च सध्वजं चापं वाजिनश्च न्यपातयत्} %||7-7-32||

\twolineshloka
{विरथस्तु गदां गृह्य माली नक्तञ्चरोत्तमः}
{आपुप्लुवे गदापाणिर्गिर्यग्रादिव केषरी} %||7-7-33||

\twolineshloka
{स तया गरुडं सङ्ख्ये ईशानमिव चान्तकः}
{ललाटदेशेऽभ्यहनद्वज्रेणेन्द्रो यथाचलम्} %||7-7-34||

\twolineshloka
{गदयाभिहतस्तेन मालिना गरुडो भृशम्}
{रणात्पराङ्मुखं देवं कृतवान्वेदनातुरः} %||7-7-35||

\twolineshloka
{पराङ्मुखे कृते देवे मालिना गरुडेन वै}
{उदतिष्ठन्महानादो रक्षसामभिनर्दताम्} %||7-7-36||

\twolineshloka
{रक्षसां नदतां नादं श्रुत्वा हरिहयानुजः}
{पराङ्मुखोऽप्युत्ससर्ज चक्रं मालिजिघांसया} %||7-7-37||

\twolineshloka
{तत्सूर्यमण्डलाभासं स्वभासा भासयन्नभः}
{कालचक्रनिभं चक्रं मालेः शीर्षमपातयत्} %||7-7-38||

\twolineshloka
{तच्छिरो राक्षसेन्द्रस्य चक्रोत्कृत्तं विभीषणम्}
{पपात रुधिरोद्गारि पुरा राहुशिरो यथा} %||7-7-39||

\twolineshloka
{ततः सुरैः सुसंहृष्टैः सर्वप्राणसमीरितः}
{सिंहनादरवो मुक्तः साधु देवेति वादिभिः} %||7-7-40||

\twolineshloka
{मालिनं निहतं दृष्ट्वा सुमाली मल्यवानपि}
{सबलौ शोकसन्तप्तौ लङ्कां प्रति विधावितौ} %||7-7-41||

\twolineshloka
{गरुडस्तु समाश्वस्तः संनिवृत्य महामनाः}
{राक्षसान्द्रावयामास पक्षवातेन कोपितः} %||7-7-42||

\fourlineindentedshloka
{नारायणोऽपीषुवराशनीभिर्-}
{ विदारयामास धनुःप्रमुक्तैः}
{नक्तञ्चरान्मुक्तविधूतकेशा}
{न्यथाशनीभिः सतडिन्महेन्द्रः} %||7-7-43||

\fourlineindentedshloka
{भिन्नातपत्रं पतमानशस्त्रम्}
{ शरैरपध्वस्तविशीर्णदेहम्}
{विनिःसृतान्त्रं भयलोलनेत्रम्}
{ बलं तदुन्मत्तनिभं बभूव} %||7-7-44||

\fourlineindentedshloka
{सिंहार्दितानामिव कुञ्जराणाम्}
{ निशाचराणां सह कुञ्जराणाम्}
{रवाश्च वेगाश्च समं बभूवुः}
{ पुराणसिंहेन विमर्दितानाम्} %||7-7-45||

\fourlineindentedshloka
{सञ्छाद्यमाना हरिबाणजालैः}
{ स्वबाणजालानि समुत्सृजन्तः}
{धावन्ति नक्तञ्चरकालमेघा}
{ वायुप्रणुन्ना इव कालमेघाः} %||7-7-46||

\fourlineindentedshloka
{चक्रप्रहारैर्विनिकृत्तशीर्षाः}
{ सञ्चूर्णिताङ्गाश्च गदाप्रहारैः}
{असिप्रहारैर्बहुधा विभक्ताः}
{ पतन्ति शैला इव राक्षसेन्द्राः} %||7-7-47||

\twolineshloka
{चक्रकृत्तास्यकमला गदासञ्चूर्णितोरसः}
{लाङ्गलग्लपितग्रीवा मुसलैर्भिन्नमस्तकाः} %||7-7-48||

\twolineshloka
{केचिच्चैवासिना छिन्नास्तथान्ये शरताडिताः}
{निपेतुरम्बरात्तूर्णं राक्षसाः सागराम्भसि} %||7-7-49||

\fourlineindentedshloka
{तदाम्बरं विगलितहारकुण्डलैर्-}
{ निशाचरैर्नीलबलाहकोपमैः}
{निपात्यमानैर्ददृशे निरन्तरम्}
{ निपात्यमानैरिव नीलपर्वतैः} %||7-8-50||


{॥इत्यार्षे श्रीमद्रामायणे वाल्मीकीये आदिकाव्ये उत्तरकाण्डे सप्तमः सर्गः॥}

\sect{अष्टमः सर्गः}

\twolineshloka
{हन्यमाने बले तस्मिन्पद्मनाभेन पृष्ठतः}
{माल्यवान्संनिवृत्तोऽथ वेलातिग इवार्णवः} %||7-8-1||

\twolineshloka
{संरक्तनयनः कोपाच्चलन्मौलिर्निशाचरः}
{पद्मनाभमिदं प्राह वचनं परुषं तदा} %||7-8-2||

\twolineshloka
{नारायण न जानीषे क्षत्रधर्मं सनातनम्}
{अयुद्धमनसो भग्नान्योऽस्मान्हंसि यथेतरः} %||7-8-3||

\twolineshloka
{पराङ्मुखवधं पापं यः करोति सुरेश्वर}
{स हन्ता न गतः स्वर्गं लभते पुण्यकर्मणाम्} %||7-8-4||

\twolineshloka
{युद्धश्रद्धाथ वा तेऽस्ति शङ्खचक्रगदाधर}
{अहं स्थितोऽस्मि पश्यामि बलं दर्शय यत्तव} %||7-8-5||

\threelineshloka
{उवाच राक्षसेन्द्रं तं देवराजानुजो बली}
{युष्मत्तो भयभीतानां देवानां वै मयाभयम्}
{राक्षसोत्सादनं दत्तं तदेतदनुपाल्यते} %||7-8-6||

\twolineshloka
{प्राणैरपि प्रियं कार्यं देवानां हि सदा मया}
{सोऽहं वो निहनिष्यामि रसातलगतानपि} %||7-8-7||

\twolineshloka
{देवमेवं ब्रुवाणं तु रक्ताम्बुरुहलोचनम्}
{शक्त्या बिभेद सङ्क्रुद्धो राक्षसेन्द्रो ररास च} %||7-8-8||

\twolineshloka
{माल्यवद्भुजनिर्मुक्ता शक्तिर्घण्टाकृतस्वना}
{हरेरुरसि बभ्राज मेघस्थेव शतह्रदा} %||7-8-9||

\twolineshloka
{ततस्तामेव चोत्कृष्य शक्तिं शक्तिधरप्रियः}
{माल्यवन्तं समुद्दिश्य चिक्षेपाम्बुरुहेक्षणः} %||7-8-10||

\twolineshloka
{स्कन्दोत्सृष्टेव सा शक्तिर्गोविन्दकरनिःसृता}
{काङ्क्षन्ती राक्षसं प्रायान्महोल्केवाञ्जनाचलम्} %||7-8-11||

\twolineshloka
{सा तस्योरसि विस्तीर्णे हारभासावभासिते}
{अपतद्राक्षसेन्द्रस्य गिरिकूट इवाशनिः} %||7-8-12||

\twolineshloka
{तया भिन्नतनुत्राणाः प्राविशद्विपुलं तमः}
{माल्यवान्पुनराश्वस्तस्तस्थौ गिरिरिवाचलः} %||7-8-13||

\twolineshloka
{ततः कार्ष्णायसं शूलं कण्टकैर्बहुभिश्चितम्}
{प्रगृह्याभ्यहनद्देवं स्तनयोरन्तरे दृढम्} %||7-8-14||

\twolineshloka
{तथैव रणरक्तस्तु मुष्टिना वासवानुजम्}
{ताडयित्वा धनुर्मात्रमपक्रान्तो निशाचरः} %||7-8-15||

\twolineshloka
{ततोऽम्बरे महाञ्शब्दः साधु साध्विति चोत्थितः}
{आहत्य राक्षसो विष्णुं गरुडं चाप्यताडयत्} %||7-8-16||

\twolineshloka
{वैनतेयस्ततः क्रुद्धः पक्षवातेन राक्षसम्}
{व्यपोहद्बलवान्वायुः शुष्कपर्णचयं यथा} %||7-8-17||

\twolineshloka
{द्विजेन्द्रपक्षवातेन द्रावितं दृश्य पूर्वजम्}
{सुमाली स्वबलैः सार्धं लङ्कामभिमुखो ययौ} %||7-8-18||

\twolineshloka
{पक्षवातबलोद्धूतो माल्यवानपि राक्षसः}
{स्वबलेन समागम्य ययौ लङ्कां ह्रिया वृतः} %||7-8-19||

\twolineshloka
{एवं ते राक्षसा राम हरिणा कमलेक्षण}
{बहुशः संयुगे भग्ना हतप्रवरनायकाः} %||7-8-20||

\twolineshloka
{अशक्नुवन्तस्ते विष्णुं प्रतियोद्धुं भयार्दिताः}
{त्यक्त्वा लङ्कां गता वस्तुं पातालं सहपत्नयः} %||7-8-21||

\twolineshloka
{सुमालिनं समासाद्य राक्षसं रघुनन्दन}
{स्थिताः प्रख्यातवीर्यास्ते वंशे सालकटङ्कटे} %||7-8-22||

\threelineshloka
{ये त्वया निहतास्ते वै पौलस्त्या नाम राक्षसाः}
{सुमाली माल्यवान्माली ये च तेषां पुरःसराः}
{सर्व एते महाभाग रावणाद्बलवत्तराः} %||7-8-23||

\twolineshloka
{न चान्यो रक्षसां हन्ता सुरेष्वपि पुरंजय}
{ऋते नारायणं देवं शङ्खचक्रगदाधरम्} %||7-8-24||

\twolineshloka
{भवान्नारायणो देवश्चतुर्बाहुः सनातनः}
{राक्षसान्हन्तुमुत्पन्नो अजेयः प्रभुरव्ययः} %||7-9-25||


{॥इत्यार्षे श्रीमद्रामायणे वाल्मीकीये आदिकाव्ये उत्तरकाण्डे अष्टमः सर्गः॥}

\sect{नवमः सर्गः}

\twolineshloka
{कस्यचित्त्वथ कालस्य सुमाली नाम राक्षसः}
{रसातलान्मर्त्यलोकं सर्वं वै विचचार ह} %||7-9-1||

\threelineshloka
{नीलजीमूतसङ्काशस्तप्तकाञ्चनकुण्डलः}
{कन्यां दुहितरं गृह्य विना पद्ममिव श्रियम्}
{अथापश्यत्स गच्छन्तं पुष्पकेण धनेश्वरम्} %||7-9-2||

\twolineshloka
{तं दृष्ट्वामरसङ्काशं गच्छन्तं पावकोपमम्}
{अथाब्ब्रवीत्सुतां रक्षः कैकसीं नाम नामतः} %||7-9-3||

\twolineshloka
{पुत्रि प्रदानकालोऽयं यौवनं तेऽतिवर्तते}
{त्वत्कृते च वयं सर्वे यन्त्रिता धर्मबुद्धयः} %||7-9-4||

\twolineshloka
{त्वं हि सर्वगुणोपेता श्रीः सपद्मेव पुत्रिके}
{प्रत्याख्यानाच्च भीतैस्त्वं न वरैः प्रतिगृह्यसे} %||7-9-5||

\twolineshloka
{कन्यापितृत्वं दुःखं हि सर्वेषां मानकाङ्क्षिणाम्}
{न ज्ञायते च कः कन्यां वरयेदिति पुत्रिके} %||7-9-6||

\twolineshloka
{मातुः कुलं पितृकुलं यत्र चैव प्रदीयते}
{कुलत्रयं सदा कन्या संशये स्थाप्य तिष्ठति} %||7-9-7||

\twolineshloka
{सा त्वं मुनिवरश्रेष्ठं प्रजापतिकुलोद्भवम्}
{गच्छ विश्रवसं पुत्रि पौलस्त्यं वरय स्वयम्} %||7-9-8||

\twolineshloka
{ईदृशास्ते भविष्यन्ति पुत्राः पुत्रि न संशयः}
{तेजसा भास्करसमा यादृशोऽयं धनेश्वरः} %||7-9-9||

\twolineshloka
{एतस्मिन्नन्तरे राम पुलस्त्यतनयो द्विजः}
{अग्निहोत्रमुपातिष्ठच्चतुर्थ इव पावकः} %||7-9-10||

\twolineshloka
{सा तु तां दारुणां वेलामचिन्त्य पितृगौरवात्}
{उपसृत्याग्रतस्तस्य चरणाधोमुखी स्थिता} %||7-9-11||

\twolineshloka
{स तु तां वीक्ष्य सुश्रोणीं पूर्णचन्द्रनिभाननाम्}
{अब्रवीत्परमोदारो दीप्यमान इवौजसा} %||7-9-12||

\twolineshloka
{भद्रे कस्यासि दुहिता कुतो वा त्वमिहागता}
{किं कार्यं कस्य वा हेतोस्तत्त्वतो ब्रूहि शोभने} %||7-9-13||

\twolineshloka
{एवमुक्ता तु सा कन्या कृताञ्जलिरथाब्रवीत्}
{आत्मप्रभावेन मुने ज्ञातुमर्हसि मे मतम्} %||7-9-14||

\twolineshloka
{किं तु विद्धि हि मां ब्रह्मञ्शासनात्पितुरागताम्}
{कैकसी नाम नाम्नाहं शेषं त्वं ज्ञातुमर्हसि} %||7-9-15||

\twolineshloka
{स तु गत्वा मुनिर्ध्यानं वाक्यमेतदुवाच ह}
{विज्ञातं ते मया भद्रे कारणं यन्मनोगतम्} %||7-9-16||

\twolineshloka
{दारुणायां तु वेलायां यस्मात्त्वं मामुपस्थिता}
{शृणु तस्मात्सुतान्भद्रे यादृशाञ्जनयिष्यसि} %||7-9-17||

\twolineshloka
{दारुणान्दारुणाकारान्दारुणाभिजनप्रियान्}
{प्रसविष्यसि सुश्रोणि राक्षसान्क्रूरकर्मणः} %||7-9-18||

\twolineshloka
{सा तु तद्वचनं श्रुत्वा प्रणिपत्याब्रवीद्वचः}
{भगवन्नेदृशाः पुत्रास्त्वत्तोऽर्हा ब्रह्मयोनितः} %||7-9-19||

\twolineshloka
{अथाब्रवीन्मुनिस्तत्र पश्चिमो यस्तवात्मजः}
{मम वंशानुरूपश्च धर्मात्मा च भविष्यति} %||7-9-20||

\twolineshloka
{एवमुक्ता तु सा कन्या राम कालेन केनचित्}
{जनयामास बीभत्सं रक्षोरूपं सुदारुणम्} %||7-9-21||

\twolineshloka
{दशशीर्षं महादंष्ट्रं नीलाञ्जनचयोपमम्}
{ताम्रौष्ठं विंशतिभुजं महास्यं दीप्तमूर्धजम्} %||7-9-22||

\twolineshloka
{जातमात्रे ततस्तस्मिन्सज्वालकवलाः शिवाः}
{क्रव्यादाश्चापसव्यानि मण्डलानि प्रचक्रिरे} %||7-9-23||

\twolineshloka
{ववर्ष रुधिरं देवो मेघाश्च खरनिस्वनाः}
{प्रबभौ न च खे सूर्यो महोल्काश्चापतन्भुवि} %||7-9-24||

\twolineshloka
{अथ नामाकरोत्तस्य पितामहसमः पिता}
{दशशीर्षः प्रसूतोऽयं दशग्रीवो भविष्यति} %||7-9-25||

\twolineshloka
{तस्य त्वनन्तरं जातः कुम्भकर्णो महाबलः}
{प्रमाणाद्यस्य विपुलं प्रमाणं नेह विद्यते} %||7-9-26||

\twolineshloka
{ततः शूर्पणखा नाम संजज्ञे विकृतानना}
{विभीषणश्च धर्मात्मा कैकस्याः पश्चिमः सुतः} %||7-9-27||

\twolineshloka
{ते तु तत्र महारण्ये ववृधुः सुमहौजसः}
{तेषां क्रूरो दशग्रीवो लोकोद्वेगकरोऽभवत्} %||7-9-28||

\twolineshloka
{कुम्भकर्णः प्रमत्तस्तु महर्षीन्धर्मसंश्रितान्}
{त्रैलोक्यं त्रासयन्दुष्टो भक्षयन्विचचार ह} %||7-9-29||

\twolineshloka
{विभीषणस्तु धर्मात्मा नित्यं धर्मपथे स्थितः}
{स्वाध्यायनियताहार उवास नियतेन्द्रियः} %||7-9-30||

\twolineshloka
{अथ वित्तेश्वरो देवस्तत्र कालेन केनचित्}
{आगच्छत्पितरं द्रष्टुं पुष्पकेण महौजसम्} %||7-9-31||

\twolineshloka
{तं दृष्ट्वा कैकसी तत्र ज्वलन्तमिव तेजसा}
{आस्थाय राक्षसीं बुद्धिं दशग्रीवमुवाच ह} %||7-9-32||

\twolineshloka
{पुत्रवैश्रवणं पश्य भ्रातरं तेजसा वृतम्}
{भ्रातृभावे समे चापि पश्यात्मानं त्वमीदृशम्} %||7-9-33||

\twolineshloka
{दशग्रीव तथा यत्नं कुरुष्वामितविक्रम}
{यथा भवसि मे पुत्र शीघ्रं वैश्वरणोपमः} %||7-9-34||

\twolineshloka
{मातुस्तद्वचनं श्रुत्वा दशग्रीवः प्रतापवान्}
{अमर्षमतुलं लेभे प्रतिज्ञां चाकरोत्तदा} %||7-9-35||

\twolineshloka
{सत्यं ते प्रतिजानामि तुल्यो भ्रात्राधिकोऽपि वा}
{भविष्याम्यचिरान्मातः सन्तापं त्यज हृद्गतम्} %||7-9-36||

\threelineshloka
{ततः क्रोधेन तेनैव दशग्रीवः सहानुजः}
{प्राप्स्यामि तपसा काममिति कृत्वाध्यवस्य च}
{आगच्छदात्मसिद्ध्यर्थं गोकर्णस्याश्रमं शुभम्} %||7-10-37||


{॥इत्यार्षे श्रीमद्रामायणे वाल्मीकीये आदिकाव्ये उत्तरकाण्डे नवमः सर्गः॥}

\sect{दशमः सर्गः}

\twolineshloka
{अथाब्रवीद्द्विजं रामः कथं ते भ्रातरो वने}
{कीदृशं तु तदा ब्रह्मंस्तपश्चेरुर्महाव्रताः} %||7-10-1||

\twolineshloka
{अगस्त्यस्त्वब्रवीत्तत्र रामं प्रयत मानसम्}
{तांस्तान्धर्मविधींस्तत्र भ्रातरस्ते समाविशन्} %||7-10-2||

\twolineshloka
{कुम्भकर्णस्तदा यत्तो नित्यं धर्मपरायणः}
{तताप ग्रैष्मिके काले पञ्चस्वग्निष्ववस्थितः} %||7-10-3||

\twolineshloka
{वर्षे मेघोदकक्लिन्नो वीरासनमसेवत}
{नित्यं च शैशिरे काले जलमध्यप्रतिश्रयः} %||7-10-4||

\twolineshloka
{एवं वर्षसहस्राणि दश तस्यातिचक्रमुः}
{धर्मे प्रयतमानस्य सत्पथे निष्ठितस्य च} %||7-10-5||

\twolineshloka
{विभीषणस्तु धर्मात्मा नित्यं धर्मपरः शुचिः}
{पञ्चवर्षसहस्राणि पादेनैकेन तस्थिवान्} %||7-10-6||

\twolineshloka
{समाप्ते नियमे तस्य ननृतुश्चाप्सरोगणाः}
{पपात पुष्पवर्षं च क्षुभिताश्चापि देवताः} %||7-10-7||

\twolineshloka
{पञ्चवर्षसहस्राणि सूर्यं चैवान्ववर्तत}
{तस्थौ चोर्ध्वशिरो बाहुः स्वाध्यायधृतमानसः} %||7-10-8||

\twolineshloka
{एवं विभीषणस्यापि गतानि नियतात्मनः}
{दशवर्षसहस्राणि स्वर्गस्थस्येव नन्दने} %||7-10-9||

\twolineshloka
{दशवर्षसहस्रं तु निराहारो दशाननः}
{पूर्णे वर्षसहस्रे तु शिरश्चाग्नौ जुहाव सः} %||7-10-10||

\twolineshloka
{एवं वर्षसहस्राणि नव तस्यातिचक्रमुः}
{शिरांसि नव चाप्यस्य प्रविष्टानि हुताशनम्} %||7-10-11||

\twolineshloka
{अथ वर्षसहस्रे तु दशमे दशमं शिरः}
{छेत्तुकामः स धर्मात्मा प्राप्तश्चात्र पितामहः} %||7-10-12||

\twolineshloka
{पितामहस्तु सुप्रीतः सार्धं देवैरुपस्थितः}
{वत्स वत्स दशग्रीव प्रीतोऽस्मीत्यभ्यभाषत} %||7-10-13||

\twolineshloka
{शीघ्रं वरय धर्मज्ञ वरो यस्तेऽभिकाङ्क्षितः}
{किं ते कामं करोम्यद्य न वृथा ते परिश्रमः} %||7-10-14||

\twolineshloka
{ततोऽब्रवीद्दशग्रीवः प्रहृष्टेनान्तरात्मना}
{प्रणम्य शिरसा देवं हर्षगद्गदया गिरा} %||7-10-15||

\twolineshloka
{भगवन्प्राणिनां नित्यं नान्यत्र मरणाद्भयम्}
{नास्ति मृत्युसमः शत्रुरमरत्वमतो वृणे} %||7-10-16||

\twolineshloka
{सुपर्णनागयक्षाणां दैत्यदानवरक्षसाम्}
{अवध्यः स्यां प्रजाध्यक्ष देवतानां च शाश्वतम्} %||7-10-17||

\twolineshloka
{न हि चिन्ता ममान्येषु प्राणिष्वमरपूजित}
{तृणभूता हि मे सर्वे प्राणिनो मानुषादयः} %||7-10-18||

\twolineshloka
{एवमुक्तस्तु धर्मात्मा दशग्रीवेण रक्षसा}
{उवाच वचनं राम सह देवैः पितामहः} %||7-10-19||

\twolineshloka
{भविष्यत्येवमेवैतत्तव राक्षसपुङ्गव}
{शृणु चापि वचो भूयः प्रीतस्येह शुभं मम} %||7-10-20||

\twolineshloka
{हुतानि यानि शीर्षाणि पूर्वमग्नौ त्वयानघ}
{पुनस्तानि भविष्यन्ति तथैव तव राक्षस} %||7-10-21||

\twolineshloka
{एवं पितामहोक्तस्य दशग्रीवस्य रक्षसः}
{अग्नौ हुतानि शीर्षाणि यानि तान्युत्थितानि वै} %||7-10-22||

\twolineshloka
{एवमुक्त्व्वा तु तं राम दशग्रीवं प्रजापतिः}
{विभीषणमथोवाच वाक्यं लोकपितामहः} %||7-10-23||

\twolineshloka
{विभीषण त्वया वत्स धर्मसंहितबुद्धिना}
{परितुष्टोऽस्मि धर्मज्ञ वरं वरय सुव्रत} %||7-10-24||

\twolineshloka
{विभीषणस्तु धर्मात्मा वचनं प्राह साञ्जलिः}
{वृतः सर्वगुणैर्नित्यं चन्द्रमा इव रश्मिभिः} %||7-10-25||

\twolineshloka
{भगवन्कृतकृत्योऽहं यन्मे लोकगुरुः स्वयम्}
{प्रीतो यदि त्वं दातव्यं वरं मे शृणु सुव्रत} %||7-10-26||

\twolineshloka
{या या मे जायते बुद्धिर्येषु येष्वाश्रमेष्विह}
{सा सा भवतु धर्मिष्ठा तं तं धर्मं च पालये} %||7-10-27||

\twolineshloka
{एष मे परमोदार वरः परमको मतः}
{न हि धर्माभिरक्तानां लोके किञ्चन दुर्लभम्} %||7-10-28||

\twolineshloka
{अथ प्रजापतिः प्रीतो विभीषणमुवाच ह}
{धर्मिष्ठस्त्वं यथा वत्स तथा चैतद्भविष्यति} %||7-10-29||

\twolineshloka
{यस्माद्राक्षसयोनौ ते जातस्यामित्रकर्षण}
{नाधर्मे जायते बुद्धिरमरत्वं ददामि ते} %||7-10-30||

\twolineshloka
{कुम्भकर्णाय तु वरं प्रयच्छन्तमरिन्दम}
{प्रजापतिं सुराः सर्वे वाक्यं प्राञ्जलयोऽब्रुवन्} %||7-10-31||

\twolineshloka
{न तावत्कुम्भकर्णाय प्रदातव्यो वरस्त्वया}
{जानीषे हि यथा लोकांस्त्रासयत्येष दुर्मतिः} %||7-10-32||

\twolineshloka
{नन्दनेऽप्सरसः सप्त महेन्द्रानुचरा दश}
{अनेन भक्षिता ब्रह्मनृषयो मानुषास्तथा} %||7-10-33||

\twolineshloka
{वरव्याजेन मोहोऽस्मै दीयताममितप्रभ}
{लोकानां स्वस्ति चैव स्याद्भवेदस्य च संनतिः} %||7-10-34||

\twolineshloka
{एवमुक्तः सुरैर्ब्रह्माचिन्तयत्पद्मसम्भवः}
{चिन्तिता चोपतस्थेऽस्य पार्श्वं देवी सरस्वती} %||7-10-35||

\twolineshloka
{प्राञ्जलिः सा तु पर्श्वस्था प्राह वाक्यं सरस्वती}
{इयमस्म्यागता देवकिं कार्यं करवाण्यहम्} %||7-10-36||

\twolineshloka
{प्रजापतिस्तु तां प्राप्तां प्राह वाक्यं सरस्वतीम्}
{वाणि त्वं राक्षसेन्द्रस्य भव या देवतेप्सिता} %||7-10-37||

\twolineshloka
{तथेत्युक्त्वा प्रविष्टा सा प्रजापतिरथाब्रवीत्}
{कुम्भकर्ण महाबाहो वरं वरय यो मतः} %||7-10-38||

\twolineshloka
{कुम्भकर्णस्तु तद्वाक्यं श्रुत्वा वचनमब्रवीत्}
{स्वप्तुं वर्षाण्यनेकानि देवदेव ममेप्सितम्} %||7-10-39||

\twolineshloka
{एवमस्त्विति तं चोक्त्वा सह देवैः पितामहः}
{देवी सरस्वती चैव मुक्त्वा तं प्रययौ दिवम्} %||7-10-40||

\twolineshloka
{कुम्भकर्णस्तु दुष्टात्मा चिन्तयामास दुःखितः}
{कीर्दृशं किं न्विदं वाक्यं ममाद्य वदनाच्च्युतम्} %||7-10-41||

\twolineshloka
{एवं लब्धवराः सर्वे भ्रातरो दीप्ततेजसः}
{श्लेष्मातकवनं गत्वा तत्र ते न्यवसन्सुखम्} %||7-11-42||


{॥इत्यार्षे श्रीमद्रामायणे वाल्मीकीये आदिकाव्ये उत्तरकाण्डे दशमः सर्गः॥}

\sect{एकादशः सर्गः}

\twolineshloka
{सुमाली वरलब्धांस्तु ज्ञात्वा तान्वै निशाचरान्}
{उदतिष्ठद्भयं त्यक्त्वा सानुगः स रसातलात्} %||7-11-1||

\twolineshloka
{मारीचश्च प्रहस्तश्च विरूपाक्षो महोदरः}
{उदतिष्ठन्सुसंरब्धाः सचिवास्तस्य रक्षसः} %||7-11-2||

\twolineshloka
{सुमाली चैव तैः सर्वैर्वृतो राक्षसपुङ्गवैः}
{अभिगम्य दशग्रीवं परिष्वज्येदमब्रवीत्} %||7-11-3||

\twolineshloka
{दिष्ट्या ते पुत्रसम्प्राप्तश्चिन्तितोऽयं मनोरथः}
{यस्त्वं त्रिभुवणश्रेष्ठाल्लब्धवान्वरमीदृशम्} %||7-11-4||

\twolineshloka
{यत्कृते च वयं लङ्कां त्यक्त्वा याता रसातलम्}
{तद्गतं नो महाबाहो महद्विष्णुकृतं भयम्} %||7-11-5||

\twolineshloka
{असकृत्तेन भग्ना हि परित्यज्य स्वमालयम्}
{विद्रुताः सहिताः सर्वे प्रविष्टाः स्म रसातलम्} %||7-11-6||

\twolineshloka
{अस्मदीया च लङ्केयं नगरी राक्षसोषिता}
{निवेशिता तव भ्रात्रा धनाध्यक्षेण धीमता} %||7-11-7||

\twolineshloka
{यदि नामात्र शक्यं स्यात्साम्ना दानेन वानघ}
{तरसा वा महाबाहो प्रत्यानेतुं कृतं भवेत्} %||7-11-8||

\twolineshloka
{त्वं च लङ्केश्वरस्तात भविष्यसि न संशयः}
{सर्वेषां नः प्रभुश्चैव भविष्यसि महाबल} %||7-11-9||

\twolineshloka
{अथाब्रवीद्दशग्रीवो मातामहमुपस्थितम्}
{वित्तेशो गुरुरस्माकं नार्हस्येवं प्रभाषितुम्} %||7-11-10||

\twolineshloka
{उक्तवन्तं तथा वाक्यं दशग्रीवं निशाचरः}
{प्रहस्तः प्रश्रितं वाक्यमिदमाह सकारणम्} %||7-11-11||

\twolineshloka
{दशग्रीव महाबाहो नार्हस्त्वं वक्तुमीदृशम्}
{सौभ्रात्रं नास्ति शूराणां शृणु चेदं वचो मम} %||7-11-12||

\twolineshloka
{अदितिश्च दितिश्चैव भगिन्यौ सहिते किल}
{भार्ये परमरूपिण्यौ कश्यपस्य प्रजापतेः} %||7-11-13||

\twolineshloka
{अदितिर्जनयामास देवांस्त्रिभुवणेश्वरान्}
{दितिस्त्वजनयद्दैत्यान्कश्यपस्यात्मसम्भवान्} %||7-11-14||

\twolineshloka
{दैत्यानां किल धर्मज्ञ पुरेयं सवनार्णवा}
{सपर्वता मही वीर तेऽभवन्प्रभविष्णवः} %||7-11-15||

\twolineshloka
{निहत्य तांस्तु समरे विष्णुना प्रभविष्णुना}
{देवानां वशमानीतं त्रैलोक्यमिदमव्ययम्} %||7-11-16||

\twolineshloka
{नैतदेको भवानेव करिष्यति विपर्ययम्}
{सुरैराचरितं पूर्वं कुरुष्वैतद्वचो मम} %||7-11-17||

\twolineshloka
{एवमुक्तो दशग्रीवः प्रहस्तेन दुरात्मना}
{चिन्तयित्वा मुहूर्तं वै बाढमित्येव सोऽब्रवीत्} %||7-11-18||

\twolineshloka
{स तु तेनैव हर्षेण तस्मिन्नहनि वीर्यवान्}
{वनं गतो दशग्रीवः सह तैः क्षणदाचरैः} %||7-11-19||

\twolineshloka
{त्रिकूटस्थः स तु तदा दशग्रीवो निशाचरः}
{प्रेषयामास दौत्येन प्रहस्तं वाक्यकोविदम्} %||7-11-20||

\twolineshloka
{प्रहस्त शीघ्रं गत्वा त्वं ब्रूहि नैरृतपुङ्गवम्}
{वचनान्मम वित्तेशं सामपूर्वमिदं वचः} %||7-11-21||

\twolineshloka
{इयं लङ्का पुरी राजन्राक्षसानां महात्मनाम्}
{त्वया निवेशिता सौम्य नैतद्युक्तं तवानघ} %||7-11-22||

\twolineshloka
{तद्भवान्यदि साम्नैतां दद्यादतुलविक्रम}
{कृता भवेन्मम प्रीतिर्धर्मश्चैवानुपालितः} %||7-11-23||

\twolineshloka
{इत्युक्तः स तदा गत्वा प्रहस्तो वाक्यकोविदः}
{दशग्रीववचः सर्वं वित्तेशाय न्यवेदयत्} %||7-11-24||

\twolineshloka
{प्रहस्तादपि संश्रुत्य देवो वैश्रवणो वचः}
{प्रत्युवाच प्रहस्तं तं वाक्यं वाक्यविशारदः} %||7-11-25||

\twolineshloka
{ब्रूहि गच्छ दशग्रीवं पुरी राज्यं च यन्मम}
{तवाप्येतन्महाबाहो भुङ्क्ष्वैतद्धतकण्टकम्} %||7-11-26||

\twolineshloka
{सर्वं कर्तास्मि भद्रं ते राक्षसेश वचोऽचिरात्}
{किं तु तावत्प्रतीक्षस्व पितुर्यावन्निवेदये} %||7-11-27||

\twolineshloka
{एवमुक्त्वा धनाध्यक्षो जगाम पितुरन्तिकम्}
{अभिवाद्य गुरुं प्राह रावणस्य यदीप्सितम्} %||7-11-28||

\threelineshloka
{एष तात दशग्रीवो दूतं प्रेषितवान्मम}
{दीयतां नगरी लङ्का पूर्वं रक्षोगणोषिता}
{मयात्र यदनुष्ठेयं तन्ममाचक्ष्व सुव्रत} %||7-11-29||

\twolineshloka
{ब्रह्मर्षिस्त्वेवमुक्तोऽसौ विश्रवा मुनिपुङ्गवः}
{उवाच धनदं वाक्यं शृणु पुत्र वचो मम} %||7-11-30||

\twolineshloka
{दशग्रीवो महाबाहुरुक्तवान्मम संनिधौ}
{मया निर्भर्त्सितश्चासीद्बहुधोक्तः सुदुर्मतिः} %||7-11-31||

\twolineshloka
{स क्रोधेन मया चोक्तो ध्वंसस्वेति पुनः पुनः}
{श्रेयोऽभियुक्तं धर्म्यं च शृणु पुत्र वचो मम} %||7-11-32||

\twolineshloka
{वरप्रदानसम्मूढो मान्यामान्यं सुदुर्मतिः}
{न वेत्ति मम शापाच्च प्रकृतिं दारुणां गतः} %||7-11-33||

\twolineshloka
{तस्माद्गच्छ महाबाहो कैलासं धरणीधरम्}
{निवेशय निवासार्थं त्यज लङ्कां सहानुगः} %||7-11-34||

\twolineshloka
{तत्र मन्दाकिनी रम्या नदीनां प्रवरा नदी}
{काञ्चनैः सूर्यसङ्काशैः पङ्कजैः संवृतोदका} %||7-11-35||

\twolineshloka
{न हि क्षमं त्वया तेन वैरं धनदरक्षसा}
{जानीषे हि यथानेन लब्धः परमको वरः} %||7-11-36||

\twolineshloka
{एवमुक्तो गृहीत्वा तु तद्वचः पितृगौरवात्}
{सदार पौरः सामात्यः सवाहनधनो गतः} %||7-11-37||

\threelineshloka
{प्रहस्तस्तु दशग्रीवं गत्वा सर्वं न्यवेदयत्}
{शून्या सा नगरी लङ्का त्रिंशद्योजनमायता}
{प्रविश्य तां सहास्माभिः स्वधर्मं तत्र पालय} %||7-11-38||

\twolineshloka
{एवमुक्तः प्रहस्तेन रावणो राक्षसस्तदा}
{विवेश नगरीं लङ्कां सभ्राता सबलानुगः} %||7-11-39||

\fourlineindentedshloka
{स चाभिषिक्तः क्षणदाचरैस्तदा}
{ निवेशयामास पुरीं दशाननः}
{निकामपूर्णा च बभूव सा पुरी}
{ निशाचरैर्नीलबलाहकोपमैः} %||7-11-40||

\fourlineindentedshloka
{धनेश्वरस्त्वथ पितृवाक्यगौरवा}
{न्न्यवेशयच्छशिविमले गिरौ पुरीम्}
{स्वलङ्कृतैर्भवनवरैर्विभूषिताम्}
{ पुरन्दरस्येव तदामरावतीम्} %||7-12-41||


{॥इत्यार्षे श्रीमद्रामायणे वाल्मीकीये आदिकाव्ये उत्तरकाण्डे एकादशः सर्गः॥}

\sect{द्वादशः सर्गः}

\twolineshloka
{राक्षसेन्द्रोऽभिषिक्तस्तु भ्रातृभ्यां सहितस्तदा}
{ततः प्रदानं राक्षस्या भगिन्याः समचिन्तयत्} %||7-12-1||

\twolineshloka
{ददौ तां कालकेयाय दानवेन्द्राय राक्षसीम्}
{स्वसां शूर्पणखां नाम विद्युज्जिह्वाय नामतः} %||7-12-2||

\twolineshloka
{अथ दत्त्वा स्वसारं स मृगयां पर्यटन्नृपः}
{तत्रापश्यत्ततो राम मयं नाम दितेः सुतम्} %||7-12-3||

\twolineshloka
{कन्यासहायं तं दृष्ट्वा दशग्रीवो निशाचरः}
{अपृच्छत्को भवनेको निर्मनुष्य मृगे वने} %||7-12-4||

\twolineshloka
{मयस्त्वथाब्रवीद्राम पृच्छन्तं तं निशाचरम्}
{श्रूयतां सर्वमाख्यास्ये यथावृत्तमिदं मम} %||7-12-5||

\twolineshloka
{हेमा नामाप्सरास्तात श्रुतपूर्वा यदि त्वया}
{दैवतैर्मम सा दत्ता पौलोमीव शतक्रतोः} %||7-12-6||

\twolineshloka
{तस्यां सक्तमनास्तात पञ्चवर्षशतान्यहम्}
{सा च दैवत कार्येण गता वर्षं चतुर्दशम्} %||7-12-7||

\twolineshloka
{तस्याः कृते च हेमायाः सर्वं हेमपुरं मया}
{वज्रवैदूर्यचित्रं च मायया निर्मितं तदा} %||7-12-8||

\twolineshloka
{तत्राहमरतिं विन्दंस्तया हीनः सुदुःखितः}
{तस्मात्पुराद्दुहितरं गृहीत्वा वनमागतः} %||7-12-9||

\twolineshloka
{इयं ममात्मजा राजंस्तस्याः कुक्षौ विवर्धिता}
{भर्तारमनया सार्धमस्याः प्राप्तोऽस्मि मार्गितुम्} %||7-12-10||

\twolineshloka
{कन्यापितृत्वं दुःखं हि नराणां मानकाङ्क्षिणाम्}
{कन्या हि द्वे कुले नित्यं संशये स्थाप्य तिष्ठति} %||7-12-11||

\twolineshloka
{द्वौ सुतौ तु मम त्वस्यां भार्यायां सम्बभूवतुः}
{मायावी प्रथमस्तात दुन्दुभिस्तदनन्तरम्} %||7-12-12||

\twolineshloka
{एतत्ते सर्वमाख्यातं याथातथ्येन पृच्छतः}
{त्वामिदानीं कथं तात जानीयां को भवानिति} %||7-12-13||

\twolineshloka
{एवमुक्तो राक्षसेन्द्रो विनीतमिदमब्रवीत्}
{अहं पौलस्त्य तनयो दशग्रीवश्च नामतः} %||7-12-14||

\twolineshloka
{ब्रह्मर्षेस्तं सुतं ज्ञात्वा मयो हर्षमुपागतः}
{दातुं दुहितरं तस्य रोचयामास तत्र वै} %||7-12-15||

\threelineshloka
{प्रहसन्प्राह दैत्येन्द्रो राक्षसेन्द्रमिदं वचः}
{इयं ममात्मजा राजन्हेमयाप्सरसा धृता}
{कन्या मन्दोदरी नाम पत्न्यर्थं प्रतिगृह्यताम्} %||7-12-16||

\twolineshloka
{बाढमित्येव तं राम दशग्रीवोऽभ्यभाषत}
{प्रज्वाल्य तत्र चैवाग्निमकरोत्पाणिसङ्ग्रहम्} %||7-12-17||

\twolineshloka
{न हि तस्य मयो राम शापाभिज्ञस्तपोधनात्}
{विदित्वा तेन सा दत्ता तस्य पैतामहं कुलम्} %||7-12-18||

\twolineshloka
{अमोघां तस्य शक्तिं च प्रददौ परमाद्भुताम्}
{परेण तपसा लब्धां जघ्निवाँल्लक्ष्मणं यया} %||7-12-19||

\twolineshloka
{एवं स कृतदारो वै लङ्कायामीश्वरः प्रभुः}
{गत्वा तु नगरं भार्ये भ्रातृभ्यां समुदावहत्} %||7-12-20||

\twolineshloka
{वैरोचनस्य दौहित्रीं वज्रज्वालेति नामतः}
{तां भार्यां कुम्भकर्णस्य रावणः समुदावहत्} %||7-12-21||

\twolineshloka
{गन्धर्वराजस्य सुतां शैलूषस्य महात्मन}
{सरमा नाम धर्मज्ञो लेभे भार्यां विभीषणः} %||7-12-22||

\twolineshloka
{तीरे तु सरसः सा वै संजज्ञे मानसस्य च}
{मानसं च सरस्तात ववृधे जलदागमे} %||7-12-23||

\twolineshloka
{मात्रा तु तस्याः कन्यायाः स्नेहनाक्रन्दितं वचः}
{सरो मा वर्धतेत्युक्तं ततः सा सरमाभवत्} %||7-12-24||

\twolineshloka
{एवं ते कृतदारा वै रेमिरे तत्र राक्षसाः}
{स्वां स्वां भार्यामुपादाय गन्धर्वा इव नन्दने} %||7-12-25||

\twolineshloka
{ततो मन्दोदरी पुत्रं मेघनादमसूयत}
{स एष इन्द्रजिन्नाम युष्माभिरभिधीयते} %||7-12-26||

\twolineshloka
{जातमात्रेण हि पुरा तेन राक्षससूनुना}
{रुदता सुमहान्मुक्तो नादो जलधरोपमः} %||7-12-27||

\twolineshloka
{जडीकृतायां लङ्कायां तेन नादेन तस्य वै}
{पिता तस्याकरोन्नाम मेघनाद इति स्वयम्} %||7-12-28||

\twolineshloka
{सोऽवर्धत तदा राम रावणान्तःपुरे शुभे}
{रक्ष्यमाणो वरस्त्रीभिश्छन्नः काष्ठैरिवानलः} %||7-13-29||


{॥इत्यार्षे श्रीमद्रामायणे वाल्मीकीये आदिकाव्ये उत्तरकाण्डे द्वादशः सर्गः॥}

\sect{त्रयोदशः सर्गः}

\twolineshloka
{अथ लोकेश्वरोत्सृष्टा तत्र कालेन केनचित्}
{निद्रा समभवत्तीव्रा कुम्भकर्णस्य रूपिणी} %||7-13-1||

\twolineshloka
{ततो भ्रातरमासीनं कुम्भकर्णोऽब्रवीद्वचः}
{निद्रा मां बाधते राजन्कारयस्व ममालयम्} %||7-13-2||

\twolineshloka
{विनियुक्तास्ततो राज्ञा शिल्पिनो विश्वकर्मवत्}
{अकुर्वन्कुम्भकर्णस्य कैलाससममालयम्} %||7-13-3||

\twolineshloka
{विस्तीर्णं योजनं शुभ्रं ततो द्विगुणमायतम्}
{दर्शनीयं निराबाधं कुम्भकर्णस्य चक्रिरे} %||7-13-4||

\twolineshloka
{स्फाटिकैः काञ्चनैश्चित्रैः स्तम्भैः सर्वत्र शोभितम्}
{वैदूर्यकृतशोभं च किङ्किणीजालकं तथा} %||7-13-5||

\twolineshloka
{दन्ततोरणविन्यस्तं वज्रस्फटिकवेदिकम्}
{सर्वर्तुसुखदं नित्यं मेरोः पुण्यां गुहामिव} %||7-13-6||

\twolineshloka
{तत्र निद्रां समाविष्टः कुम्भकर्णो निशाचरः}
{बहून्यब्दसहस्राणि शयानो नावबुध्यते} %||7-13-7||

\twolineshloka
{निद्राभिभूते तु तदा कुम्भकर्णे दशाननः}
{देवर्षियक्षगन्धर्वान्बाधते स्म स नित्यशः} %||7-13-8||

\twolineshloka
{उद्यानानि विचित्राणि नन्दनादीनि यानि च}
{तानि गत्वा सुसङ्क्रुद्धो भिनत्ति स्म दशाननः} %||7-13-9||

\twolineshloka
{नदीं गज इव क्रीडन्वृक्षान्वायुरिव क्षिपन्}
{नगान्वज्र इव सृष्टो विध्वंसयति नित्यशः} %||7-13-10||

\twolineshloka
{तथा वृत्तं तु विज्ञाय दशग्रीवं धनेश्वरः}
{कुलानुरूपं धर्मज्ञ वृत्तं संस्मृत्य चात्मनः} %||7-13-11||

\twolineshloka
{सौभ्रात्रदर्शनार्थं तु दूतं वैश्वरणस्तदा}
{लङ्कां सम्प्रेषयामास दशग्रीवस्य वै हितम्} %||7-13-12||

\twolineshloka
{स गत्वा नगरीं लङ्कामाससाद विभीषणम्}
{मानितस्तेन धर्मेण पृष्ठश्चागमनं प्रति} %||7-13-13||

\twolineshloka
{पृष्ट्वा च कुशलं राज्ञो ज्ञातीनपि च बान्धवान्}
{सभायां दर्शयामास तमासीनं दशाननम्} %||7-13-14||

\twolineshloka
{स दृष्ट्वा तत्र राजानं दीप्यमानं स्वतेजसा}
{जयेन चाभिसम्पूज्य तूष्णीमासीन्मुहूर्तकम्} %||7-13-15||

\twolineshloka
{तस्योपनीते पर्यङ्के वरास्तरणसंवृते}
{उपविश्य दशग्रीवं दूतो वाक्यमथाब्रवीत्} %||7-13-16||

\twolineshloka
{राजन्वदामि ते सर्वं भ्राता तव यदब्रवीत्}
{उभयोः सदृशं सौम्य वृत्तस्य च कुलस्य च} %||7-13-17||

\twolineshloka
{साधु पर्याप्तमेतावत्कृतश्चारित्रसङ्ग्रहः}
{साधु धर्मे व्यवस्थानं क्रियतां यदि शक्यते} %||7-13-18||

\twolineshloka
{दृष्टं मे नन्दनं भग्नमृषयो निहताः श्रुताः}
{देवानां तु समुद्योगस्त्वत्तो राजञ्श्रुतश्च मे} %||7-13-19||

\twolineshloka
{निराकृतश्च बहुशस्त्वयाहं राक्षसाधिप}
{अपराद्धा हि बाल्याच्च रक्षणीयाः स्वबान्धवाः} %||7-13-20||

\twolineshloka
{अहं तु हिमवत्पृष्ठं गतो धर्ममुपासितुम्}
{रौद्रं व्रतं समास्थाय नियतो नियतेन्द्रियः} %||7-13-21||

\twolineshloka
{तत्र देवो मया दृष्टः सह देव्योमया प्रभुः}
{सव्यं चक्षुर्मया चैव तत्र देव्यां निपातितम्} %||7-13-22||

\twolineshloka
{का न्वियं स्यादिति शुभा न खल्वन्येन हेतुना}
{रूपं ह्यनुपमं कृत्वा तत्र क्रीडति पार्वती} %||7-13-23||

\twolineshloka
{ततो देव्याः प्रभावेन दग्धं सव्यं ममेक्षणम्}
{रेणुध्वस्तमिव ज्योतिः पिङ्गलत्वमुपागतम्} %||7-13-24||

\twolineshloka
{ततोऽहमन्यद्विस्तीर्णं गत्वा तस्य गिरेस्तटम्}
{पूर्णं वर्षशतान्यष्टौ समवाप महाव्रतम्} %||7-13-25||

\twolineshloka
{समाप्ते नियमे तस्मिंस्तत्र देवो महेश्वरः}
{प्रीतः प्रीतेन मनसा प्राह वाक्यमिदं प्रभुः} %||7-13-26||

\twolineshloka
{प्रीतोऽस्मि तव धर्मज्ञ तपसानेन सुव्रत}
{मया चैतद्व्रतं चीर्णं त्वया चैव धनाधिप} %||7-13-27||

\twolineshloka
{तृतीयः पुरुषो नास्ति यश्चरेद्व्रतमीदृशम्}
{व्रतं सुदुश्चरं ह्येतन्मयैवोत्पादितं पुरा} %||7-13-28||

\twolineshloka
{तत्सखित्वं मया सार्धं रोचयस्व धनेश्वर}
{तपसा निर्जितत्वाद्धि सखा भव ममानघ} %||7-13-29||

\twolineshloka
{देव्या दग्धं प्रभावेन यच्च साव्यं तवेक्षणम्}
{एकाक्षि पिङ्गलेत्येव नाम स्थास्यति शाश्वतम्} %||7-13-30||

\twolineshloka
{एवं तेन सखित्वं च प्राप्यानुज्ञां च शङ्करात्}
{आगम्य च श्रुतोऽयं मे तव पापविनिश्चयः} %||7-13-31||

\twolineshloka
{तदधर्मिष्ठसंयोगान्निवर्त कुलदूषण}
{चिन्त्यते हि वधोपायः सर्षिसङ्घैः सुरैस्तव} %||7-13-32||

\twolineshloka
{एवमुक्तो दशग्रीवः क्रुद्धः संरक्तलोचनः}
{हस्तान्दन्तांश सम्पीड्य वाक्यमेतदुवाच ह} %||7-13-33||

\twolineshloka
{विज्ञातं ते मया दूत वाक्यं यत्त्वं प्रभाषसे}
{नैव त्वमसि नैवासौ भ्रात्रा येनासि प्रेषितः} %||7-13-34||

\twolineshloka
{हितं न स ममैतद्धि ब्रवीति धनरक्षकः}
{महेश्वरसखित्वं तु मूढ श्रावयसे किल} %||7-13-35||

\twolineshloka
{न हन्तव्यो गुरुर्ज्येष्ठो ममायमिति मन्यते}
{तस्य त्विदानीं श्रुत्वा मे वाक्यमेषा कृता मतिः} %||7-13-36||

\threelineshloka
{त्रीँल्लोकानपि जेष्यामि बाहुवीर्यमुपाश्रितः}
{एतन्मुहूर्तमेषोऽहं तस्यैकस्य कृते च वै}
{चतुरो लोकपालांस्तान्नयिष्यामि यमक्षयम्} %||7-13-37||

\twolineshloka
{एवमुक्त्वा तु लङ्केशो दूतं खड्गेन जघ्निवान्}
{ददौ भक्षयितुं ह्येनं राक्षसानां दुरात्मनाम्} %||7-13-38||

\twolineshloka
{ततः कृतस्वस्त्ययनो रथमारुह्य रावणः}
{त्रैलोक्यविजयाकाङ्क्षी ययौ तत्र धनेश्वरः} %||7-14-39||


{॥इत्यार्षे श्रीमद्रामायणे वाल्मीकीये आदिकाव्ये उत्तरकाण्डे त्रयोदशः सर्गः॥}

\sect{चतुर्दशः सर्गः}

\twolineshloka
{ततः स सचिवैः सार्धं षड्भिर्नित्यं बलोत्कटैः}
{महोदरप्रहस्ताभ्यां मारीचशुकसारणैः} %||7-14-1||

\twolineshloka
{धूम्राक्षेण च वीरेण नित्यं समरगृध्नुना}
{वृतः सम्प्रययौ श्रीमान्क्रोधाल्लोकान्दहन्निव} %||7-14-2||

\twolineshloka
{पुराणि स नदीः शैलान्वनान्युपवनानि च}
{अतिक्रम्य मुहूर्तेन कैलासं गिरिमाविशत्} %||7-14-3||

\twolineshloka
{तं निविष्टं गिरौ तस्मिन्राक्षसेन्द्रं निशम्य तु}
{राज्ञो भ्रातायमित्युक्त्वा गता यत्र धनेश्वरः} %||7-14-4||

\twolineshloka
{गत्वा तु सर्वमाचख्युर्भ्रातुस्तस्य विनिश्चयम्}
{अनुज्ञाता ययुश्चैव युद्धाय धनदेन ते} %||7-14-5||

\twolineshloka
{ततो बलस्य सङ्क्षोभः सागरस्येव वर्धतः}
{अभून्नैरृतराजस्य गिरिं सञ्चालयन्निव} %||7-14-6||

\twolineshloka
{ततो युद्धं समभवद्यक्षराक्षससङ्कुलम्}
{व्यथिताश्चाभवंस्तत्र सचिवास्तस्य रक्षसः} %||7-14-7||

\twolineshloka
{तं दृष्ट्वा तादृशं सैन्यं दशग्रीवो निशाचरः}
{हर्षान्नादं ततः कृत्वा रोषात्समभिवर्तत} %||7-14-8||

\twolineshloka
{ये तु ते राक्षसेन्द्रस्य सचिवा घोरविक्रमः}
{ते सहस्रं सहस्राणामेकैकं समयोधयन्} %||7-14-9||

\twolineshloka
{ततो गदाभिः परिघैरसिभिः शक्तितोमरैः}
{वध्यमानो दशग्रीवस्तत्सैन्यं समगाहत} %||7-14-10||

\twolineshloka
{तैर्निरुच्छ्वासवत्तत्र वध्यमानो दशाननः}
{वर्षमाणैरिव घनैर्यक्षेन्द्रैः संनिरुध्यत} %||7-14-11||

\twolineshloka
{स दुरात्मा समुद्यम्य कालदण्डोपमां गदाम्}
{प्रविवेश ततः सैन्यं नयन्यक्षान्यमक्षयम्} %||7-14-12||

\twolineshloka
{स कक्षमिव विस्तीर्णं शुष्केन्धनसमाकुलम्}
{वातेनाग्निरिवायत्तोऽदहत्सैन्यं सुदारुणम्} %||7-14-13||

\twolineshloka
{तैस्तु तस्य मृधेऽमात्यैर्महोदरशुकादिभिः}
{अल्पावशिष्टास्ते यक्षाः कृता वातैरिवाम्बुदाः} %||7-14-14||

\twolineshloka
{केचित्त्वायुधभग्नाङ्गाः पतिताः समरक्षितौ}
{ओष्ठान्स्वदशनैस्तीक्ष्णैर्दंशन्तो भुवि पातिताः} %||7-14-15||

\twolineshloka
{भयादन्योन्यमालिङ्ग्य भ्रष्टशस्त्रा रणाजिरे}
{निषेदुस्ते तदा यक्षाः कूला जलहता इव} %||7-14-16||

\twolineshloka
{हतानां स्वर्गसंस्थानां युध्यतां पृथिवीतले}
{प्रेक्षतामृषिसङ्घानां न बभूवान्तरं दिवि} %||7-14-17||

\twolineshloka
{एतस्मिन्नन्तरे राम विस्तीर्णबलवाहनः}
{अगमत्सुमहान्यक्षो नाम्ना संयोधकण्टकः} %||7-14-18||

\twolineshloka
{तेन यक्षेण मारीचो विष्णुनेव समाहतः}
{पतितः पृथिवीं भेजे क्षीणपुण्य इवाम्बरात्} %||7-14-19||

\twolineshloka
{प्राप्तसंज्ञो मुहूर्तेन विश्रम्य च निशाचरः}
{तं यक्षं योधयामास स च भग्नः प्रदुद्रुवे} %||7-14-20||

\twolineshloka
{ततः काञ्चनचित्राङ्गं वैदूर्यरजतोक्षितम्}
{मर्यादां द्वारपालानां तोरणं तत्समाविशत्} %||7-14-21||

\twolineshloka
{ततो राम दशग्रीवं प्रविशन्तं निशाचरम्}
{सूर्यभानुरिति ख्यातो द्वारपालो न्यवारयत्} %||7-14-22||

\threelineshloka
{ततस्तोरणमुत्पाट्य तेन यक्षेण ताडितः}
{राक्षसो यक्षसृष्टेन तोरणेन समाहतः}
{न क्षितिं प्रययौ राम वरात्सलिलयोनिनः} %||7-14-23||

\twolineshloka
{स तु तेनैव तं यक्षं तोरणेन समाहनत्}
{नादृश्यत तदा यक्षो भस्म तेन कृतस्तु सः} %||7-14-24||

\twolineshloka
{ततः प्रदुद्रुवुः सर्वे यक्षा दृष्ट्वा पराक्रमम्}
{ततो नदीर्गुहाश्चैव विविशुर्भयपीडिताः} %||7-15-25||


{॥इत्यार्षे श्रीमद्रामायणे वाल्मीकीये आदिकाव्ये उत्तरकाण्डे चतुर्दशः सर्गः॥}

\sect{पञ्चदशः सर्गः}

\twolineshloka
{ततस्तान्विद्रुतान्दृष्ट्वा यक्षाञ्शतसहस्रशः}
{स्वयमेव धनाध्यक्षो निर्जगाम रणं प्रति} %||7-15-1||

\twolineshloka
{तत्र माणिचारो नाम यक्षः परमदुर्जयः}
{वृतो यक्षसहस्रैः स चतुर्भिः समयोधयत्} %||7-15-2||

\twolineshloka
{ते गदामुसलप्रासशक्तितोमरमुद्गरैः}
{अभिघ्नन्तो रणे यक्षा राक्षसानभिदुद्रुवुः} %||7-15-3||

\twolineshloka
{ततः प्रहस्तेन तदा सहस्रं निहतं रणे}
{महोदरेण गदया सहस्रमपरं हतम्} %||7-15-4||

\twolineshloka
{क्रुद्धेन च तदा राम मारीचेन दुरात्मना}
{निमेषान्तरमात्रेण द्वे सहस्रे निपातिते} %||7-15-5||

\twolineshloka
{धूम्राक्षेण समागम्य माणिभद्रो महारणे}
{मुसलेनोरसि क्रोधात्ताडितो न च कम्पितः} %||7-15-6||

\twolineshloka
{ततो गदां समाविध्य माणिभद्रेण राक्षसः}
{धूम्राक्षस्ताडितो मूर्ध्नि विह्वलो निपपात ह} %||7-15-7||

\twolineshloka
{धूम्राक्षं ताडितं दृष्ट्वा पतितं शोणितोक्षितम्}
{अभ्यधावत्सुसङ्क्रुद्धो माणिभद्रं दशाननः} %||7-15-8||

\twolineshloka
{तं क्रुद्धमभिधावन्तं युगान्ताग्निमिवोत्थितम्}
{शक्तिभिस्ताडयामास तिसृभिर्यक्षपुङ्गवः} %||7-15-9||

\threelineshloka
{ततो राक्षसराजेन ताडितो गदया रणे}
{तस्य तेन प्रहारेण मुकुटः पार्श्वमागतः}
{तदा प्रभृति यक्षोऽसौ पार्श्वमौलिरिति स्मृतः} %||7-15-10||

\twolineshloka
{तस्मिंस्तु विमुखे यक्षे माणिभद्रे महात्मनि}
{संनादः सुमहान्राम तस्मिञ्शैले व्यवर्धत} %||7-15-11||

\twolineshloka
{ततो दूरात्प्रददृशे धनाध्यक्षो गदाधरः}
{शुक्रप्रोष्टःपदाभ्यां च शङ्खपद्मसमावृतः} %||7-15-12||

\twolineshloka
{स दृष्ट्वा भ्रातरं सङ्ख्ये शापाद्विभ्रष्टगौरवम्}
{उवाच वचनं धीमान्युक्तं पैतामहे कुले} %||7-15-13||

\twolineshloka
{मया त्वं वार्यमाणोऽपि नावगच्छसि दुर्मते}
{पश्चादस्य फलं प्राप्य ज्ञास्यसे निरयं गतः} %||7-15-14||

\twolineshloka
{यो हि मोहाद्विषं पीत्वा नावगच्छति मानवः}
{परिणामे स वि मूढो जानीते कर्मणः फलम्} %||7-15-15||

\twolineshloka
{दैवतानि हि नन्दन्ति धर्मयुक्तेन केनचित्}
{येन त्वमीदृशं भावं नीतस्तच्च न बुध्यसे} %||7-15-16||

\twolineshloka
{यो हि मातॄह्पितॄन्भ्रातॄनाचर्यांश्चावमन्यते}
{स पश्यति फलं तस्य प्रेतराजवशं गतः} %||7-15-17||

\twolineshloka
{अध्रुवे हि शरीरे यो न करोति तपोऽर्जनम्}
{स पश्चात्तप्यते मूढो मृतो दृष्ट्वात्मनो गतिम्} %||7-15-18||

\twolineshloka
{कस्यचिन्न हि दुर्बुधेश्छन्दतो जायते मतिः}
{यादृशं कुरुते कर्म तादृशं फलमश्नुते} %||7-15-19||

\twolineshloka
{बुद्धिं रूपं बलं वित्तं पुत्रान्माहात्म्यमेव च}
{प्रप्नुवन्ति नराः सर्वं स्वकृतैः पूर्वकर्मभिः} %||7-15-20||

\twolineshloka
{एवं निरयगामी त्वं यस्य ते मतिरीदृशी}
{न त्वां समभिभाषिष्ये दुर्वृत्तस्यैष निर्णयः} %||7-15-21||

\twolineshloka
{एवमुक्त्वा ततस्तेन तस्यामात्याः समाहताः}
{मारीचप्रमुखाः सर्वे विमुखा विप्रदुद्रुवुः} %||7-15-22||

\twolineshloka
{ततस्तेन दशग्रीवो यक्षेन्द्रेण महात्मना}
{गदयाभिहतो मूर्ध्नि न च स्थानाद्व्यकम्पत} %||7-15-23||

\twolineshloka
{ततस्तौ राम निघ्नन्तावन्योन्यं परमाहवे}
{न विह्वलौ न च श्रान्तौ बभूवतुरमर्षणैः} %||7-15-24||

\twolineshloka
{आग्नेयमस्त्रं स ततो मुमोच धनदो रणे}
{वारुणेन दशग्रीवस्तदस्त्रं प्रत्यवारयत्} %||7-15-25||

\twolineshloka
{ततो मायां प्रविष्टः स राक्षसीं राक्षसेश्वरः}
{जघान मूर्ध्नि धनदं व्याविध्य महतीं गदाम्} %||7-15-26||

\twolineshloka
{एवं स तेनाभिहतो विह्वलः शोणितोक्षितः}
{कृत्तमूल इवाशोको निपपात धनाधिपः} %||7-15-27||

\twolineshloka
{ततः पद्मादिभिस्तत्र निधिभिः स धनाधिपः}
{नन्दनं वनमानीय धनदो श्वासितस्तदा} %||7-15-28||

\twolineshloka
{ततो निर्जित्य तं राम धनदं राक्षसाधिपः}
{पुष्पकं तस्य जग्राह विमानं जयलक्षणम्} %||7-15-29||

\twolineshloka
{काञ्चनस्तम्भसंवीतं वैदूर्यमणितोरणम्}
{मुक्ताजालप्रतिच्छन्नं सर्वकामफलद्रुमम्} %||7-15-30||

\twolineshloka
{तत्तु राजा समारुह्य कामगं वीर्यनिर्जितम्}
{जित्वा वैश्रवणं देवं कैलासादवरोहत} %||7-16-31||


{॥इत्यार्षे श्रीमद्रामायणे वाल्मीकीये आदिकाव्ये उत्तरकाण्डे पञ्चदशः सर्गः॥}

\sect{षोडशः सर्गः}

\twolineshloka
{स जित्वा भ्रातरं राम धनदं राक्षसाधिपः}
{महासेनप्रसूतिं तु ययौ शरवणं ततः} %||7-16-1||

\twolineshloka
{अथापश्यद्दशग्रीवो रौक्मं शरवणं तदा}
{गभस्तिजालसंवीतं द्वितीयमिव भास्करम्} %||7-16-2||

\twolineshloka
{पर्वतं स समासाद्य किञ्चिद्रम्यवनान्तरम्}
{अपश्यत्पुष्पकं तत्र राम विष्टम्भितं दिवि} %||7-16-3||

\twolineshloka
{विष्टब्धं पुष्पकं दृष्ट्वा कामगं ह्यगमं कृतम्}
{राक्षसश्चिन्तयामास सचिवैस्तैः समावृतः} %||7-16-4||

\twolineshloka
{किमिदं यन्निमित्तं मे न च गच्छति पुष्पकम्}
{पर्वतस्योपरिस्थस्य कस्य कर्म त्विदं भवेत्} %||7-16-5||

\twolineshloka
{ततोऽब्रवीद्दशग्रीवं मारीचो बुद्धिकोविदः}
{नैतन्निष्कारणं राजन्पुष्पकोऽयं न गच्छति} %||7-16-6||

\twolineshloka
{ततः पार्श्वमुपागम्य भवस्यानुचरो बली}
{नन्दीश्वर उवाचेदं राक्षसेन्द्रमशङ्कितः} %||7-16-7||

{निवर्तस्व दशग्रीव शैले क्रीडति शङ्करः॥\devanumber{8}॥} %||7-16-8|| (Check)

\addtocounter{shlokacount}{1}

\twolineshloka
{सुपर्णनागयक्षाणां दैत्यदानवरक्षसाम्}
{प्राणिनामेव सर्वेषामगम्यः पर्वतः कृतः} %||7-16-9||

\twolineshloka
{स रोषात्ताम्रनयनः पुष्पकादवरुह्य च}
{कोऽयं शङ्कर इत्युक्त्वा शैलमूलमुपागमत्} %||7-16-10||

\twolineshloka
{नन्दीश्वरमथापश्यदविदूरस्थितं प्रभुम्}
{दीप्तं शूलमवष्टभ्य द्वितीयमिव शङ्करम्} %||7-16-11||

\twolineshloka
{स वानरमुखं दृष्ट्वा तमवज्ञाय राक्षसः}
{प्रहासं मुमुचे मौर्ख्यात्सतोय इव तोयदः} %||7-16-12||

\twolineshloka
{सङ्क्रुद्धो भगवान्नन्दी शङ्करस्यापरा तनुः}
{अब्रवीद्राक्षसं तत्र दशग्रीवमुपस्थितम्} %||7-16-13||

\twolineshloka
{यस्माद्वानरमूर्तिं मां दृष्ट्वा राक्षसदुर्मते}
{मौर्ख्यात्त्वमवजानीषे परिहासं च मुञ्चसि} %||7-16-14||

\twolineshloka
{तस्मान्मद्रूपसंयुक्ता मद्वीर्यसमतेजसः}
{उत्पत्स्यन्ते वधार्थं हि कुलस्य तव वानराः} %||7-16-15||

\twolineshloka
{किं त्विदानीं मया शक्यं कर्तुं यत्त्वां निशाचर}
{न हन्तव्यो हतस्त्वं हि पूर्वमेव स्वकर्मभिः} %||7-16-16||

\twolineshloka
{अचिन्तयित्वा स तदा नन्दिवाक्यं निशाचरः}
{पर्वतं तं समासाद्य वाक्यमेतदुवाच ह} %||7-16-17||

\twolineshloka
{पुष्पकस्य गतिश्छिन्ना यत्कृते मम गच्छतः}
{तदेतच्छैलमुन्मूलं करोमि तव गोपते} %||7-16-18||

\twolineshloka
{केन प्रभावेन भवस्तत्र क्रीडति राजवत्}
{विज्ञातव्यं न जानीषे भयस्थानमुपस्थितम्} %||7-16-19||

\twolineshloka
{एवमुक्त्वा ततो राजन्भुजान्प्रक्षिप्य पर्वते}
{तोलयामास तं शैलं समृगव्यालपादपम्} %||7-16-20||

\twolineshloka
{ततो राम महादेवः प्रहसन्वीक्ष्य तत्कृतम्}
{पादाङ्गुष्ठेन तं शैलं पीडयामास लीलया} %||7-16-21||

\twolineshloka
{ततस्ते पीडितास्तस्य शैलस्याधो गता भुजाः}
{विस्मिताश्चाभवंस्तत्र सचिवास्तस्य रक्षसः} %||7-16-22||

\twolineshloka
{रक्षसा तेन रोषाच्च भुजानां पीडनात्तथा}
{मुक्तो विरावः सुमहांस्त्रैलोक्यं येन पूरितम्} %||7-16-23||

\twolineshloka
{मानुषाः शब्दवित्रस्ता मेनिरे लोकसङ्क्षयम्}
{देवताश्चापि सङ्क्षुब्धाश्चलिताः स्वेषु कर्मसु} %||7-16-24||

\twolineshloka
{ततः प्रीतो महादेवः शैलाग्रे विष्ठितस्तदा}
{मुक्त्वा तस्य भुजान्राजन्प्राह वाक्यं दशाननम्} %||7-16-25||

\twolineshloka
{प्रीतोऽस्मि तव वीर्याच्च शौण्डीर्याच्च निशाचर}
{रवतो वेदना मुक्तः स्वरः परमदारुणः} %||7-16-26||

\twolineshloka
{यस्माल्लोकत्रयं त्वेतद्रावितं भयमागतम्}
{तस्मात्त्वं रावणो नाम नाम्ना तेन भविष्यसि} %||7-16-27||

\twolineshloka
{देवता मानुषा यक्षा ये चान्ये जगतीतले}
{एवं त्वामभिधास्यन्ति रावणं लोकरावणम्} %||7-16-28||

\twolineshloka
{गच्छ पौलस्त्य विस्रब्धः पथा येन त्वमिच्छसि}
{मया त्वमभ्यनुज्ञातो राक्षसाधिप गम्यताम्} %||7-16-29||

\twolineshloka
{साक्षान्महेश्वरेणैवं कृतनामा स रावणः}
{अभिवाद्य महादेवं विमानं तत्समारुहत्} %||7-16-30||

\twolineshloka
{ततो महीतले राम परिचक्राम रावणः}
{क्षत्रियान्सुमहावीर्यान्बाधमानस्ततस्ततः} %||7-17-31||


{॥इत्यार्षे श्रीमद्रामायणे वाल्मीकीये आदिकाव्ये उत्तरकाण्डे षोडशः सर्गः॥}

\sect{सप्तदशः सर्गः}

\twolineshloka
{अथ राजन्महाबाहुर्विचरन्स महीतलम्}
{हिमवद्वनमासाद्य परिचक्राम रावणः} %||7-17-1||

\twolineshloka
{तत्रापश्यत वै कन्यां कृष्टाजिनजटाधराम्}
{आर्षेण विधिना युक्तां तपन्तीं देवतामिव} %||7-17-2||

\twolineshloka
{स दृष्ट्वा रूपसम्पन्नां कन्यां तां सुमहाव्रताम्}
{काममोहपरीतात्मा पप्रच्छ प्रहसन्निव} %||7-17-3||

\twolineshloka
{किमिदं वर्तसे भद्रे विरुद्धं यौवनस्य ते}
{न हि युक्ता तवैतस्य रूपस्येयं प्रतिक्रिया} %||7-17-4||

\twolineshloka
{कस्यासि दुहिता भद्रे को वा भर्ता तवानघे}
{पृच्छतः शंस मे शीघ्रं को वा हेतुस्तपोऽर्जने} %||7-17-5||

\twolineshloka
{एवमुक्ता तु सा कन्या तेनानार्येण रक्षसा}
{अब्रवीद्विधिवत्कृत्वा तस्यातिथ्यं तपोधना} %||7-17-6||

\twolineshloka
{कुशध्वजो नाम पिता ब्रह्मर्षिर्मम धार्मिकः}
{बृहस्पतिसुतः श्रीमान्बुद्ध्या तुल्यो बृहस्पतेः} %||7-17-7||

\twolineshloka
{तस्याहं कुर्वतो नित्यं वेदाभ्यासं महात्मनः}
{सम्भूता वान्मयी कन्या नाम्ना वेदवती स्मृता} %||7-17-8||

\twolineshloka
{ततो देवाः सगन्धर्वा यक्षराक्षसपन्नगाः}
{ते चापि गत्वा पितरं वरणं रोचयन्ति मे} %||7-17-9||

\twolineshloka
{न च मां स पिता तेभ्यो दत्तवान्राक्षसेश्वर}
{कारणं तद्वदिष्यामि निशामय महाभुज} %||7-17-10||

\twolineshloka
{पितुस्तु मम जामाता विष्णुः किल सुरोत्तमः}
{अभिप्रेतस्त्रिलोकेशस्तस्मान्नान्यस्य मे पिताः} %||7-17-11||

\threelineshloka
{दातुमिच्छति धर्मात्मा तच्छ्रुत्वा बलदर्पितः}
{शम्भुर्नाम ततो राजा दैत्यानां कुपितोऽभवत्}
{तेन रात्रौ प्रसुप्तो मे पिता पापेन हिंसितः} %||7-17-12||

\twolineshloka
{ततो मे जननी दीना तच्छरीरं पितुर्मम}
{परिष्वज्य महाभागा प्रविष्टा दहनं सह} %||7-17-13||

\twolineshloka
{ततो मनोरथं सत्यं पितुर्नारायणं प्रति}
{करोमीति ममेच्छा च हृदये साधु विष्ठिता} %||7-17-14||

\twolineshloka
{अहं प्रेतगतस्यापि करिष्ये काङ्क्षितं पितुः}
{इति प्रतिज्ञामारुह्य चरामि विपुलं तपः} %||7-17-15||

\twolineshloka
{एतत्ते सर्वमाख्यातं मया राक्षसपुङ्गव}
{आश्रितां विद्धि मां धर्मं नारायणपतीच्छया} %||7-17-16||

\twolineshloka
{विज्ञातस्त्वं हि मे राजन्गच्छ पौलस्त्यनन्दन}
{जानामि तपसा सर्वं त्रैलोक्ये यद्धि वर्तते} %||7-17-17||

\twolineshloka
{सोऽब्रवीद्रावणस्तत्र तां कन्यां सुमहाव्रताम्}
{अवरुह्य विमानाग्रात्कन्दर्पशरपीडितः} %||7-17-18||

\twolineshloka
{अवलिप्तासि सुश्रोणि यस्यास्ते मतिरीदृशी}
{वृद्धानां मृगशावाक्षि भ्राजते धर्मसञ्चयः} %||7-17-19||

\twolineshloka
{त्वं सर्वगुणसम्पन्ना नार्हसे कर्तुमीदृशम्}
{त्रैलोक्यसुन्दरी भीरु यौवने वार्धकं विधिम्} %||7-17-20||

\threelineshloka
{कश्च तावदसौ यं त्वं विष्णुरित्यभिभाषसे}
{वीर्येण तपसा चैव भोगेन च बलेन च}
{न मयासौ समो भद्रे यं त्वं कामयसेऽङ्गने} %||7-17-21||

\twolineshloka
{म मैवमिति सा कन्या तमुवाच निशाचरम्}
{मूर्धजेषु च तां रक्षः कराग्रेण परामृशत्} %||7-17-22||

\twolineshloka
{ततो वेदवती क्रुद्धा केशान्हस्तेन साच्छिनत्}
{उवाचाग्निं समाधाय मरणाय कृतत्वरा} %||7-17-23||

\twolineshloka
{धर्षितायास्त्वयानार्य नेदानीं मम जीवितम्}
{रक्षस्तस्मात्प्रवेक्ष्यामि पश्यतस्ते हुताशनम्} %||7-17-24||

\twolineshloka
{यस्मात्तु धर्षिता चाहमपापा चाप्यनाथवत्}
{तस्मात्तव वधार्थं वै समुत्पत्स्याम्यहं पुनः} %||7-17-25||

\twolineshloka
{न हि शक्यः स्त्रिया पाप हन्तुं त्वं तु विशेषतः}
{शापे त्वयि मयोत्सृष्टे तपसश्च व्ययो भवेत्} %||7-17-26||

\twolineshloka
{यदि त्वस्ति मया किञ्चित्कृतं दत्तं हुतं तथा}
{तेन ह्ययोनिजा साध्वी भवेयं धर्मिणः सुता} %||7-17-27||

\twolineshloka
{एवमुक्त्वा प्रविष्टा सा ज्वलन्तं वै हुताशनम्}
{पपात च दिवो दिव्या पुष्पवृष्टिः समन्ततः} %||7-17-28||

\twolineshloka
{पूर्वं क्रोधहतः शत्रुर्ययासौ निहतस्त्वया}
{समुपाश्रित्य शैलाभं तव वीर्यममानुषम्} %||7-17-29||

\twolineshloka
{एवमेषा महाभागा मर्त्येषूत्पद्यते पुनः}
{क्षेत्रे हलमुखग्रस्ते वेद्यामग्निशिखोपमा} %||7-17-30||

\threelineshloka
{एषा वेदवती नाम पूर्वमासीत्कृते युगे}
{त्रेतायुगमनुप्राप्य वधार्थं तस्य रक्षसः}
{सीतोत्पन्नेति सीतैषा मानुषैः पुनरुच्यते} %||7-18-31||


{॥इत्यार्षे श्रीमद्रामायणे वाल्मीकीये आदिकाव्ये उत्तरकाण्डे सप्तदशः सर्गः॥}

\sect{अष्टादशः सर्गः}

\twolineshloka
{प्रविष्टायां हुताशं तु वेदवत्यां स रावणः}
{पुष्पकं तत्समारुह्य परिचक्राम मेदिनीम्} %||7-18-1||

\twolineshloka
{ततो मरुत्तं नृपतिं यजन्तं सह दैवतैः}
{उशीरबीजमासाद्य ददर्श स तु राक्षसः} %||7-18-2||

\twolineshloka
{संवर्तो नाम ब्रह्मर्षिर्भ्राता साक्षाद्बृहस्पतेः}
{याजयामास धर्मज्ञः सर्वैर्ब्रह्मगणैर्वृतः} %||7-18-3||

\twolineshloka
{दृष्ट्वा देवास्तु तद्रक्षो वरदानेन दुर्जयम्}
{तां तां योनिं समापन्नास्तस्य धर्षणभीरवः} %||7-18-4||

\twolineshloka
{इन्द्रो मयूरः संवृत्तो धर्मराजस्तु वायसः}
{कृकलासो धनाध्यक्षो हंसो वै वरुणोऽभवत्} %||7-18-5||

\twolineshloka
{तं च राजानमासाद्य रावणो राक्षसाधिपः}
{प्राह युद्धं प्रयच्चेति निर्जितोऽस्मीति वा वद} %||7-18-6||

\twolineshloka
{ततो मरुत्तो नृपतिः को भवानित्युवाच तम्}
{अवहासं ततो मुक्त्वा राक्षसो वाक्यमब्रवीत्} %||7-18-7||

\twolineshloka
{अकुतूहलभावेन प्रीतोऽस्मि तव पार्थिव}
{धनदस्यानुजं यो मां नावगच्छसि रावणम्} %||7-18-8||

\twolineshloka
{त्रिषु लोकेषु कः सोऽस्ति यो न जानाति मे बलम्}
{भ्रातरं येन निर्जित्य विमानमिदमाहृतम्} %||7-18-9||

\twolineshloka
{ततो मरुत्तो नृपतिस्तं राक्षसमथाब्रवीत्}
{धन्यः खलु भवान्येन ज्येष्ठो भ्राता रणे जितः} %||7-18-10||

\twolineshloka
{नाधर्मसहितं श्लाघ्यं न लोकप्रतिसंहितम्}
{कर्म दौरात्म्यकं कृत्वा श्लाघसे भ्रातृनिर्जयात्} %||7-18-11||

\twolineshloka
{किं त्वं प्राक्केवलं धर्मं चरित्वा लब्धवान्वरम्}
{श्रुतपूर्वं हि न मया यादृशं भाषसे स्वयम्} %||7-18-12||

\twolineshloka
{ततः शरासनं गृह्य सायकांश्च स पार्थिवः}
{रणाय निर्ययौ क्रुद्धः संवर्तो मार्गमावृणोत्} %||7-18-13||

\twolineshloka
{सोऽब्रवीत्स्नेहसंयुक्तं मरुत्तं तं महानृषिः}
{श्रोतव्यं यदि मद्वाक्यं सम्प्रहारो न ते क्षमः} %||7-18-14||

\twolineshloka
{माहेश्वरमिदं सत्रमसमाप्तं कुलं दहेत्}
{दीक्षितस्य कुतो युद्धं क्रूरत्वं दीक्षिते कुतः} %||7-18-15||

\threelineshloka
{संशयश्च रणे नित्यं राक्षसश्चैष दुर्जयः}
{स निवृत्तो गुरोर्वाक्यान्मरुत्तः पृथिवीपतिः}
{विसृज्य सशरं चापं स्वस्थो मखमुखोऽभवत्} %||7-18-16||

\twolineshloka
{ततस्तं निर्जितं मत्वा घोषयामास वै शुकः}
{रावणो जितवांश्चेति हर्षान्नादं च मुक्तवान्} %||7-18-17||

\twolineshloka
{तान्भक्षयित्वा तत्रस्थान्महर्षीन्यज्ञमागतान्}
{वितृप्तो रुधिरैस्तेषां पुनः सम्प्रययौ महीम्} %||7-18-18||

\twolineshloka
{रावणे तु गते देवाः सेन्द्राश्चैव दिवौकसः}
{ततः स्वां योनिमासाद्य तानि सत्त्वान्यथाब्रुवन्} %||7-18-19||

\twolineshloka
{हर्षात्तदाब्रवीदिन्द्रो मयूरं नीलबर्हिणम्}
{प्रीतोऽस्मि तव धर्मज्ञ उपकाराद्विहङ्गम} %||7-18-20||

\twolineshloka
{मम नेत्रसहस्रं यत्तत्ते बर्हे भविष्यति}
{वर्षमाणे मयि मुदं प्राप्स्यसे प्रीतिलक्षणम्} %||7-18-21||

\twolineshloka
{नीलाः किल पुरा बर्हा मयूराणां नराधिप}
{सुराधिपाद्वरं प्राप्य गताः सर्वे विचित्रताम्} %||7-18-22||

\twolineshloka
{धर्मराजोऽब्रवीद्राम प्राग्वंशे वायसं स्थितम्}
{पक्षिंस्तवास्मि सुप्रीतः प्रीतस्य च वचः शृणु} %||7-18-23||

\twolineshloka
{यथान्ये विविधै रोगैः पीड्यन्ते प्राणिनो मया}
{ते न ते प्रभविष्यन्ति मयि प्रीते न संशयः} %||7-18-24||

\twolineshloka
{मृत्युतस्ते भयं नास्ति वरान्मम विहङ्गम}
{यावत्त्वां न वधिष्यन्ति नरास्तावद्भविष्यसि} %||7-18-25||

\twolineshloka
{ये च मद्विषयस्थास्तु मानवाः क्षुधयार्दिताः}
{त्वयि भुक्ते तु तृप्तास्ते भविष्यन्ति सबान्धवाः} %||7-18-26||

\twolineshloka
{वरुणस्त्वब्रवीद्धंसं गङ्गातोयविचारिणम्}
{श्रूयतां प्रीतिसंयुक्तं वचः पत्ररथेश्वर} %||7-18-27||

\twolineshloka
{वर्णो मनोहरः सौम्यश्चन्द्रमण्डलसंनिभः}
{भविष्यति तवोदग्रः शुक्लफेनसमप्रभः} %||7-18-28||

\twolineshloka
{मच्छरीरं समासाद्य कान्तो नित्यं भविष्यसि}
{प्राप्स्यसे चातुलां प्रीतिमेतन्मे प्रीतिलक्षणम्} %||7-18-29||

\twolineshloka
{हंसानां हि पुरा राम न वर्णः सर्वपाण्डुरः}
{पक्षा नीलाग्रसंवीताः क्रोडाः शष्पाग्रनिर्मलाः} %||7-18-30||

\twolineshloka
{अथाब्रवीद्वैश्रवणः कृकलासं गिरौ स्थितम्}
{हैरण्यं सम्प्रयच्छामि वर्णं प्रीतिस्तवाप्यहम्} %||7-18-31||

\twolineshloka
{सद्रव्यं च शिरो नित्यं भविष्यति तवाक्षयम्}
{एष काञ्चनको वर्णो मत्प्रीत्या ते भविष्यति} %||7-18-32||

\twolineshloka
{एवं दत्त्वा वरांस्तेभ्यस्तस्मिन्यज्ञोत्सवे सुराः}
{निवृत्ते सह राज्ञा वै पुनः स्वभवनं गताः} %||7-19-33||


{॥इत्यार्षे श्रीमद्रामायणे वाल्मीकीये आदिकाव्ये उत्तरकाण्डे अष्टादशः सर्गः॥}

\sect{एकोनविंशः सर्गः}

\twolineshloka
{अथ जित्वा मरुत्तं स प्रययौ राक्षसाधिपः}
{नगराणि नरेन्द्राणां युद्धकाङ्क्षी दशाननः} %||7-19-1||

\twolineshloka
{स समासाद्य राजेन्द्रान्महेन्द्रवरुणोपमान्}
{अब्रवीद्राक्षसेन्द्रस्तु युद्धं मे दीयतामिति} %||7-19-2||

\twolineshloka
{निर्जिताः स्मेति वा ब्रूत एषो हि मम निश्चयः}
{अन्यथा कुर्वतामेवं मोक्षो वो नोपपद्यते} %||7-19-3||

\twolineshloka
{ततस्तु बहवः प्राज्ञाः पार्थिवा धर्मणिश्चयाः}
{निर्जिताः स्मेत्यभाषन्त ज्ञात्वा वरबलं रिपोः} %||7-19-4||

\twolineshloka
{दुष्यन्तः सुरथो गाधिर्गयो राजा पुरूरवाः}
{एते सर्वेऽब्रुवंस्तात निर्जिताः स्मेति पार्थिवाः} %||7-19-5||

\twolineshloka
{अथायोध्यां समासाद्य रावणो राक्षसाधिपः}
{सुगुप्तामनरण्येन शक्रेणेवामरावतीम्} %||7-19-6||

\twolineshloka
{प्राह राजानमासाद्य युद्धं मे सम्प्रदीयताम्}
{निर्जितोऽस्मीति वा ब्रूहि ममैतदिह शासनम्} %||7-19-7||

\twolineshloka
{अनरण्यः सुसङ्क्रुद्धो राक्षसेन्द्रमथाब्रवीत्}
{दीयते द्वन्द्वयुद्धं ते राक्षसाधिपते मया} %||7-19-8||

\twolineshloka
{अथ पूर्वं श्रुतार्थेन सज्जितं सुमहद्धि यत्}
{निष्क्रामत्तन्नरेन्द्रस्य बलं रक्षोवधोद्यतम्} %||7-19-9||

\twolineshloka
{नागानां बहुसाहस्रं वाजिनामयुतं तथा}
{महीं सञ्छाद्य निष्क्रान्तं सपदातिरथं क्षणात्} %||7-19-10||

\twolineshloka
{तद्रावणबलं प्राप्य बलं तस्य महीपतेः}
{प्राणश्यत तदा राजन्हव्यं हुतमिवानले} %||7-19-11||

\twolineshloka
{सोऽपश्यत नरेन्द्रस्तु नश्यमानं महद्बलम्}
{महार्णवं समासाद्य यथा पञ्चापगा जलम्} %||7-19-12||

\twolineshloka
{ततः शक्रधनुःप्रख्यं धनुर्विस्फारयन्स्वयम्}
{आसदाद नरेन्द्रास्तं रावणं क्रोधमूर्छितः} %||7-19-13||

\twolineshloka
{ततो बाणशतान्यष्टौ पातयामास मूर्धनि}
{तस्य राक्षसराजस्य इक्ष्वाकुकुलनन्दनः} %||7-19-14||

\twolineshloka
{तस्य बाणाः पतन्तस्ते चक्रिरे न क्षतं क्वचित्}
{वारिधारा इवाभ्रेभ्यः पतन्त्यो नगमूर्धनि} %||7-19-15||

\twolineshloka
{ततो राक्षसराजेन क्रुद्धेन नृपतिस्तदा}
{तलेन भिहतो मूर्ध्नि स रथान्निपपात ह} %||7-19-16||

\twolineshloka
{स राजा पतितो भूमौ विह्वलाङ्गः प्रवेपितः}
{वज्रदग्ध इवारण्ये सालो निपतितो महान्} %||7-19-17||

\twolineshloka
{तं प्रहस्याब्रवीद्रक्ष इक्ष्वाकुं पृथिवीपतिम्}
{किमिदानीं त्वया प्राप्तं फलं मां प्रति युध्यता} %||7-19-18||

\twolineshloka
{त्रैलोक्ये नास्ति यो द्वन्द्वं मम दद्यान्नराधिप}
{शङ्के प्रमत्तो भोगेषु न शृणोषि बलं मम} %||7-19-19||

\twolineshloka
{तस्यैवं ब्रुवतो राजा मन्दासुर्वाक्यमब्रवीत्}
{किं शक्यमिह कर्तुं वै यत्कालो दुरतिक्रमः} %||7-19-20||

\twolineshloka
{न ह्यहं निर्जितो रक्षस्त्वया चात्मप्रशंसिना}
{कालेनेह विपन्नोऽहं हेतुभूतस्तु मे भवान्} %||7-19-21||

\twolineshloka
{किं त्विदानीं मया शक्यं कर्तुं प्राणपरिक्षये}
{इक्ष्वाकुपरिभावित्वाद्वचो वक्ष्यामि राक्षस} %||7-19-22||

\twolineshloka
{यदि दत्तं यदि हुतं यदि मे सुकृतं तपः}
{यदि गुप्ताः प्रजाः सम्यक्तथा सत्यं वचोऽस्तु मे} %||7-19-23||

\twolineshloka
{उत्पत्स्यते कुले ह्यस्मिन्निक्ष्वाकूणां महात्मनाम्}
{राजा परमतेजस्वी यस्ते प्राणान्हरिष्यति} %||7-19-24||

\twolineshloka
{ततो जलधरोदग्रस्ताडितो देवदुन्दुभिः}
{तस्मिन्नुदाहृते शापे पुष्पवृष्टिश्च खाच्च्युता} %||7-19-25||

\twolineshloka
{ततः स राजा राजेन्द्र गतः स्थानं त्रिविष्टपम्}
{स्वर्गते च नृपे राम राक्षसः स न्यवर्तत} %||7-20-26||


{॥इत्यार्षे श्रीमद्रामायणे वाल्मीकीये आदिकाव्ये उत्तरकाण्डे एकोनविंशः सर्गः॥}

\sect{विंशः सर्गः}

\twolineshloka
{ततो वित्रासयन्मर्त्यान्पृथिव्यां राक्षसाधिपः}
{आससाद घने तस्मिन्नारदं मुनिसत्तमम्} %||7-20-1||

\twolineshloka
{नारदस्तु महातेजा देवर्षिरमितप्रभः}
{अब्रवीन्मेघपृष्ठस्थो रावणं पुष्पके स्थितम्} %||7-20-2||

\twolineshloka
{राक्षसाधिपते सौम्य तिष्ठ विश्रवसः सुत}
{प्रीतोऽस्म्यभिजनोपेत विक्रमैरूर्जितैस्तव} %||7-20-3||

\twolineshloka
{विष्णुना दैत्यघातैश्च तार्क्ष्यस्योरगधर्षणैः}
{त्वया समरमर्दैश्च भृशं हि परितोषितः} %||7-20-4||

\twolineshloka
{किञ्चिद्वक्ष्यामि तावत्ते श्रोतव्यं श्रोष्यसे यदि}
{श्रुत्वा चानन्तरं कार्यं त्वया राक्षसपुङ्गव} %||7-20-5||

\twolineshloka
{किमयं वध्यते लोकस्त्वयावध्येन दैवतैः}
{हत एव ह्ययं लोको यदा मृत्युवशं गतः} %||7-20-6||

\twolineshloka
{पश्य तावन्महाबाहो राक्षसेश्वरमानुषम्}
{लोकमेनं विचित्रार्थं यस्य न ज्ञायते गतिः} %||7-20-7||

\twolineshloka
{क्वचिद्वादित्रनृत्तानि सेव्यन्ते मुदितैर्जनैः}
{रुद्यते चापरैरार्तैर्धाराश्रुनयनाननैः} %||7-20-8||

\twolineshloka
{माता पितृसुतस्नेहैर्भार्या बन्धुमनोरमैः}
{मोहेनायं जनो ध्वस्तः क्लेशं स्वं नावबुध्यते} %||7-20-9||

\twolineshloka
{तत्किमेवं परिक्लिश्य लोकं मोहनिराकृतम्}
{जित एव त्वया सौम्य मर्त्यलोको न संशयः} %||7-20-10||

\twolineshloka
{एवमुक्तस्तु लङ्केशो दीप्यमान इवौजसा}
{अब्रवीन्नारदं तत्र सम्प्रहस्याभिवाद्य च} %||7-20-11||

\twolineshloka
{महर्षे देवगन्धर्वविहार समरप्रिय}
{अहं खलूद्यतो गन्तुं विजयार्थी रसातलम्} %||7-20-12||

\twolineshloka
{ततो लोकत्रयं जित्वा स्थाप्य नागान्सुरान्वशे}
{समुद्रममृतार्थं वै मथिष्यामि रसालयम्} %||7-20-13||

\twolineshloka
{अथाब्रवीद्दशग्रीवं नारदो भगवानृषिः}
{क्व खल्विदानीं मार्गेण त्वयानेन गमिष्यते} %||7-20-14||

\twolineshloka
{अयं खलु सुदुर्गम्यः पितृराज्ञः पुरं प्रति}
{मार्गो गच्छति दुर्धर्षो यमस्यामित्रकर्शन} %||7-20-15||

\twolineshloka
{स तु शारदमेघाभं मुक्त्वा हासं दशाननः}
{उवाच कृतमित्येव वचनं चेदमब्रवीत्} %||7-20-16||

\twolineshloka
{तस्मादेष महाब्रह्मन्वैवस्वतवधोद्यतः}
{गच्छामि दक्षिणामाशां यत्र सूर्यात्मजो नृपः} %||7-20-17||

\twolineshloka
{मया हि भगवन्क्रोधात्प्रतिज्ञातं रणार्थिना}
{अवजेष्यामि चतुरो लोकपालानिति प्रभो} %||7-20-18||

\twolineshloka
{तेनैष प्रस्थितोऽहं वै पितृराजपुरं प्रति}
{प्राणिसङ्क्लेशकर्तारं योजयिष्यामि मृत्युना} %||7-20-19||

\twolineshloka
{एवमुक्त्वा दशग्रीवो मुनिं तमभिवाद्य च}
{प्रययौ दक्षिणामाशां प्रहृष्टैः सह मन्त्रिभिः} %||7-20-20||

\twolineshloka
{नारदस्तु महातेजा मुहूर्तं ध्यानमास्थितः}
{चिन्तयामास विप्रेन्द्रो विधूम इव पावकः} %||7-20-21||

\twolineshloka
{येन लोकास्त्रयः सेन्द्राः क्लिश्यन्ते सचराचराः}
{क्षीणे चायुषि धर्मे च स कालो हिंस्यते कथम्} %||7-20-22||

\twolineshloka
{यस्य नित्यं त्रयो लोका विद्रवन्ति भयार्दिताः}
{तं कथं राक्षसेन्द्रोऽसौ स्वयमेवाभिगच्छति} %||7-20-23||

\twolineshloka
{यो विधाता च धाता च सुकृते दुष्कृते यथा}
{त्रैलोक्यं विजितं येन तं कथं नु विजेष्यति} %||7-20-24||

\twolineshloka
{अपरं किं नु कृत्वैवं विधानं संविधास्यति}
{कौतूहलसमुत्पन्नो यास्यामि यमसादनम्} %||7-21-25||


{॥इत्यार्षे श्रीमद्रामायणे वाल्मीकीये आदिकाव्ये उत्तरकाण्डे विंशः सर्गः॥}

\sect{एकविंशः सर्गः}

\twolineshloka
{एवं सञ्चिन्त्य विप्रेन्द्रो जगाम लघुविक्रमः}
{आख्यातुं तद्यथावृत्तं यमस्य सदनं प्रति} %||7-21-1||

\twolineshloka
{अपश्यत्स यमं तत्र देवमग्निपुरस्कृतम्}
{विधानमुपतिष्ठन्तं प्राणिनो यस्य यादृशम्} %||7-21-2||

\twolineshloka
{स तु दृष्ट्वा यमः प्राप्तं महर्षिं तत्र नारदम्}
{अब्रवीत्सुखमासीनमर्घ्यमावेद्य धर्मतः} %||7-21-3||

\twolineshloka
{कच्चित्क्षेमं नु देवर्षे कच्चिद्धर्मो न नश्यति}
{किमागमनकृत्यं ते देवगन्धर्वसेवित} %||7-21-4||

\twolineshloka
{अब्रवीत्तु तदा वाक्यं नारदो भगवानृषिः}
{श्रूयतामभिधास्यामि विधानं च विधीयताम्} %||7-21-5||

\twolineshloka
{एष नाम्ना दशग्रीवः पितृराज निशाचरः}
{उपयाति वशं नेतुं विक्रमैस्त्वां सुदुर्जयम्} %||7-21-6||

\twolineshloka
{एतेन कारणेनाहं त्वरितोऽस्म्यागतः प्रभो}
{दण्डप्रहरणस्याद्य तव किं नु करिष्यति} %||7-21-7||

\twolineshloka
{एतस्मिन्नन्तरे दूरादंशुमन्तमिवोदितम्}
{ददृशे दिव्यमायान्तं विमानं तस्य रक्षसः} %||7-21-8||

\twolineshloka
{तं देशं प्रभया तस्य पुष्पकस्य महाबलः}
{कृत्वा वितिमिरं सर्वं समीपं समवर्तत} %||7-21-9||

\twolineshloka
{स त्वपश्यन्महाबाहुर्दशग्रीवस्ततस्ततः}
{प्राणिनः सुकृतं कर्म भुञ्जानांश्चैव दुष्कृतम्} %||7-21-10||

\twolineshloka
{ततस्तान्वध्यमानांस्तु कर्मभिर्दुष्कृतैः स्वकैः}
{रावणो मोचयामास विक्रमेण बलाद्बली} %||7-21-11||

\twolineshloka
{प्रेतेषु मुच्यमानेषु राक्षसेन बलीयसा}
{प्रेतगोपाः सुसंरब्धा राक्षसेन्द्रमभिद्रवन्} %||7-21-12||

\twolineshloka
{ते प्रासैः परिघैः शूलैर्मुद्गरैः शक्तितोमरैः}
{पुष्पकं समवर्षन्त शूराः शतसहस्रशः} %||7-21-13||

\twolineshloka
{तस्यासनानि प्रासादान्वेदिकास्तरणानि च}
{पुष्पकस्य बभञ्जुस्ते शीघ्रं मधुकरा इव} %||7-21-14||

\twolineshloka
{देवनिष्ठानभूतं तद्विमानं पुष्पकं मृधे}
{भज्यमानं तथैवासीदक्षयं ब्रह्मतेजसा} %||7-21-15||

\twolineshloka
{ततस्ते रावणामात्या यथाकामं यथाबलम्}
{अयुध्यन्त महावीर्याः स च राजा दशाननः} %||7-21-16||

\twolineshloka
{ते तु शोणितदिग्धाङ्गाः सर्वशस्त्रसमाहताः}
{अमात्या राक्षसेन्द्रस्य चक्रुरायोधनं महत्} %||7-21-17||

\twolineshloka
{अन्योन्यं च महाभागा जघ्नुः प्रहरणैर्युधि}
{यमस्य च महत्सैन्यं राक्षसस्य च मन्त्रिणः} %||7-21-18||

\twolineshloka
{अमात्यांस्तांस्तु सन्त्यज्य राक्षसस्य महौजसः}
{तमेव समधावन्त शूलवर्षैर्दशाननम्} %||7-21-19||

\twolineshloka
{ततः शोणितदिग्धाङ्गः प्रहारैर्जर्जरीकृतः}
{विमाने राक्षसश्रेष्ठः फुल्लाशोक इवाबभौ} %||7-21-20||

\twolineshloka
{स शूलानि गदाः प्रासाञ्शक्तितोमरसायकान्}
{मुसलानि शिलावृक्षान्मुमोचास्त्रबलाद्बली} %||7-21-21||

\twolineshloka
{तांस्तु सर्वान्समाक्षिप्य तदस्त्रमपहत्य च}
{जघ्नुस्ते राक्षसं घोरमेकं शतसहस्रकः} %||7-21-22||

\twolineshloka
{परिवार्य च तं सर्वे शैलं मेघोत्करा इव}
{भिन्दिपालैश्च शूलैश्च निरुच्छ्वासमकारयन्} %||7-21-23||

\twolineshloka
{विमुक्तकवचः क्रुद्धो सिक्तः शोणितविस्रवैः}
{स पुष्पकं परित्यज्य पृथिव्यामवतिष्ठत} %||7-21-24||

\twolineshloka
{ततः स कार्मुकी बाणी पृथिव्यां राक्षसाधिपः}
{लब्धसंज्ञो मुहूर्तेन क्रुद्धस्तस्थौ यथान्तकः} %||7-21-25||

\twolineshloka
{ततः पाशुपतं दिव्यमस्त्रं सन्धाय कार्मुके}
{तिष्ठ तिष्ठेति तानुक्त्वा तच्चापं व्यपकर्षत} %||7-21-26||

\twolineshloka
{ज्वालामाली स तु शरः क्रव्यादानुगतो रणे}
{मुक्तो गुल्मान्द्रुमांश्चैव भस्मकृत्वा प्रधावति} %||7-21-27||

\twolineshloka
{ते तस्य तेजसा दग्धाः सैन्या वैवस्वतस्य तु}
{रणे तस्मिन्निपतिता दावदग्धा नगा इव} %||7-21-28||

\twolineshloka
{ततः स सचिवैः सार्धं राक्षसो भीमविक्रमः}
{ननाद सुमहानादं कम्पयन्निव मेदिनीम्} %||7-22-29||


{॥इत्यार्षे श्रीमद्रामायणे वाल्मीकीये आदिकाव्ये उत्तरकाण्डे एकविंशः सर्गः॥}

\sect{द्वाविंशः सर्गः}

\twolineshloka
{स तु तस्य महानादं श्रुत्वा वैवस्वतो यमः}
{शत्रुं विजयिनं मेने स्वबलस्य च सङ्क्षयम्} %||7-22-1||

\twolineshloka
{स तु योधान्हतान्मत्वा क्रोधपर्याकुलेक्षणः}
{अब्रवीत्त्वरितं सूतं रथः समुपनीयताम्} %||7-22-2||

\twolineshloka
{तस्य सूतो रथं दिव्यमुपस्थाप्य महास्वनम्}
{स्थितः स च महातेजा आरुरोह महारथम्} %||7-22-3||

\twolineshloka
{पाशमुद्गरहस्तश्च मृत्युस्तस्याग्रतो स्थितः}
{येन सङ्क्षिप्यते सर्वं त्रैलोक्यं सचराचरम्} %||7-22-4||

\twolineshloka
{कालदण्डश्च पार्श्वस्थो मूर्तिमान्स्यन्दने स्थितः}
{यमप्रहरणं दिव्यं प्रज्वलन्निव तेजसा} %||7-22-5||

\twolineshloka
{ततो लोकास्त्रयस्त्रस्ताः कम्पन्ते च दिवौकसः}
{कालं क्रुद्धं तदा दृष्ट्वा लोकत्रयभयावहम्} %||7-22-6||

\twolineshloka
{दृष्ट्वा तु ते तं विकृतं रथं मृत्युसमन्वितम्}
{सचिवा राक्षसेन्द्रस्य सर्वलोकभयावहम्} %||7-22-7||

\twolineshloka
{लघुसत्त्वतया सर्वे नष्टसंज्ञा भयार्दिताः}
{नात्र योद्धुं समर्थाः स्म इत्युक्त्वा विप्रदुद्रुवुः} %||7-22-8||

\twolineshloka
{स तु तं तादृशं दृष्ट्वा रथं लोकभयावहम्}
{नाक्षुभ्यत तदा रक्षो व्यथा चैवास्य नाभवत्} %||7-22-9||

\twolineshloka
{स तु रावणमासाद्य विसृजञ्शक्तितोमरान्}
{यमो मर्माणि सङ्क्रुद्धो राक्षसस्य न्यकृन्तत} %||7-22-10||

\twolineshloka
{रावणस्तु स्थितः स्वस्थः शरवर्षं मुमोच ह}
{तस्मिन्वैवस्वतरथे तोयवर्षमिवाम्बुदः} %||7-22-11||

\twolineshloka
{ततो महाशक्तिशतैः पात्यमानैर्महोरसि}
{प्रतिकर्तुं स नाशक्नोद्राक्षसः शल्यपीडितः} %||7-22-12||

\twolineshloka
{नानाप्रहरणैरेवं यमेनामित्रकर्शिना}
{सप्तरात्रं कृते सङ्ख्ये न भग्नो विजितोऽपि वा} %||7-22-13||

\twolineshloka
{ततोऽभवत्पुनर्युद्धं यमराक्षसयोस्तदा}
{विजयाकाङ्क्षिणोस्तत्र समरेष्वनिवर्तिनोः} %||7-22-14||

\twolineshloka
{ततो देवाः सगन्धर्वाः सिद्धाश्च परमर्षयः}
{प्रजापतिं पुरस्कृत्य ददृशुस्तद्रणाजिरम्} %||7-22-15||

\twolineshloka
{संवर्त इव लोकानामभवद्युध्यतोस्तयोः}
{राक्षसानां च मुख्यस्य प्रेतानामीश्वरस्य च} %||7-22-16||

\twolineshloka
{राक्षसेन्द्रस्ततः क्रुद्धश्चापमायम्य संयुगे}
{निरन्तरमिवाकाशं कुर्वन्बाणान्मुमोच ह} %||7-22-17||

\twolineshloka
{मृत्युं चतुर्भिर्विशिखैः सूतं सप्तभिरर्दयत्}
{यमं शरसहस्रेण शीघ्रं मर्मस्वताडयत्} %||7-22-18||

\twolineshloka
{ततः क्रुद्धस्य सहसा यमस्याभिविनिःसृतः}
{ज्वालामालो विनिश्वासो वदनात्क्रोधपावकः} %||7-22-19||

\twolineshloka
{ततोऽपश्यंस्तदाश्चर्यं देवदानवराक्षसाः}
{क्रोधजं पावकं दीप्तं दिधक्षन्तं रिपोर्बलम्} %||7-22-20||

\twolineshloka
{मृत्युस्तु परमक्रुद्धो वैवस्वतमथाब्रवीत्}
{मुञ्च मां देव शीघ्रं त्वं निहन्मि समरे रिपुम्} %||7-22-21||

\twolineshloka
{नरकः शम्बरो वृत्रः शम्भुः कार्तस्वरो बली}
{नमुचिर्विरोचनश्चैव तावुभौ मधुकैटभौ} %||7-22-22||

\twolineshloka
{एते चान्ये च बहवो बलवन्तो दुरासदाः}
{विनिपन्ना मया दृष्टाः का चिन्तास्मिन्निशाचरे} %||7-22-23||

\twolineshloka
{मुञ्च मां साधु धर्मज्ञ यावदेनं निहन्म्यहम्}
{न हि कश्चिन्मया दृष्टो मुहूर्तमपि जीवति} %||7-22-24||

\twolineshloka
{बलं मम न खल्वेतन्मर्यादैषा निसर्गतः}
{संस्पृष्टो हि मया कश्चिन्न जीवेदिति निश्चयः} %||7-22-25||

\twolineshloka
{एतत्तु वचनं श्रुत्वा धर्मराजः प्रतापवान्}
{अब्रवीत्तत्र तं मृत्युमयमेनं निहन्म्यहम्} %||7-22-26||

\twolineshloka
{ततः संरक्तनयनः क्रुद्धो वैवस्वतः प्रभुः}
{कालदण्डममोघं तं तोलयामास पाणिना} %||7-22-27||

\twolineshloka
{यस्य पार्श्वेषु निश्छिद्राः कालपाशाः प्रतिष्ठिताः}
{पावकस्पर्शसङ्काशो मुद्गरो मूर्तिमान्स्थितः} %||7-22-28||

\twolineshloka
{दर्शनादेव यः प्राणान्प्राणिनामुपरुध्यति}
{किं पुनस्ताडनाद्वापि पीडनाद्वापि देहिनः} %||7-22-29||

\twolineshloka
{स ज्वालापरिवारस्तु पिबन्निव निशाचरम्}
{करस्पृष्टो बलवता दण्डः क्रुद्धः सुदारुणः} %||7-22-30||

\twolineshloka
{ततो विदुद्रुवुः सर्वे सत्त्वास्तस्माद्रणाजिरात्}
{सुराश्च क्षुभिता दृष्ट्वा कालदण्डोद्यतं यमम्} %||7-22-31||

\twolineshloka
{तस्मिन्प्रहर्तुकामे तु दण्डमुद्यम्य रावणम्}
{यमं पितामहः साक्षाद्दर्शयित्वेदमब्रवीत्} %||7-22-32||

\twolineshloka
{वैवस्वत महाबाहो न खल्वतुलविक्रम}
{प्रहर्तव्यं त्वयैतेन दण्डेनास्मिन्निशाचरे} %||7-22-33||

\twolineshloka
{वरः खलु मया दत्तस्तस्य त्रिदशपुङ्गव}
{तत्त्वया नानृतं कार्यं यन्मया व्याहृतं वचः} %||7-22-34||

\twolineshloka
{अमोघो ह्येष सर्वासां प्रजानां विनिपातने}
{कालदण्डो मया सृष्टः पूर्वं मृत्युपुरस्कृतः} %||7-22-35||

\twolineshloka
{तन्न खल्वेष ते सौम्य पात्यो राक्षसमूर्धनि}
{न ह्यस्मिन्पतिते कश्चिन्मुहूर्तमपि जीवति} %||7-22-36||

\twolineshloka
{यदि ह्यस्मिन्निपतिते न म्रियेतैष राक्षसः}
{म्रियेत वा दशग्रीवस्तथाप्युभयतोऽनृतम्} %||7-22-37||

\twolineshloka
{राक्षसेन्द्रान्नियच्छाद्य दण्डमेनं वधोद्यतम्}
{सत्यं मम कुरुष्वेदं लोकांस्त्वं समवेक्ष्य च} %||7-22-38||

\twolineshloka
{एवमुक्तस्तु धर्मात्मा प्रत्युवाच यमस्तदा}
{एष व्यावर्तितो दण्डः प्रभविष्णुर्भवान्हि नः} %||7-22-39||

\twolineshloka
{किं त्विदानीं मया शक्यं कर्तुं रणगतेन हि}
{यन्मया यन्न हन्तव्यो राक्षसो वरदर्पितः} %||7-22-40||

\twolineshloka
{एष तस्मात्प्रणश्यामि दर्शनादस्य रक्षसः}
{इत्युक्त्वा सरथः साश्वस्तत्रैवान्तरधीयत} %||7-22-41||

\twolineshloka
{दशग्रीवस्तु तं जित्वा नाम विश्राव्य चात्मनः}
{पुष्पकेण तु संहृष्टो निष्क्रान्तो यमसादनात्} %||7-22-42||

\twolineshloka
{ततो वैवस्वतो देवैः सह ब्रह्मपुरोगमैः}
{जगाम त्रिदिवं हृष्टो नारदश्च महामुनिः} %||7-23-43||


{॥इत्यार्षे श्रीमद्रामायणे वाल्मीकीये आदिकाव्ये उत्तरकाण्डे द्वाविंशः सर्गः॥}

\sect{त्रयोविंशः सर्गः}

\twolineshloka
{स तु जित्वा दशग्रीवो यमं त्रिदशपुङ्गवम्}
{रावणस्तु जयश्लाघी स्वसहायान्ददर्श ह} %||7-23-1||

\twolineshloka
{जयेन वर्धयित्वा च मारीचप्रमुखास्ततः}
{पुष्पकं भेजिरे सर्वे सान्त्विता रवणेन ह} %||7-23-2||

\twolineshloka
{ततो रसातलं हृष्टः प्रविष्टः पयसो निधिम्}
{दैत्योरग गणाध्युष्टं वरुणेन सुरक्षितम्} %||7-23-3||

\twolineshloka
{स तु भोगवतीं गत्वा पुरीं वासुकिपालिताम्}
{स्थाप्य नागान्वशे कृत्वा ययौ मणिमतीं पुरीम्} %||7-23-4||

\twolineshloka
{निवातकवचास्तत्र दैत्या लब्धवरा वसन्}
{राक्षसस्तान्समासाद्य युद्धेन समुपाह्वयत्} %||7-23-5||

\twolineshloka
{ते तु सर्वे सुविक्रान्ता दैतेया बलशालिनः}
{नानाप्रहरणास्तत्र प्रयुद्धा युद्धदुर्मदाः} %||7-23-6||

\twolineshloka
{तेषां तु युध्यमानानां साग्रः संवत्सरो गतः}
{न चान्यतरयोस्तत्र विजयो वा क्षयोऽपि वा} %||7-23-7||

\twolineshloka
{ततः पितामहस्तत्र त्रैलोक्यगतिरव्ययः}
{आजगाम द्रुतं देवो विमानवरमास्थितः} %||7-23-8||

\twolineshloka
{निवातकवचानां तु निवार्य रणकर्म तत्}
{वृद्धः पितामहो वाक्यमुवाच विदितार्थवत्} %||7-23-9||

\twolineshloka
{न ह्ययं रावणो युद्धे शक्यो जेतुं सुरासुरैः}
{न भवन्तः क्षयं नेतुं शक्याः सेन्द्रैः सुरासुरैः} %||7-23-10||

\twolineshloka
{राक्षसस्य सखित्वं वै भवद्भिः सह रोचते}
{अविभक्ता हि सर्वार्थाः सुहृदां नात्र संशयः} %||7-23-11||

\twolineshloka
{ततोऽग्निसाक्षिकं सख्यं कृतवांस्तत्र रावणः}
{निवातकवचैः सार्धं प्रीतिमानभवत्तदा} %||7-23-12||

\twolineshloka
{अर्चितस्तैर्यथान्यायं संवत्सरसुखोषितः}
{स्वपुरान्निर्विशेषं च पूजां प्राप्तो दशाननः} %||7-23-13||

\twolineshloka
{स तूपधार्य मायानां शतमेकोनमात्मवान्}
{सलिलेन्द्रपुरान्वेषी स बभ्राम रसातलम्} %||7-23-14||

\twolineshloka
{ततोऽश्मनगरं नाम कालकेयाभिरक्षितम्}
{तं विजित्य मुहूर्तेन जघ्ने दैत्यांश्चतुःशतम्} %||7-23-15||

\twolineshloka
{ततः पाण्डुरमेघाभं कैलासमिव संस्थितम्}
{वरुणस्यालयं दिव्यमपश्यद्राक्षसाधिपः} %||7-23-16||

\twolineshloka
{क्षरन्तीं च पयो नित्यं सुरभिं गामवस्थिताम्}
{यस्याः पयोविनिष्यन्दात्क्षीरोदो नाम सागरः} %||7-23-17||

\threelineshloka
{यस्माच्चन्द्रः प्रभवति शीतरश्मिः प्रजाहितः}
{यं समासाद्य जीवन्ति फेनपाः परमर्षयः}
{अमृतं यत्र चोत्पन्नं सुरा चापि सुराशिनाम्} %||7-23-18||

\threelineshloka
{यां ब्रुवन्ति नरा लोके सुरभिं नाम नामतः}
{प्रदक्षिणं तु तां कृत्वा रावणः परमाद्भुताम्}
{प्रविवेश महाघोरं गुप्तं बहुविधैर्बलैः} %||7-23-19||

\twolineshloka
{ततो धाराशताकीर्णं शारदाभ्रनिभं तदा}
{नित्यप्रहृष्टं ददृशे वरुणस्य गृहोत्तमम्} %||7-23-20||

\twolineshloka
{ततो हत्वा बलाध्यक्षान्समरे तैश्च ताडितः}
{अब्रवीत्क्व गतो यो वो राजा शीघ्रं निवेद्यताम्} %||7-23-21||

\twolineshloka
{युद्धार्थी रावणः प्राप्तस्तस्य युद्धं प्रदीयताम्}
{वद वा न भयं तेऽस्ति निर्जितोऽस्मीति साञ्जलिः} %||7-23-22||

\twolineshloka
{एतस्मिन्नन्तरे क्रुद्धा वरुणस्य महात्मनः}
{पुत्राः पौत्राश्च निष्क्रामन्गौश्च पुष्कर एव च} %||7-23-23||

\twolineshloka
{ते तु वीर्यगुणोपेता बलैः परिवृताः स्वकैः}
{युक्त्वा रथान्कामगमानुद्यद्भास्करवर्चसः} %||7-23-24||

\twolineshloka
{ततो युद्धं समभवद्दारुणं लोमहर्षणम्}
{सलिलेन्द्रस्य पुत्राणां रावणस्य च रक्षसः} %||7-23-25||

\twolineshloka
{अमात्यैस्तु महावीर्यैर्दशग्रीवस्य रक्षसः}
{वारुणं तद्बलं कृत्स्नं क्षणेन विनिपातितम्} %||7-23-26||

\twolineshloka
{समीक्ष्य स्वबलं सङ्ख्ये वरुणस्या सुतास्तदा}
{अर्दिताः शरजालेन निवृत्ता रणकर्मणः} %||7-23-27||

\twolineshloka
{महीतलगतास्ते तु रावणं दृश्य पुष्पके}
{आकाशमाशु विविशुः स्यन्दनैः शीघ्रगामिभिः} %||7-23-28||

\twolineshloka
{महदासीत्ततस्तेषां तुल्यं स्थानमवाप्य तत्}
{आकाशयुद्धं तुमुलं देवदानवयोरिव} %||7-23-29||

\twolineshloka
{ततस्ते रावणं युद्धे शरैः पावकसंनिभैः}
{विमुखीकृत्य संहृष्टा विनेदुर्विविधान्रवान्} %||7-23-30||

\twolineshloka
{ततो महोदरः क्रुद्धो राजानं दृश्य धर्षितम्}
{त्यक्त्वा मृत्युभयं वीरो युद्धकाङ्क्षी व्यलोकयत्} %||7-23-31||

\twolineshloka
{तेन तेषां हया ये च कामगाः पवनोपमाः}
{महोदरेण गदया हतास्ते प्रययुः क्षितिम्} %||7-23-32||

\twolineshloka
{तेषां वरुणसूनूनां हत्वा योधान्हयांश्च तान्}
{मुमोचाशु महानादं विरथान्प्रेक्ष्य तान्स्थितान्} %||7-23-33||

\twolineshloka
{ते तु तेषां रथाः साश्वाः सह सारथिभिर्वरैः}
{महोदरेण निहताः पतिताः पृथिवीतले} %||7-23-34||

\twolineshloka
{ते तु त्यक्त्वा रथान्पुत्रा वरुणस्य महात्मनः}
{आकाशे विष्ठिताः शूराः स्वप्रभावान्न विव्यथुः} %||7-23-35||

\twolineshloka
{धनूंषि कृत्वा सज्यानि विनिर्भिद्य महोदरम्}
{रावणं समरे क्रुद्धाः सहिताः समभिद्रवन्} %||7-23-36||

\twolineshloka
{ततः क्रुद्धो दशग्रीवः कालाग्निरिव विष्ठितः}
{शरवर्षं महावेगं तेषां मर्मस्वपातयत्} %||7-23-37||

\threelineshloka
{मुसलानि विचित्राणि ततो भल्लशतानि च}
{पट्टसांश्चैव शक्तीश्च शतघ्नीस्तोमरांस्तथा}
{पातयामास दुर्धर्षस्तेषामुपरि विष्ठितः} %||7-23-38||

{अथ विद्धास्तु ते वीरा विनिष्पेतुः पदातयः॥\devanumber{39}॥} %||7-23-39|| (Check)

\addtocounter{shlokacount}{1}

\twolineshloka
{ततो रक्षो महानादं मुक्त्वा हन्ति स्म वारुणान्}
{नानाप्रहरणैर्घोरैर्धारापातैरिवाम्बुदः} %||7-23-40||

\twolineshloka
{ततस्ते विमुखाः सर्वे पतिता धरणीतले}
{रणात्स्वपुरुषैः शीघ्रं गृहाण्येव प्रवेशिताः} %||7-23-41||

\twolineshloka
{तानब्रवीत्ततो रक्षो वरुणाय निवेद्यताम्}
{रावणं चाब्रवीन्मन्त्री प्रभासो नाम वारुणः} %||7-23-42||

\twolineshloka
{गतः खलु महातेजा ब्रह्मलोकं जलेश्वरः}
{गान्धर्वं वरुणः श्रोतुं यं त्वमाह्वयसे युधि} %||7-23-43||

\twolineshloka
{तत्किं तव वृथा वीर परिश्राम्य गते नृपे}
{ये तु संनिहिता वीराः कुमारास्ते पराजिताः} %||7-23-44||

\twolineshloka
{राक्षसेन्द्रस्तु तच्छ्रुत्वा नाम विश्राव्य चात्मनः}
{हर्षान्नादं विमुञ्चन्वै निष्क्रान्तो वरुणालयात्} %||7-23-45||

\twolineshloka
{आगतस्तु पथा येन तेनैव विनिवृत्य सः}
{लङ्कामभिमुखो रक्षो नभस्तलगतो ययौ} %||7-24-46||


{॥इत्यार्षे श्रीमद्रामायणे वाल्मीकीये आदिकाव्ये उत्तरकाण्डे त्रयोविंशः सर्गः॥}

\sect{चतुर्विंशः सर्गः}

\twolineshloka
{निवर्तमानः संहृष्टो रावणः स दुरात्मवान्}
{जह्रे पथि नरेन्द्रर्षिदेवगन्धर्वकन्यकाः} %||7-24-1||

\twolineshloka
{दर्शनीयां हि यां रक्षः कन्यां स्त्रीं वाथ पश्यति}
{हत्वा बन्धुजनं तस्या विमाने संन्यवेशयत्} %||7-24-2||

\twolineshloka
{तत्र पन्नगयक्षाणां मानुषाणां च रक्षसाम्}
{दैत्यानां दानवानां च कन्या जग्राह रावणः} %||7-24-3||

\twolineshloka
{दीर्घकेश्यः सुचार्वङ्ग्यः पूर्णचन्द्रनिभाननाः}
{शोकायत्तास्तरुण्यश्च समस्ता स्तननम्रिताः} %||7-24-4||

\twolineshloka
{तुल्यमग्न्यर्चिषां तत्र शोकाग्निभयसम्भवम्}
{प्रवेपमाना दुःखार्ता मुमुचुर्बाष्पजं जलम्} %||7-24-5||

\twolineshloka
{तासां निश्वसमानानां निश्वसैः सम्प्रदीपितम्}
{अग्निहोत्रमिवाभाति संनिरुद्धाग्निपुष्पकम्} %||7-24-6||

\threelineshloka
{काचिद्दध्यौ सुदुःखार्ता हन्यादपि हि मामयम्}
{स्मृत्वा मातॄह्पितॄन्भ्रातॄन्पुत्रान्वै श्वशुरानपि}
{दुःखशोकसमाविष्टा विलेपुः सहिताः स्त्रियः} %||7-24-7||

\twolineshloka
{कथं नु खलु मे पुत्रः करिष्यति मया विना}
{कथं माता कथं भ्राता निमग्नाः शोकसागरे} %||7-24-8||

\twolineshloka
{हा कथं नु करिष्यामि भर्तारं दैवतं विना}
{मृत्यो प्रसीद याचे त्वां नय मां यमसादनम्} %||7-24-9||

\twolineshloka
{किं नु मे दुष्कृतं कर्म कृतं देहान्तरे पुरा}
{ततोऽस्मि धर्षितानेन पतिता शोकसागरे} %||7-24-10||

\twolineshloka
{न खल्विदानीं पश्यामि दुःखस्यान्तमिहात्मनः}
{अहो धिन्मानुषाँल्लोकान्नास्ति खल्वधमः परः} %||7-24-11||

\twolineshloka
{यद्दुर्बला बलवता बान्धवा रावणेन मे}
{उदितेनैव सूर्येण तारका इव नाशिताः} %||7-24-12||

\twolineshloka
{अहो सुबलवद्रक्षो वधोपायेषु रज्यते}
{अहो दुर्वृत्तमात्मानं स्वयमेव न बुध्यते} %||7-24-13||

\twolineshloka
{सर्वथा सदृशस्तावद्विक्रमोऽस्य दुरात्मनः}
{इदं त्वसदृशं कर्म परदाराभिमर्शनम्} %||7-24-14||

\twolineshloka
{यस्मादेष परख्यासु स्त्रीषु रज्यति दुर्मतिः}
{तस्माद्धि स्त्रीकृतेनैव वधं प्राप्स्यति वारणः} %||7-24-15||

\twolineshloka
{शप्तः स्त्रीभिः स तु तदा हततेजाः सुनिष्प्रभ}
{पतिव्रताभिः साध्वीभिः स्थिताभिः साधु वर्त्मनि} %||7-24-16||

\twolineshloka
{एवं विलपमानासु रावणो राक्षसाधिपः}
{प्रविवेश पुरीं लङ्कां पूज्यमानो निशाचरैः} %||7-24-17||

\twolineshloka
{ततो राक्षसराजस्य स्वसा परमदुःखिता}
{पादयोः पतिता तस्य वक्तुमेवोपचक्रमे} %||7-24-18||

\twolineshloka
{ततः स्वसारमुत्थाप्य रावणः परिसान्त्वयन्}
{अब्रवीत्किमिदं भद्रे वक्तुमर्हसि मे द्रुतम्} %||7-24-19||

\twolineshloka
{सा बाष्पपरिरुद्धाक्षी राक्षसी वाक्यमब्रवीत्}
{हतास्मि विधवा राजंस्त्वया बलवता कृता} %||7-24-20||

\twolineshloka
{एते विर्यात्त्वया राजन्दैत्या विनिहता रणे}
{कालकेया इति ख्याता महाबलपराक्रमाः} %||7-24-21||

\twolineshloka
{तत्र मे निहतो भर्ता गरीयाञ्जीवितादपि}
{स त्वया दयितस्तत्र भ्रात्रा शत्रुसमेन वै} %||7-24-22||

\twolineshloka
{या त्वयास्मि हता राजन्स्वयमेवेह बन्धुना}
{दुःखं वैधव्यशब्दं च दत्तं भोक्ष्याम्यहं त्वया} %||7-24-23||

\twolineshloka
{ननु नाम त्वया रक्ष्यो जामाता समरेष्वपि}
{तं निहत्य रणे राजन्स्वयमेव न लज्जसे} %||7-24-24||

\twolineshloka
{एवमुक्तस्तया रक्षो भगिन्या क्रोशमानया}
{अब्रवीत्सान्त्वयित्वा तां सामपूर्वमिदं वचः} %||7-24-25||

\twolineshloka
{अलं वत्से विषादेन न भेतव्यं च सर्वशः}
{मानदानविशेषैस्त्वां तोषयिष्यामि नित्यशः} %||7-24-26||

\threelineshloka
{युद्धे प्रमत्तो व्याक्षिप्तो जयकाङ्क्षी क्षिपञ्शरान्}
{नावगच्छामि युद्धेषु स्वान्परान्वाप्यहं शुभे}
{तेनासौ निहतः सङ्ख्ये मया भर्ता तव स्वसः} %||7-24-27||

\twolineshloka
{अस्मिन्काले तु यत्प्राप्तं तत्करिष्यामि ते हितम्}
{भ्रातुरैश्वर्यसंस्थस्य खरस्य भव पार्श्वतः} %||7-24-28||

\twolineshloka
{चतुर्दशानां भ्राता ते सहस्राणां भविष्यति}
{प्रभुः प्रयाणे दाने च राक्षसानां महौजसाम्} %||7-24-29||

\twolineshloka
{तत्र मातृष्वसुः पुत्रो भ्राता तव खरः प्रभुः}
{भविष्यति सदा कुर्वन्यद्वक्ष्यसि वचः स्वयम्} %||7-24-30||

\twolineshloka
{शीघ्रं गच्छत्वयं शूरो दण्डकान्परिरक्षितुम्}
{दूषणोऽस्य बलाध्यक्षो भविष्यति महाबलः} %||7-24-31||

\twolineshloka
{स हि शप्तो वनोद्देशः क्रुद्धेनोशनसा पुरा}
{राक्षसानामयं वासो भविष्यति न संशयः} %||7-24-32||

\twolineshloka
{एवमुक्त्वा दशग्रीवः सैन्यं तस्यादिदेश ह}
{चतुर्दश सहस्राणि रक्षसां कामरूपिणाम्} %||7-24-33||

\twolineshloka
{स तैः सर्वैः परिवृतो राक्षसैर्घोरदर्शनैः}
{खरः सम्प्रययौ शीघ्रं दण्डकानकुतोभयः} %||7-24-34||

\twolineshloka
{स तत्र कारयामास राज्यं निहतकण्टकम्}
{सा च शूर्पणखा प्रीता न्यवसद्दण्डकावने} %||7-25-35||


{॥इत्यार्षे श्रीमद्रामायणे वाल्मीकीये आदिकाव्ये उत्तरकाण्डे चतुर्विंशः सर्गः॥}

\sect{पञ्चविंशः सर्गः}

\twolineshloka
{स तु दत्त्वा दशग्रीवो वनं घोरं खरस्य तत्}
{भगिनीं च समाश्वास्य हृष्टः स्वस्थतरोऽभवत्} %||7-25-1||

\twolineshloka
{ततो निकुम्भिला नाम लङ्कायाः काननं महत्}
{महात्मा राक्षसेन्द्रस्तत्प्रविवेश सहानुगः} %||7-25-2||

\twolineshloka
{तत्र यूपशताकीर्णं सौम्यचैत्योपशोभितम्}
{ददर्श विष्ठितं यज्ञं सम्प्रदीप्तमिव श्रिया} %||7-25-3||

\twolineshloka
{ततः कृष्णाजिनधरं कमण्डलुशिखाध्वजम्}
{ददर्श स्वसुतं तत्र मेघनादमरिन्दमम्} %||7-25-4||

\twolineshloka
{रक्षःपतिः समासाद्य समाश्लिष्य च बाहुभिः}
{अब्रवीत्किमिदं वत्स वर्तते तद्ब्रवीहि मे} %||7-25-5||

\twolineshloka
{उशना त्वब्रवीत्तत्र गुरुर्यज्ञसमृद्धये}
{रावणं राक्षसश्रेष्ठं द्विजश्रेष्ठो महातपाः} %||7-25-6||

\twolineshloka
{अहमाख्यामि ते राजञ्श्रूयतां सर्वमेव च}
{यज्ञास्ते सप्त पुत्रेण प्राप्ताः सुबहुविस्तराः} %||7-25-7||

\twolineshloka
{अग्निष्टोमोऽश्वमेधश्च यज्ञो बहुसुवर्णकः}
{राजसूयस्तथा यज्ञो गोमेधो वैष्णवस्तथा} %||7-25-8||

\twolineshloka
{माहेश्वरे प्रवृत्ते तु यज्ञे पुम्भिः सुदुर्लभे}
{वरांस्ते लब्धवान्पुत्रः साक्षात्पशुपतेरिह} %||7-25-9||

\twolineshloka
{कामगं स्यन्दनं दिव्यमन्तरिक्षचरं ध्रुवम्}
{मायां च तामसीं नाम यया सम्पद्यते तमः} %||7-25-10||

\twolineshloka
{एतया किल सङ्ग्रामे मायया राक्षसेश्वर}
{प्रयुद्धस्य गतिः शक्या न हि ज्ञातुं सुरासुरैः} %||7-25-11||

\twolineshloka
{अक्षयाविषुधी बाणैश्चापं चापि सुदुर्जयम्}
{अस्त्रं च बलवत्सौम्य शत्रुविध्वंसनं रणे} %||7-25-12||

\twolineshloka
{एतान्सर्वान्वराँल्लब्ध्वा पुत्रस्तेऽयं दशानन}
{अद्य यज्ञसमाप्तौ च त्वत्प्रतीक्षः स्थितो अहम्} %||7-25-13||

\twolineshloka
{ततोऽब्रवीद्दशग्रीवो न शोभनमिदं कृतम्}
{पूजिताः शत्रवो यस्माद्द्रव्यैरिन्द्रपुरोगमाः} %||7-25-14||

\twolineshloka
{एहीदानीं कृतं यद्धि तदकर्तुं न शक्यते}
{आगच्छ सौम्य गच्छामः स्वमेव भवनं प्रति} %||7-25-15||

\twolineshloka
{ततो गत्वा दशग्रीवः सपुत्रः सविभीषणः}
{स्त्रियोऽवतारयामास सर्वास्ता बाष्पविक्लवाः} %||7-25-16||

\twolineshloka
{लक्षिण्यो रत्नबूताश्च देवदानवरक्षसाम्}
{नानाभूषणसम्पन्ना ज्वलन्त्यः स्वेन तेजसा} %||7-25-17||

\twolineshloka
{विभीषणस्तु ता नारीर्दृष्ट्वा शोकसमाकुलाः}
{तस्य तां च मतिं ज्ञात्वा धर्मात्मा वाक्यमब्रवीत्} %||7-25-18||

\twolineshloka
{ईदृशैस्तैः समाचारैर्यशोऽर्थकुलनाशनैः}
{धर्षणं प्राणिनां दत्त्वा स्वमतेन विचेष्टसे} %||7-25-19||

\twolineshloka
{ज्ञातीन्वै धर्षयित्वेमास्त्वयानीता वराङ्गनाः}
{त्वामतिक्रम्य मधुना राजन्कुम्भीनसी हृता} %||7-25-20||

\twolineshloka
{रावणस्त्वब्रवीद्वाक्यं नावगच्छामि किं त्विदम्}
{को वायं यस्त्वयाख्यातो मधुरित्येव नामतः} %||7-25-21||

\twolineshloka
{विभीषणस्तु सङ्क्रुद्धो भ्रातरं वाक्यमब्रवीत्}
{श्रूयतामस्य पापस्य कर्मणः फलमागतम्} %||7-25-22||

\twolineshloka
{मातामहस्य योऽस्माकं ज्येष्ठो भ्राता सुमालिनः}
{माल्यवानिति विख्यातो वृद्धप्राज्ञो निशाचरः} %||7-25-23||

\twolineshloka
{पितुर्ज्येष्ठो जनन्याश्च अस्माकं त्वार्यकोऽभवत्}
{तस्य कुम्भीनसी नाम दुहितुर्दुहिताभवत्} %||7-25-24||

\twolineshloka
{मातृष्वसुरथास्माकं सा कन्या चानलोद्भवा}
{भवत्यस्माकमेषा वै भ्रातॄणां धर्मतः स्वसा} %||7-25-25||

\twolineshloka
{सा हृता मधुना राजन्राक्षसेन बलीयसा}
{यज्ञप्रवृत्ते पुत्रे ते मयि चान्तर्जलोषिते} %||7-25-26||

\twolineshloka
{निहत्य राक्षसश्रेष्ठानमात्यांस्तव सम्मतान्}
{धर्षयित्वा हृता राजन्गुप्ता ह्यन्तःपुरे तव} %||7-25-27||

\threelineshloka
{श्रुत्वा त्वेतन्महाराज क्षान्तमेव हतो न सः}
{यस्मादवश्यं दातव्या कन्या भर्त्रे हि दातृभिः}
{अस्मिन्नेवाभिसम्प्राप्तं लोके विदितमस्तु ते} %||7-25-28||

\twolineshloka
{ततोऽब्रवीद्दशग्रीवः क्रुद्धः संरक्तलोचनः}
{कल्प्यतां मे रथः शीघ्रं शूराः सज्जीभवन्तु च} %||7-25-29||

\twolineshloka
{भ्राता मे कुम्भकर्णश्च ये च मुख्या निशाचराः}
{वाहनान्यधिरोहन्तु नानाप्रहरणायुधाः} %||7-25-30||

\twolineshloka
{अद्य तं समरे हत्वा मधुं रावणनिर्भयम्}
{इन्द्रलोकं गमिष्यामि युद्धकाङ्क्षी सुहृद्वृतः} %||7-25-31||

\twolineshloka
{ततो विजित्य त्रिदिवं वशे स्थाप्य पुरन्दरम्}
{निर्वृतो विहरिष्यामि त्रैलोक्यैश्वर्यशोभितः} %||7-25-32||

\twolineshloka
{अक्षौहिणीसहस्राणि चत्वार्युग्राणि रक्षसाम्}
{नानाप्रहरणान्याशु निर्ययुर्युद्धकाङ्क्षिणाम्} %||7-25-33||

\twolineshloka
{इन्द्रजित्त्वग्रतः सैन्यं सैनिकान्परिगृह्य च}
{रावणो मध्यतः शूरः कुम्भकर्णश्च पृष्ठतः} %||7-25-34||

\twolineshloka
{विभीषणस्तु धर्मात्मा लङ्कायां धर्ममाचरत्}
{ते तु सर्वे महाभागा ययुर्मधुपुरं प्रति} %||7-25-35||

\twolineshloka
{रथैर्नागैः खरैरुष्ट्रैर्हयैर्दीप्तैर्महोरगैः}
{राक्षसाः प्रययुः सर्वे कृत्वाकाशं निरन्तरम्} %||7-25-36||

\twolineshloka
{दैत्याश्च शतशस्तत्र कृतवैराः सुरैः सह}
{रावणं प्रेक्ष्य गच्छन्तमन्वगच्छन्त पृष्ठतः} %||7-25-37||

\twolineshloka
{स तु गत्वा मधुपुरं प्रविश्य च दशाननः}
{न ददर्श मधुं तत्र भगिनीं तत्र दृष्टवान्} %||7-25-38||

\twolineshloka
{सा प्रह्वा प्राञ्जलिर्भूत्वा शिरसा पादयोर्गता}
{तस्य राक्षसराजस्य त्रस्ता कुम्भीनसी स्वसा} %||7-25-39||

\twolineshloka
{तां समुत्थापयामास न भेतव्यमिति ब्रुवन्}
{रावणो राक्षसश्रेष्ठः किं चापि करवाणि ते} %||7-25-40||

\twolineshloka
{साब्रवीद्यदि मे राजन्प्रसन्नस्त्वं महाबल}
{भर्तारं न ममेहाद्य हन्तुमर्हसि मानद} %||7-25-41||

\twolineshloka
{सत्यवाग्भव राजेन्द्र मामवेक्षस्व याचतीम्}
{त्वया ह्युक्तं महाबाहो न भेतव्यमिति स्वयम्} %||7-25-42||

\twolineshloka
{रावणस्त्वब्रवीद्धृष्टः स्वसारं तत्र संस्थिताम्}
{क्व चासौ तव भर्ता वै मम शीघ्रं निवेद्यताम्} %||7-25-43||

\twolineshloka
{सह तेन गमिष्यामि सुरलोकं जयाय वै}
{तव कारुण्यसौहार्दान्निवृत्तोऽस्मि मधोर्वधात्} %||7-25-44||

\twolineshloka
{इत्युक्ता सा प्रसुप्तं तं समुत्थाप्य निशाचरम्}
{अब्रवीत्सम्प्रहृष्टेव राक्षसी सुविपश्चितम्} %||7-25-45||

\twolineshloka
{एष प्राप्तो दशग्रीवो मम भ्राता निशाचरः}
{सुरलोकजयाकाङ्क्षी साहाय्ये त्वां वृणोति च} %||7-25-46||

\twolineshloka
{तदस्य त्वं सहायार्थं सबन्धुर्गच्छ राक्षस}
{स्निग्धस्य भजमानस्य युक्तमर्थाय कल्पितुम्} %||7-25-47||

\twolineshloka
{तस्यास्तद्वचनं श्रुत्वा तथेत्याह मधुर्वचः}
{ददर्श राक्षसश्रेष्ठं यथान्यायमुपेत्य सः} %||7-25-48||

\threelineshloka
{पूजयामास धर्मेण रावणं राक्षसाधिपम्}
{प्राप्तपूजो दशग्रीवो मधुवेश्मनि वीर्यवान्}
{तत्र चैकां निशामुष्य गमनायोपचक्रमे} %||7-25-49||

\twolineshloka
{ततः कैलासमासाद्य शैलं वैश्रवणालयम्}
{राक्षसेन्द्रो महेन्द्राभः सेनामुपनिवेशयत्} %||7-26-50||


{॥इत्यार्षे श्रीमद्रामायणे वाल्मीकीये आदिकाव्ये उत्तरकाण्डे पञ्चविंशः सर्गः॥}

\sect{षड्विंशः सर्गः}

\twolineshloka
{स तु तत्र दशग्रीवः सह सैन्येन वीर्यवान्}
{अस्तं प्राप्ते दिनकरे निवासं समरोचयत्} %||7-26-1||

\twolineshloka
{उदिते विमले चन्द्रे तुल्यपर्वतवर्चसि}
{स ददर्श गुणांस्तत्र चन्द्रपादोपशोभितान्} %||7-26-2||

\twolineshloka
{कर्णिकारवनैर्दिव्यैः कदम्बगहनैस्तथा}
{पद्मिनीभिश्च फुल्लाभिर्मन्दाकिन्या जलैरपि} %||7-26-3||

\twolineshloka
{घण्टानामिव संनादः शुश्रुवे मधुरस्वनः}
{अप्सरोगणसङ्घनां गायतां धनदालये} %||7-26-4||

\twolineshloka
{पुष्पवर्षाणि मुञ्चन्तो नगाः पवनताडिताः}
{शैलं तं वासयन्तीव मधुमाधवगन्धिनः} %||7-26-5||

\twolineshloka
{मधुपुष्परजःपृक्तं गन्धमादाय पुष्कलम्}
{प्रववौ वर्धयन्कामं रावणस्य सुखोऽनिलः} %||7-26-6||

\twolineshloka
{गेयात्पुष्पसमृद्ध्या च शैत्याद्वायोर्गुणैर्गिरेः}
{प्रवृत्तायां रजन्यां च चन्द्रस्योदयनेन च} %||7-26-7||

\twolineshloka
{रावणः सुमहावीर्यः कामबाणवशं गतः}
{विनिश्वस्य विनिश्वस्य शशिनं समवैक्षत} %||7-26-8||

\twolineshloka
{एतस्मिन्नन्तरे तत्र दिव्यपुष्पविभूषिता}
{सर्वाप्सरोवरा रम्भा पूर्णचन्द्रनिभानना} %||7-26-9||

\twolineshloka
{कृतैर्विशेषकैरार्द्रैः षडर्तुकुसुमोत्सवैः}
{नीलं सतोयमेघाभं वस्त्रं समवगुण्ठिता} %||7-26-10||

\threelineshloka
{यस्य वक्त्रं शशिनिभं भ्रुवौ चापनिभे शुभे}
{ऊरू करिकराकारौ करौ पल्लवकोमलौ}
{सैन्यमध्येन गच्छन्ती रावणेनोपलक्षिता} %||7-26-11||

\twolineshloka
{तां समुत्थाय रक्षेन्द्रः कामबाणबलार्दितः}
{करे गृहीत्वा गच्छन्तीं स्मयमानोऽभ्यभाषत} %||7-26-12||

\twolineshloka
{क्व गच्छसि वरारोहे कां सिद्धिं भजसे स्वयम्}
{कस्याभ्युदयकालोऽयं यस्त्वां समुपभोक्ष्यते} %||7-26-13||

\twolineshloka
{तवाननरसस्याद्य पद्मोत्पलसुगन्धिनः}
{सुधामृतरसस्येव कोऽद्य तृप्तिं गमिष्यति} %||7-26-14||

\twolineshloka
{स्वर्णकुम्भनिभौ पीनौ शुभौ भीरु निरन्तरौ}
{कस्योरस्थलसंस्पर्शं दास्यतस्ते कुचाविमौ} %||7-26-15||

\twolineshloka
{सुवर्णचक्रप्रतिमं स्वर्णदामचितं पृथु}
{अध्यारोक्ष्यति कस्तेऽद्य स्वर्गं जघनरूपिणम्} %||7-26-16||

\twolineshloka
{मद्विशिष्टः पुमान्कोऽन्यः शक्रो विष्णुरथाश्विनौ}
{मामतीत्य हि यस्य त्वं यासि भीरु न शोभनम्} %||7-26-17||

\twolineshloka
{विश्रम त्वं पृथुश्रोणि शिलातलमिदं शुभम्}
{त्रैलोक्ये यः प्रभुश्चैव तुल्यो मम न विद्यते} %||7-26-18||

\twolineshloka
{तदेष प्राञ्जलिः प्रह्वो याचते त्वां दशाननः}
{यः प्रभुश्चापि भर्ता च त्रैलोक्यस्य भजस्व माम्} %||7-26-19||

\twolineshloka
{एवमुक्ताब्रवीद्रम्भा वेपमाना कृताञ्जलिः}
{प्रसीद नार्हसे वक्तुमीदृशं त्वं हि मे गुरुः} %||7-26-20||

\twolineshloka
{अन्येभ्योऽपि त्वया रक्ष्या प्राप्नुयां धर्षणं यदि}
{धर्मतश्च स्नुषा तेऽहं तत्त्वमेतद्ब्रवीमि ते} %||7-26-21||

\twolineshloka
{अब्रवीत्तां दशग्रीवश्चरणाधोमुखीं स्थिताम्}
{सुतस्य यदि मे भार्या ततस्त्वं मे स्नुषा भवेः} %||7-26-22||

\twolineshloka
{बाढमित्येव सा रम्भा प्राह रावणमुत्तरम्}
{धर्मतस्ते सुतस्याहं भार्या राक्षसपुङ्गव} %||7-26-23||

\twolineshloka
{पुत्रः प्रियतरः प्राणैर्भ्रातुर्वैश्रवणस्य ते}
{ख्यातो यस्त्रिषु लोकेषु नलकूबर इत्यसौ} %||7-26-24||

\twolineshloka
{धर्मतो यो भवेद्विप्रः क्षत्रियो वीर्यतो भवेत्}
{क्रोधाद्यश्च भवेदग्निः क्षान्त्या च वसुधासमः} %||7-26-25||

\twolineshloka
{तस्यास्मि कृतसङ्केता लोकपालसुतस्य वै}
{तमुद्दिश्य च मे सर्वं विभूषणमिदं कृतम्} %||7-26-26||

\twolineshloka
{यस्य तस्य हि नान्यस्य भावो मां प्रति तिष्ठति}
{तेन सत्येन मां राजन्मोक्तुमर्हस्यरिन्दम} %||7-26-27||

\twolineshloka
{स हि तिष्ठति धर्मात्मा साम्प्रतं मत्समुत्सुकः}
{तन्न विघ्नं सुतस्येह कर्तुमर्हसि मुञ्च माम्} %||7-26-28||

\twolineshloka
{सद्भिराचरितं मार्गं गच्छ राक्षसपुङ्गव}
{माननीयो मया हि त्वं लालनीया तथास्मि ते} %||7-26-29||

\threelineshloka
{एवं ब्रुवाणां रम्भां तां धर्मार्थसहितं वचः}
{निर्भर्त्स्य राक्षसो मोहात्प्रतिगृह्य बलाद्बली}
{काममोहाभिसंरब्धो मैथुनायोपचक्रमे} %||7-26-30||

\twolineshloka
{सा विमुक्ता ततो रम्भा भ्रष्टमाल्यविभूषणा}
{गजेन्द्राक्रीडमथिता नदीवाकुलतां गता} %||7-26-31||

\twolineshloka
{सा वेपमाना लज्जन्ती भीता करकृताञ्जलिः}
{नलकूबरमासाद्य पादयोर्निपपात ह} %||7-26-32||

\twolineshloka
{तदवस्थां च तां दृष्ट्वा महात्मा नलकूबरः}
{अब्रवीत्किमिदं भद्रे पादयोः पतितासि मे} %||7-26-33||

\twolineshloka
{सा तु निश्वसमाना च वेपमानाथ साञ्जलिः}
{तस्मै सर्वं यथातथ्यमाख्यातुमुपचक्रमे} %||7-26-34||

\twolineshloka
{एष देव दशग्रीवः प्राप्तो गन्तुं त्रिविष्टपम्}
{तेन सैन्यसहायेन निशेह परिणाम्यते} %||7-26-35||

\twolineshloka
{आयान्ती तेन दृष्टास्मि त्वत्सकाशमरिन्दम}
{गृहीत्वा तेन पृष्टास्मि कस्य त्वमिति रक्षसा} %||7-26-36||

\twolineshloka
{मया तु सर्वं यत्सत्यं तद्धि तस्मै निवेदितम्}
{काममोहाभिभूतात्मा नाश्रौषीत्तद्वचो मम} %||7-26-37||

\twolineshloka
{याच्यमानो मया देव स्नुषा तेऽहमिति प्रभो}
{तत्सर्वं पृष्ठतः कृत्वा बलात्तेनास्मि धर्षिता} %||7-26-38||

\twolineshloka
{एवं त्वमपराधं मे क्षन्तुमर्हसि मानद}
{न हि तुल्यं बलं सौम्य स्त्रियाश्च पुरुषस्य च} %||7-26-39||

\twolineshloka
{एवं श्रुत्वा तु सङ्क्रुद्धस्तदा वैश्वरणात्मजः}
{धर्षणां तां परां श्रुत्वा ध्यानं सम्प्रविवेश ह} %||7-26-40||

\twolineshloka
{तस्य तत्कर्म विज्ञाय तदा वैश्रवणात्मजः}
{मुहूर्ताद्रोषताम्राक्षस्तोयं जग्राह पाणिना} %||7-26-41||

\twolineshloka
{गृहीत्वा सलिलं दिव्यमुपस्पृश्य यथाविधि}
{उत्ससर्ज तदा शापं राक्षसेन्द्राय दारुणम्} %||7-26-42||

\twolineshloka
{अकामा तेन यस्मात्त्वं बलाद्भद्रे प्रधर्षिता}
{तस्मात्स युवतीमन्यां नाकामामुपयास्यति} %||7-26-43||

\twolineshloka
{यदा त्वकामां कामार्तो धर्षयिष्यति योषितम्}
{मूर्धा तु सप्तधा तस्य शकलीभविता तदा} %||7-26-44||

\twolineshloka
{तस्मिन्नुदाहृते शापे ज्वलिताग्निसमप्रभे}
{देवदुन्दुभयो नेदुः पुष्पवृष्टिश्च खाच्च्युता} %||7-26-45||

\twolineshloka
{प्रजापतिमुखाश्चापि सर्वे देवाः प्रहर्षिताः}
{ज्ञात्वा लोकगतिं सर्वां तस्य मृत्युं च रक्षसः} %||7-26-46||

\twolineshloka
{श्रुत्वा तु स दशग्रीवस्तं शापं रोमहर्षणम्}
{नारीषु मैथुनं भावं नाकामास्वभ्यरोचयत्} %||7-27-47||


{॥इत्यार्षे श्रीमद्रामायणे वाल्मीकीये आदिकाव्ये उत्तरकाण्डे षड्विंशः सर्गः॥}

\sect{सप्तविंशः सर्गः}

\twolineshloka
{कैलासं लङ्घयित्वाथ दशग्रीवः सराक्षसः}
{आससाद महातेजा इन्द्रलोकं निशाचरः} %||7-27-1||

\twolineshloka
{तस्य राक्षससैन्यस्य समन्तादुपयास्यतः}
{देवलोकं ययौ शब्दो भिद्यमानार्णवोपमः} %||7-27-2||

\twolineshloka
{श्रुत्वा तु रावणं प्राप्तमिन्द्रः सञ्चलितासनः}
{अब्रवीत्तत्र तान्देवान्सर्वानेव समागतान्} %||7-27-3||

\twolineshloka
{आदित्यान्सवसून्रुद्रान्विश्वान्साध्यान्मरुद्गणान्}
{सज्जीभवत युद्धार्थं रावणस्य दुरात्मनः} %||7-27-4||

\twolineshloka
{एवमुक्तास्तु शक्रेण देवाः शक्रसमा युधि}
{संनह्यन्त महासत्त्वा युद्धश्रद्धासमन्विताः} %||7-27-5||

\twolineshloka
{स तु दीनः परित्रस्तो महेन्द्रो रावणं प्रति}
{विष्णोः समीपमागत्य वाक्यमेतदुवाच ह} %||7-27-6||

\twolineshloka
{विष्णो कथं करिष्यामो महावीर्यपराक्रम}
{असौ हि बलवान्रक्षो युद्धार्थमभिवर्तते} %||7-27-7||

\twolineshloka
{वरप्रदानाद्बलवान्न खल्वन्येन हेतुना}
{तच्च सत्यं हि कर्तव्यं वाक्यं देव प्रजापतेः} %||7-27-8||

\twolineshloka
{तद्यथा नमुचिर्वृत्रो बलिर्नरकशम्बरौ}
{त्वन्मतं समवष्टभ्य यथा दग्धास्तथा कुरु} %||7-27-9||

\twolineshloka
{न ह्यन्यो देव देवानामापत्सु सुमहाबल}
{गतिः परायणं वास्ति त्वामृते पुरुषोत्तम} %||7-27-10||

\twolineshloka
{त्वं हि नारायणः श्रीमान्पद्मनाभः सनातनः}
{त्वयाहं स्थापितश्चैव देवराज्ये सनातने} %||7-27-11||

\twolineshloka
{तदाख्याहि यथातत्त्वं देवदेव मम स्वयम्}
{असिचक्रसहायस्त्वं युध्यसे संयुगे रिपुम्} %||7-27-12||

\twolineshloka
{एवमुक्तः स शक्रेण देवो नारायणः प्रभुः}
{अब्रवीन्न परित्रासः कार्यस्ते श्रूयतां च मे} %||7-27-13||

\twolineshloka
{न तावदेष दुर्वृत्तः शक्यो दैवतदानवैः}
{हन्तुं युधि समासाद्य वरदानेन दुर्जयः} %||7-27-14||

\twolineshloka
{सर्वथा तु महत्कर्म करिष्यति बलोत्कटः}
{रक्षः पुत्रसहायोऽसौ दृष्टमेतन्निसर्गतः} %||7-27-15||

\twolineshloka
{ब्रवीषि यत्तु मां शक्र संयुगे योत्स्यसीति ह}
{नैवाहं प्रतियोत्स्ये तं रावणं राक्षसाधिपम्} %||7-27-16||

\twolineshloka
{अनिहत्य रिपुं विष्णुर्न हि प्रतिनिवर्तते}
{दुर्लभश्चैष कामोऽद्य वरमासाद्य राक्षसे} %||7-27-17||

\twolineshloka
{प्रतिजानामि देवेन्द्र त्वत्समीपं शतक्रतो}
{राक्षसस्याहमेवास्य भविता मृत्युकारणम्} %||7-27-18||

\twolineshloka
{अहमेनं वधिष्यामि रावणं ससुतं युधि}
{देवतास्तोषयिष्यामि ज्ञात्वा कालमुपस्थितम्} %||7-27-19||

\twolineshloka
{एतस्मिन्नन्तरे नादः शुश्रुवे रजनीक्षये}
{तस्य रावणसैन्यस्य प्रयुद्धस्य समन्ततः} %||7-27-20||

\twolineshloka
{अथ युद्धं समभवद्देवराक्षसयोस्तदा}
{घोरं तुमुलनिर्ह्रादं नानाप्रहरणायुधम्} %||7-27-21||

\twolineshloka
{एतस्मिन्नन्तरे शूरा राक्षसा घोरदर्शनाः}
{युद्धार्थमभ्यधावन्त सचिवा रावणाज्ञया} %||7-27-22||

\twolineshloka
{मारीचश्च प्रहस्तश्च महापार्श्वमहोदरौ}
{अकम्पनो निकुम्भश्च शुकः सारण एव च} %||7-27-23||

\twolineshloka
{संह्रादिर्धूमकेतुश्च महादंष्ट्रो महामुखः}
{जम्बुमाली महामाली विरूपाक्षश्च राक्षसः} %||7-27-24||

\twolineshloka
{एतैः सर्वैर्महावीर्यैर्वृतो राक्षसपुङ्गवः}
{रावणस्यार्यकः सैन्यं सुमाली प्रविवेश ह} %||7-27-25||

\twolineshloka
{स हि देवगणान्सर्वान्नानाप्रहरणैः शितैः}
{विध्वंसयति सङ्क्रुद्धः सह तैः क्षणदाचरैः} %||7-27-26||

\twolineshloka
{एतस्मिन्नन्तरे शूरो वसूनामष्टमो वसुः}
{सावित्र इति विख्यातः प्रविवेश महारणम्} %||7-27-27||

\twolineshloka
{ततो युद्धं समभवत्सुराणां राक्षसैः सह}
{क्रुद्धानां रक्षसां कीर्तिं समरेष्वनिवर्तिनाम्} %||7-27-28||

\twolineshloka
{ततस्ते राक्षसाः शूरा देवांस्तान्समरे स्थितान्}
{नानाप्रहरणैर्घोरैर्जघ्नुः शतसहस्रशः} %||7-27-29||

\twolineshloka
{सुरास्तु राक्षसान्घोरान्महावीर्यान्स्वतेजसा}
{समरे विविधैः शस्त्रैरनयन्यमसादनम्} %||7-27-30||

\twolineshloka
{एतस्मिन्नन्तरे शूरः सुमाली नाम राक्षसः}
{नानाप्रहरणैः क्रुद्धो रणमेवाभ्यवर्तत} %||7-27-31||

\twolineshloka
{देवानां तद्बलं सर्वं नानाप्रहरणैः शितैः}
{विध्वंसयति सङ्क्रुद्धो वायुर्जलधरानिव} %||7-27-32||

\twolineshloka
{ते महाबाणवर्षैश्च शूलैः प्रासैश्च दारुणैः}
{पीड्यमानाः सुराः सर्वे न व्यतिष्ठन्समाहिताः} %||7-27-33||

\twolineshloka
{ततो विद्राव्यमाणेषु त्रिदशेषु सुमालिना}
{वसूनामष्टमो देवः सावित्रो व्यवतिष्ठत} %||7-27-34||

\twolineshloka
{संवृतः स्वैरनीकैस्तु प्रहरन्तं निशाचरम्}
{विक्रमेण महातेजा वारयामास संयुगे} %||7-27-35||

\twolineshloka
{सुमत्तयोस्तयोरासीद्युद्धं लोके सुदारुणम्}
{सुमालिनो वसोश्चैव समरेष्वनिवर्तिनोः} %||7-27-36||

\twolineshloka
{ततस्तस्य महाबाणैर्वसुना सुमहात्मना}
{महान्स पन्नगरथः क्षणेन विनिपातितः} %||7-27-37||

\twolineshloka
{हत्वा तु संयुगे तस्य रथं बाणशतैः शितैः}
{गदां तस्य वधार्थाय वसुर्जग्राह पाणिना} %||7-27-38||

\twolineshloka
{तां प्रदीप्तां प्रगृह्याशु कालदण्डनिभां शुभाम्}
{तस्य मूर्धनि सावित्रः सुमालेर्विनिपातयत्} %||7-27-39||

\twolineshloka
{तस्य मूर्धनि सोल्काभा पतन्ती च तदा बभौ}
{सहस्राक्षसमुत्सृष्टा गिराविव महाशनिः} %||7-27-40||

\twolineshloka
{तस्य नैवास्थि कायो वा न मांसं ददृशे तदा}
{गदया भस्मसाद्भूतो रणे तस्मिन्निपातितः} %||7-27-41||

\twolineshloka
{तं दृष्ट्वा निहतं सङ्ख्ये राक्षसास्ते समन्ततः}
{दुद्रुवुः सहिताः सर्वे क्रोशमाना महास्वनम्} %||7-28-42||


{॥इत्यार्षे श्रीमद्रामायणे वाल्मीकीये आदिकाव्ये उत्तरकाण्डे सप्तविंशः सर्गः॥}

\sect{अष्टाविंशः सर्गः}

\twolineshloka
{सुमालिनं हतं दृष्ट्वा वसुना भस्मसात्कृतम्}
{विद्रुतं चापि स्वं सैन्यं लक्षयित्वार्दितं शरैः} %||7-28-1||

\twolineshloka
{ततः स बलवान्क्रुद्धो रावणस्य सुतो युधि}
{निवर्त्य राक्षसान्सर्वान्मेघनादो व्यतिष्ठत} %||7-28-2||

\twolineshloka
{स रथेनाग्निवर्णेन कामगेन महारथः}
{अभिदुद्राव सेनां तां वनान्यग्निरिव ज्वलन्} %||7-28-3||

\twolineshloka
{ततः प्रविशतस्तस्य विविधायुधधारिणः}
{विदुद्रुवुर्दिशः सर्वा देवास्तस्य च दर्शनात्} %||7-28-4||

\twolineshloka
{न तत्रावस्थितः कश्चिद्रणे तस्य युयुत्सतः}
{सर्वानाविध्य वित्रस्तान्दृष्ट्वा शक्रोऽभ्यभाषत} %||7-28-5||

\twolineshloka
{न भेतव्यं न गन्तव्यं निवर्तध्वं रणं प्रति}
{एष गच्छति मे पुत्रो युद्धार्थमपराजितः} %||7-28-6||

\twolineshloka
{ततः शक्रसुतो देवो जयन्त इति विश्रुतः}
{रथेनाद्भुतकल्पेन सङ्ग्राममभिवर्तत} %||7-28-7||

\twolineshloka
{ततस्ते त्रिदशाः सर्वे परिवार्य शचीसुतम्}
{रावणस्य सुतं युद्धे समासाद्य व्यवस्थिताः} %||7-28-8||

\twolineshloka
{तेषां युद्धं महदभूत्सदृशं देवरक्षसाम्}
{कृते महेन्द्रपुत्रस्य राक्षसेन्द्रसुतस्य च} %||7-28-9||

\twolineshloka
{ततो मातलिपुत्रे तु गोमुखे राक्षसात्मजः}
{सारथौ पातयामास शरान्काञ्चनभूषणान्} %||7-28-10||

\twolineshloka
{शचीसुतस्त्वपि तथा जयन्तस्तस्य सारथिम्}
{तं चैव रावणिं क्रुद्धः प्रत्यविध्यद्रणाजिरे} %||7-28-11||

\twolineshloka
{ततः क्रुद्धो महातेजा रक्षो विस्फारितेक्षणः}
{रावणिः शक्रपुत्रं तं शरवर्षैरवाकिरत्} %||7-28-12||

\threelineshloka
{ततः प्रगृह्य शस्त्राणि सारवन्ति महान्ति च}
{शतघ्नीस्तोमरान्प्रासान्गदाखड्गपरश्वधान्}
{सुमहान्त्यद्रिशृङ्गाणि पातयामास रावणिः} %||7-28-13||

\twolineshloka
{ततः प्रव्यथिता लोकाः संजज्ञे च तमो महत्}
{तस्य रावणपुत्रस्य तदा शत्रूनभिघ्नतः} %||7-28-14||

\twolineshloka
{ततस्तद्दैवतबलं समन्तात्तं शचीसुतम्}
{बहुप्रकारमस्वस्थं तत्र तत्र स्म धावति} %||7-28-15||

\twolineshloka
{नाभ्यजानंस्तदान्योन्यं शत्रून्वा दैवतानि वा}
{तत्र तत्र विपर्यस्तं समन्तात्परिधावितम्} %||7-28-16||

\twolineshloka
{एतस्मिन्नन्तरे शूरः पुलोमा नाम वीर्यवान्}
{दैतेयस्तेन सङ्गृह्य शचीपुत्रोऽपवाहितः} %||7-28-17||

\twolineshloka
{गृहीत्वा तं तु नप्तारं प्रविष्टः स महोदधिम्}
{मातामहोऽर्यकस्तस्य पौलोमी येन सा शची} %||7-28-18||

\twolineshloka
{प्रणाशं दृश्य तु सुरा जयन्तस्यातिदारुणम्}
{व्यथिताश्चाप्रहृष्टाश्च समन्ताद्विप्रदुद्रुवुः} %||7-28-19||

\twolineshloka
{रावणिस्त्वथ संहृष्टो बलैः परिवृतः स्वकैः}
{अभ्यधावत देवांस्तान्मुमोच च महास्वनम्} %||7-28-20||

\twolineshloka
{दृष्ट्वा प्रणाशं पुत्रस्य रावणेश्चापि विक्रमम्}
{मातलिं प्राह देवेन्द्रो रथः समुपनीयताम्} %||7-28-21||

\twolineshloka
{स तु दिव्यो महाभीमः सज्ज एव महारथः}
{उपस्थितो मातलिना वाह्यमानो मनोजवः} %||7-28-22||

\twolineshloka
{ततो मेघा रथे तस्मिंस्तडिद्वन्तो महास्वनाः}
{अग्रतो वायुचपला गच्छन्तो व्यनदंस्तदा} %||7-28-23||

\twolineshloka
{नानावाद्यानि वाद्यन्त स्तुतयश्च समाहिताः}
{ननृतुश्चाप्सरःसङ्घाः प्रयाते वासवे रणम्} %||7-28-24||

\twolineshloka
{रुद्रैर्वसुभिरादित्यैः साध्यैश्च समरुद्गणैः}
{वृतो नानाप्रहरणैर्निर्ययौ त्रिदशाधिपः} %||7-28-25||

\twolineshloka
{निर्गच्छतस्तु शक्रस्य परुषं पवनो ववौ}
{भास्करो निष्प्रभश्चासीन्महोल्काश्च प्रपेदिरे} %||7-28-26||

\twolineshloka
{एतस्मिन्नन्तरे शूरो दशग्रीवः प्रतापवान्}
{आरुरोह रथं दिव्यं निर्मितं विश्वकर्मणा} %||7-28-27||

\twolineshloka
{पन्नगैः सुमहाकायैर्वेष्टितं लोमहर्षणैः}
{येषां निश्वासवातेन प्रदीप्तमिव संयुगम्} %||7-28-28||

\twolineshloka
{दैत्यैर्निशाचरैः शूरै रथः सम्परिवारितः}
{समराभिमुखो दिव्यो महेन्द्रमभिवर्तत} %||7-28-29||

\twolineshloka
{पुत्रं तं वारयित्वासौ स्वयमेव व्यवस्थितः}
{सोऽपि युद्धाद्विनिष्क्रम्य रावणिः समुपाविशत्} %||7-28-30||

\twolineshloka
{ततो युद्धं प्रवृत्तं तु सुराणां राक्षसैः सह}
{शस्त्राभिवर्षणं घोरं मेघानामिव संयुगे} %||7-28-31||

\twolineshloka
{कुम्भकर्णस्तु दुष्टात्मा नानाप्रहरणोद्यतः}
{नाज्ञायत तदा युद्धे सह केनाप्ययुध्यत} %||7-28-32||

\twolineshloka
{दन्तैर्भुजाभ्यां पद्भ्यां च शक्तितोमरसायकैः}
{येन केनैव संरब्धस्ताडयामास वै सुरान्} %||7-28-33||

\twolineshloka
{ततो रुद्रैर्महाभागैः सहादित्यैर्निशाचरः}
{प्रयुद्धस्तैश्च सङ्ग्रामे कृत्तः शस्त्रैर्निरन्तरम्} %||7-28-34||

\twolineshloka
{ततस्तद्राक्षसं सैन्यं त्रिदशैः समरुद्गणैः}
{रणे विद्रावितं सर्वं नानाप्रहरणैः शितैः} %||7-28-35||

\twolineshloka
{केचिद्विनिहताः शस्त्रैर्वेष्टन्ति स्म महीतले}
{वाहनेष्ववसक्ताश्च स्थिता एवापरे रणे} %||7-28-36||

\twolineshloka
{रथान्नागान्खरानुष्ट्रान्पन्नगांस्तुरगांस्तथा}
{शिंशुमारान्वराहांश्च पिशाचवदनांस्तथा} %||7-28-37||

\twolineshloka
{तान्समालिङ्ग्य बाहुभ्यां विष्टब्धाः केचिदुच्छ्रिताः}
{देवैस्तु शस्त्रसंविद्धा मम्रिरे च निशाचराः} %||7-28-38||

\twolineshloka
{चित्रकर्म इवाभाति स तेषां रणसम्प्लवः}
{निहतानां प्रमत्तानां राक्षसानां महीतले} %||7-28-39||

\twolineshloka
{शोणितोदक निष्यन्दाकङ्कगृध्रसमाकुला}
{प्रवृत्ता संयुगमुखे शस्त्रग्राहवती नदी} %||7-28-40||

\twolineshloka
{एतस्मिन्नन्तरे क्रुद्धो दशग्रीवः प्रतापवान्}
{निरीक्ष्य तद्बलं सर्वं दैवतैर्विनिपातितम्} %||7-28-41||

\twolineshloka
{स तं प्रतिविगाह्याशु प्रवृद्धं सैन्यसागरम्}
{त्रिदशान्समरे निघ्नञ्शक्रमेवाभ्यवर्तत} %||7-28-42||

\twolineshloka
{ततः शक्रो महच्चापं विस्फार्य सुमहास्वनम्}
{यस्य विस्फारघोषेण स्वनन्ति स्म दिशो दश} %||7-28-43||

\twolineshloka
{तद्विकृष्य महच्चापमिन्द्रो रावणमूर्धनि}
{निपातयामास शरान्पावकादित्यवर्चसः} %||7-28-44||

\twolineshloka
{तथैव च महाबाहुर्दशग्रीवो व्यवस्थितः}
{शक्रं कार्मुकविभ्रष्टैः शरवर्षैरवाकिरत्} %||7-28-45||

\twolineshloka
{प्रयुध्यतोरथ तयोर्बाणवर्षैः समन्ततः}
{नाज्ञायत तदा किञ्चित्सर्वं हि तमसा वृतम्} %||7-29-46||


{॥इत्यार्षे श्रीमद्रामायणे वाल्मीकीये आदिकाव्ये उत्तरकाण्डे अष्टाविंशः सर्गः॥}

\sect{एकोनत्रिंशः सर्गः}

\twolineshloka
{ततस्तमसि संजाते राक्षसा दैवतैः सह}
{अयुध्यन्त बलोन्मत्ताः सूदयन्तः परस्परम्} %||7-29-1||

\twolineshloka
{ततस्तु देवसैन्येन राक्षसानां महद्बलम्}
{दशांशं स्थापितं युद्धे शेषं नीतं यमक्षयम्} %||7-29-2||

\twolineshloka
{तस्मिंस्तु तमसा नद्धे सर्वे ते देवराक्षसाः}
{अन्योन्यं नाभ्यजानन्त युध्यमानाः परस्परम्} %||7-29-3||

\twolineshloka
{इन्द्रश्च रावणश्चैव रावणिश्च महाबलः}
{तस्मिंस्तमोजालवृते मोहमीयुर्न ते त्रयः} %||7-29-4||

\twolineshloka
{स तु दृष्ट्वा बलं सर्वं निहतं रावणो रणे}
{क्रोधमभ्यागमत्तीव्रं महानादं च मुक्तवान्} %||7-29-5||

\twolineshloka
{क्रोधात्सूतं च दुर्धर्षः स्यन्दनस्थमुवाच ह}
{परसैन्यस्य मध्येन यावदन्तं नयस्व माम्} %||7-29-6||

\twolineshloka
{अद्यैतांस्त्रिदशान्सर्वान्विक्रमैः समरे स्वयम्}
{नानाशस्त्रैर्महासारैर्नाशयामि नभस्तलात्} %||7-29-7||

\twolineshloka
{अहमिन्द्रं वधिष्यामि वरुणं धनदं यमम्}
{त्रिदशान्विनिहत्याशु स्वयं स्थास्याम्यथोपरि} %||7-29-8||

\twolineshloka
{विषादो न च कर्तव्यः शीघ्रं वाहय मे रथम्}
{द्विः खलु त्वां ब्रवीम्यद्य यावदन्तं नयस्व माम्} %||7-29-9||

\twolineshloka
{अयं स नन्दनोद्देशो यत्र वर्तामहे वयम्}
{नय मामद्य तत्र त्वमुदयो यत्र पर्वतः} %||7-29-10||

\twolineshloka
{तस्य तद्वचनं श्रुत्वा तुरगान्स मनोजवान्}
{आदिदेशाथ शत्रूणां मध्येनैव च सारथिः} %||7-29-11||

\twolineshloka
{तस्य तं निश्चयं ज्ञात्वा शक्रो देवेश्वरस्तदा}
{रथस्थः समरस्थांस्तान्देवान्वाक्यमथाब्रवीत्} %||7-29-12||

\twolineshloka
{सुराः शृणुत मद्वाक्यं यत्तावन्मम रोचते}
{जीवन्नेव दशग्रीवः साधु रक्षो निगृह्यताम्} %||7-29-13||

\twolineshloka
{एष ह्यतिबलः सैन्ये रथेन पवनौजसा}
{गमिष्यति प्रवृद्धोर्मिः समुद्र इव पर्वणि} %||7-29-14||

\twolineshloka
{न ह्येष हन्तुं शक्योऽद्य वरदानात्सुनिर्भयः}
{तद्ग्रहीष्यामहे रक्षो यत्ता भवत संयुगे} %||7-29-15||

\twolineshloka
{यथा बलिं निगृह्यैतत्त्रैलोक्यं भुज्यते मया}
{एवमेतस्य पापस्य निग्रहो मम रोचते} %||7-29-16||

\twolineshloka
{ततोऽन्यं देशमास्थाय शक्रः सन्त्यज्य रावणम्}
{अयुध्यत महातेजा राक्षसान्नाशयन्रणे} %||7-29-17||

\twolineshloka
{उत्तरेण दशग्रीवः प्रविवेशानिवर्तितः}
{दक्षिणेन तु पार्श्वेन प्रविवेश शतक्रतुः} %||7-29-18||

\twolineshloka
{ततः स योजनशतं प्रविष्टो राक्षसाधिपः}
{देवतानां बलं कृत्स्नं शरवर्षैरवाकिरत्} %||7-29-19||

\twolineshloka
{ततः शक्रो निरीक्ष्याथ प्रविष्टं तं बलं स्वकम्}
{न्यवर्तयदसम्भ्रान्तः समावृत्य दशाननम्} %||7-29-20||

\twolineshloka
{एतस्मिन्नन्तरे नादो मुक्तो दानवराक्षसैः}
{हा हताः स्मेति तं दृष्ट्वा ग्रस्तं शक्रेण रावणम्} %||7-29-21||

\twolineshloka
{ततो रथं समारुह्य रावणिः क्रोधमूर्छितः}
{तत्सैन्यमतिसङ्क्रुद्धः प्रविवेश सुदारुणम्} %||7-29-22||

\twolineshloka
{स तां प्रविश्य मायां तु दत्तां गोपतिना पुरा}
{अदृश्यः सर्वभूतानां तत्सैन्यं समवाकिरत्} %||7-29-23||

\twolineshloka
{ततः स देवान्सन्त्यज्य शक्रमेवाभ्ययाद्द्रुतम्}
{महेन्द्रश्च महातेजा न ददर्श सुतं रिपोः} %||7-29-24||

\twolineshloka
{स मातलिं हयांश्चैव ताडयित्वा शरोत्तमैः}
{महेन्द्रं बाणवर्षेण शीघ्रहस्तो ह्यवाकिरत्} %||7-29-25||

\twolineshloka
{ततः शक्रो रथं त्यक्त्व विसृज्य च स मातलिम्}
{ऐरावतं समारुह्य मृगयामास रावणिम्} %||7-29-26||

\twolineshloka
{स तु माया बलाद्रक्षः सङ्ग्रामे नाभ्यदृश्यत}
{किरमाणः शरौघेन महेन्द्रममितौजसम्} %||7-29-27||

\twolineshloka
{स तं यदा परिश्रान्तमिन्द्रं मेनेऽथ रावणिः}
{तदैनं मायया बद्ध्वा स्वसैन्यमभितोऽनयत्} %||7-29-28||

\threelineshloka
{तं दृष्ट्वाथ बलात्तस्मिन्माययापहृतं रणे}
{महेन्द्रममराः सर्वे किं न्वेतदिति चुक्रुशुः}
{न हि दृश्यति विद्यावान्मायया येन नीयते} %||7-29-29||

\twolineshloka
{एतस्मिन्नन्तरे चापि सर्वे सुरगणास्तदा}
{अभ्यद्रवन्सुसङ्क्रुद्धा रावणं शस्त्रवृष्टिभिः} %||7-29-30||

\twolineshloka
{रावणस्तु समासाद्य वस्वादित्यमरुद्गणान्}
{न शशाक रणे स्थातुं न योद्धुं शस्त्रपीडितः} %||7-29-31||

\twolineshloka
{तं तु दृष्ट्वा परिश्रान्तं प्रहारैर्जर्जरच्छविम्}
{रावणिः पितरं युद्धेऽदर्शनस्थोऽब्रवीदिदम्} %||7-29-32||

\twolineshloka
{आगच्छ तात गच्छावो निवृत्तं रणकर्म तत्}
{जितं ते विदितं भोऽस्तु स्वस्थो भव गतज्वरः} %||7-29-33||

\twolineshloka
{अयं हि सुरसैन्यस्य त्रैलोक्यस्य च यः प्रभुः}
{स गृहीतो मया शक्रो भग्नमानाः सुराः कृताः} %||7-29-34||

\twolineshloka
{यथेष्टं भुङ्क्ष्व त्रैलोक्यं निगृह्य रिपुमोजसा}
{वृथा ते किं श्रमं कृत्वा युद्धं हि तव निष्फलम्} %||7-29-35||

\twolineshloka
{स दैवतबलात्तस्मान्निवृत्तो रणकर्मणः}
{तच्छ्रुत्वा रावणेर्वाक्यं स्वस्थचेता दशाननः} %||7-29-36||

\fourlineindentedshloka
{अथ रणविगतज्वरः प्रभुर्-}
{ विजयमवाप्य निशाचराधिपः}
{भवनमभि ततो जगाम हृष्टः}
{ स्वसुतमवाप्य च वाक्यमब्रवीत्} %||7-29-37||

\fourlineindentedshloka
{अतिबलसदृशैः पराक्रमैस्तैर्-}
{ मम कुलमानविवर्धनं कृतम्}
{यदमरसमविक्रम त्वया}
{ त्रिदशपतिस्त्रिदशाश्च निर्जिताः} %||7-29-38||

\fourlineindentedshloka
{त्वरितमुपनयस्व वासवम्}
{ नगरमितो व्रज सैन्यसंवृतः}
{अहमपि तव गच्छतो द्रुतम्}
{ सह सचिवैरनुयामि पृष्ठतः} %||7-29-39||

\fourlineindentedshloka
{अथ स बलवृतः सवाहन}
{स्त्रिदशपतिं परिगृह्य रावणिः}
{स्वभवनमुपगम्य राक्षसो}
{ मुदितमना विससर्ज राक्षसान्} %||7-30-40||


{॥इत्यार्षे श्रीमद्रामायणे वाल्मीकीये आदिकाव्ये उत्तरकाण्डे एकोनत्रिंशः सर्गः॥}

\sect{त्रिंशः सर्गः}

\twolineshloka
{जिते महेन्द्रेऽतिबले रावणस्य सुतेन वै}
{प्रजापतिं पुरस्कृत्य गता लङ्कां सुरास्तदा} %||7-30-1||

\twolineshloka
{तं रावणं समासाद्य पुत्रभ्रातृभिरावृतम्}
{अब्रवीद्गगने तिष्ठन्सान्त्वपूर्वं प्रजापतिः} %||7-30-2||

\twolineshloka
{वत्स रावण तुष्टोऽस्मि तव पुत्रस्य संयुगे}
{अहोऽस्य विक्रमौदार्यं तव तुल्योऽधिकोऽपि वा} %||7-30-3||

\twolineshloka
{जितं हि भवता सर्वं त्रैलोक्यं स्वेन तेजसा}
{कृता प्रतिज्ञा सफला प्रीतोऽस्मि स्वसुतेन वै} %||7-30-4||

\twolineshloka
{अयं च पुत्रोऽतिबलस्तव रावणरावणिः}
{इन्द्रजित्त्विति विख्यातो जगत्येष भविष्यति} %||7-30-5||

\twolineshloka
{बलवाञ्शत्रुनिर्जेता भविष्यत्येष राक्षसः}
{यमाश्रित्य त्वया राजन्स्थापितास्त्रिदशा वशे} %||7-30-6||

\twolineshloka
{तन्मुच्यतां महाबाहो महेन्द्रः पाकशासनः}
{किं चास्य मोक्षणार्थाय प्रयच्छन्ति दिवौकसः} %||7-30-7||

\twolineshloka
{अथाब्रवीन्महातेजा इन्द्रजित्समितिंजयः}
{अमरत्वमहं देव वृणोमीहास्य मोक्षणे} %||7-30-8||

\twolineshloka
{अब्रवीत्तु तदा देवो रावणिं कमलोद्भवः}
{नास्ति सर्वामरत्वं हि केषाञ्चित्प्राणिनां भुवि} %||7-30-9||

\twolineshloka
{अथाब्रवीत्स तत्रस्थमिन्द्रजित्पद्मसम्भवम्}
{श्रूयतां या भवेत्सिद्धिः शतक्रतुविमोक्षणे} %||7-30-10||

\twolineshloka
{ममेष्टं नित्यशो देव हव्यैः सम्पूज्य पावकम्}
{सङ्ग्राममवतर्तुं वै शत्रुनिर्जयकाङ्क्षिणः} %||7-30-11||

\twolineshloka
{तस्मिंश्चेदसमाप्ते तु जप्यहोमे विभावसोः}
{युध्येयं देव सङ्ग्रामे तदा मे स्याद्विनाशनम्} %||7-30-12||

\twolineshloka
{सर्वो हि तपसा चैव वृणोत्यमरतां पुमान्}
{विक्रमेण मया त्वेतदमरत्वं प्रवर्तितम्} %||7-30-13||

\twolineshloka
{एवमस्त्विति तं प्राह वाक्यं देवः प्रजापतिः}
{मुक्तश्चेन्द्रजिता शक्रो गताश्च त्रिदिवं सुराः} %||7-30-14||

\twolineshloka
{एतस्मिन्नन्तरे शक्रो दीनो भ्रष्टाम्बरस्रजः}
{राम चिन्तापरीतात्मा ध्यानतत्परतां गतः} %||7-30-15||

\twolineshloka
{तं तु दृष्ट्वा तथाभूतं प्राह देवः प्रजापतिः}
{शक्रक्रतो किमुत्कण्ठां करोषि स्मर दुष्कृतम्} %||7-30-16||

\twolineshloka
{अमरेन्द्र मया बह्व्यः प्रजाः सृष्टाः पुरा प्रभो}
{एकवर्णाः समाभाषा एकरूपाश्च सर्वशः} %||7-30-17||

\twolineshloka
{तासां नास्ति विशेषो हि दर्शने लक्षणेऽपि वा}
{ततोऽहमेकाग्रमनास्ताः प्रजाः पर्यचिन्तयम्} %||7-30-18||

\twolineshloka
{सोऽहं तासां विशेषार्थं स्त्रियमेकां विनिर्ममे}
{यद्यत्प्रजानां प्रत्यङ्गं विशिष्टं तत्तदुद्धृतम्} %||7-30-19||

\twolineshloka
{ततो मया रूपगुणैरहल्या स्त्री विनिर्मिता}
{अहल्येत्येव च मया तस्या नाम प्रवर्तितम्} %||7-30-20||

\twolineshloka
{निर्मितायां तु देवेन्द्र तस्यां नार्यां सुरर्षभ}
{भविष्यतीति कस्यैषा मम चिन्ता ततोऽभवत्} %||7-30-21||

\twolineshloka
{त्वं तु शक्र तदा नारीं जानीषे मनसा प्रभो}
{स्थानाधिकतया पत्नी ममैषेति पुरन्दर} %||7-30-22||

\twolineshloka
{सा मया न्यासभूता तु गौतमस्य महात्मनः}
{न्यस्ता बहूनि वर्षाणि तेन निर्यातिता च सा} %||7-30-23||

\twolineshloka
{ततस्तस्य परिज्ञाय मया स्थैर्यं महामुनेः}
{ज्ञात्वा तपसि सिद्धिं च पत्न्यर्थं स्पर्शिता तदा} %||7-30-24||

\twolineshloka
{स तया सह धर्मात्मा रमते स्म महामुनिः}
{आसन्निराशा देवास्तु गौतमे दत्तया तया} %||7-30-25||

\twolineshloka
{त्वं क्रुद्धस्त्विह कामात्मा गत्वा तस्याश्रमं मुनेः}
{दृष्टवांश्च तदा तां स्त्रीं दीप्तामग्निशिखामिव} %||7-30-26||

\twolineshloka
{सा त्वया धर्षिता शक्र कामार्तेन समन्युना}
{दृष्टस्त्वं च तदा तेन आश्रमे परमर्षिणा} %||7-30-27||

\twolineshloka
{ततः क्रुद्धेन तेनासि शप्तः परमतेजसा}
{गतोऽसि येन देवेन्द्र दशाभागविपर्ययम्} %||7-30-28||

\twolineshloka
{यस्मान्मे धर्षिता पत्नी त्वया वासव निर्भयम्}
{तस्मात्त्वं समरे राजञ्शत्रुहस्तं गमिष्यसि} %||7-30-29||

\twolineshloka
{अयं तु भावो दुर्बुद्धे यस्त्वयेह प्रवर्तितः}
{मानुषेष्वपि सर्वेषु भविष्यति न संशयः} %||7-30-30||

\twolineshloka
{तत्राधर्मः सुबलवान्समुत्थास्यति यो महान्}
{तत्रार्धं तस्य यः कर्ता त्वय्यर्धं निपतिष्यति} %||7-30-31||

\twolineshloka
{न च ते स्थावरं स्थानं भविष्यति पुरन्दर}
{एतेनाधर्मयोगेन यस्त्वयेह प्रवर्तितः} %||7-30-32||

\twolineshloka
{यश्च यश्च सुरेन्द्रः स्याद्ध्रुवः स न भविष्यति}
{एष शापो मया मुक्त इत्यसौ त्वां तदाब्रवीत्} %||7-30-33||

\twolineshloka
{तां तु भार्यां विनिर्भर्त्स्य सोऽब्रवीत्सुमहातपाः}
{दुर्विनीते विनिध्वंस ममाश्रमसमीपतः} %||7-30-34||

\twolineshloka
{रूपयौवनसम्पन्ना यस्मात्त्वमनवस्थिता}
{तस्माद्रूपवती लोके न त्वमेका भविष्यसि} %||7-30-35||

\twolineshloka
{रूपं च तत्प्रजाः सर्वा गमिष्यन्ति सुदुर्लभम्}
{यत्तवेदं समाश्रित्य विभ्रमेऽयमुपस्थितः} %||7-30-36||

\twolineshloka
{तदा प्रभृति भूयिष्ठं प्रजा रूपसमन्विताः}
{शापोत्सर्गाद्धि तस्येदं मुनेः सर्वमुपागतम्} %||7-30-37||

\twolineshloka
{तत्स्मर त्वं महाबाहो दुष्कृतं यत्त्वया कृतम्}
{येन त्वं ग्रहणं शत्रोर्गतो नान्येन वासव} %||7-30-38||

\twolineshloka
{शीघ्रं यजस्व यज्ञं त्वं वैष्णवं सुसमाहितः}
{पावितस्तेन यज्ञेन यास्यसि त्रिदिवं ततः} %||7-30-39||

\twolineshloka
{पुत्रश्च तव देवेन्द्र न विनष्टो महारणे}
{नीतः संनिहितश्चैव अर्यकेण महोदधौ} %||7-30-40||

\twolineshloka
{एतच्छ्रुत्वा महेन्द्रस्तु यज्ञमिष्ट्वा च वैष्णवम्}
{पुनस्त्रिदिवमाक्रामदन्वशासच्च देवताः} %||7-30-41||

\twolineshloka
{एतदिन्द्रजितो राम बलं यत्कीर्तितं मया}
{निर्जितस्तेन देवेन्द्रः प्राणिनोऽन्ये च किं पुनः} %||7-31-42||


{॥इत्यार्षे श्रीमद्रामायणे वाल्मीकीये आदिकाव्ये उत्तरकाण्डे त्रिंशः सर्गः॥}

\sect{एकत्रिंशः सर्गः}

\twolineshloka
{ततो रामो महातेजा विस्मयात्पुनरेव हि}
{उवाच प्रणतो वाक्यमगस्त्यमृषिसत्तमम्} %||7-31-1||

\twolineshloka
{भगवन्किं तदा लोकाः शून्या आसन्द्विजोत्तम}
{धर्षणां यत्र न प्राप्तो रावणो राक्षसेश्वरः} %||7-31-2||

\twolineshloka
{उताहो हीनवीर्यास्ते बभुवुः पृथिवीक्षितः}
{बहिष्कृता वरास्त्रैश्च बहवो निर्जिता नृपाः} %||7-31-3||

\twolineshloka
{राघवस्य वचः श्रुत्वा अगस्त्यो भगवानृषिः}
{उवाच रामं प्रहसन्पितामह इवेश्वरम्} %||7-31-4||

\twolineshloka
{स एवं बाधमानस्तु पार्थिवान्पार्थिवर्षभ}
{चचार रावणो राम पृथिव्यां पृथिवीपते} %||7-31-5||

\twolineshloka
{ततो माहिष्मतीं नाम पुरीं स्वर्गपुरीप्रभाम्}
{सम्प्राप्तो यत्र साम्निध्यं परमं वसुरेतसः} %||7-31-6||

\twolineshloka
{तुल्य आसीन्नृपस्तस्य प्रतापाद्वसुरेतसः}
{अर्जुनो नाम यस्याग्निः शरकुण्डे शयः सदा} %||7-31-7||

\twolineshloka
{तमेव दिवसं सोऽथ हैहयाधिपतिर्बली}
{अर्जुनो नर्मदां रन्तुं गतः स्त्रीभिः सहेश्वरः} %||7-31-8||

\twolineshloka
{रावणो राक्षसेन्द्रस्तु तस्यामात्यानपृच्छत}
{क्वार्जुनो वो नृपः सोऽद्य शीघ्रमाख्यातुमर्हथ} %||7-31-9||

\twolineshloka
{रावणोऽहमनुप्राप्तो युद्धेप्सुर्नृवरेण तु}
{ममागमनमव्यग्रैर्युष्माभिः संनिवेद्यताम्} %||7-31-10||

\twolineshloka
{इत्येवं रावणेनोक्तास्तेऽमात्याः सुविपश्चितः}
{अब्रुवन्राक्षसपतिमसाम्निध्यं महीपतेः} %||7-31-11||

\twolineshloka
{श्रुत्वा विश्रवसः पुत्रः पौराणामर्जुनं गतम्}
{अपसृत्यागतो विन्ध्यं हिमवत्संनिभं गिरिम्} %||7-31-12||

\twolineshloka
{स तमभ्रमिवाविष्टमुद्भ्रान्तमिव मेदिनीम्}
{अपश्यद्रावणो विन्ध्यमालिखन्तमिवाम्बरम्} %||7-31-13||

\twolineshloka
{सहस्रशिखरोपेतं सिंहाध्युषितकन्दरम्}
{प्रपात पतितैः शीतैः साट्टहासमिवाम्बुभिः} %||7-31-14||

\twolineshloka
{देवदानवगन्धर्वैः साप्सरोगणकिंनरैः}
{साह स्त्रीभिः क्रीडमानैः स्वर्गभूतं महोच्छ्रयम्} %||7-31-15||

\twolineshloka
{नदीभिः स्यन्दमानाभिरगतिप्रतिमं जलम्}
{स्फुटीभिश्चलजिह्वाभिर्वमन्तमिव विष्ठितम्} %||7-31-16||

\twolineshloka
{उल्कावन्तं दरीवन्तं हिमवत्संनिभं गिरिम्}
{पश्यमानस्ततो विन्ध्यं रावणो नर्मदां ययौ} %||7-31-17||

\threelineshloka
{चलोपलजलां पुण्यां पश्चिमोदधिगामिनीम्}
{महिषैः सृमरैः सिंहैः शार्दूलर्क्षगजोत्तमैः}
{उष्णाभितप्तैस्तृषितैः सङ्क्षोभितजलाशयाम्} %||7-31-18||

\twolineshloka
{चक्रवाकैः सकारण्डैः सहंसजलकुक्कुटैः}
{सारसैश्च सदामत्तैः कोकूजद्भिः समावृताम्} %||7-31-19||

\twolineshloka
{फुल्लद्रुमकृतोत्तंसां चक्रवाकयुगस्तनीम्}
{विस्तीर्णपुलिनश्रोणीं हंसावलिसुमेखलाम्} %||7-31-20||

\twolineshloka
{पुष्परेण्वनुलिप्ताङ्गीं जलफेनामलांशुकाम्}
{जलावगाहसंस्पर्शां फुल्लोत्पलशुभेक्षणाम्} %||7-31-21||

\twolineshloka
{पुष्पकादवरुह्याशु नर्मदां सरितां वराम्}
{इष्टामिव वरां नारीमवगाह्य दशाननः} %||7-31-22||

\threelineshloka
{स तस्याः पुलिने रम्ये नानाकुसुमशोभिते}
{उपोपविष्टः सचिवैः सार्धं राक्षसपुङ्गवः}
{नर्मदा दर्शजं हर्षमाप्तवान्राक्षसेश्वरः} %||7-31-23||

\twolineshloka
{ततः सलीलं प्रहसान्रावणो राक्षसाधिपः}
{उवाच सचिवांस्तत्र मारीचशुकसारणान्} %||7-31-24||

\threelineshloka
{एष रश्मिसहस्रेण जगत्कृत्वेव काञ्चनम्}
{तीक्ष्णतापकरः सूर्यो नभसो मध्यमास्थितः}
{मामासीनं विदित्वेह चन्द्रायाति दिवाकरः} %||7-31-25||

\twolineshloka
{नर्मदा जलशीतश्च सुगन्धिः श्रमनाशनः}
{मद्भयादनिलो ह्येष वात्यसौ सुसमाहितः} %||7-31-26||

\twolineshloka
{इयं चापि सरिच्छ्रेष्ठा नर्मदा नर्म वर्धिनी}
{लीनमीनविहङ्गोर्मिः सभयेवाङ्गना स्थिता} %||7-31-27||

\twolineshloka
{तद्भवन्तः क्षताः शस्त्रैर्नृपैरिन्द्रसमैर्युधि}
{चन्दनस्य रसेनेव रुधिरेण समुक्षिताः} %||7-31-28||

\twolineshloka
{ते यूयमवगाहध्वं नर्मदां शर्मदां नृणाम्}
{महापद्ममुखा मत्ता गङ्गामिव महागजाः} %||7-31-29||

{अस्यां स्नात्वा महानद्यां पाप्मानं विप्रमोक्ष्यथ॥\devanumber{30}॥} %||7-31-30|| (Check)

\addtocounter{shlokacount}{1}

\twolineshloka
{अहमप्यत्र पुलिने शरदिन्दुसमप्रभे}
{पुष्पोपहरं शनकैः करिष्यामि उमापतेः} %||7-31-31||

\twolineshloka
{रावणेनैवमुक्तास्तु मारीचशुकसारणाः}
{समहोदरधूम्राक्षा नर्मदामवगाहिरे} %||7-31-32||

\twolineshloka
{राक्षसेन्द्रगजैस्तैस्तु क्षोभ्यते नर्मदा नदी}
{वामनाञ्जनपद्माद्यैर्गङ्गा इव महागजैः} %||7-31-33||

\twolineshloka
{ततस्ते राक्षसाः स्नात्वा नर्मदाया वराम्भसि}
{उत्तीर्य पुष्पाण्याजह्रुर्बल्यर्थं रावणस्य तु} %||7-31-34||

\twolineshloka
{नर्मदा पुलिने रम्ये शुभ्राभ्रसदृशप्रभे}
{राक्षसेन्द्रैर्मुहूर्तेन कृतः पुष्पमयो गिरिः} %||7-31-35||

\twolineshloka
{पुष्पेषूपहृतेष्वेव रावणो राक्षसेश्वरः}
{अवतीर्णो नदीं स्नातुं गङ्गामिव महागजः} %||7-31-36||

\twolineshloka
{तत्र स्नात्वा च विधिवज्जप्त्वा जप्यमनुत्तमम्}
{नर्मदा सलिलात्तस्मादुत्ततार स रावणः} %||7-31-37||

\threelineshloka
{रावणं प्राञ्जलिं यान्तमन्वयुः सप्तराक्षसाः}
{यत्र यत्र स याति स्म रावणो राक्षसाधिपः}
{जाम्बूनदमयं लिङ्गं तत्र तत्र स्म नीयते} %||7-31-38||

\twolineshloka
{वालुकवेदिमध्ये तु तल्लिङ्गं स्थाप्य रावणः}
{अर्चयामास गन्धैश्च पुष्पैश्चामृतगन्धिभिः} %||7-31-39||

\fourlineindentedshloka
{ततः सतामार्तिहरं हरं परम्}
{ वरप्रदं चन्द्रमयूखभूषणम्}
{समर्चयित्वा स निशाचरो जगौ}
{ प्रसार्य हस्तान्प्रणनर्त चायतान्} %||7-32-40||


{॥इत्यार्षे श्रीमद्रामायणे वाल्मीकीये आदिकाव्ये उत्तरकाण्डे एकत्रिंशः सर्गः॥}

\sect{द्वात्रिंशः सर्गः}

\twolineshloka
{नर्मदा पुलिने यत्र राक्षसेन्द्रः स रावणः}
{पुष्पोपहारं कुरुते तस्माद्देशाददूरतः} %||7-32-1||

\twolineshloka
{अर्जुनो जयतां श्रेष्ठो माहिष्मत्याः पतिः प्रभुः}
{क्रीडिते सह नारीभिर्नर्मदातोयमाश्रितः} %||7-32-2||

\twolineshloka
{तासां मध्यगतो राज रराज स ततोऽर्जुनः}
{करेणूनां सहस्रस्य मध्यस्थ इव कुञ्जरः} %||7-32-3||

\twolineshloka
{जिज्ञासुः स तु बाहूनां सहस्रस्योत्तमं बलम्}
{रुरोध नर्मदा वेगं बाहुभिः स तदार्जुनः} %||7-32-4||

\twolineshloka
{कार्तवीर्यभुजासेतुं तज्जलं प्राप्य निर्मलम्}
{कूलापहारं कुर्वाणं प्रतिस्रोतः प्रधावति} %||7-32-5||

\twolineshloka
{समीननक्रमकरः सपुष्पकुशसंस्तरः}
{स नर्मदाम्भसो वेगः प्रावृट्काल इवाबभौ} %||7-32-6||

\twolineshloka
{स वेगः कार्तवीर्येण सम्प्रेषिट इवाम्भसः}
{पुष्पोपहारं तत्सर्वं रावणस्य जहार ह} %||7-32-7||

\twolineshloka
{रावणोऽर्धसमाप्तं तु उत्सृज्य नियमं तदा}
{नर्मदां पश्यते कान्तां प्रतिकूलां यथा प्रियाम्} %||7-32-8||

\twolineshloka
{पश्चिमेन तु तं दृष्ट्वा सागरोद्गारसंनिभम्}
{वर्धन्तमम्भसो वेगं पूर्वामाशां प्रविश्य तु} %||7-32-9||

\twolineshloka
{ततोऽनुद्भ्रान्तशकुनां स्वाभाव्ये परमे स्थिताम्}
{निर्विकाराङ्गनाभासां पश्यते रावणो नदीम्} %||7-32-10||

\twolineshloka
{सव्येतरकराङ्गुल्या सशब्दं च दशाननः}
{वेगप्रभवमन्वेष्टुं सोऽदिशच्छुकसारणौ} %||7-32-11||

\twolineshloka
{तौ तु रावणसन्दिष्टौ भ्रातरौ शुकसारणौ}
{व्योमान्तरचरौ वीरौ प्रस्थितौ पश्चिमोन्मुखौ} %||7-32-12||

\twolineshloka
{अर्धयोजनमात्रं तु गत्वा तौ तु निशाचरौ}
{पश्येतां पुरुषं तोये क्रीडन्तं सहयोषितम्} %||7-32-13||

\twolineshloka
{बृहत्सालप्रतीकशं तोयव्याकुलमूर्धजम्}
{मदरक्तान्तनयनं मदनाकारवर्चसम्} %||7-32-14||

\twolineshloka
{नदीं बाहुसहस्रेण रुन्धन्तमरिमर्दनम्}
{गिरिं पादसहस्रेण रुन्धन्तमिव मेदिनीम्} %||7-32-15||

\twolineshloka
{बालानां वरनारीणां सहस्रेणाभिसंवृतम्}
{समदानां करेणूनां सहस्रेणेव कुञ्जरम्} %||7-32-16||

\twolineshloka
{तमद्भुततमं दृष्ट्वा राक्षसौ शुकसारणौ}
{संनिवृत्तावुपागम्य रावणं तमथोचतुः} %||7-32-17||

\twolineshloka
{बृहत्सालप्रतीकाशः कोऽप्यसौ राक्षसेश्वर}
{नर्मदां रोधवद्रुद्ध्वा क्रीडापयति योषितः} %||7-32-18||

\twolineshloka
{तेन बाहुसहस्रेण संनिरुद्धजला नदी}
{सागरोद्गारसङ्काशानुद्गारान्सृजते मुहुः} %||7-32-19||

\twolineshloka
{इत्येवं भाषमाणौ तौ निशम्य शुकसारणौ}
{रावणोऽर्जुन इत्युक्त्वा उत्तस्थौ युद्धलालसः} %||7-32-20||

\twolineshloka
{अर्जुनाभिमुखे तस्मिन्प्रस्थिते राक्षसेश्वरे}
{सकृदेव कृतो रावः सरक्तः प्रेषितो घनैः} %||7-32-21||

\twolineshloka
{महोदरमहापार्श्वधूम्राक्षशुकसारणैः}
{संवृतो राक्षसेन्द्रस्तु तत्रागाद्यत्र सोऽर्जुनः} %||7-32-22||

\twolineshloka
{नातिदीर्घेण कालेन स ततो राक्षसो बली}
{तं नर्मदा ह्रदं भीममाजगामाञ्जनप्रभः} %||7-32-23||

\twolineshloka
{स तत्र स्त्रीपरिवृतं वाशिताभिरिव द्विपम्}
{नरेन्द्रं पश्यते राजा राक्षसानां तदार्जुनम्} %||7-32-24||

\twolineshloka
{स रोषाद्रक्तनयनो राक्षसेन्द्रो बलोद्धतः}
{इत्येवमर्जुनामात्यानाह गम्भीरया गिरा} %||7-32-25||

\twolineshloka
{अमात्याः क्षिप्रमाख्यात हैहयस्य नृपस्य वै}
{युद्धार्थं समनुप्राप्तो रावणो नाम नामतः} %||7-32-26||

\twolineshloka
{रावणस्य वचः श्रुत्वा मन्त्रिणोऽथार्जुनस्य ते}
{उत्तस्थुः सायुधास्तं च रावणं वाक्यमब्रुवन्} %||7-32-27||

\threelineshloka
{युद्धस्य कालो विज्ञातः साधु भोः साधु रावण}
{यः क्षीबं स्त्रीवृतं चैव योद्धुमिच्छसि नो नृपम्}
{वाशितामध्यगं मत्तं शार्दूल इव कुञ्जरम्} %||7-32-28||

\twolineshloka
{क्षमस्वाद्य दशग्रीव उष्यतां रजनी त्वया}
{युद्धश्रद्धा तु यद्यस्ति श्वस्तात समरेऽर्जुनम्} %||7-32-29||

\twolineshloka
{यदि वापि त्वरा तुभ्यं युद्धतृष्णासमावृता}
{निहत्यास्मांस्ततो युद्धमर्जुनेनोपयास्यसि} %||7-32-30||

\twolineshloka
{ततस्ते रावणामात्यैरमात्याः पार्थिवस्य तु}
{सूदिताश्चापि ते युद्धे भक्षिताश्च बुभुक्षितैः} %||7-32-31||

\twolineshloka
{ततो हलहलाशब्दो नर्मदा तिर आबभौ}
{अर्जुनस्यानुयात्राणां रावणस्य च मन्त्रिणाम्} %||7-32-32||

\twolineshloka
{इषुभिस्तोमरैः शूलैर्वज्रकल्पैः सकर्षणैः}
{सरावणानर्दयन्तः समन्तात्समभिद्रुताः} %||7-32-33||

\twolineshloka
{हैहयाधिपयोधानां वेग आसीत्सुदारुणः}
{सनक्रमीनमकरसमुद्रस्येव निस्वनः} %||7-32-34||

\twolineshloka
{रावणस्य तु तेऽमात्याः प्रहस्तशुकसारणाः}
{कार्तवीर्यबलं क्रुद्धा निर्दहन्त्यग्नितेजसः} %||7-32-35||

\twolineshloka
{अर्जुनाय तु तत्कर्म रावणस्य समन्त्रिणः}
{क्रीडमानाय कथितं पुरुषैर्द्वाररक्षिभिः} %||7-32-36||

\twolineshloka
{उक्त्वा न भेतव्यमिति स्त्रीजनं स ततोऽर्जुनः}
{उत्ततार जलात्तस्माद्गङ्गातोयादिवाञ्जनः} %||7-32-37||

\twolineshloka
{क्रोधदूषितनेत्रस्तु स ततोऽर्जुन पावकः}
{प्रजज्वाल महाघोरो युगान्त इव पावकः} %||7-32-38||

\twolineshloka
{स तूर्णतरमादाय वरहेमाङ्गदो गदाम्}
{अभिद्रवति रक्षांसि तमांसीव दिवाकरः} %||7-32-39||

\twolineshloka
{बाहुविक्षेपकरणां समुद्यम्य महागदाम्}
{गारुडं वेगमास्थाय आपपातैव सोऽर्जुनः} %||7-32-40||

\twolineshloka
{तस्य मर्गं समावृत्य विन्ध्योऽर्कस्येव पर्वतः}
{स्थितो विन्ध्य इवाकम्प्यः प्रहस्तो मुसलायुधः} %||7-32-41||

\twolineshloka
{ततोऽस्य मुसलं घोरं लोहबद्धं मदोद्धतः}
{प्रहस्तः प्रेषयन्क्रुद्धो ररास च यथाम्बुदः} %||7-32-42||

\twolineshloka
{तस्याग्रे मुसलस्याग्निरशोकापीडसंनिभः}
{प्रहस्तकरमुक्तस्य बभूव प्रदहन्निव} %||7-32-43||

\twolineshloka
{आधावमानं मुसलं कार्तवीर्यस्तदार्जुनः}
{निपुणं वञ्चयामास सगदो गजविक्रमः} %||7-32-44||

\twolineshloka
{ततस्तमभिदुद्राव प्रहस्तं हैहयाधिपः}
{भ्रामयाणो गदां गुर्वीं पञ्चबाहुशतोच्छ्रयाम्} %||7-32-45||

\twolineshloka
{तेनाहतोऽतिवेगेन प्रहस्तो गदया तदा}
{निपपात स्थितः शैलो वज्रिवज्रहतो यथा} %||7-32-46||

\twolineshloka
{प्रहस्तं पतितं दृष्ट्वा मारीचशुकसारणाः}
{समहोदरधूम्राक्षा अपसृप्ता रणाजिरात्} %||7-32-47||

\twolineshloka
{अपक्रान्तेष्वमात्येषु प्रहस्ते च निपातिते}
{रावणोऽभ्यद्रवत्तूर्णमर्जुनं नृपसत्तमम्} %||7-32-48||

\twolineshloka
{सहस्रबाहोस्तद्युद्धं विंशद्बाहोश्च दारुणम्}
{नृपराक्षसयोस्तत्र आरब्धं लोमहर्षणम्} %||7-32-49||

\twolineshloka
{सागराविव सङ्क्षुब्धौ चलमूलाविवाचलौ}
{तेजोयुक्ताविवादित्यौ प्रदहन्ताविवानलौ} %||7-32-50||

\twolineshloka
{बलोद्धतौ यथा नागौ वाशितार्थे यथा वृषौ}
{मेघाविव विनर्दन्तौ सिंहाविव बलोत्कटौ} %||7-32-51||

\twolineshloka
{रुद्रकालाविव क्रुद्धौ तौ तथा राक्षसार्जुनौ}
{परस्परं गदाभ्यां तौ ताडयामासतुर्भृशम्} %||7-32-52||

\twolineshloka
{वज्रप्रहारानचला यथा घोरान्विषेहिरे}
{गदाप्रहारांस्तद्वत्तौ सहेते नरराक्षसौ} %||7-32-53||

\twolineshloka
{यथाशनिरवेभ्यस्तु जायते वै प्रतिश्रुतिः}
{तथा ताभ्यां गदापातैर्दिशः सर्वाः प्रतिश्रुताः} %||7-32-54||

\twolineshloka
{अर्जुनस्य गदा सा तु पात्यमानाहितोरसि}
{काञ्चनाभं नभश्चक्रे विद्युत्सौदामनी यथा} %||7-32-55||

\twolineshloka
{तथैव रावणेनापि पात्यमाना मुहुर्मुहुः}
{अर्जुनोरसि निर्भाति गदोल्केव महागिरौ} %||7-32-56||

\twolineshloka
{नार्जुनः खेदमाप्नोति न राक्षसगणेश्वरः}
{सममासीत्तयोर्युद्धं यथा पूर्वं बलीन्द्रयोः} %||7-32-57||

\twolineshloka
{शृङ्गैर्महर्षभौ यद्वद्दन्ताग्रैरिव कुञ्जरौ}
{परस्परं विनिघ्नन्तौ नरराक्षससत्तमौ} %||7-32-58||

\twolineshloka
{ततोऽर्जुनेन क्रुद्धेन सर्वप्राणेन सा गदा}
{स्तनयोरन्तरे मुक्ता रावणस्य महाहवे} %||7-32-59||

\twolineshloka
{वरदानकृतत्राणे सा गदा रावणोरसि}
{दुर्बलेव यथा सेना द्विधाभूतापतत्क्षितौ} %||7-32-60||

\twolineshloka
{स त्वर्जुनप्रमुक्तेन गदापातेन रावणः}
{अपासर्पद्धनुर्मात्रं निषसाद च निष्टनन्} %||7-32-61||

\twolineshloka
{स विह्वलं तदालक्ष्य दशग्रीवं ततोऽर्जुनः}
{सहसा प्रतिजग्राह गरुत्मानिव पन्नगम्} %||7-32-62||

\twolineshloka
{स तं बाहुसहस्रेण बलाद्गृह्य दशाननम्}
{बबन्ध बलवान्राजा बलिं नारायणो यथा} %||7-32-63||

\twolineshloka
{बध्यमाने दशग्रीवे सिद्धचारणदेवताः}
{साध्वीति वादिनः पुष्पैः किरन्त्यर्जुनमूर्धनि} %||7-32-64||

\twolineshloka
{व्याघ्रो मृगमिवादाय सिंहराडिव दन्तिनम्}
{ररास हैहयो राजा हर्षादम्बुदवन्मुहुः} %||7-32-65||

\twolineshloka
{प्रहस्तस्तु समाश्वस्तो दृष्ट्वा बद्धं दशाननम्}
{सह तै राकसैः क्रुद्ध अभिदुद्राव पार्थिवम्} %||7-32-66||

\twolineshloka
{नक्तञ्चराणां वेगस्तु तेषामापततां बभौ}
{उद्धृत आतपापाये समुद्राणामिवाद्भुतः} %||7-32-67||

\twolineshloka
{मुञ्च मुञ्चेति भाषन्तस्तिष्ठ तिष्ठेति चासकृत्}
{मुसलानि च शूलानि उत्ससर्जुस्तदार्जुने} %||7-32-68||

\twolineshloka
{अप्राप्तान्येव तान्याशु असम्भ्रान्तस्तदार्जुनः}
{आयुधान्यमरारीणां जग्राह रिपुसूदनः} %||7-32-69||

\twolineshloka
{ततस्तैरेव रक्षांसि दुर्धरैः प्रवरायुधैः}
{भित्त्वा विद्रावयामास वायुरम्बुधरानिव} %||7-32-70||

\twolineshloka
{राक्षसांस्त्रासयित्वा तु कार्तवीर्यार्जुनस्तदा}
{रावणं गृह्य नगरं प्रविवेश सुहृद्वृतः} %||7-32-71||

\fourlineindentedshloka
{स कीर्यमाणः कुसुमाक्षतोत्करैर्-}
{ द्विजैः सपौरैः पुरुहूतसंनिभः}
{तदार्जुनः सम्प्रविवेश तां पुरीम्}
{ बलिं निगृह्यैव सहस्रलोचनः} %||7-33-72||


{॥इत्यार्षे श्रीमद्रामायणे वाल्मीकीये आदिकाव्ये उत्तरकाण्डे द्वात्रिंशः सर्गः॥}

\sect{त्रयस्त्रिंशः सर्गः}

\twolineshloka
{रावणग्रहणं तत्तु वायुग्रहणसंनिभम्}
{ऋषिः पुलस्त्यः शुश्राव कथितं दिवि दैवतैः} %||7-33-1||

\twolineshloka
{ततः पुत्रसुतस्नेहात्कम्प्यमानो महाधृतिः}
{माहिष्मतीपतिं द्रष्टुमाजगाम महानृषिः} %||7-33-2||

\twolineshloka
{स वायुमार्गमास्थाय वायुतुल्यगतिर्द्विजः}
{पुरीं माहिष्मतीं प्राप्तो मनःसन्तापविक्रमः} %||7-33-3||

\twolineshloka
{सोऽमरावतिसङ्काशां हृष्टपुष्टजनावृताम्}
{प्रविवेश पुरीं ब्रह्मा इन्द्रस्येवामरावतीम्} %||7-33-4||

\twolineshloka
{पादचारमिवादित्यं निष्पतन्तं सुदुर्दृशम्}
{ततस्ते प्रत्यभिज्ञाय अर्जुनाय न्यवेदयन्} %||7-33-5||

\twolineshloka
{पुलस्त्य इति तं श्रुत्वा वचनं हैहयाधिपः}
{शिरस्यञ्जलिमुद्धृत्य प्रत्युद्गच्छद्द्विजोत्तमम्} %||7-33-6||

\twolineshloka
{पुरोहितोऽस्य गृह्यार्घ्यं मधुपर्कं तथैव च}
{पुरस्तात्प्रययौ राज्ञ इन्द्रस्येव बृहस्पतिः} %||7-33-7||

\twolineshloka
{ततस्तमृषिमायान्तमुद्यन्तमिव भास्करम्}
{अर्जुनो दृश्य सम्प्राप्तं ववन्देन्द्र इवेश्वरम्} %||7-33-8||

\twolineshloka
{स तस्य मधुपर्कं च पाद्यमर्घ्यं च दापयन्}
{पुलस्त्यमाह राजेन्द्रो हर्षगद्गदया गिरा} %||7-33-9||

\twolineshloka
{अद्येयममरावत्या तुल्या माहिष्मती कृता}
{अद्याहं तु द्विजेन्द्रेन्द्र यस्मात्पश्यामि दुर्दृशम्} %||7-33-10||

\twolineshloka
{अद्य मे कुशलं देव अद्य मे कुलमुद्धृतम्}
{यत्ते देवगणैर्वन्द्यौ वन्देऽहं चरणाविमौ} %||7-33-11||

\twolineshloka
{इदं राज्यमिमे पुत्रा इमे दारा इमे वयम्}
{ब्रह्मन्किं कुर्म किं कार्यमाज्ञापयतु नो भवान्} %||7-33-12||

\twolineshloka
{तं धर्मेऽग्निषु भृत्येषु शिवं पृष्ट्वाथ पार्थिवम्}
{पुलस्त्योवाच राजानं हैहयानां तदार्जुनम्} %||7-33-13||

\twolineshloka
{राजेन्द्रामलपद्माक्षपूर्णचन्द्रनिभानन}
{अतुलं ते बलं येन दशग्रीवस्त्वया जितः} %||7-33-14||

\twolineshloka
{भयाद्यस्यावतिष्ठेतां निष्पन्दौ सागरानिलौ}
{सोऽयमद्य त्वया बद्धः पौत्रो मेऽतीवदुर्जयः} %||7-33-15||

\twolineshloka
{तत्पुत्रक यशः स्फीतं नाम विश्रावितं त्वया}
{मद्वाक्याद्याच्यमानोऽद्य मुञ्च वत्स दशाननम्} %||7-33-16||

\twolineshloka
{पुलस्त्याज्ञां स गृह्याथ अकिञ्चनवचोऽर्जुनः}
{मुमोच पार्थिवेन्द्रेन्द्रो राक्षसेन्द्रं प्रहृष्टवत्} %||7-33-17||

\fourlineindentedshloka
{स तं प्रमुक्त्वा त्रिदशारिमर्जुनः}
{ प्रपूज्य दिव्याभरणस्रगम्बरैः}
{अहिंसाकं सख्यमुपेत्य साग्निकम्}
{ प्रणम्य स ब्रह्मसुतं गृहं ययौ} %||7-33-18||

\twolineshloka
{पुलस्त्येनापि सङ्गम्य राक्षसेन्द्रः प्रतापवान्}
{परिष्वङ्गकृतातिथ्यो लज्जमानो विसर्जितः} %||7-33-19||

\twolineshloka
{पितामहसुतश्चापि पुलस्त्यो मुनिसत्तमः}
{मोचयित्वा दशग्रीवं ब्रह्मलोकं जगाम सः} %||7-33-20||

\twolineshloka
{एवं स रावणः प्राप्तः कार्तवीर्यात्तु धर्षणात्}
{पुलस्त्यवचनाच्चापि पुनर्मोक्षमवाप्तवान्} %||7-33-21||

\twolineshloka
{एवं बलिभ्यो बलिनः सन्ति राघवनन्दन}
{नावज्ञा परतः कार्या य इच्छेच्छ्रेय आत्मनः} %||7-33-22||

\fourlineindentedshloka
{ततः स राजा पिशिताशनानाम्}
{ सहस्रबाहोरुपलभ्य मैत्रीम्}
{पुनर्नराणां कदनं चकार}
{ चचार सर्वां पृथिवीं च दर्पात्} %||7-34-23||


{॥इत्यार्षे श्रीमद्रामायणे वाल्मीकीये आदिकाव्ये उत्तरकाण्डे त्रयस्त्रिंशः सर्गः॥}

\sect{चतुस्त्रिंशः सर्गः}

\twolineshloka
{अर्जुनेन विमुक्तस्तु रावणो राक्षसाधिपः}
{चचार पृथिवीं सर्वामनिर्विण्णस्तथा कृतः} %||7-34-1||

\twolineshloka
{राक्षसं वा मनुष्यं वा शृणुते यं बलाधिकम्}
{रावणस्तं समासाद्य युद्धे ह्वयति दर्पितः} %||7-34-2||

\twolineshloka
{ततः कदाचित्किष्किन्धां नगरीं वालिपालिताम्}
{गत्वाह्वयति युद्धाय वालिनं हेममालिनम्} %||7-34-3||

\twolineshloka
{ततस्तं वानरामात्यस्तारस्तारापिता प्रभुः}
{उवाच रावणं वाक्यं युद्धप्रेप्सुमुपागतम्} %||7-34-4||

\twolineshloka
{राक्षसेन्द्र गतो वाली यस्ते प्रतिबलो भवेत्}
{नान्यः प्रमुखतः स्थातुं तव शक्तः प्लवङ्गमः} %||7-34-5||

\twolineshloka
{चतुर्भ्योऽपि समुद्रेभ्यः सन्ध्यामन्वास्य रावण}
{इमं मुहूर्तमायाति वाली तिष्ठ मुहूर्तकम्} %||7-34-6||

\twolineshloka
{एतानस्थिचयान्पश्य य एते शङ्खपाण्डुराः}
{युद्धार्थिनामिमे राजन्वानराधिपतेजसा} %||7-34-7||

\twolineshloka
{यद्वामृतरसः पीतस्त्वया रावणराक्षस}
{तथा वालिनमासाद्य तदन्तं तव जीवितम्} %||7-34-8||

\twolineshloka
{अथ वा त्वरसे मर्तुं गच्छ दक्षिणसागरम्}
{वालिनं द्रक्ष्यसे तत्र भूमिष्ठमिव भास्करम्} %||7-34-9||

\twolineshloka
{स तु तारं विनिर्भर्त्स्य रावणो राक्षसेश्वरः}
{पुष्पकं तत्समारुह्य प्रययौ दक्षिणार्णवम्} %||7-34-10||

\twolineshloka
{तत्र हेमगिरिप्रख्यं तरुणार्कनिभाननम्}
{रावणो वालिनं दृष्ट्वा सन्ध्योपासनतत्परम्} %||7-34-11||

\twolineshloka
{पुष्पकादवरुह्याथ रावणोऽञ्जनसंनिभः}
{ग्रहीतुं वालिनं तूर्णं निःशब्दपदमाद्रवत्} %||7-34-12||

\twolineshloka
{यदृच्छयोन्मीलयता वालिनापि स रावणः}
{पापाभिप्रायवान्दृष्टश्चकार न च सम्भ्रमम्} %||7-34-13||

\twolineshloka
{शशमालक्ष्य सिंहो वा पन्नगं गरुडो यथा}
{न चिन्तयति तं वाली रावणं पापनिश्चयम्} %||7-34-14||

\twolineshloka
{जिघृक्षमाणमद्यैनं रावणं पापबुद्धिनम्}
{कक्षावलम्बिनं कृत्वा गमिष्यामि महार्णवान्} %||7-34-15||

\twolineshloka
{द्रक्ष्यन्त्यरिं ममाङ्कस्थं स्रंसितोरुकराम्बरम्}
{लम्बमानं दशग्रीवं गरुडस्येव पन्नगम्} %||7-34-16||

\twolineshloka
{इत्येवं मतिमास्थाय वाली कर्णमुपाश्रितः}
{जपन्वै नैगमान्मन्त्रांस्तस्थौ पर्वतराडिव} %||7-34-17||

\twolineshloka
{तावन्योन्यं जिघृक्षन्तौ हरिराक्षसपार्थिवौ}
{प्रयत्नवन्तौ तत्कर्म ईहतुर्बलदर्पितौ} %||7-34-18||

\twolineshloka
{हस्तग्राह्यं तु तं मत्वा पादशब्देन रावणम्}
{पराङ्मुखोऽपि जग्राह वाली सर्पमिवाण्डजः} %||7-34-19||

\twolineshloka
{ग्रहीतुकामं तं गृह्य रक्षसामीश्वरं हरिः}
{खमुत्पपात वेगेन कृत्वा कक्षावलम्बिनम्} %||7-34-20||

\twolineshloka
{स तं पीड्दयमानस्तु वितुदन्तं नखैर्मुहुः}
{जहार रावणं वाली पवनस्तोयदं यथा} %||7-34-21||

\twolineshloka
{अथ ते राक्षसामात्या ह्रियमाणे दशानने}
{मुमोक्षयिषवो घोरा रवमाणा ह्यभिद्रवन्} %||7-34-22||

\twolineshloka
{अन्वीयमानस्तैर्वाली भ्राजतेऽम्बरमध्यगः}
{अन्वीयमानो मेघौघैरम्बरस्थ इवांशुमान्} %||7-34-23||

\twolineshloka
{तेऽशक्नुवन्तः सम्प्राप्तं वालिनं राक्षसोत्तमाः}
{तस्य बाहूरुवेगेन परिश्रान्तः पतन्ति च} %||7-34-24||

{वालिमार्गादपाक्रामन्पर्वतेन्द्रा हि गच्छतः॥\devanumber{25}॥} %||7-34-25|| (Check)

\addtocounter{shlokacount}{1}

\twolineshloka
{अपक्षिगणसम्पातो वानरेन्द्रो महाजवः}
{क्रमशः सागरान्सर्वान्सन्ध्याकालमवन्दत} %||7-34-26||

\twolineshloka
{सभाज्यमानो भूतैस्तु खेचरैः खेचरो हरिः}
{पश्चिमं सागरं वाली आजगाम सरावणः} %||7-34-27||

\twolineshloka
{तत्र सन्ध्यामुपासित्वा स्नात्वा जप्त्वा च वानरः}
{उत्तरं सागरं प्रायाद्वहमानो दशाननम्} %||7-34-28||

\twolineshloka
{उत्तरे सागरे सन्ध्यामुपासित्वा दशाननम्}
{वहमानोऽगमद्वाली पूर्वमम्बुमहानिधिम्} %||7-34-29||

\twolineshloka
{तत्रापि सन्ध्यामन्वास्य वासविः स हरीश्वरः}
{किष्किन्धाभिमुखो गृह्य रावणं पुनरागमत्} %||7-34-30||

\twolineshloka
{चतुर्ष्वपि समुद्रेषु सन्ध्यामन्वास्य वानरः}
{रावणोद्वहनश्रान्तः किष्किन्धोपवनेऽपतत्} %||7-34-31||

\twolineshloka
{रावणं तु मुमोचाथ स्वकक्षात्कपिसत्तमः}
{कुतस्त्वमिति चोवाच प्रहसन्रावणं प्रति} %||7-34-32||

\twolineshloka
{विस्मयं तु महद्गत्वा श्रमलोकनिरीक्षणः}
{राक्षसेशो हरीशं तमिदं वचनमब्रवीत्} %||7-34-33||

\twolineshloka
{वानरेन्द्र महेन्द्राभ राक्षसेन्द्रोऽस्मि रावणः}
{युद्धेप्सुरहं सम्प्राप्तः स चाद्यासादितस्त्वया} %||7-34-34||

\twolineshloka
{अहो बलमहो वीर्यमहो गम्भीरता च ते}
{येनाहं पशुवद्गृह्य भ्रामितश्चतुरोऽर्णवान्} %||7-34-35||

\twolineshloka
{एवमश्रान्तवद्वीर शीघ्रमेव च वानर}
{मां चैवोद्वहमानस्तु कोऽन्यो वीरः क्रमिष्यति} %||7-34-36||

\twolineshloka
{त्रयाणामेव भूतानां गतिरेषा प्लवङ्गम}
{मनोऽनिलसुपर्णानां तव वा नात्र संशयः} %||7-34-37||

\twolineshloka
{सोऽहं दृष्टबलस्तुभ्यमिच्छामि हरिपुङ्गव}
{त्वया सह चिरं सख्यं सुस्निग्धं पावकाग्रतः} %||7-34-38||

\twolineshloka
{दाराः पुत्राः पुरं राष्ट्रं भोगाच्छादनभोजनम्}
{सर्वमेवाविभक्तं नौ भविष्यति हरीश्वर} %||7-34-39||

\twolineshloka
{ततः प्रज्वालयित्वाग्निं तावुभौ हरिराक्षसौ}
{भ्रातृत्वमुपसम्पन्नौ परिष्वज्य परस्परम्} %||7-34-40||

\twolineshloka
{अन्योन्यं लम्बितकरौ ततस्तौ हरिराक्षसौ}
{किष्किन्धां विशतुर्हृष्टौ सिंहौ गिरिगुहामिव} %||7-34-41||

\twolineshloka
{स तत्र मासमुषितः सुग्रीव इव रावणः}
{अमात्यैरागतैर्नीचस्त्रैलोक्योत्सादनार्थिभिः} %||7-34-42||

\twolineshloka
{एवमेतत्पुरावृत्तं वालिना रावणः प्रभो}
{धर्षितश्च कृतश्चापि भ्राता पावकसंनिधौ} %||7-34-43||

\twolineshloka
{बलमप्रतिमं राम वालिनोऽभवदुत्तमम्}
{सोऽपि तया विनिर्दग्धः शलभो वह्निना यथा} %||7-35-44||


{॥इत्यार्षे श्रीमद्रामायणे वाल्मीकीये आदिकाव्ये उत्तरकाण्डे चतुस्त्रिंशः सर्गः॥}

\sect{पञ्चत्रिंशः सर्गः}

\twolineshloka
{अपृच्छत ततो रामो दक्षिणाशालयं मुनिम्}
{प्राञ्जलिर्विनयोपेत इदमाह वचोऽर्थवत्} %||7-35-1||

\twolineshloka
{अतुलं बलमेताभ्यां वालिनो रावणस्य च}
{न त्वेतौ हनुमद्वीर्यैः समाविति मतिर्मम} %||7-35-2||

\twolineshloka
{शौर्यं दाक्ष्यं बलं धैर्यं प्राज्ञता नयसाधनम्}
{विक्रमश्च प्रभावश्च हनूमति कृतालयाः} %||7-35-3||

\twolineshloka
{दृष्ट्वोदधिं विषीदन्तीं तदैष कपिवाहिनीम्}
{समाश्वास्य कपीन्भूयो योजनानां शतं प्लुतः} %||7-35-4||

\twolineshloka
{धर्षयित्वा पुरीं लङ्कां रावणान्तःपुरं तथा}
{दृष्ट्वा सम्भाषिता चापि सीता विश्वासिता तथा} %||7-35-5||

\twolineshloka
{सेनाग्रगा मन्त्रिसुताः किङ्करा रावणात्मजः}
{एते हनुमता तत्र एकेन विनिपातिताः} %||7-35-6||

\twolineshloka
{भूयो बन्धाद्विमुक्तेन सम्भाषित्वा दशाननम्}
{लङ्का भस्मीकृता तेन पावकेनेव मेदिनी} %||7-35-7||

\twolineshloka
{न कालस्य न शक्रस्य न विष्णोर्वित्तपस्य च}
{कर्माणि तानि श्रूयन्ते यानि युद्धे हनूमतः} %||7-35-8||

\twolineshloka
{एतस्य बाहुवीर्येण लङ्का सीता च लक्ष्मणः}
{प्राप्तो मया जयश्चैव राज्यं मित्राणि बान्धवाः} %||7-35-9||

\twolineshloka
{हनूमान्यदि मे न स्याद्वानराधिपतेः सखा}
{प्रवृत्तमपि को वेत्तुं जानक्याः शक्तिमान्भवेत्} %||7-35-10||

\twolineshloka
{किमर्थं वाली चैतेन सुग्रीवप्रियकाम्यया}
{तदा वैरे समुत्पन्ने न दग्धो वीरुधो यथा} %||7-35-11||

\twolineshloka
{न हि वेदितवान्मन्ये हनूमानात्मनो बलम्}
{यद्दृष्टवाञ्जीवितेष्टं क्लिश्यन्तं वानराधिपम्} %||7-35-12||

\twolineshloka
{एतन्मे भगवन्सर्वं हनूमति महामुने}
{विस्तरेण यथातत्त्वं कथयामरपूजित} %||7-35-13||

\twolineshloka
{राघवस्य वचः श्रुत्वा हेतुयुक्तमृषिस्ततः}
{हनूमतः समक्षं तमिदं वचनमब्रवीत्} %||7-35-14||

\twolineshloka
{सत्यमेतद्रघुश्रेष्ठ यद्ब्रवीषि हनूमतः}
{न बले विद्यते तुल्यो न गतौ न मतौ परः} %||7-35-15||

\twolineshloka
{अमोघशापैः शापस्तु दत्तोऽस्य ऋषिभिः पुरा}
{न वेदिता बलं येन बली सन्नरिमर्दनः} %||7-35-16||

\twolineshloka
{बाल्येऽप्येतेन यत्कर्म कृतं राम महाबल}
{तन्न वर्णयितुं शक्यमतिबालतयास्य ते} %||7-35-17||

\twolineshloka
{यदि वास्ति त्वभिप्रायस्तच्छ्रोतुं तव राघव}
{समाधाय मतिं राम निशामय वदाम्यहम्} %||7-35-18||

\twolineshloka
{सूर्यदत्तवरस्वर्णः सुमेरुर्नाम पर्वतः}
{यत्र राज्यं प्रशास्त्यस्य केषरी नाम वै पिता} %||7-35-19||

\twolineshloka
{तस्य भार्या बभूवेष्टा ह्यञ्जनेति परिश्रुता}
{जनयामास तस्यां वै वायुरात्मजमुत्तमम्} %||7-35-20||

\twolineshloka
{शालिशूकसमाभासं प्रासूतेमं तदाञ्जना}
{फलान्याहर्तुकामा वै निष्क्रान्ता गहने चरा} %||7-35-21||

\twolineshloka
{एष मातुर्वियोगाच्च क्षुधया च भृशार्दितः}
{रुरोद शिशुरत्यर्थं शिशुः शरभराडिव} %||7-35-22||

\twolineshloka
{ततोद्यन्तं विवस्वन्तं जपापुष्पोत्करोपमम्}
{ददृशे फललोभाच्च उत्पपात रविं प्रति} %||7-35-23||

\twolineshloka
{बालार्काभिमुखो बालो बालार्क इव मूर्तिमान्}
{ग्रहीतुकामो बालार्कं प्लवतेऽम्बरमध्यगः} %||7-35-24||

\twolineshloka
{एतस्मिन्प्लवमाने तु शिशुभावे हनूमति}
{देवदानवसिद्धानां विस्मयः सुमहानभूत्} %||7-35-25||

\twolineshloka
{नाप्येवं वेगवान्वायुर्गरुडो न मनस्तथा}
{यथायं वायुपुत्रस्तु क्रमतेऽम्बरमुत्तमम्} %||7-35-26||

\twolineshloka
{यदि तावच्छिशोरस्य ईदृशौ गतिविक्रमौ}
{यौवनं बलमासाद्य कथं वेगो भविष्यति} %||7-35-27||

\twolineshloka
{तमनुप्लवते वायुः प्लवन्तं पुत्रमात्मनः}
{सूर्यदाहभयाद्रक्षंस्तुषारचयशीतलः} %||7-35-28||

\twolineshloka
{बहुयोजनसाहस्रं क्रमत्येष ततोऽम्बरम्}
{पितुर्बलाच्च बाल्याच्च भास्कराभ्याशमागतः} %||7-35-29||

\twolineshloka
{शिशुरेष त्वदोषज्ञ इति मत्वा दिवाकरः}
{कार्यं चात्र समायत्तमित्येवं न ददाह सः} %||7-35-30||

\twolineshloka
{यमेव दिवसं ह्येष ग्रहीतुं भास्करं प्लुतः}
{तमेव दिवसं राहुर्जिघृक्षति दिवाकरम्} %||7-35-31||

\twolineshloka
{अनेन च परामृष्टो राम सूर्यरथोपरि}
{अपक्रान्तस्ततस्त्रस्तो राहुश्चन्द्रार्कमर्दनः} %||7-35-32||

\twolineshloka
{स इन्द्रभवनं गत्वा सरोषः सिंहिकासुतः}
{अब्रवीद्भ्रुकुटीं कृत्वा देवं देवगणैर्वृतम्} %||7-35-33||

\twolineshloka
{बुभुक्षापनयं दत्त्वा चन्द्रार्कौ मम वासव}
{किमिदं तत्त्वया दत्तमन्यस्य बलवृत्रहन्} %||7-35-34||

\twolineshloka
{अद्याहं पर्वकाले तु जिघृक्षुः सूर्यमागतः}
{अथान्यो राहुरासाद्य जग्राह सहसा रविम्} %||7-35-35||

\twolineshloka
{स राहोर्वचनं श्रुत्वा वासवः सम्भ्रमान्वितः}
{उत्पपातासनं हित्वा उद्वहन्काञ्चनस्रजम्} %||7-35-36||

\twolineshloka
{ततः कैलासकूटाभं चतुर्दन्तं मदस्रवम्}
{शृङ्गारकारिणं प्रांशुं स्वर्णघण्टाट्टहासिनम्} %||7-35-37||

\twolineshloka
{इन्द्रः करीन्द्रमारुह्य राहुं कृत्वा पुरःसरम्}
{प्रायाद्यत्राभवत्सूर्यः सहानेन हनूमता} %||7-35-38||

\twolineshloka
{अथातिरभसेनागाद्राहुरुत्सृज्य वासवम्}
{अनेन च स वै दृष्ट आधावञ्शैलकूटवत्} %||7-35-39||

\twolineshloka
{ततः सूर्यं समुत्सृज्य राहुमेवमवेक्ष्य च}
{उत्पपात पुनर्व्योम ग्रहीतुं सिंहिकासुतम्} %||7-35-40||

\twolineshloka
{उत्सृज्यार्कमिमं राम आधावन्तं प्लवङ्गमम्}
{दृष्ट्वा राहुः परावृत्य मुखशेषः पराङ्मुखः} %||7-35-41||

\twolineshloka
{इन्द्रमाशंसमानस्तु त्रातारं सिंहिकासुतः}
{इन्द्र इन्द्रेति सन्त्रासान्मुहुर्मुहुरभाषत} %||7-35-42||

\twolineshloka
{राहोर्विक्रोशमानस्य प्रागेवालक्षितः स्वरः}
{श्रुत्वेन्द्रोवाच मा भैषीरयमेनं निहन्म्यहम्} %||7-35-43||

\twolineshloka
{ऐरावतं ततो दृष्ट्वा महत्तदिदमित्यपि}
{फलं तं हस्तिराजानमभिदुद्राव मारुतिः} %||7-35-44||

\twolineshloka
{तदास्य धावतो रूपमैरावतजिघृक्षया}
{मुहूर्तमभवद्घोरमिन्द्राग्न्योरिव भास्वरम्} %||7-35-45||

\twolineshloka
{एवमाधावमानं तु नातिक्रुद्धः शचीपतिः}
{हस्तान्तेनातिमुक्तेन कुलिशेनाभ्यताडयत्} %||7-35-46||

\twolineshloka
{ततो गिरौ पपातैष इन्द्रवज्राभिताडितः}
{पतमानस्य चैतस्य वामो हनुरभज्यत} %||7-35-47||

\twolineshloka
{तस्मिंस्तु पतिते बाले वज्रताडनविह्वले}
{चुक्रोधेन्द्राय पवनः प्रजानामशिवाय च} %||7-35-48||

\twolineshloka
{विण्मूत्राशयमावृत्य प्रजास्वन्तर्गतः प्रभुः}
{रुरोध सर्वभूतानि यथा वर्षाणि वासवः} %||7-35-49||

\twolineshloka
{वायुप्रकोपाद्भूतानि निरुच्छ्वासानि सर्वतः}
{सन्धिभिर्भज्यमानानि काष्ठभूतानि जज्ञिरे} %||7-35-50||

\twolineshloka
{निःस्वधं निर्वषट्कारं निष्क्रियं धर्मवर्जितम्}
{वायुप्रकोपात्त्रैलोक्यं निरयस्थमिवाबभौ} %||7-35-51||

\twolineshloka
{ततः प्रजाः सगन्धर्वाः सदेवासुरमानुषाः}
{प्रजापतिं समाधावन्नसुखार्ताः सुखैषिणः} %||7-35-52||

\twolineshloka
{ऊचुः प्राञ्जलयो देवा दरोदरनिभोदराः}
{त्वया स्म भगवन्सृष्टाः प्रजानाथ चतुर्विधाः} %||7-35-53||

\twolineshloka
{त्वया दत्तोऽयमस्माकमायुषः पवनः पतिः}
{सोऽस्मान्प्राणेश्वरो भूत्वा कस्मादेषोऽद्य सत्तम} %||7-35-54||

\twolineshloka
{रुरोध दुःखं जनयन्नन्तःपुर इव स्त्रियः}
{तस्मात्त्वां शरणं प्राप्ता वायुनोपहता विभो} %||7-35-55||

{वायुसंरोधजं दुःखमिदं नो नुद शत्रुहन्॥\devanumber{56}॥} %||7-35-56|| (Check)

\addtocounter{shlokacount}{1}

\twolineshloka
{एतत्प्रजानां श्रुत्वा तु प्रजानाथः प्रजापतिः}
{कारणादिति तानुक्त्वा प्रजाः पुनरभाषत} %||7-35-57||

\twolineshloka
{यस्मिन्वः कारणे वायुश्चुक्रोध च रुरोध च}
{प्रजाः शृणुध्वं तत्सर्वं श्रोतव्यं चात्मनः क्षमम्} %||7-35-58||

\twolineshloka
{पुत्रस्तस्यामरेशेन इन्द्रेणाद्य निपातितः}
{राहोर्वचनमाज्ञाय राज्ञा वः कोपितोऽनिलः} %||7-35-59||

\twolineshloka
{अशरीरः शरीरेषु वायुश्चरति पालयन्}
{शरीरं हि विना वायुं समतां याति रेणुभिः} %||7-35-60||

\twolineshloka
{वायुः प्राणाः सुखं वायुर्वायुः सर्वमिदं जगत्}
{वायुना सम्परित्यक्तं न सुखं विन्दते जगत्} %||7-35-61||

\twolineshloka
{अद्यैव च परित्यक्तं वायुना जगदायुषा}
{अद्यैवेमे निरुच्छ्वासाः काष्ठकुड्योपमाः स्थिताः} %||7-35-62||

\twolineshloka
{तद्यामस्तत्र यत्रास्ते मारुतो रुक्प्रदो हि वः}
{मा विनाशं गमिष्याम अप्रसाद्यादितेः सुतम्} %||7-35-63||

\fourlineindentedshloka
{ततः प्रजाभिः सहितः प्रजापतिः}
{ सदेवगन्धर्वभुजङ्गगुह्यकः}
{जगाम तत्रास्यति यत्र मारुतः}
{ सुतं सुरेन्द्राभिहतं प्रगृह्य सः} %||7-35-64||

\fourlineindentedshloka
{ततोऽर्कवैश्वानरकाञ्चनप्रभम्}
{ सुतं तदोत्सङ्गगतं सदागतेः}
{चतुर्मुखो वीक्ष्य कृपामथाकरोत्}
{ सदेवसिद्धर्षिभुजङ्गराक्षसः} %||7-36-65||


{॥इत्यार्षे श्रीमद्रामायणे वाल्मीकीये आदिकाव्ये उत्तरकाण्डे पञ्चत्रिंशः सर्गः॥}

\sect{षट्त्रिंशः सर्गः}

\twolineshloka
{ततः पितामहं दृष्ट्वा वायुः पुत्रवधार्दितः}
{शिशुकं तं समादाय उत्तस्थौ धातुरग्रतः} %||7-36-1||

\twolineshloka
{चलत्कुण्डलमौलिस्रक्तपनीयविभूषणः}
{पादयोर्न्यपतद्वायुस्तिस्रोऽवस्थाय वेधसे} %||7-36-2||

\twolineshloka
{तं तु वेदविदाद्यस्तु लम्बाभरणशोभिना}
{वायुमुत्थाप्य हस्तेन शिशुं तं परिमृष्टवान्} %||7-36-3||

\twolineshloka
{स्पृष्टमात्रस्ततः सोऽथ सलीलं पद्मजन्मना}
{जलसिक्तं यथा सस्यं पुनर्जीवितमाप्तवान्} %||7-36-4||

\twolineshloka
{प्राणवन्तमिमं दृष्ट्वा प्राणो गन्धवहो मुदा}
{चचार सर्वभूतेषु संनिरुद्धं यथापुरा} %||7-36-5||

\twolineshloka
{मरुद्रोगविनिर्मुक्ताः प्रजा वै मुदिताभवन्}
{शीतवातविनिर्मुक्ताः पद्मिन्य इव साम्बुजाः} %||7-36-6||

\twolineshloka
{ततस्त्रियुग्मस्त्रिककुत्त्रिधामा त्रिदशार्चितः}
{उवाच देवता ब्रह्मा मारुतप्रियकाम्यया} %||7-36-7||

\twolineshloka
{भो महेन्द्राग्निवरुणधनेश्वरमहेश्वराः}
{जानतामपि तत्सर्वं हितं वक्ष्यामि श्रूयताम्} %||7-36-8||

\twolineshloka
{अनेन शिशुना कार्यं कर्तव्यं वो भविष्यति}
{ददतास्य वरान्सर्वे मारुतस्यास्य तुष्टिदान्} %||7-36-9||

\twolineshloka
{ततः सहस्रनयनः प्रीतिरक्तः शुभाननः}
{कुशे शयमयीं मालां समुत्क्षिप्येदमब्रवीत्} %||7-36-10||

\twolineshloka
{मत्करोत्सृष्टवज्रेण हनुरस्य यथा क्षतः}
{नाम्नैष कपिशार्दूलो भविता हनुमानिति} %||7-36-11||

\twolineshloka
{अहमेवास्य दास्यामि परमं वरमुत्तमम्}
{अतः प्रभृति वज्रस्य ममावध्यो भविष्यति} %||7-36-12||

\twolineshloka
{मार्ताण्डस्त्वब्रवीत्तत्र भगवांस्तिमिरापहः}
{तेजसोऽस्य मदीयस्य ददामि शतिकां कलाम्} %||7-36-13||

\twolineshloka
{यदा तु शास्त्राण्यध्येतुं शक्तिरस्य भविष्यति}
{तदास्य शास्त्रं दास्यामि येन वाग्मी भविष्यति} %||7-36-14||

\twolineshloka
{वरुणश्च वरं प्रादान्नास्य मृत्युर्भविष्यति}
{वर्षायुतशतेनापि मत्पाशादुदकादपि} %||7-36-15||

\twolineshloka
{यमोऽपि दण्डावध्यत्वमरोगत्वं च नित्यशः}
{दिशतेऽस्य वरं तुष्ट अविषादं च संयुगे} %||7-36-16||

\twolineshloka
{गदेयं मामिका नैनं संयुगेषु वधिष्यति}
{इत्येवं वरदः प्राह तदा ह्येकाक्षिपिङ्गलः} %||7-36-17||

\twolineshloka
{मत्तो मदायुधानां च न वध्योऽयं भविष्यति}
{इत्येवं शङ्करेणापि दत्तोऽस्य परमो वरः} %||7-36-18||

\twolineshloka
{सर्वेषां ब्रह्मदण्डानामवध्योऽयं भविष्यति}
{दीर्घायुश्च महात्मा च इति ब्रह्माब्रवीद्वचः} %||7-36-19||

\twolineshloka
{विश्वकर्मा तु दृष्ट्वैनं बालसूर्योपमं शिशुम्}
{शिल्पिनां प्रवरः प्राह वरमस्य महामतिः} %||7-36-20||

\twolineshloka
{विनिर्मितानि देवानामायुधानीह यानि तु}
{तेषां सङ्ग्रामकाले तु अवध्योऽयं भविष्यति} %||7-36-21||

\twolineshloka
{ततः सुराणां तु वरैर्दृष्ट्वा ह्येनमलङ्कृतम्}
{चतुर्मुखस्तुष्टमुखो वायुमाह जगद्गुरुः} %||7-36-22||

\twolineshloka
{अमित्राणां भयकरो मित्राणामभयङ्करः}
{अजेयो भविता तेऽत्र पुत्रो मारुतमारुतिः} %||7-36-23||

\twolineshloka
{रावणोत्सादनार्थानि रामप्रीतिकराणि च}
{रोमहर्षकराण्येष कर्ता कर्माणि संयुगे} %||7-36-24||

\twolineshloka
{एवमुक्त्वा तमामन्त्र्य मारुतं तेऽमरैः सह}
{यथागतं ययुः सर्वे पितामहपुरोगमाः} %||7-36-25||

\twolineshloka
{सोऽपि गन्धवहः पुत्रं प्रगृह्य गृहमानयत्}
{अञ्जनायास्तमाख्याय वरं दत्तं विनिःसृतः} %||7-36-26||

\twolineshloka
{प्राप्य राम वरानेष वरदानबलान्वितः}
{बलेनात्मनि संस्थेन सोऽपूर्यत यथार्णवः} %||7-36-27||

\twolineshloka
{बलेनापूर्यमाणो हि एष वानरपुङ्गवः}
{आश्रमेषु महर्षीणामपराध्यति निर्भयः} %||7-36-28||

\twolineshloka
{स्रुग्भाण्डानग्निहोत्रं च वल्कलानां च सञ्चयान्}
{भग्नविच्छिन्नविध्वस्तान्सुशान्तानां करोत्ययम्} %||7-36-29||

\twolineshloka
{सर्वेषां ब्रह्मदण्डानामवध्यं ब्रह्मणा कृतम्}
{जानन्त ऋषयस्तं वै क्षमन्ते तस्य नित्यशः} %||7-36-30||

\twolineshloka
{यदा केषरिणा त्वेष वायुना साञ्जनेन च}
{प्रतिषिद्धोऽपि मर्यादां लङ्घयत्येव वानरः} %||7-36-31||

\twolineshloka
{ततो महर्षयः क्रुद्धा भृग्वङ्गिरसवंशजाः}
{शेपुरेनं रघुश्रेष्ठ नातिक्रुद्धातिमन्यवः} %||7-36-32||

\twolineshloka
{बाधसे यत्समाश्रित्य बलमस्मान्प्लवङ्गम}
{तद्दीर्घकालं वेत्तासि नास्माकं शापमोहितः} %||7-36-33||

\twolineshloka
{ततस्तु हृततेजौजा महर्षिवचनौजसा}
{एषो श्रमाणि नात्येति मृदुभावगतश्चरन्} %||7-36-34||

\twolineshloka
{अथ ऋक्षरजा नाम वालिसुग्रीवयोः पिता}
{सर्ववानरराजासीत्तेजसा इव भास्करः} %||7-36-35||

\twolineshloka
{स तु राज्यं चिरं कृत्वा वानराणां हरीश्वरः}
{ततस्त्वर्क्षरजा नाम कालधर्मेण सङ्गतः} %||7-36-36||

\twolineshloka
{तस्मिन्नस्तमिते वाली मन्त्रिभिर्मन्त्रकोविदैः}
{पित्र्ये पदे कृतो राजा सुग्रीवो वालिनः पदे} %||7-36-37||

\twolineshloka
{सुग्रीवेण समं त्वस्य अद्वैधं छिद्रवर्जितम्}
{अहार्यं सख्यमभवदनिलस्य यथाग्निना} %||7-36-38||

\twolineshloka
{एष शापवशादेव न वेदबलमात्मनः}
{वालिसुग्रीवयोर्वैरं यदा राम समुत्थितम्} %||7-36-39||

\twolineshloka
{न ह्येष राम सुग्रीवो भ्राम्यमाणोऽपि वालिना}
{वेदयानो न च ह्येष बलमात्मनि मारुतिः} %||7-36-40||

\fourlineindentedshloka
{पराक्रमोत्साहमतिप्रतापैः}
{ सौशील्यमाधुर्यनयानयैश्च}
{गाम्भीर्यचातुर्यसुवीर्यधैर्यैर्-}
{ हनूमतः कोऽप्यधिकोऽस्ति लोके} %||7-36-41||

\fourlineindentedshloka
{असौ पुरा व्याकरणं ग्रहीष्य}
{न्सूर्योन्मुखः पृष्ठगमः कपीन्द्रः}
{उद्यद्गिरेरस्तगिरिं जगाम}
{ ग्रन्थं महद्धारयदप्रमेयः} %||7-36-42||

\fourlineindentedshloka
{प्रवीविविक्षोरिव सागरस्य}
{ लोकान्दिधक्षोरिव पावकस्य}
{लोकक्षयेष्वेव यथान्तकस्य}
{ हनूमतः स्थास्यति कः पुरस्तात्} %||7-36-43||

\fourlineindentedshloka
{एषोऽपि चान्ये च महाकपीन्द्राः}
{ सुग्रीवमैन्दद्विविदाः सनीलाः}
{सतारतारेयनलाः सरम्भा}
{स्त्वत्कारणाद्राम सुरैर्हि सृष्टाः} %||7-36-44||

\twolineshloka
{तदेतत्कथितं सर्वं यन्मां त्वं परिपृच्छसि}
{हनूमतो बालभावे कर्मैतत्कथितं मया} %||7-36-45||

\twolineshloka
{दृष्टः सम्भाषितश्चासि राम गच्छमहे वयम्}
{एवमुक्त्वा गताः सर्वे ऋषयस्ते यथागतम्} %||7-37-46||


{॥इत्यार्षे श्रीमद्रामायणे वाल्मीकीये आदिकाव्ये उत्तरकाण्डे षट्त्रिंशः सर्गः॥}

\sect{सप्तत्रिंशः सर्गः}

\twolineshloka
{विमृश्य च ततो रामो वयस्यमकुतोभयम्}
{प्रतर्दनं काशिपतिं परिष्वज्येदमब्रवीत्} %||7-37-1||

\twolineshloka
{दर्शिता भवता प्रीतिर्दर्शितं सौहृदं परम्}
{उद्योगश्च कृतो राजन्भरतेन त्वया सह} %||7-37-2||

\twolineshloka
{तद्भवानद्य काशेयीं पुरीं वाराणसीं व्रज}
{रमणीयां त्वया गुप्तां सुप्राकारां सुतोरणाम्} %||7-37-3||

\twolineshloka
{एतावदुक्त्वा उत्थाय काकुत्स्थः परमासनात्}
{पर्यष्वजत धर्मात्मा निरन्तरमुरोगतम्} %||7-37-4||

\twolineshloka
{विसृज्य तं वयस्यं स स्वागतान्पृथिवीपतीन्}
{प्रहसन्राघवो वाक्यमुवाच मधुराक्षरम्} %||7-37-5||

\twolineshloka
{भवतां प्रीतिरव्यग्रा तेजसा परिरक्षिता}
{धर्मश्च नियतो नित्यं सत्यं च भवतां सदा} %||7-37-6||

\twolineshloka
{युष्माकं च प्रभावेन तेजसा च महात्मनाम्}
{हतो दुरात्मा दुर्बुद्धी रावणो राक्षसाधिपः} %||7-37-7||

\twolineshloka
{हेतुमात्रमहं तत्र भवतां तेजसां हतः}
{रावणः सगणो युद्धे सपुत्रः सहबान्धवः} %||7-37-8||

\twolineshloka
{भवन्तश्च समानीता भरतेन महात्मना}
{श्रुत्वा जनकराजस्य कानने तनयां हृताम्} %||7-37-9||

\twolineshloka
{उद्युक्तानां च सर्वेषां पार्थिवानां महात्मनाम्}
{कालो ह्यतीतः सुमहान्गमने रोचतां मतिः} %||7-37-10||

\twolineshloka
{प्रत्यूचुस्तं च राजानो हर्षेण महतान्विताः}
{दिष्ट्या त्वं विजयी राम राज्यं चापि प्रतिष्ठितम्} %||7-37-11||

\twolineshloka
{दिष्ट्या प्रत्याहृता सीता दिष्ट्या शत्रुः पराजितः}
{एष नः परमः काम एषा नः कीर्तिरुत्तमा} %||7-37-12||

\twolineshloka
{यत्त्वां विजयिनं राम पश्यामो हतशात्रवम्}
{उपपन्नं च काकुत्स्थ यत्त्वमस्मान्प्रशंससि} %||7-37-13||

\twolineshloka
{प्रशंसार्हा हि जानन्ति प्रशंसां वक्तुमीदृशीम्}
{आपृच्छामो गमिष्यामो हृदिस्थो नः सदा भवान्} %||7-37-14||

{भवेच्च ते महाराज प्रीतिरस्मासु नित्यदा॥\devanumber{15}॥} %||7-38-15|| (Check)

\addtocounter{shlokacount}{1}


{॥इत्यार्षे श्रीमद्रामायणे वाल्मीकीये आदिकाव्ये उत्तरकाण्डे सप्तत्रिंशः सर्गः॥}

\sect{अष्टात्रिंशः सर्गः}

\twolineshloka
{ते प्रयाता महात्मानः पार्थिवाः सर्वतो दिशम्}
{कम्पयन्तो महीं वीराः स्वपुराणि प्रहृष्टवत्} %||7-38-1||

\twolineshloka
{अक्षौहिणी सहस्रैस्ते समवेतास्त्वनेकशः}
{हृष्टाः प्रतिगताः सर्वे राघवार्थे समागताः} %||7-38-2||

\twolineshloka
{ऊचुश्चैव महीपाला बलदर्पसमन्विताः}
{न नाम रावणं युद्धे पश्यामः पुरतः स्थितम्} %||7-38-3||

\twolineshloka
{भरतेन वयं पश्चात्समानीता निरर्थकम्}
{हता हि राक्षसास्तत्र पार्थिवैः स्युर्न संशयः} %||7-38-4||

\twolineshloka
{रामस्य बाहुवीर्येण पालिता लक्ष्मणस्य च}
{सुखं पारे समुद्रस्य युध्येम विगतज्वराः} %||7-38-5||

\twolineshloka
{एताश्चान्याश्च राजानः कथास्तत्र सहस्रशः}
{कथयन्तः स्वराष्ट्राणि विविशुस्ते महारथाः} %||7-38-6||

\twolineshloka
{यथापुराणि ते गत्वा रत्नानि विविधानि च}
{रामाय प्रियकामार्थमुपहारान्नृपा ददुः} %||7-38-7||

\twolineshloka
{अश्वान्रत्नानि वस्त्राणि हस्तिनश्च मदोत्कटान्}
{चन्दनानि च दिव्यानि दिव्यान्याभरणानि च} %||7-38-8||

\twolineshloka
{भरतो लक्ष्मणश्चैव शत्रुघ्नश्च महारथः}
{आदाय तानि रत्नानि अयोध्यामगमन्पुनः} %||7-38-9||

\twolineshloka
{आगताश्च पुरीं रम्यामयोध्यां पुरुषर्षभाः}
{ददुः सर्वाणि रत्नानि राघवाय महात्मने} %||7-38-10||

\twolineshloka
{प्रतिगृह्य च तत्सर्वं प्रीतियुक्तः स राघवः}
{सर्वाणि तानि प्रददौ सुग्रीवाय महात्मने} %||7-38-11||

\twolineshloka
{विभीषणाय च ददौ ये चान्ये ऋक्षवानराः}
{हनूमत्प्रमुखा वीरा राक्षसाश्च महाबलाः} %||7-38-12||

\twolineshloka
{ते सर्वे हृष्टमनसो रामदत्तानि तान्यथ}
{शिरोभिर्धारयामासुर्बाहुभिश्च महाबलाः} %||7-38-13||

\twolineshloka
{पपुश्चैव सुगन्धीनि मधूनि विविधानि च}
{मांसानि च सुमृष्टानि फलान्यास्वादयन्ति च} %||7-38-14||

\twolineshloka
{एवं तेषां निवसतां मासः साग्रो गतस्तदा}
{मुहूर्तमिव तत्सर्वं रामभक्त्या समर्थयन्} %||7-38-15||

\twolineshloka
{रेमे रामः स तैः सार्धं वानरैः कामरूपिभिः}
{राजभिश्च महावीर्यै राक्षसैश्च महाबलैः} %||7-38-16||

\twolineshloka
{एवं तेषां ययौ मासो द्वितीयः शैशिरः सुखम्}
{वानराणां प्रहृष्टानां राक्षसानां च सर्वशः} %||7-39-17||


{॥इत्यार्षे श्रीमद्रामायणे वाल्मीकीये आदिकाव्ये उत्तरकाण्डे अष्टात्रिंशः सर्गः॥}

\sect{एकोनचत्वारिंशः सर्गः}

\twolineshloka
{तथा स्म तेषां वसतामृक्षवानररक्षसाम्}
{राघवस्तु महातेजाः सुग्रीवमिदमब्रवीत्} %||7-39-1||

\twolineshloka
{गम्यतां सौम्य किष्किन्धां दुराधर्षं सुरासुरैः}
{पालयस्व सहामात्यै राज्यं निहतकण्टकम्} %||7-39-2||

\twolineshloka
{अङ्गदं च महाबाहो प्रीत्या परमयान्वितः}
{पश्य त्वं हनुमन्तं च नलं च सुमहाबलम्} %||7-39-3||

\twolineshloka
{सुषेणं श्वशुरं शूरं तारं च बलिनां वरम्}
{कुमुदं चैव दुर्धर्षं नीलं च सुमहाबलम्} %||7-39-4||

\twolineshloka
{वीरं शतबलिं चैव मैन्दं द्विविदमेव च}
{गजं गवाक्षं गवयं शरभं च महाबलम्} %||7-39-5||

\twolineshloka
{ऋक्षराजं च दुर्धर्षं जाम्बवन्तं महाबलम्}
{पश्य प्रीतिसमायुक्तो गन्धमादनमेव च} %||7-39-6||

\twolineshloka
{ये चान्ये सुमहात्मानो मदर्थे त्यक्तजीविताः}
{पश्य त्वं प्रीतिसंयुक्तो मा चैषां विप्रियं कृथाः} %||7-39-7||

\twolineshloka
{एवमुक्त्वा च सुग्रीवं प्रशस्य च पुनः पुनः}
{विभीषणमथोवाच रामो मधुरया गिरा} %||7-39-8||

\twolineshloka
{तङ्कां प्रशाधि धर्मेण सम्मतो ह्यसि पार्थिव}
{पुरस्य राक्षसानां च भ्रातुर्वैश्वरणस्य च} %||7-39-9||

\twolineshloka
{मा च बुद्धिमधर्मे त्वं कुर्या राजन्कथञ्चन}
{बुद्धिमन्तो हि राजानो ध्रुवमश्नन्ति मेदिनीम्} %||7-39-10||

\twolineshloka
{अहं च नित्यशो राजन्सुग्रीवसहितस्त्वया}
{स्मर्तव्यः परया प्रीत्या गच्छ त्वं विगतज्वरः} %||7-39-11||

\twolineshloka
{रामस्य भाषितं श्रुत्वा ऋष्कवानरराक्षसाः}
{साधु साध्विति काकुत्स्थं प्रशशंसुः पुनः पुनः} %||7-39-12||

\twolineshloka
{तव बुद्धिर्महाबाहो वीर्यमद्भुतमेव च}
{माधुर्यं परमं राम स्वयम्भोरिव नित्यदा} %||7-39-13||

\twolineshloka
{तेषामेवं ब्रुवाणानां वानराणां च रक्षसाम्}
{हनूमत्प्रणतो भूत्वा राघवं वाक्यमब्रवीत्} %||7-39-14||

\twolineshloka
{स्नेहो मे परमो राजंस्त्वयि नित्यं प्रतिष्ठितः}
{भक्तिश्च नियता वीर भावो नान्यत्र गच्छति} %||7-39-15||

\twolineshloka
{यावद्रामकथां वीर श्रोष्येऽहं पृथिवीतले}
{तावच्छरीरे वत्स्यन्तु मम प्राणा न संशयः} %||7-39-16||

\twolineshloka
{एवं ब्रुवाणं राजेन्द्रो हनूमन्तमथासनात्}
{उत्थाय च परिष्वज्य वाक्यमेतदुवाच ह} %||7-39-17||

\twolineshloka
{एवमेतत्कपिश्रेष्ठ भविता नात्र संशयः}
{लोका हि यावत्स्थास्यन्ति तावत्स्थास्यति मे कथा} %||7-39-18||

\twolineshloka
{चरिष्यति कथा यावल्लोकानेषा हि मामिका}
{तावच्छरीरे वत्स्यन्ति प्राणास्तव न संशयः} %||7-39-19||

\twolineshloka
{ततोऽस्य हारं चन्द्राभं मुच्य कण्ठात्स राघवः}
{वैदूर्यतरलं स्नेहादाबबन्धे हनूमति} %||7-39-20||

\twolineshloka
{तेनोरसि निबद्धेन हारेण स महाकपिः}
{रराज हेमशैलेन्द्रश्चन्द्रेणाक्रान्तमस्तकः} %||7-39-21||

\twolineshloka
{श्रुत्वा तु राघवस्यैतदुत्थायोत्थाय वानराः}
{प्रणम्य शिरसा पादौ प्रजग्मुस्ते महाबलाः} %||7-39-22||

\twolineshloka
{सुग्रीवश्चैव रामेण परिष्वक्तो महाभुजः}
{विभीषणश्च धर्मात्मा निरन्तरमुरोगतः} %||7-39-23||

\twolineshloka
{सर्वे च ते बाष्पगलाः साश्रुनेत्रा विचेतसः}
{सम्मूढा इव दुःखेन त्यजन्ते राघवं तदा} %||7-40-24||


{॥इत्यार्षे श्रीमद्रामायणे वाल्मीकीये आदिकाव्ये उत्तरकाण्डे एकोनचत्वारिंशः सर्गः॥}

\sect{चत्वारिंशः सर्गः}

\twolineshloka
{विसृज्य च महाबाहुरृक्षवानरराक्षसान्}
{भ्रातृभिः सहितो रामः प्रमुमोद सुखी सुखम्} %||7-40-1||

\twolineshloka
{अथापराह्णसमये भ्रातृभिः सह राघवः}
{शुश्राव मधुरां वाणीमन्तरिक्षात्प्रभाषिताम्} %||7-40-2||

\twolineshloka
{सौम्य राम निरीक्षस्व सौम्येन वदनेन माम्}
{कैलासशिखरात्प्राप्तं विद्धि मां पुष्करं प्रभो} %||7-40-3||

\twolineshloka
{तव शासनमाज्ञाय गतोऽस्मि धनदं प्रति}
{उपस्थातुं नरश्रेष्ठ स च मां प्रत्यभाषत} %||7-40-4||

\twolineshloka
{निर्जितस्त्वं नरेन्द्रेण राघवेण महात्मना}
{निहत्य युधि दुर्धर्षं रावणं राक्षसाधिपम्} %||7-40-5||

\twolineshloka
{ममापि परमा प्रीतिर्हते तस्मिन्दुरात्मनि}
{रावणे सगणे सौम्य सपुत्रामात्यबान्धवे} %||7-40-6||

\twolineshloka
{स त्वं रामेण लङ्कायां निर्जितः परमात्मना}
{वह सौम्य तमेव त्वमहमाज्ञापयामि ते} %||7-40-7||

\twolineshloka
{एष मे परमः कामो यत्त्वं राघवनन्दनम्}
{वहेर्लोकस्य संयानं गच्छस्व विगतज्वरः} %||7-40-8||

\twolineshloka
{तच्छासनमहं ज्ञात्वा धनदस्य महात्मनः}
{त्वत्सकाशं पुनः प्राप्तः स एवं प्रतिगृह्ण माम्} %||7-40-9||

\twolineshloka
{बाढमित्येव काकुत्स्थः पुष्पकं समपूजयत्}
{लाजाक्षतैश्च पुष्पैश्च गन्धैश्च सुसुगन्धिभिः} %||7-40-10||

\threelineshloka
{गम्यतां च यथाकाममागच्छेस्त्वं यदा स्मरे}
{एवमस्त्विति रामेण विसृष्टः पुष्पकः पुनः}
{अभिप्रेतां दिशं प्रायात्पुष्पकः पुष्पभूषितः} %||7-40-11||

\twolineshloka
{एवमन्तर्हिते तस्मिन्पुष्पके विविधात्मनि}
{भरतः प्राञ्जलिर्वाक्यमुवाच रघुनन्दनम्} %||7-40-12||

\twolineshloka
{अत्यद्भुतानि दृश्यन्ते त्वयि राज्यं प्रशासति}
{अमानुषाणां सत्त्वानां व्याहृतानि मुहुर्मुहुः} %||7-40-13||

\twolineshloka
{अनामयाच्च मर्त्यानां साग्रो मासो गतो ह्ययम्}
{जीर्णानामपि सत्त्वानां मृत्युर्नायाति राघव} %||7-40-14||

\twolineshloka
{पुत्रान्नार्यः प्रसूयन्ते वपुष्मन्तश्च मानवाः}
{हर्षश्चाभ्यधिको राजञ्जनस्य पुरवासिनः} %||7-40-15||

\twolineshloka
{काले च वासवो वर्षं पातयत्यमृतोपमम्}
{वायवश्चापि वायन्ते स्पर्शवन्तः सुखप्रदाः} %||7-40-16||

\twolineshloka
{ईदृशो नश्चिरं राजा भवत्विति नरेश्वर}
{कथयन्ति पुरे पौरा जना जनपदेषु च} %||7-40-17||

\twolineshloka
{एता वाचः सुमधुरा भरतेन समीरिताः}
{श्रुत्वा रामो मुदा युक्तः प्रमुमोद सुखी सुखम्} %||7-41-18||


{॥इत्यार्षे श्रीमद्रामायणे वाल्मीकीये आदिकाव्ये उत्तरकाण्डे चत्वारिंशः सर्गः॥}

\sect{एकचत्वारिंशः सर्गः}

\twolineshloka
{स विसृज्य ततो रामः पुष्पकं हेमभूषितम्}
{प्रविवेश महाबाहुरशोकवनिकां तदा} %||7-41-1||

\twolineshloka
{चन्दनागरुचूतैश्च तुङ्गकालेयकैरपि}
{देवदारुवनैश्चापि समन्तादुपशोभिताम्} %||7-41-2||

\twolineshloka
{प्रियङ्गुभिः कदम्बैश्च तथा कुरबकैरपि}
{जम्बूभिः पाटलीभिश्च कोविदारैश्च संवृताम्} %||7-41-3||

\twolineshloka
{सर्वदा कुसुमै रम्यैः फलवद्भिर्मनोरमैः}
{चारुपल्लवपुष्पाढ्यैर्मत्तभ्रमरसङ्कुलैः} %||7-41-4||

\twolineshloka
{कोकिलैर्भृङ्गराजैश्च नानावर्णैश्च पक्षिभिः}
{शोभितां शतशश्चित्रैश्चूतवृक्षावतंसकैः} %||7-41-5||

\twolineshloka
{शातकुम्भनिभाः केचित्केचिदग्निशिखोपमाः}
{नीलाञ्जननिभाश्चान्ये भान्ति तत्र स्म पादपाः} %||7-41-6||

\twolineshloka
{दीर्घिका विविधाकाराः पूर्णाः परमवारिणा}
{महार्हमणिसोपानस्फटिकान्तरकुट्टिमाः} %||7-41-7||

\twolineshloka
{फुल्लपद्मोत्पलवनाश्चक्रवाकोपशोभिताः}
{प्राकारैर्विविधाकारैः शोभिताश्च शिलातलैः} %||7-41-8||

\twolineshloka
{तत्र तत्र वनोद्देशे वैदूर्यमणिसंनिभैः}
{शाद्वलैः परमोपेताः पुष्पितद्रुमसंयुताः} %||7-41-9||

\twolineshloka
{नन्दनं हि यथेन्द्रस्य ब्राह्मं चैत्ररथं यथा}
{तथारूपं हि रामस्य काननं तन्निवेशितम्} %||7-41-10||

\twolineshloka
{बह्वासनगृहोपेतां लतागृहसमावृताम्}
{अशोकवनिकां स्फीतां प्रविश्य रघुनन्दनः} %||7-41-11||

\twolineshloka
{आसने तु शुभाकारे पुष्पस्तबकभूषिते}
{कुथास्तरणसंवीते रामः संनिषसाद ह} %||7-41-12||

\twolineshloka
{सीतां सङ्गृह्य बाहुभ्यां मधुमैरेयमुत्तमम्}
{पाययामास काकुत्स्थः शचीमिन्द्रो यथामृतम्} %||7-41-13||

\twolineshloka
{मांसानि च विचित्राणि फलानि विविधानि च}
{रामस्याभ्यवहारार्थं किङ्करास्तूर्णमाहरन्} %||7-41-14||

\twolineshloka
{उपनृत्यन्ति राजानं नृत्यगीतविशारदाः}
{बालाश्च रूपवत्यश्च स्त्रियः पानवशं गताः} %||7-41-15||

\twolineshloka
{एवं रामो मुदा युक्तः सीतां सुरुचिराननाम्}
{रमयामास वैदेहीमहन्यहनि देववत्} %||7-41-16||

\twolineshloka
{तथा तु रममाणस्य तस्यैवं शिशिरः शुभः}
{अत्यक्रामन्नरेन्द्रस्य राघवस्य महात्मनः} %||7-41-17||

\twolineshloka
{पूर्वाह्णे पौरकृत्यानि कृत्वा धर्मेण धर्मवित्}
{शेषं दिवसभागार्धमन्तःपुरगतोऽभवत्} %||7-41-18||

\twolineshloka
{सीता च देवकार्याणि कृत्वा पौर्वाह्णिकानि तु}
{श्वश्रूणामविशेषेण सर्वासां प्राञ्जलिः स्थिता} %||7-41-19||

\twolineshloka
{ततो राममुपागच्छद्विचित्रबहुभूषणा}
{त्रिविष्टपे सहस्राक्षमुपविष्टं यथा शची} %||7-41-20||

\twolineshloka
{दृष्ट्वा तु राघवः पत्नीं कल्याणेन समन्विताम्}
{प्रहर्षमतुलं लेभे साधु साध्विति चाब्रवीत्} %||7-41-21||

\twolineshloka
{अपत्यलाभो वैदेहि ममायं समुपस्थितः}
{किमिच्छसि हि तद्ब्रूहि कः कामः क्रियतां तव} %||7-41-22||

\twolineshloka
{प्रहसन्ती तु वैदेही रामं वाक्यमथाब्रवीत्}
{तपोवनानि पुण्यानि द्रष्टुमिच्छामि राघव} %||7-41-23||

\twolineshloka
{गङ्गातीरे निविष्टानि ऋषीणां पुण्यकर्मणाम्}
{फलमूलाशिनां वीर पादमूलेषु वर्तितुम्} %||7-41-24||

\twolineshloka
{एष मे परमः कामो यन्मूलफलभोजिषु}
{अप्येकरात्रं काकुत्स्थ वसेयं पुण्यशालिषु} %||7-41-25||

\twolineshloka
{तथेति च प्रतिज्ञातं रामेणाक्लिष्टकर्मणा}
{विस्रब्धा भव वैदेहि श्वो गमिष्यस्यसंशयम्} %||7-41-26||

\twolineshloka
{एवमुक्त्वा तु काकुत्स्थो मैथिलीं जनकात्मजाम्}
{मध्यकक्षान्तरं रामो निर्जगाम सुहृद्वृतः} %||7-42-27||


{॥इत्यार्षे श्रीमद्रामायणे वाल्मीकीये आदिकाव्ये उत्तरकाण्डे एकचत्वारिंशः सर्गः॥}

\sect{द्विचत्वारिंशः सर्गः}

\twolineshloka
{तत्रोपविष्टं राजानमुपासन्ते विचक्षणाः}
{कथानां बहुरूपाणां हास्यकाराः समन्ततः} %||7-42-1||

\twolineshloka
{विजयो मधुमत्तश्च काश्यपः पिङ्गलः कुशः}
{सुराजिः कालियो भद्रो दन्तवक्रः समागधः} %||7-42-2||

\twolineshloka
{एते कथा बहुविधा परिहाससमन्विताः}
{कथयन्ति स्म संहृष्टा राघवस्य महात्मनः} %||7-42-3||

\twolineshloka
{ततः कथायां कस्याञ्चिद्राघवः समभाषत}
{काः कथा नगरे भद्र वर्तन्ते विषयेषु च} %||7-42-4||

\twolineshloka
{मामाश्रितानि कान्याहुः पौरजानपदा जनाः}
{किं च सीतां समाश्रित्य भरतं किं नु लक्ष्मणम्} %||7-42-5||

\twolineshloka
{किं नु शत्रुघ्नमाश्रित्य कैकेयीं मातरं च मे}
{वक्तव्यतां च राजानो नवे राज्ये व्रजन्ति हि} %||7-42-6||

\twolineshloka
{एवमुक्ते तु रामेण भद्रः प्राञ्जलिरब्रवीत्}
{स्थिताः कथाः शुभा राजन्वर्तन्ते पुरवासिनाम्} %||7-42-7||

\twolineshloka
{अयं तु विजयः सौम्य दशग्रीववधाश्रितः}
{भूयिष्ठं स्वपुरे पौरैः कथ्यते पुरुषर्षभ} %||7-42-8||

\twolineshloka
{एवमुक्तस्तु भद्रेण राघवो वाक्यमब्रवीत्}
{कथयस्व यथातथ्यं सर्वं निरवशेषतः} %||7-42-9||

\twolineshloka
{शुभाशुभानि वाक्यानि यान्याहुः पुरवासिनः}
{श्रुत्वेदानीं शुभं कुर्यां न कुर्यामशुभानि च} %||7-42-10||

\twolineshloka
{कथयस्व च विस्रब्धो निर्भयो विगतज्वरः}
{कथयन्ते यथा पौरा जना जनपदेषु च} %||7-42-11||

\twolineshloka
{राघवेणैवमुक्तस्तु भद्रः सुरुचिरं वचः}
{प्रत्युवाच महाबाहुं प्राञ्जलिः सुसमाहितः} %||7-42-12||

\twolineshloka
{शृणु राजन्यथा पौराः कथयन्ति शुभाशुभम्}
{चत्वरापणरथ्यासु वनेषूपवनेषु च} %||7-42-13||

\twolineshloka
{दुष्करं कृतवान्रामः समुद्रे सेतुबन्धनम्}
{अकृतं पूर्वकैः कैश्चिद्देवैरपि सदानवैः} %||7-42-14||

\twolineshloka
{रावणश्च दुराधर्षो हतः सबलवाहनः}
{वानराश्च वशं नीता ऋक्षाश्च सह राक्षसैः} %||7-42-15||

\twolineshloka
{हत्वा च रावणं युद्धे सीतामाहृत्य राघवः}
{अमर्षं पृष्ठतः कृत्वा स्ववेश्म पुनरानयत्} %||7-42-16||

\twolineshloka
{कीदृशं हृदये तस्य सीतासम्भोगजं सुखम्}
{अङ्कमारोप्य हि पुरा रावणेन बलाद्धृताम्} %||7-42-17||

\twolineshloka
{लङ्कामपि पुनर्नीतामशोकवनिकां गताम्}
{रक्षसां वशमापन्नां कथं रामो न कुत्सते} %||7-42-18||

\twolineshloka
{अस्माकमपि दारेषु सहनीयं भविष्यति}
{यथा हि कुरुते राजा प्रजा तमनुवर्तते} %||7-42-19||

\twolineshloka
{एवं बहुविधा वाचो वदन्ति पुरवासिनः}
{नगरेषु च सर्वेषु राजञ्जनपदेषु च} %||7-42-20||

\twolineshloka
{तस्यैतद्भाषितं श्रुत्वा राघवः परमार्तवत्}
{उवाच सर्वान्सुहृदः कथमेतन्निवेद्यताम्} %||7-42-21||

\twolineshloka
{सर्वे तु शिरसा भूमावभिवाद्य प्रणम्य च}
{प्रत्यूचू राघवं दीनमेवमेतन्न संशयः} %||7-42-22||

\twolineshloka
{श्रुत्वा तु वाक्यं काकुत्स्थः सर्वेषां समुदीरितम्}
{विसर्जयामास तदा सर्वांस्ताञ्शत्रुतापनः} %||7-43-23||


{॥इत्यार्षे श्रीमद्रामायणे वाल्मीकीये आदिकाव्ये उत्तरकाण्डे द्विचत्वारिंशः सर्गः॥}

\sect{त्रिचत्वारिंशः सर्गः}

\twolineshloka
{विसृज्य तु सुहृद्वर्गं बुद्ध्या निश्चित्य राघवः}
{समीपे द्वाःस्थमासीनमिदं वचनमब्रवीत्} %||7-43-1||

\twolineshloka
{शीघ्रमानय सौमित्रिं लक्ष्मणं शुभलक्षणम्}
{भरतं च महाबाहुं शत्रुघ्नं चापराजितम्} %||7-43-2||

\twolineshloka
{रामस्य भाषितं श्रुत्वा द्वाःस्थो मूर्ध्नि कृताञ्जलिः}
{लक्ष्मणस्य गृहं गत्वा प्रविवेशानिवारितः} %||7-43-3||

\twolineshloka
{उवाच च तदा वाक्यं वर्धयित्वा कृताञ्जलिः}
{द्रष्टुमिच्छति राजा त्वां गम्यतां तत्र मा चिरम्} %||7-43-4||

\twolineshloka
{बाढमित्येव सौमित्रिः श्रुत्वा राघवशासनम्}
{प्राद्रवद्रथमारुह्य राघवस्य निवेशनम्} %||7-43-5||

\twolineshloka
{प्रयान्तं लक्ष्मणं दृष्ट्वा द्वाःस्थो भरतमन्तिकात्}
{उवाच प्राञ्जलिर्वाक्यं राजा त्वां द्रष्टुमिच्छति} %||7-43-6||

\twolineshloka
{भरतस्तु वचः श्रुत्वा द्वाःस्थाद्रामसमीरितम्}
{उत्पपातासनात्तूर्णं पद्भ्यामेव ततोऽगमत्} %||7-43-7||

\twolineshloka
{दृष्ट्वा प्रयान्तं भरतं त्वरमाणः कृताञ्जलिः}
{शत्रुघ्नभवनं गत्वा ततो वाक्यं जगाद ह} %||7-43-8||

\twolineshloka
{एह्यागच्छ रघुश्रेष्ठ राजा त्वां द्रष्टुमिच्छति}
{गतो हि लक्ष्मणः पूर्वं भरतश्च महायशाः} %||7-43-9||

\twolineshloka
{श्रुत्वा तु वचनं तस्य शत्रुघ्नो रामशासनम्}
{शिरसा वन्द्य धरणीं प्रययौ यत्र राघवः} %||7-43-10||

\twolineshloka
{कुमारानागताञ्श्रुत्वा चिन्ताव्याकुलितेन्द्रियः}
{अवाक्शिरा दीनमना द्वाःस्थं वचनमब्रवीत्} %||7-43-11||

\twolineshloka
{प्रवेशय कुमारांस्त्वं मत्समीपं त्वरान्वितः}
{एतेषु जीवितं मह्यमेते प्राणा बहिश्चराः} %||7-43-12||

\twolineshloka
{आज्ञप्तास्तु नरेन्द्रेण कुमाराः शुक्लवाससः}
{प्रह्वाः प्राञ्जलयो भूत्वा विविशुस्ते समाहिताः} %||7-43-13||

\twolineshloka
{ते तु दृष्ट्वा मुखं तस्य सग्रहं शशिनं यथा}
{सन्ध्यागतमिवादित्यं प्रभया परिवर्जितम्} %||7-43-14||

\twolineshloka
{बाष्पपूर्णे च नयने दृष्ट्वा रामस्य धीमतः}
{हतशोभं यथा पद्मं मुखं वीक्ष्य च तस्य ते} %||7-43-15||

\twolineshloka
{ततोऽभिवाद्य त्वरिताः पादौ रामस्य मूर्धभिः}
{तस्थुः समाहिताः सर्वे रामश्चाश्रूण्यवर्तयत्} %||7-43-16||

\twolineshloka
{तान्परिष्वज्य बाहुभ्यामुत्थाप्य च महाभुजः}
{आसनेष्वाध्वमित्युक्त्वा ततो वाक्यं जगाद ह} %||7-43-17||

\twolineshloka
{भवन्तो मम सर्वस्वं भवन्तो मम जीवितम्}
{भवद्भिश्च कृतं राज्यं पालयामि नरेश्वराः} %||7-43-18||

\twolineshloka
{भवन्तः कृतशास्त्रार्था बुद्धौ च परिनिष्ठिताः}
{सम्भूय च मदर्थोऽयमन्वेष्टव्यो नरेश्वराः} %||7-44-19||


{॥इत्यार्षे श्रीमद्रामायणे वाल्मीकीये आदिकाव्ये उत्तरकाण्डे त्रिचत्वारिंशः सर्गः॥}

\sect{चतुश्चत्वारिंशः सर्गः}

\twolineshloka
{तेषां समुपविष्टानां सर्वेषां दीनचेतसाम्}
{उवाच वाक्यं काकुत्स्थो मुखेन परिशुष्यता} %||7-44-1||

\twolineshloka
{सर्वे शृणुत भद्रं वो मा कुरुध्वं मनोऽन्यथा}
{पौराणां मम सीतायां यादृशी वर्तते कथा} %||7-44-2||

\twolineshloka
{पौरापवादः सुमहांस्तथा जनपदस्य च}
{वर्तते मयि बीभत्सः स मे मर्माणि कृन्तति} %||7-44-3||

\twolineshloka
{अहं किल कुले जात इक्ष्वाकूणां महात्मनाम्}
{सीतां पापसमाचारामानयेयं कथं पुरे} %||7-44-4||

\twolineshloka
{जानासि हि यथा सौम्य दण्डके विजने वने}
{रावणेन हृता सीता स च विध्वंसितो मया} %||7-44-5||

\twolineshloka
{प्रत्यक्षं तव सौमित्रे देवानां हव्यवाहनः}
{अपापां मैथिलीमाह वायुश्चाकाशगोचरः} %||7-44-6||

\twolineshloka
{चन्द्रादित्यौ च शंसेते सुराणां संनिधौ पुरा}
{ऋषीणां चैव सर्वेषामपापां जनकात्मजाम्} %||7-44-7||

\twolineshloka
{एवं शुद्धसमाचारा देवगन्धर्वसंनिधौ}
{लङ्काद्वीपे महेन्द्रेण मम हस्ते निवेशिता} %||7-44-8||

\twolineshloka
{अन्तरात्मा च मे वेत्ति सीतां शुद्धां यशस्विनीम्}
{ततो गृहीत्वा वैदेहीमयोध्यामहमागतः} %||7-44-9||

\twolineshloka
{अयं तु मे महान्वादः शोकश्च हृदि वर्तते}
{पौरापवादः सुमहांस्तथा जनपदस्य च} %||7-44-10||

\twolineshloka
{अकीर्तिर्यस्य गीयेत लोके भूतस्य कस्यचित्}
{पतत्येवाधमाँल्लोकान्यावच्छब्दः स कीर्त्यते} %||7-44-11||

\twolineshloka
{अकीर्तिर्निन्द्यते दैवैः कीर्तिर्देवेषु पूज्यते}
{कीर्त्यर्थं च समारम्भः सर्व एव महात्मनाम्} %||7-44-12||

\twolineshloka
{अप्यहं जीवितं जह्यां युष्मान्वा पुरुषर्षभाः}
{अपवादभयाद्भीतः किं पुनर्जनकात्मजाम्} %||7-44-13||

\twolineshloka
{तस्माद्भवन्तः पश्यन्तु पतितं शोकसागरे}
{न हि पश्याम्यहं भूयः किञ्चिद्दुःखमतोऽधिकम्} %||7-44-14||

\twolineshloka
{श्वस्त्वं प्रभाते सौमित्रे सुमन्त्राधिष्ठितं रथम्}
{आरुह्य सीतामारोप्य विषयान्ते समुत्सृज} %||7-44-15||

\twolineshloka
{गङ्गायास्तु परे पारे वाल्मीकेः सुमहात्मनः}
{आश्रमो दिव्यसङ्काशस्तमसातीरमाश्रितः} %||7-44-16||

\twolineshloka
{तत्रैनां विजने कक्षे विसृज्य रघुनन्दन}
{शीघ्रमागच्छ सौमित्रे कुरुष्व वचनं मम} %||7-44-17||

\twolineshloka
{न चास्मि प्रतिवक्तव्यः सीतां प्रति कथञ्चन}
{अप्रीतिः परमा मह्यं भवेत्तु प्रतिवारिते} %||7-44-18||

\twolineshloka
{शापिताश्च मया यूयं भुजाभ्यां जीवितेन च}
{ये मां वाक्यान्तरे ब्रूयुरनुनेतुं कथञ्चन} %||7-44-19||

\twolineshloka
{मानयन्तु भवन्तो मां यदि मच्छासने स्थिताः}
{इतोऽद्य नीयतां सीता कुरुष्व वचनं मम} %||7-44-20||

\twolineshloka
{पूर्वमुक्तोऽहमनया गङ्गातीरे महाश्रमान्}
{पश्येयमिति तस्याश्च कामः संवर्त्यतामयम्} %||7-44-21||

\twolineshloka
{एवमुक्त्वा तु काकुत्स्थो बाष्पेण पिहितेक्षणः}
{प्रविवेश स धर्मात्मा भ्रातृभिः परिवारितः} %||7-45-22||


{॥इत्यार्षे श्रीमद्रामायणे वाल्मीकीये आदिकाव्ये उत्तरकाण्डे चतुश्चत्वारिंशः सर्गः॥}

\sect{पञ्चचत्वारिंशः सर्गः}

\twolineshloka
{ततो रजन्यां व्युष्टायां लक्ष्मणो दीनचेतनः}
{सुमन्त्रमब्रवीद्वाक्यं मुखेन परिशुष्यता} %||7-45-1||

\twolineshloka
{सारथे तुरगाञ्शीघ्रं योजयस्व रथोत्तमे}
{स्वास्तीर्णं राजभवनात्सीतायाश्चासनं शुभम्} %||7-45-2||

\twolineshloka
{सीता हि राजभवनादाश्रमं पुण्यकर्मणाम्}
{मया नेया महर्षीणां शीघ्रमानीयतां रथः} %||7-45-3||

\twolineshloka
{सुमन्त्रस्तु तथेत्युक्त्वा युक्तं परमवाजिभिः}
{रथं सुरुचिरप्रख्यं स्वास्तीर्णं सुखशय्यया} %||7-45-4||

\twolineshloka
{आदायोवाच सौमित्रिं मित्राणां हर्षवर्धनम्}
{रथोऽयं समनुप्राप्तो यत्कार्यं क्रियतां प्रभो} %||7-45-5||

\twolineshloka
{एवमुक्तः सुमन्त्रेण राजवेश्म स लक्ष्मणः}
{प्रविश्य सीतामासाद्य व्याजहार नरर्षभः} %||7-45-6||

\twolineshloka
{गङ्गातीरे मया देवि मुनीनामाश्रमे शुभे}
{शीघ्रं गत्वोपनेयासि शासनात्पार्थिवस्य नः} %||7-45-7||

\twolineshloka
{एवमुक्ता तु वैदेही लक्ष्मणेन महात्मना}
{प्रहर्षमतुलं लेभे गमनं चाभ्यरोचयत्} %||7-45-8||

\twolineshloka
{वासांसि च महार्हाणि रत्नानि विविधानि च}
{गृहीत्वा तानि वैदेही गमनायोपचक्रमे} %||7-45-9||

\threelineshloka
{इमानि मुनिपत्नीनां दास्याम्याभरणान्यहम्}
{सौमित्रिस्तु तथेत्युक्त्वा रथमारोप्य मैथिलीम्}
{प्रययौ शीघ्रतुरगो रामस्याज्ञामनुस्मरन्} %||7-45-10||

\twolineshloka
{अब्रवीच्च तदा सीता लक्ष्मणं लक्ष्मिवर्धनम्}
{अशुभानि बहून्यद्य पश्यामि रघुनन्दन} %||7-45-11||

\twolineshloka
{नयनं मे स्फुरत्यद्य गात्रोत्कम्पश्च जायते}
{हृदयं चैव सौमित्रे अस्वस्थमिव लक्षये} %||7-45-12||

\twolineshloka
{औत्सुक्यं परमं चापि अधृतिश्च परा मम}
{शून्यामिव च पश्यामि पृथिवीं पृथुलोचन} %||7-45-13||

\twolineshloka
{अपि स्वस्ति भवेत्तस्य भ्रातुस्ते भ्रातृभिः सह}
{श्वश्रूणां चैव मे वीर सर्वासामविशेषतः} %||7-45-14||

\twolineshloka
{पुरे जनपदे चैव कुशलं प्राणिनामपि}
{इत्यञ्जलिकृता सीता देवता अभ्ययाचत} %||7-45-15||

\twolineshloka
{लक्ष्मणोऽर्थं तु तं श्रुत्वा शिरसा वन्द्य मैथिलीम्}
{शिवमित्यब्रवीद्धृष्टो हृदयेन विशुष्यता} %||7-45-16||

\twolineshloka
{ततो वासमुपागम्य गोमतीतीर आश्रमे}
{प्रभाते पुनरुत्थाय सौमित्रिः सूतमब्रवीत्} %||7-45-17||

\twolineshloka
{योजयस्व रथं शीघ्रमद्य भागीरथीजलम्}
{शिरसा धारयिष्यामि त्र्यम्बकः पर्वते यथा} %||7-45-18||

\twolineshloka
{सोऽश्वान्विचारयित्वाशु रथे युक्त्वा मनोजवान्}
{आरोहस्वेति वैदेहीं सूतः प्राञ्जलिरब्रवीत्} %||7-45-19||

\twolineshloka
{सा तु सूतस्य वचनादारुरोह रथोत्तमम्}
{सीता सौमित्रिणा सार्धं सुमन्त्रेण च धीमता} %||7-45-20||

\twolineshloka
{अथार्धदिवसं गत्वा भागीरथ्या जलाशयम्}
{निरीक्ष्य लक्ष्मणो दीनः प्ररुरोद महास्वनम्} %||7-45-21||

\twolineshloka
{सीता तु परमायत्ता दृष्ट्वा लक्ष्मणमातुरम्}
{उवाच वाक्यं धर्मज्ञ किमिदं रुद्यते त्वया} %||7-45-22||

\twolineshloka
{जाह्नवीतीरमासाद्य चिराभिलषितं मम}
{हर्षकाले किमर्थं मां विषादयसि लक्ष्मण} %||7-45-23||

\twolineshloka
{नित्यं त्वं रामपादेषु वर्तसे पुरुषर्षभ}
{कच्चिद्विनाकृतस्तेन द्विरात्रे शोकमागतः} %||7-45-24||

\twolineshloka
{ममापि दयितो रामो जीवितेनापि लक्ष्मण}
{न चाहमेवं शोचामि मैवं त्वं बालिशो भव} %||7-45-25||

\twolineshloka
{तारयस्व च मां गङ्गां दर्शयस्व च तापसान्}
{ततो धनानि वासांसि दास्याम्याभरणानि च} %||7-45-26||

\twolineshloka
{ततः कृत्वा महर्षीणां यथार्हमभिवादनम्}
{तत्र चैकां निशामुष्य यास्यामस्तां पुरीं पुनः} %||7-45-27||

\twolineshloka
{तस्यास्तद्वचनं श्रुत्वा प्रमृज्य नयने शुभे}
{तितीर्षुर्लक्ष्मणो गङ्गां शुभां नावमुपाहरत्} %||7-46-28||


{॥इत्यार्षे श्रीमद्रामायणे वाल्मीकीये आदिकाव्ये उत्तरकाण्डे पञ्चचत्वारिंशः सर्गः॥}

\sect{षट्चत्वारिंशः सर्गः}

\twolineshloka
{अथ नावं सुविस्तीर्णां नैषादीं राघवानुजः}
{आरुरोह समायुक्तां पूर्वमारोप्य मैथिलीम्} %||7-46-1||

\twolineshloka
{सुमन्त्रं चैव सरथं स्थीयतामिति लक्ष्मणः}
{उवाच शोकसन्तप्तः प्रयाहीति च नाविकम्} %||7-46-2||

\twolineshloka
{ततस्तीरमुपागम्य भागीरथ्याः स लक्ष्मणः}
{उवाच मैथिलीं वाक्यं प्राञ्जलिर्बाष्पगद्गदः} %||7-46-3||

\twolineshloka
{हृद्गतं मे महच्छल्यं यदस्म्यार्येण धीमता}
{अस्मिन्निमित्ते वैदेहि लोकस्य वचनीकृतः} %||7-46-4||

\twolineshloka
{श्रेयो हि मरणं मेऽद्य मृत्योर्वा यत्परं भवेत्}
{न चास्मिन्नीदृशे कार्ये नियोज्यो लोकनिन्दिते} %||7-46-5||

\twolineshloka
{प्रसीद न च मे रोषं कर्तुमर्हसि सुव्रते}
{इत्यञ्जलिकृतो भूमौ निपपात स लक्ष्मणः} %||7-46-6||

\twolineshloka
{रुदन्तं प्राञ्जलिं दृष्ट्वा काङ्क्षन्तं मृत्युमात्मनः}
{मैथिली भृशसंविग्ना लक्ष्मणं वाक्यमब्रवीत्} %||7-46-7||

\twolineshloka
{किमिदं नावगच्छामि ब्रूहि तत्त्वेन लक्ष्मण}
{पश्यामि त्वां च न स्वस्थमपि क्षेमं महीपतेः} %||7-46-8||

\twolineshloka
{शापितोऽसि नरेन्द्रेण यत्त्वं सन्तापमात्मनः}
{तद्ब्रूयाः संनिधौ मह्यमहमाज्ञापयामि ते} %||7-46-9||

\twolineshloka
{वैदेह्या चोद्यमानस्तु लक्ष्मणो दीनचेतनः}
{अवाङ्मुखो बाष्पगलो वाक्यमेतदुवाच ह} %||7-46-10||

\twolineshloka
{श्रुत्वा परिषदो मध्ये अपवादं सुदारुणम्}
{पुरे जनपदे चैव त्वत्कृते जनकात्मजे} %||7-46-11||

\twolineshloka
{न तानि वचनीयानि मया देवि तवाग्रतः}
{यानि राज्ञा हृदि न्यस्तान्यमर्षः पृष्ठतः कृतः} %||7-46-12||

\twolineshloka
{सा त्वं त्यक्ता नृपतिना निर्दोषा मम संनिधौ}
{पौरापवादभीतेन ग्राह्यं देवि न तेऽन्यथा} %||7-46-13||

\twolineshloka
{आश्रमान्तेषु च मया त्यक्तव्या त्वं भविष्यसि}
{राज्ञः शासनमाज्ञाय तवैवं किल दौर्हृदम्} %||7-46-14||

\twolineshloka
{तदेतज्जाह्नवीतीरे ब्रह्मर्षीणां तपोवनम्}
{पुण्यं च रमणीयं च मा विषादं कृथाः शुभे} %||7-46-15||

\twolineshloka
{राज्ञो दशरथस्यैष पितुर्मे मुनिपुङ्गवः}
{सखा परमको विप्रो वाल्मीकिः सुमहायशाः} %||7-46-16||

\twolineshloka
{पादच्छायामुपागम्य सुखमस्य महात्मनः}
{उपवासपरैकाग्रा वस त्वं जनकात्मजे} %||7-46-17||

\twolineshloka
{पतिव्रतात्वमास्थाय रामं कृत्वा सदा हृदि}
{श्रेयस्ते परमं देवि तथा कृत्वा भविष्यति} %||7-47-18||


{॥इत्यार्षे श्रीमद्रामायणे वाल्मीकीये आदिकाव्ये उत्तरकाण्डे षट्चत्वारिंशः सर्गः॥}

\sect{सप्तचत्वारिंशः सर्गः}

\twolineshloka
{लक्ष्मणस्य वचः श्रुत्वा दारुणं जनकात्मजा}
{परं विषादमागम्य वैदेही निपपात ह} %||7-47-1||

\twolineshloka
{सा मुहूर्तमिवासंज्ञा बाष्पव्याकुलितेक्षणा}
{लक्ष्मणं दीनया वाचा उवाच जनकात्मजा} %||7-47-2||

\twolineshloka
{मामिकेयं तनुर्नूनं सृष्टा दुःखाय लक्ष्मण}
{धात्रा यस्यास्तथा मेऽद्य दुःखमूर्तिः प्रदृश्यते} %||7-47-3||

\twolineshloka
{किं नु पापं कृतं पूर्वं को वा दारैर्वियोजितः}
{याहं शुद्धसमाचारा त्यक्ता नृपतिना सती} %||7-47-4||

\twolineshloka
{पुराहमाश्रमे वासं रामपादानुवर्तिनी}
{अनुरुध्यापि सौमित्रे दुःखे विपरिवर्तिनी} %||7-47-5||

\twolineshloka
{सा कथं ह्याश्रमे सौम्य वत्स्यामि विजनीकृता}
{आख्यास्यामि च कस्याहं दुःखं दुःखपरायणा} %||7-47-6||

\twolineshloka
{किं च वक्ष्यामि मुनिषु किं मयापकृतं नृपे}
{कस्मिन्वा कारणे त्यक्ता राघवेण महात्मना} %||7-47-7||

\twolineshloka
{न खल्वद्यैव सौमित्रे जीवितं जाह्नवीजले}
{त्यजेयं राजवंशस्तु भर्तुर्मे परिहास्यते} %||7-47-8||

\twolineshloka
{यथाज्ञां कुरु सौमित्रे त्यज मां दुःखभागिनीम्}
{निदेशे स्थीयतां राज्ञः शृणु चेदं वचो मम} %||7-47-9||

\twolineshloka
{श्वश्रूणामविशेषेण प्राञ्जलिः प्रग्रहेण च}
{शिरसा वन्द्य चरणौ कुशलं ब्रूहि पार्थिवम्} %||7-47-10||

\twolineshloka
{यथा भ्रातृषु वर्तेथास्तथा पौरेषु नित्यदा}
{परमो ह्येष धर्मः स्यादेषा कीर्तिरनुत्तमा} %||7-47-11||

\threelineshloka
{यत्त्वं पौरजनं राजन्धर्मेण समवाप्नुयाः}
{अहं तु नानुशोचामि स्वशरीरं नरर्षभ}
{यथापवादं पौराणां तथैव रघुनन्दन} %||7-47-12||

\twolineshloka
{एवं ब्रुवन्त्यां सीतायां लक्ष्मणो दीनचेतनः}
{शिरसा धरणीं गत्वा व्याहर्तुं न शशाक ह} %||7-47-13||

\twolineshloka
{प्रदक्षिणं च कृत्वा स रुदन्नेव महास्वनम्}
{आरुरोह पुनर्नावं नाविकं चाभ्यचोदयत्} %||7-47-14||

\twolineshloka
{स गत्वा चोत्तरं कूलं शोकभारसमन्वितः}
{सम्मूढ इव दुःखेन रथमध्यारुहद्द्रुतम्} %||7-47-15||

\twolineshloka
{मुहुर्मुहुरपावृत्य दृष्ट्वा सीतामनाथवत्}
{वेष्टन्तीं परतीरस्थां लक्ष्मणः प्रययावथ} %||7-47-16||

\twolineshloka
{दूरस्थं रथमालोक्य लक्ष्मणं च मुहुर्मुहुः}
{निरीक्षमाणामुद्विग्नां सीतां शोकः समाविशत्} %||7-47-17||

\fourlineindentedshloka
{सा दुःखभारावनता तपस्विनी}
{ यशोधरा नाथमपश्यती सती}
{रुरोद सा बर्हिणनादिते वने}
{ महास्वनं दुःखपरायणा सती} %||7-48-18||


{॥इत्यार्षे श्रीमद्रामायणे वाल्मीकीये आदिकाव्ये उत्तरकाण्डे सप्तचत्वारिंशः सर्गः॥}

\sect{अष्टचत्वारिंशः सर्गः}

\twolineshloka
{सीतां तु रुदतीं दृष्ट्वा ये तत्र मुनिदारकाः}
{प्राद्रवन्यत्र भगवानास्ते वाल्मीकिरग्र्यधीः} %||7-48-1||

\twolineshloka
{अभिवाद्य मुनेः पादौ मुनिपुत्रा महर्षये}
{सर्वे निवेदयामासुस्तस्यास्तु रुदितस्वनम्} %||7-48-2||

\twolineshloka
{अदृष्टपूर्वा भगवन्कस्याप्येषा महात्मनः}
{पत्नी श्रीरिव सम्मोहाद्विरौति विकृतस्वरा} %||7-48-3||

\twolineshloka
{भगवन्साधु पश्येमां देवतामिव खाच्च्युताम्}
{न ह्येनां मानुषीं विद्मः सत्क्रियास्याः प्रयुज्यताम्} %||7-48-4||

\twolineshloka
{तेषां तद्वचनं श्रुत्वा बुद्ध्या निश्चित्य धर्मवित्}
{तपसा लब्धचक्षुष्मान्प्राद्रवद्यत्र मैथिली} %||7-48-5||

\threelineshloka
{तं तु देशमभिप्रेत्य किञ्चित्पद्भ्यां महामुनिः}
{अर्घ्यमादाय रुचिरं जाह्वनीतीरमाश्रितः}
{ददर्श राघवस्येष्टां पत्नीं सीतामनाथवत्} %||7-48-6||

\twolineshloka
{तां सीतां शोकभारार्तां वाल्मीकिर्मुनिपुङ्गवः}
{उवाच मधुरां वाणीं ह्लादयन्निव तेजसा} %||7-48-7||

\twolineshloka
{स्नुषा दशरथस्य त्वं रामस्य महिषी सती}
{जनकस्य सुता राज्ञः स्वागतं ते पतिव्रते} %||7-48-8||

\twolineshloka
{आयान्त्येवासि विज्ञाता मया धर्मसमाधिना}
{कारणं चैव सर्वं मे हृदयेनोपलक्षितम्} %||7-48-9||

\twolineshloka
{अपापां वेद्मि सीते त्वां तपोलब्धेन चक्षुषा}
{विशुद्धभावा वैदेहि साम्प्रतं मयि वर्तसे} %||7-48-10||

\twolineshloka
{आश्रमस्याविदूरे मे तापस्यस्तपसि स्थिताः}
{तास्त्वां वत्से यथा वत्सं पालयिष्यन्ति नित्यशः} %||7-48-11||

\twolineshloka
{इदमर्घ्यं प्रतीच्छ त्वं विस्रब्धा विगतज्वरा}
{यथा स्वगृहमभ्येत्य विषादं चैव मा कृथाः} %||7-48-12||

\twolineshloka
{श्रुत्वा तु भाषितं सीता मुनेः परममद्भुतम्}
{शिरसा वन्द्य चरणौ तथेत्याह कृताञ्जलिः} %||7-48-13||

\twolineshloka
{तं प्रयान्तं मुनिं सीता प्राञ्जलिः पृष्ठतोऽन्वगात्}
{अन्वयाद्यत्र तापस्यो धर्मनित्याः समाहिताः} %||7-48-14||

\twolineshloka
{तं दृष्ट्वा मुनिमायान्तं वैदेह्यानुगतं तदा}
{उपाजग्मुर्मुदा युक्ता वचनं चेदमब्रुवन्} %||7-48-15||

\twolineshloka
{स्वागतं ते मुनिश्रेष्ठ चिरस्यागमनं प्रभो}
{अभिवादयामः सर्वास्त्वामुच्यतां किं च कुर्महे} %||7-48-16||

\twolineshloka
{तासां तद्वचनं श्रुत्वा वाल्मीकिरिदमब्रवीत्}
{सीतेयं समनुप्राप्ता पत्नी रामस्य धीमतः} %||7-48-17||

\twolineshloka
{स्नुषा दशरथस्यैषा जनकस्य सुता सती}
{अपापा पतिना त्यक्ता परिपाल्या मया सदा} %||7-48-18||

\twolineshloka
{इमां भवत्यः पश्यन्तु स्नेहेन परमेण ह}
{गौरवान्मम वाक्यस्य पूज्या वोऽस्तु विशेषतः} %||7-48-19||

\twolineshloka
{मुहुर्मुहुश्च वैदेहीं परिसान्त्व्य महायशाः}
{स्वमाश्रमं शिष्यवृतः पुनरायान्महातपाः} %||7-49-20||


{॥इत्यार्षे श्रीमद्रामायणे वाल्मीकीये आदिकाव्ये उत्तरकाण्डे अष्टचत्वारिंशः सर्गः॥}

\sect{एकोनपञ्चाशत्तमः सर्गः}

\twolineshloka
{दृष्ट्वा तु मैथिलीं सीतामाश्रमं सम्प्रवेशिताम्}
{सन्तापमकरोद्घोरं लक्ष्मणो दीनचेतनः} %||7-49-1||

\twolineshloka
{अब्रवीच्च महातेजाः सुमन्त्रं मन्त्रसारथिम्}
{सीतासन्तापजं दुःखं पश्य रामस्य धीमतः} %||7-49-2||

\twolineshloka
{अतो दुःखतरं किं नु राघवस्य भविष्यति}
{पत्नीं शुद्धसमाचारां विसृज्य जनकात्मजाम्} %||7-49-3||

\twolineshloka
{व्यक्तं दैवादहं मन्ये राघवस्य विनाभवम्}
{वैदेह्या सारथे सार्धं दैवं हि दुरतिक्रमम्} %||7-49-4||

\twolineshloka
{यो हि देवान्सगन्धर्वानसुरान्सह राक्षसैः}
{निहन्याद्राघवः क्रुद्धः स दैवमनुवर्तते} %||7-49-5||

\twolineshloka
{पुरा मम पितुर्वाक्यैर्दण्डके विजने वने}
{उषितो नववर्षाणि पञ्च चैव सुदारुणे} %||7-49-6||

\twolineshloka
{ततो दुःखतरं भूयः सीताया विप्रवासनम्}
{पौराणां वचनं श्रुत्वा नृशंसं प्रतिभाति मे} %||7-49-7||

\twolineshloka
{को नु धर्माश्रयः सूत कर्मण्यस्मिन्यशोहरे}
{मैथिलीं प्रति सम्प्राप्तः पौरैर्हीनार्थवादिभिः} %||7-49-8||

\twolineshloka
{एता बहुविधा वाचः श्रुत्वा लक्ष्मणभाषिताः}
{सुमन्त्रः प्राञ्जलिर्भूत्वा वाक्यमेतदुवाच ह} %||7-49-9||

\twolineshloka
{न सन्तापस्त्वया कार्यः सौमित्रे मैथिलीं प्रति}
{दृष्टमेतत्पुरा विप्रैः पितुस्ते लक्ष्मणाग्रतः} %||7-49-10||

\threelineshloka
{भविष्यति दृढं रामो दुःखप्रायोऽल्पसौख्यवान्}
{त्वां चैव मैथिलीं चैव शत्रुघ्नभरतौ तथा}
{सन्त्यजिष्यति धर्मात्मा कालेन महता महान्} %||7-49-11||

\twolineshloka
{न त्विदं त्वयि वक्तव्यं सौमित्रे भरतेऽपि वा}
{राज्ञा वोऽव्याहृतं वाक्यं दुर्वासा यदुवाच ह} %||7-49-12||

\twolineshloka
{महाराजसमीपे च मम चैव नरर्षभ}
{ऋषिणा व्याहृतं वाक्यं वसिष्ठस्य च संनिधौ} %||7-49-13||

\twolineshloka
{ऋषेस्तु वचनं श्रुत्वा मामाह पुरुषर्षभः}
{सूत न क्वचिदेवं ते वक्तव्यं जनसंनिधौ} %||7-49-14||

\twolineshloka
{तस्याहं लोकपालस्य वाक्यं तत्सुसमाहितः}
{नैव जात्वनृतं कुर्यामिति मे सौम्य दर्शनम्} %||7-49-15||

\twolineshloka
{सर्वथा नास्त्यवक्तव्यं मया सौम्य तवाग्रतः}
{यदि ते श्रवणे श्रद्धा श्रूयतां रघुनन्दन} %||7-49-16||

\twolineshloka
{यद्यप्यहं नरेन्द्रेण रहस्यं श्रावितः पुरा}
{तच्चाप्युदाहरिष्यामि दैवं हि दुरतिक्रमम्} %||7-49-17||

\twolineshloka
{तच्छ्रुत्वा भाषितं तस्य गम्भीरार्थपदं महत्}
{तथ्यं ब्रूहीति सौमित्रिः सूतं वाक्यमथाब्रवीत्} %||7-50-18||


{॥इत्यार्षे श्रीमद्रामायणे वाल्मीकीये आदिकाव्ये उत्तरकाण्डे एकोनपञ्चाशत्तमः सर्गः॥}

\sect{पञ्चाशत्तमः सर्गः}

\twolineshloka
{तथा सञ्चोदितः सूतो लक्ष्मणेन महात्मना}
{तद्वाक्यमृषिणा प्रोक्तं व्याहर्तुमुपचक्रमे} %||7-50-1||

\twolineshloka
{पुरा नाम्ना हि दुर्वासा अत्रेः पुत्रो महामुनिः}
{वसिष्ठस्याश्रमे पुण्ये स वार्षिक्यमुवास ह} %||7-50-2||

\twolineshloka
{तमाश्रमं महातेजाः पिता ते सुमहायशाः}
{पुरोधसं महात्मानं दिदृक्षुरगमत्स्वयम्} %||7-50-3||

\threelineshloka
{स दृष्ट्वा सूर्यसङ्काशं ज्वलन्तमिव तेजसा}
{उपविष्टं वसिष्ठस्य सव्ये पार्श्वे महामुनिम्}
{तौ मुनी तापसश्रेष्ठौ विनीतस्त्वभ्यवादयत्} %||7-50-4||

\twolineshloka
{स ताभ्यां पूजितो राजा स्वागतेनासनेन च}
{पाद्येन फलमूलैश्च सोऽप्यास्ते मुनिभिः सह} %||7-50-5||

\twolineshloka
{तेषां तत्रोपविष्टानां तास्ताः सुमधुराः कथाः}
{बभूवुः परमर्षीणां मध्यादित्यगतेऽहनि} %||7-50-6||

\twolineshloka
{ततः कथायां कस्याञ्चित्प्राञ्जलिः प्रग्रहो नृपः}
{उवाच तं महात्मानमत्रेः पुत्रं तपोधनम्} %||7-50-7||

\twolineshloka
{भगवन्किम्प्रमाणेन मम वंशो भविष्यति}
{किमायुश्च हि मे रामः पुत्राश्चान्ये किमायुषः} %||7-50-8||

\twolineshloka
{रामस्य च सुता ये स्युस्तेषामायुः कियद्भवेत्}
{काम्यया भगवन्ब्रूहि वंशस्यास्य गतिं मम} %||7-50-9||

\twolineshloka
{तच्छ्रुत्वा व्याहृतं वाक्यं राज्ञो दशरथस्य तु}
{दुर्वासाः सुमहातेजा व्याहर्तुमुपचक्रमे} %||7-50-10||

\twolineshloka
{अयोध्यायाः पती रामो दीर्घकालं भविष्यति}
{सुखिनश्च समृद्धाश्च भविष्यन्त्यस्य चानुजाः} %||7-50-11||

\twolineshloka
{कस्मिंश्चित्करणे त्वां च मैथिलीं च यशस्विनीम्}
{सन्त्यजिष्यति धर्मात्मा कालेन महता किल} %||7-50-12||

\twolineshloka
{दशवर्षसहस्रणि दशवर्षशतानि च}
{रामो राज्यमुपासित्वा ब्रह्मलोकं गमिष्यति} %||7-50-13||

\twolineshloka
{समृद्धैर्हयमेधैश्च इष्ट्वा परपुरंजयः}
{राजवंशांश्च काकुत्स्थो बहून्संस्थापयिष्यति} %||7-50-14||

\twolineshloka
{स सर्वमखिलं राज्ञो वंशस्यास्य गतागतम्}
{आख्याय सुमहातेजास्तूष्णीमासीन्महाद्युतिः} %||7-50-15||

\twolineshloka
{तूष्णीम्भूते मुनौ तस्मिन्राजा दशरथस्तदा}
{अभिवाद्य महात्मानौ पुनरायात्पुरोत्तमम्} %||7-50-16||

\twolineshloka
{एतद्वचो मया तत्र मुनिना व्याहृतं पुरा}
{श्रुतं हृदि च निक्षिप्तं नान्यथा तद्भविष्यति} %||7-50-17||

\twolineshloka
{एवं गते न सन्तापं गन्तुमर्हसि राघव}
{सीतार्थे राघवार्थे वा दृढो भव नरोत्तम} %||7-50-18||

\twolineshloka
{तच्छ्रुत्वा व्याहृतं वाक्यं सूतस्य परमाद्भुतम्}
{प्रहर्षमतुलं लेभे साधु साध्विति चाब्रवीत्} %||7-50-19||

\twolineshloka
{तयोः संवदतोरेवं सूतलक्ष्मणयोः पथि}
{अस्तमर्को गतो वासं गोमत्यां तावथोषतुः} %||7-51-20||


{॥इत्यार्षे श्रीमद्रामायणे वाल्मीकीये आदिकाव्ये उत्तरकाण्डे पञ्चाशत्तमः सर्गः॥}

\sect{एकपञ्चाशत्तमः सर्गः}

\twolineshloka
{तत्र तां रजनीमुष्य गोमत्यां रघुनन्दनः}
{प्रभाते पुनरुत्थाय लक्ष्मणः प्रययौ तदा} %||7-51-1||

\twolineshloka
{ततोऽर्धदिवसे प्राप्ते प्रविवेश महारथः}
{अयोध्यां रत्नसम्पूर्णां हृष्टपुष्टजनावृताम्} %||7-51-2||

\twolineshloka
{सौमित्रिस्तु परं दैन्यं जगाम सुमहामतिः}
{रामपादौ समासाद्य वक्ष्यामि किमहं गतः} %||7-51-3||

\twolineshloka
{तस्यैवं चिन्तयानस्य भवनं शशिसंनिभम्}
{रामस्य परमोदारं पुरस्तात्समदृश्यत} %||7-51-4||

\twolineshloka
{राज्ञस्तु भवनद्वारि सोऽवतीर्य नरोत्तमः}
{अवाङ्मुखो दीनमनाः प्राविवेशानिवारितः} %||7-51-5||

\twolineshloka
{स दृष्ट्वा राघवं दीनमासीनं परमासने}
{नेत्राभ्यामश्रुपूर्णाभ्यां ददर्शाग्रजमग्रतः} %||7-51-6||

\twolineshloka
{जग्राह चरणौ तस्य लक्ष्मणो दीनचेतनः}
{उवाच दीनया वाचा प्राञ्जलिः सुसमाहितः} %||7-51-7||

\threelineshloka
{आर्यस्याज्ञां पुरस्कृत्य विसृज्य जनकात्मजाम्}
{गङ्गातीरे यथोद्दिष्टे वाल्मीकेराश्रमे शुभे}
{पुनरस्म्यागतो वीर पादमूलमुपासितुम्} %||7-51-8||

\twolineshloka
{मा शुचः पुरुषव्याघ्र कालस्य गतिरीदृशी}
{त्वद्विधा न हि शोचन्ति सत्त्ववन्तो मनस्विनः} %||7-51-9||

\twolineshloka
{सर्वे क्षयान्ता निचयाः पतनान्ताः समुच्छ्रयाः}
{संयोगा विप्रयोगान्ता मरणान्तं च जीवितम्} %||7-51-10||

\twolineshloka
{शक्तस्त्वमात्मनात्मानं विजेतुं मनसैव हि}
{लोकान्सर्वांश्च काकुत्स्थ किं पुनर्दुःखमीदृशम्} %||7-51-11||

\twolineshloka
{नेदृशेषु विमुह्यन्ति त्वद्विधाः पुरुषर्षभाः}
{यदर्थं मैथिली त्यक्ता अपवादभयान्नृप} %||7-51-12||

\twolineshloka
{स त्वं पुरुषशार्दूल धैर्येण सुसमाहितः}
{त्यजेमां दुर्बलां बुद्धिं सन्तापं मा कुरुष्व ह} %||7-51-13||

\twolineshloka
{एवमुक्तस्तु काकुत्स्थो लक्ष्मणेन महात्मना}
{उवाच परया प्रीत्या सौमित्रिं मित्रवत्सलम्} %||7-51-14||

\twolineshloka
{एवमेतन्नरश्रेष्ठ यथा वदसि लक्ष्मण}
{परितोषश्च मे वीर मम कार्यानुशासने} %||7-51-15||

\twolineshloka
{निर्वृतिश्च कृता सौम्य सन्तापश्च निराकृतः}
{भवद्वाक्यैः सुमधुरैरनुनीतोऽस्मि लक्ष्मण} %||7-52-16||


{॥इत्यार्षे श्रीमद्रामायणे वाल्मीकीये आदिकाव्ये उत्तरकाण्डे एकपञ्चाशत्तमः सर्गः॥}

\sect{द्विपञ्चाशत्तमः सर्गः}

\twolineshloka
{ततः सुमन्त्रस्त्वागम्य राघवं वाक्यमब्रवीत्}
{एते निवारिता राजन्द्वारि तिष्ठन्ति तापसाः} %||7-52-1||

\threelineshloka
{भार्गवं च्यवनं नाम पुरस्कृत्य महर्षयः}
{दर्शनं ते महाराज चोदयन्ति कृतत्वराः}
{प्रीयमाणा नरव्याघ्र यमुनातीरवासिनः} %||7-52-2||

\twolineshloka
{तस्य तद्वचनं श्रुत्वा रामः प्रोवाच धर्मवित्}
{प्रवेश्यन्तां महात्मानो भार्गवप्रमुखा द्विजाः} %||7-52-3||

\twolineshloka
{राज्ञस्त्वाज्ञां पुरस्कृत्य द्वाःस्थो मूर्ध्नि कृताञ्जलिः}
{प्रवेशयामास ततस्तापसान्सम्मतान्बहून्} %||7-52-4||

\twolineshloka
{शतं समधिकं तत्र दीप्यमानं स्वतेजसा}
{प्रविष्टं राजभवनं तापसानां महात्मनाम्} %||7-52-5||

\twolineshloka
{ते द्विजाः पूर्णकलशैः सर्वतीर्थाम्बु सत्कृतम्}
{गृहीत्वा फलमूलं च रामस्याभ्याहरन्बहु} %||7-52-6||

\twolineshloka
{प्रतिगृह्य तु तत्सर्वं रामः प्रीतिपुरस्कृतः}
{तीर्थोदकानि सर्वाणि फलानि विविधानि च} %||7-52-7||

\twolineshloka
{उवाच च महाबाहुः सर्वानेव महामुनीन्}
{इमान्यासनमुख्यानि यथार्हमुपविश्यताम्} %||7-52-8||

\twolineshloka
{रामस्य भाषितं श्रुत्वा सर्व एव महर्षयः}
{बृसीषु रुचिराख्यासु निषेदुः काञ्चनीषु ते} %||7-52-9||

\twolineshloka
{उपविष्टानृषींस्तत्र दृष्ट्वा परपुरंजयः}
{प्रयतः प्राञ्जलिर्भूत्वा राघवो वाक्यमब्रवीत्} %||7-52-10||

\twolineshloka
{किमागमनकार्यं वः किं करोमि तपोधनाः}
{आज्ञाप्योऽहं महर्षीणां सर्वकामकरः सुखम्} %||7-52-11||

\twolineshloka
{इदं राज्यं च सकलं जीवितं च हृदि स्थितम्}
{सर्वमेतद्द्विजार्थं मे सत्यमेतद्ब्रवीमि वः} %||7-52-12||

\twolineshloka
{तस्य तद्वचनं श्रुत्वा साधुवादो महानभूत्}
{ऋषीणामुग्रतपसां यमुनातीरवासिनाम्} %||7-52-13||

\twolineshloka
{ऊचुश्च ते महात्मानो हर्षेण महतान्विताः}
{उपपन्नं नरश्रेष्ठ तवैव भुवि नान्यतः} %||7-52-14||

\twolineshloka
{बहवः पार्थिवा राजन्नतिक्रान्ता महाबलाः}
{कार्यगौरवमश्रुत्वा प्रतिज्ञां नाभ्यरोचयन्} %||7-52-15||

\fourlineindentedshloka
{त्वया पुनर्ब्राह्मणगौरवादियम्}
{ कृता प्रतिज्ञा ह्यनवेक्ष्य कारणम्}
{कुरुष्व कर्ता ह्यसि नात्र संशयो}
{ महाभयात्त्रातुमृषींस्त्वमर्हसि} %||7-53-16||


{॥इत्यार्षे श्रीमद्रामायणे वाल्मीकीये आदिकाव्ये उत्तरकाण्डे द्विपञ्चाशत्तमः सर्गः॥}

\sect{त्रिपञ्चाशत्तमः सर्गः}

\twolineshloka
{ब्रुवद्भिरेवमृषिभिः काकुत्स्थो वाक्यमब्रवीत्}
{किं कार्यं ब्रूत भवतां भयं नाशयितास्मि वः} %||7-53-1||

\twolineshloka
{तथा वदति काकुत्स्थे भार्गवो वाक्यमब्रवीत्}
{भयं नः शृणु यन्मूलं देशस्य च नरेश्वर} %||7-53-2||

\twolineshloka
{पूर्वं कृतयुगे राम दैतेयः सुमहाबलः}
{लोलापुत्रोऽभवज्ज्येष्ठो मधुर्नाम महासुरः} %||7-53-3||

\twolineshloka
{ब्रह्मण्यश्च शरण्यश्च बुद्ध्या च परिनिष्ठितः}
{सुरैश्च परमोदारैः प्रीतिस्तस्यातुलाभवत्} %||7-53-4||

\twolineshloka
{स मधुर्वीर्यसम्पन्नो धर्मे च सुसमाहितः}
{बहुमानाच्च रुद्रेण दत्तस्तस्याद्भुतो वरः} %||7-53-5||

\twolineshloka
{शूलं शूलाद्विनिष्कृष्य महावीर्यं महाप्रभम्}
{ददौ महात्मा सुप्रीतो वाक्यं चैतदुवाच ह} %||7-53-6||

\twolineshloka
{त्वयायमतुलो धर्मो मत्प्रसादात्कृतः शुभः}
{प्रीत्या परमया युक्तो ददाम्यायुधमुत्तमम्} %||7-53-7||

\twolineshloka
{यावत्सुरैश्च विप्रैश्च न विरुध्येर्महासुर}
{तावच्छूलं तवेदं स्यादन्यथा नाशमाप्नुयात्} %||7-53-8||

\twolineshloka
{यश्च त्वामभियुञ्जीत युद्धाय विगतज्वरः}
{तं शूलं भस्मसात्कृत्वा पुनरेष्यति ते करम्} %||7-53-9||

\twolineshloka
{एवं रुद्राद्वरं लब्ध्वा भूय एव महासुरः}
{प्रणिपत्य महादेवं वाक्यमेतदुवाच ह} %||7-53-10||

\twolineshloka
{भगवन्मम वंशस्य शूलमेतदनुत्तमम्}
{भवेत्तु सततं देव सुराणामीश्वरो ह्यसि} %||7-53-11||

\twolineshloka
{तं ब्रुवाणं मधुं देवः सर्वभूतपतिः शिवः}
{प्रत्युवाच महादेवो नैतदेवं भविष्यति} %||7-53-12||

\twolineshloka
{मा भूत्ते विफला वाणी मत्प्रसादकृता शुभा}
{भवतः पुत्रमेकं तु शूलमेतद्गमिष्यति} %||7-53-13||

\twolineshloka
{यावत्करस्थः शूलोऽयं भविष्यति सुतस्य ते}
{अवध्यः सर्वभूतानां शूलहस्तो भविष्यति} %||7-53-14||

\twolineshloka
{एवं मधुर्वरं लब्ध्वा देवात्सुमहदद्भुतम्}
{भवनं चासुरश्रेष्ठः कारयामास सुप्रभम्} %||7-53-15||

\twolineshloka
{तस्य पत्नी महाभागा प्रिया कुम्भीनसी हि या}
{विश्वावसोरपत्यं सा ह्यनलायां महाप्रभा} %||7-53-16||

\twolineshloka
{तस्याः पुत्रो महावीर्यो लवणो नाम दारुणः}
{बाल्यात्प्रभृति दुष्टात्मा पापान्येव समाचरत्} %||7-53-17||

\twolineshloka
{तं पुत्रं दुर्विनीतं तु दृष्ट्वा दुःखसमन्वितः}
{मधुः स शोकमापेदे न चैनं किञ्चिदब्रवीत्} %||7-53-18||

\twolineshloka
{स विहाय इमं लोकं प्रविष्टो वरुणालयम्}
{शूलं निवेश्य लवणे वरं तस्मै न्यवेदयत्} %||7-53-19||

\twolineshloka
{स प्रभावेन शूलस्य दौरात्म्येनात्मनस्तथा}
{सन्तापयति लोकांस्त्रीन्विशेषेण तु तापसान्} %||7-53-20||

\twolineshloka
{एवम्प्रभावो लवणः शूलं चैव तथाविधम्}
{श्रुत्वा प्रमाणं काकुत्स्थ त्वं हि नः परमा गतिः} %||7-53-21||

\twolineshloka
{बहवः पार्थिवा राम भयार्तैरृषिभिः पुरा}
{अभयं याचिता वीर त्रातारं न च विद्महे} %||7-53-22||

\threelineshloka
{ते वयं रावणं श्रुत्वा हतं सबलवाहनम्}
{त्रातारं विद्महे राम नान्यं भुवि नराधिपम्}
{तत्परित्रातुमिच्छामो लवणाद्भयपीडिताः} %||7-54-23||


{॥इत्यार्षे श्रीमद्रामायणे वाल्मीकीये आदिकाव्ये उत्तरकाण्डे त्रिपञ्चाशत्तमः सर्गः॥}

\sect{चतुःपञ्चाशत्तमः सर्गः}

\twolineshloka
{तथोक्ते तानृषीन्रामः प्रत्युवाच कृताञ्जलिः}
{किमाहारः किमाचारो लवणः क्व च वर्तते} %||7-54-1||

\twolineshloka
{राघवस्य वचः श्रुत्वा ऋषयः सर्व एव ते}
{ततो निवेदयामासुर्लवणो ववृधे यथा} %||7-54-2||

\twolineshloka
{आहारः सर्वसत्त्वानि विशेषेण च तापसाः}
{आचारो रौद्रता नित्यं वासो मधुवने सदा} %||7-54-3||

\twolineshloka
{हत्वा दशसहस्राणि सिंहव्याघ्रमृगद्विपान्}
{मानुषांश्चैव कुरुते नित्यमाहारमाह्निकम्} %||7-54-4||

\twolineshloka
{ततोऽपराणि सत्त्वानि खादते स महाबलः}
{संहारे समनुप्राप्ते व्यादितास्य इवान्तकः} %||7-54-5||

\twolineshloka
{तच्छ्रुत्वा राघवो वाक्यमुवाच स महामुनीन्}
{घातयिष्यामि तद्रक्षो व्यपगच्छतु वो भयम्} %||7-54-6||

\twolineshloka
{तथा तेषां प्रतिज्ञाय मुनीनामुग्रतेजसाम्}
{स भ्रातॄन्सहितान्सर्वानुवाच रघुनन्दनः} %||7-54-7||

\twolineshloka
{को हन्ता लवणं वीराः कस्यांशः स विधीयताम्}
{भरतस्य महाबाहोः शत्रुघ्नस्याथवा पुनः} %||7-54-8||

\twolineshloka
{राघवेणैवमुक्तस्तु भरतो वाक्यमब्रवीत्}
{अहमेनं वधिष्यामि ममांशः स विधीयताम्} %||7-54-9||

\twolineshloka
{भरतस्य वचः श्रुत्वा शौर्यवीर्यसमन्वितम्}
{लक्ष्मणावरजस्तस्थौ हित्वा सौवर्णमासनम्} %||7-54-10||

\twolineshloka
{शत्रुघ्नस्त्वब्रवीद्वाक्यं प्रणिपत्य नराधिपम्}
{कृतकर्मा महाबाहुर्मध्यमो रघुनन्दनः} %||7-54-11||

\twolineshloka
{आर्येण हि पुरा शून्या अयोध्या रक्षिता पुरी}
{सन्तापं हृदये कृत्वा आर्यस्यागमनं प्रति} %||7-54-12||

\twolineshloka
{दुःखानि च बहूनीह अनुभूतानि पार्थिव}
{शयानो दुःखशय्यासु नन्दिग्रामे महात्मना} %||7-54-13||

\threelineshloka
{फलमूलाशनो भूत्वा जटाचीरधरस्तथा}
{अनुभूयेदृशं दुःखमेष राघवनन्दनः}
{प्रेष्ये मयि स्थिते राजन्न भूयः क्लेशमाप्नुयात्} %||7-54-14||

\twolineshloka
{तथा ब्रुवति शत्रुघ्ने राघवः पुनरब्रवीत्}
{एवं भवतु काकुत्स्थ क्रियतां मम शासनम्} %||7-54-15||

\twolineshloka
{राज्ये त्वामभिषेक्ष्यामि मधोस्तु नगरे शुभे}
{निवेशय महाबाहो भरतं यद्यवेक्षसे} %||7-54-16||

\twolineshloka
{शूरस्त्वं कृतविद्यश्च समर्थः संनिवेशने}
{नगरं मधुना जुष्टं तथा जनपदाञ्शुभान्} %||7-54-17||

\twolineshloka
{यो हि वंशं समुत्पाट्य पार्थिवस्य पुनः क्षये}
{न विधत्ते नृपं तत्र नरकं स निगच्छति} %||7-54-18||

\twolineshloka
{स त्वं हत्वा मधुसुतं लवणं पापनिश्चयम्}
{राज्यं प्रशाधि धर्मेण वाक्यं मे यद्यवेक्षसे} %||7-54-19||

\twolineshloka
{उत्तरं च न वक्तव्यं शूर वाक्यान्तरे मम}
{बालेन पूर्वजस्याज्ञा कर्तव्या नात्र संशयः} %||7-54-20||

\twolineshloka
{अभिषेकं च काकुत्स्थ प्रतीच्छस्व मयोद्यतम्}
{वसिष्ठप्रमुखैर्विप्रैर्विधिमन्त्रपुरस्कृतम्} %||7-55-21||


{॥इत्यार्षे श्रीमद्रामायणे वाल्मीकीये आदिकाव्ये उत्तरकाण्डे चतुःपञ्चाशत्तमः सर्गः॥}

\sect{पञ्चपञ्चाशत्तमः सर्गः}

\twolineshloka
{एवमुक्तस्तु रामेण परां व्रीडामुपागतः}
{शत्रुघ्नो वीर्यसम्पन्नो मन्दं मन्दमुवाच ह} %||7-55-1||

\threelineshloka
{अवश्यं करणीयं च शासनं पुरुषर्षभ}
{तव चैव महाभाग शासनं दुरतिक्रमम्}
{अयं कामकरो राजंस्तवास्मि पुरुषर्षभ} %||7-55-2||

\twolineshloka
{एवमुक्ते तु शूरेण शत्रुघ्नेन महात्मना}
{उवाच रामः संहृष्टो लक्ष्मणं भरतं तथा} %||7-55-3||

\twolineshloka
{सम्भारानभिषेकस्य आनयध्वं समाहिताः}
{अद्यैव पुरुषव्याघ्रमभिषेक्ष्यामि दुर्जयम्} %||7-55-4||

\twolineshloka
{पुरोधसं च काकुत्स्थौ नैगमानृत्विजस्तथा}
{मन्त्रिणश्चैव मे सर्वानानयध्वं ममाज्ञया} %||7-55-5||

\threelineshloka
{राज्ञः शासनमाज्ञाय तथाकुर्वन्महारथाः}
{अभिषेकसमारम्भं पुरस्कृत्य पुरोधसम्}
{प्रविष्टा राजभवनं पुरन्दरगृहोपमम्} %||7-55-6||

\twolineshloka
{ततोऽभिषेको ववृधे शत्रुघ्नस्य महात्मनः}
{सम्प्रहर्षकरः श्रीमान्राघवस्य पुरस्य च} %||7-55-7||

\twolineshloka
{ततोऽभिषिक्तं शत्रुघ्नमङ्कमारोप्य राघवः}
{उवाच मधुरां वाणीं तेजस्तस्याभिपूरयन्} %||7-55-8||

\twolineshloka
{अयं शरस्त्वमोघस्ते दिव्यः परपुरंजयः}
{अनेन लवणं सौम्य हन्तासि रघुनन्दन} %||7-55-9||

\twolineshloka
{सृष्टः शरोऽयं काकुत्स्थ यदा शेते महार्णवे}
{स्वयम्भूरजितो देवो यं नापश्यन्सुरासुराः} %||7-55-10||

\threelineshloka
{अदृश्यः सर्वभूतानां तेनायं हि शरोत्तमः}
{सृष्टः क्रोधाभिभूतेन विनाशार्थं दुरात्मनोः}
{मधुकैटभयोर्वीर विघाते वर्तमानयोः} %||7-55-11||

\twolineshloka
{स्रष्टुकामेन लोकांस्त्रींस्तौ चानेन हतौ युधि}
{अनेन शरमुख्येन ततो लोकांश्चकार सः} %||7-55-12||

\twolineshloka
{नायं मया शरः पूर्वं रावणस्य वधार्थिना}
{मुक्तः शत्रुघ्न भूतानां महांस्त्रासो भवेदिति} %||7-55-13||

\twolineshloka
{यच्च तस्य महच्छूलं त्र्यम्बकेण महात्मना}
{दत्तं शत्रुविनाशाय मधोरायुधमुत्तमम्} %||7-55-14||

\twolineshloka
{तत्संनिक्षिप्य भवने पूज्यमानं पुनः पुनः}
{दिशः सर्वाः समालोक्य प्राप्नोत्याहारमात्मनः} %||7-55-15||

\twolineshloka
{यदा तु युद्धमाकाङ्क्षन्कश्चिदेनं समाह्वयेत्}
{तदा शूलं गृहीत्वा तद्भस्म रक्षः करोति तम्} %||7-55-16||

\twolineshloka
{स त्वं पुरुषशार्दूल तमायुधविवर्जितम्}
{अप्रविष्टपुरं पूर्वं द्वारि तिष्ठ धृतायुधः} %||7-55-17||

\twolineshloka
{अप्रविष्टं च भवनं युद्धाय पुरुषर्षभ}
{आह्वयेथा महाबाहो ततो हन्तासि राक्षसम्} %||7-55-18||

\twolineshloka
{अन्यथा क्रियमाणे तु अवध्यः स भविष्यति}
{यदि त्वेवं कृते वीर विनाशमुपयास्यति} %||7-55-19||

\twolineshloka
{एतत्ते सर्वमाख्यातं शूलस्य च विपर्ययम्}
{श्रीमतः शितिकण्ठस्य कृत्यं हि दुरतिक्रमम्} %||7-56-20||


{॥इत्यार्षे श्रीमद्रामायणे वाल्मीकीये आदिकाव्ये उत्तरकाण्डे पञ्चपञ्चाशत्तमः सर्गः॥}

\sect{षट्पञ्चाशत्तमः सर्गः}

\twolineshloka
{एवमुक्त्वा तु काकुत्स्थं प्रशस्य च पुनः पुनः}
{पुनरेवापरं वाक्यमुवाच रघुनन्दनः} %||7-56-1||

\twolineshloka
{इमान्यश्वसहस्राणि चत्वारि पुरुषर्षभ}
{रथानां च सहस्रे द्वे गजानां शतमेव च} %||7-56-2||

\twolineshloka
{अन्तरापणवीथ्यश्च नानापण्योपशोभिताः}
{अनुगच्छन्तु शत्रुघ्न तथैव नटनर्तकाः} %||7-56-3||

\twolineshloka
{हिरण्यस्य सुवर्णस्य अयुतं पुरुषर्षभ}
{गृहीत्वा गच्छ शत्रुघ्न पर्याप्तधनवाहनः} %||7-56-4||

\twolineshloka
{बलं च सुभृतं वीर हृष्टपुष्टमनुत्तमम्}
{सम्भाष्य सम्प्रदानेन रञ्जयस्व नरोत्तम} %||7-56-5||

\twolineshloka
{न ह्यर्थास्तत्र तिष्ठन्ति न दारा न च बान्धवाः}
{सुप्रीतो भृत्यवर्गस्तु यत्र तिष्ठति राघव} %||7-56-6||

\twolineshloka
{अतो हृष्टजनाकीर्णां प्रस्थाप्य महतीं चमूम्}
{एक एव धनुष्पानिस्तद्गच्छ त्वं मधोर्वनम्} %||7-56-7||

\twolineshloka
{यथा त्वां न प्रजानाति गच्छन्तं युद्धकाङ्क्षिणम्}
{लवणस्तु मधोः पुत्रस्तथा गच्छेरशङ्कितः} %||7-56-8||

\twolineshloka
{न तस्य मृत्युरन्योऽस्ति कश्चिद्धि पुरुषर्षभ}
{दर्शनं योऽभिगच्छेत स वध्यो लवणेन हि} %||7-56-9||

\twolineshloka
{स ग्रीष्मे व्यपयाते तु वर्षरात्र उपस्थिते}
{हन्यास्त्वं लवणं सौम्य स हि कालोऽस्य दुर्मतेः} %||7-56-10||

\twolineshloka
{महर्षींस्तु पुरस्कृत्य प्रयान्तु तव सैनिकाः}
{यथा ग्रीष्मावशेषेण तरेयुर्जाह्नवीजलम्} %||7-56-11||

\twolineshloka
{ततः स्थाप्य बलं सर्वं नदीतीरे समाहितः}
{अग्रतो धनुषा सार्धं गच्छ त्वं लघुविक्रम} %||7-56-12||

\twolineshloka
{एवमुक्तस्तु रामेण शत्रुघ्नस्तान्महाबलान्}
{सेनामुख्यान्समानीय ततो वाक्यमुवाच ह} %||7-56-13||

\twolineshloka
{एते वो गणिता वासा यत्र यत्र निवत्स्यथ}
{स्थातव्यं चाविरोधेन यथा बाधा न कस्यचित्} %||7-56-14||

\twolineshloka
{तथा तांस्तु समाज्ञाप्य निर्याप्य च महद्बलम्}
{कौसल्यां च सुमित्रां च कैकेयीं चाभ्यवादयत्} %||7-56-15||

\twolineshloka
{रामं प्रदक्षिणं कृत्वा शिरसाभिप्रणम्य च}
{राणेण चाभ्यनुज्ञातः शत्रुघ्नः शत्रुतापनः} %||7-56-16||

\threelineshloka
{लक्ष्मणं भरतं चैव प्रणिपत्य कृताञ्जलिः}
{पुरोधसं वसिष्ठं च शत्रुघ्नः प्रयतात्मवान्}
{प्रदक्षिणमथो कृत्वा निर्जगाम महाबलः} %||7-57-17||


{॥इत्यार्षे श्रीमद्रामायणे वाल्मीकीये आदिकाव्ये उत्तरकाण्डे षट्पञ्चाशत्तमः सर्गः॥}

\sect{सप्तपञ्चाशत्तमः सर्गः}

\twolineshloka
{प्रस्थाप्य तद्बलं सर्वं मासमात्रोषितः पथि}
{एक एवाशु शत्रुघ्नो जगाम त्वरितस्तदा} %||7-57-1||

\twolineshloka
{द्विरात्रमन्तरे शूर उष्य राघवनन्दनः}
{वाल्मीकेराश्रमं पुण्यमगच्छद्वासमुत्तमम्} %||7-57-2||

\twolineshloka
{सोऽभिवाद्य महात्मानं वाल्मीकिं मुनिसत्तमम्}
{कृताञ्जलिरथो भूत्वा वाक्यमेतदुवाच ह} %||7-57-3||

\twolineshloka
{भगवन्वस्तुमिच्छामि गुरोः कृत्यादिहागतः}
{श्वः प्रभाते गमिष्यामि प्रतीचीं वारुणीं दिशम्} %||7-57-4||

\twolineshloka
{शत्रुघ्नस्य वचः श्रुत्वा प्रहस्य मुनिपुङ्गवः}
{प्रत्युवाच महात्मानं स्वागतं ते महायशः} %||7-57-5||

\twolineshloka
{स्वमाश्रममिदं सौम्य राघवाणां कुलस्य ह}
{आसनं पाद्यमर्घ्यं च निर्विशङ्कः प्रतीच्छ मे} %||7-57-6||

\twolineshloka
{प्रतिगृह्य ततः पूजां फलमूलं च भोजनम्}
{भक्षयामास काकुत्स्थस्तृप्तिं च परमां गतः} %||7-57-7||

\twolineshloka
{स तु भुक्त्वा महाबाहुर्महर्षिं तमुवाच ह}
{पूर्वं यज्ञविभूतीयं कस्याश्रमसमीपतः} %||7-57-8||

\twolineshloka
{तस्य तद्भाषितं श्रुत्वा वाल्मीकिर्वाक्यमब्रवीत्}
{शत्रुघ्न शृणु यस्येदं बभूवायतनं पुरा} %||7-57-9||

\twolineshloka
{युष्माकं पूर्वको राजा सुदासस्य महात्मनः}
{पुत्रो मित्रसहो नाम वीर्यवानतिधार्मिकः} %||7-57-10||

\twolineshloka
{स बाल एव सौदासो मृगयामुपचक्रमे}
{चञ्चूर्यमाणं ददृशे स शूरो राक्षसद्वयम्} %||7-57-11||

\twolineshloka
{शार्दूलरूपिणौ घोरौ मृगान्बहुसहस्रशः}
{भक्षयाणावसन्तुष्टौ पर्याप्तिं च न जग्मतुः} %||7-57-12||

\twolineshloka
{स तु तौ राक्षसौ दृष्ट्वा निर्मृगं च वनं कृतम्}
{क्रोधेन महताविष्टो जघानैकं महेषुणा} %||7-57-13||

\twolineshloka
{विनिपात्य तमेकं तु सौदासः पुरुषर्षभः}
{विज्वरो विगतामर्षो हतं रक्षोऽभ्यवैक्षत} %||7-57-14||

\twolineshloka
{निरीक्षमाणं तं दृष्ट्वा सहायस्तस्य रक्षसः}
{सन्तापमकरोद्घोरं सौदासं चेदमब्रवीत्} %||7-57-15||

\twolineshloka
{यस्मादनपराद्धं त्वं सहायं मम जघ्निवान्}
{तस्मात्तवापि पापिष्ठ प्रदास्यामि प्रतिक्रियाम्} %||7-57-16||

\twolineshloka
{एवमुक्त्वा तु तं रक्षस्तत्रैवान्तरधीयत}
{कालपर्याययोगेन राजा मित्रसहोऽभवत्} %||7-57-17||

\twolineshloka
{राजापि यजते यज्ञं तस्याश्रमसमीपतः}
{अश्वमेधं महायज्ञं तं वसिष्ठोऽभ्यपालयत्} %||7-57-18||

\twolineshloka
{तत्र यज्ञो महानासीद्बहुवर्षगणायुतान्}
{समृद्धः परया लक्ष्म्या देवयज्ञसमोऽभवत्} %||7-57-19||

\twolineshloka
{अथावसाने यज्ञस्य पूर्ववैरमनुस्मरन्}
{वसिष्ठरूपी राजानमिति होवाच राक्षसः} %||7-57-20||

\twolineshloka
{अद्य यज्ञावसानान्ते सामिषं भोजनं मम}
{दीयतामिति शीघ्रं वै नात्र कार्या विचारणा} %||7-57-21||

\twolineshloka
{तच्छ्रुत्वा व्याहृतं वाक्यं रक्षसा कामरूपिणा}
{भक्षसंस्कारकुशलमुवाच पृथिवीपतिः} %||7-57-22||

\twolineshloka
{हविष्यं सामिषं स्वादु यथा भवति भोजनम्}
{तथा कुरुष्व शीघ्रं वै परितुष्येद्यथा गुरुः} %||7-57-23||

\twolineshloka
{शासनात्पार्थिवेन्द्रस्य सूदः सम्भ्रान्तमानसः}
{स च रक्षः पुनस्तत्र सूदवेषमथाकरोत्} %||7-57-24||

\twolineshloka
{स मानुषमथो मांसं पार्थिवाय न्यवेदयत्}
{इदं स्वादुहविष्यं च सामिषं चान्नमाहृतम्} %||7-57-25||

\twolineshloka
{स भोजनं वसिष्ठाय पत्न्या सार्धमुपाहरत्}
{मदयन्त्या नरव्याघ्र सामिषं रक्षसा हृतम्} %||7-57-26||

\twolineshloka
{ज्ञात्वा तदामिषं विप्रो मानुषं भोजनाहृतम्}
{क्रोधेन महताविष्टो व्याहर्तुमुपचक्रमे} %||7-57-27||

\twolineshloka
{यस्मात्त्वं भोजनं राजन्ममैतद्दातुमिच्छसि}
{तस्माद्भोजनमेतत्ते भविष्यति न संशयः} %||7-57-28||

\twolineshloka
{स राजा सह पत्न्या वै प्रणिपत्य मुहुर्मुहुः}
{पुनर्वसिष्ठं प्रोवाच यदुक्तं ब्रह्मरूपिणा} %||7-57-29||

\twolineshloka
{तच्छ्रुता पार्थिवेन्द्रस्य रक्षसा विकृतं च तत्}
{पुनः प्रोवाच राजानं वसिष्ठः पुरुषर्षभम्} %||7-57-30||

\twolineshloka
{मया रोषपरीतेन यदिदं व्याहृतं वचः}
{नैतच्छक्यं वृथा कर्तुं प्रदास्यामि च ते वरम्} %||7-57-31||

\twolineshloka
{कालो द्वादश वर्षाणि शापस्यास्य भविष्यति}
{मत्प्रसादाच्च राजेन्द्र अतीतं न स्मरिष्यसि} %||7-57-32||

\twolineshloka
{एवं स राजा तं शापमुपभुज्यारिमर्दनः}
{प्रतिलेभे पुना राज्यं प्रजाश्चैवान्वपालयत्} %||7-57-33||

\twolineshloka
{तस्य कल्माषपादस्य यज्ञस्यायतनं शुभम्}
{आश्रमस्य समीपेऽस्मिन्यस्मिन्पृच्छसि राघव} %||7-57-34||

\twolineshloka
{तस्य तां पार्थिवेन्द्रस्य कथां श्रुत्वा सुदारुणाम्}
{विवेश पर्णशालायां महर्षिमभिवाद्य च} %||7-58-35||


{॥इत्यार्षे श्रीमद्रामायणे वाल्मीकीये आदिकाव्ये उत्तरकाण्डे सप्तपञ्चाशत्तमः सर्गः॥}

\sect{अष्टपञ्चाशत्तमः सर्गः}

\twolineshloka
{यामेव रात्रिं शत्रुघ्न पर्णशालां समाविशत्}
{तामेव रात्रिं सीतापि प्रसूता दारकद्वयम्} %||7-58-1||

\threelineshloka
{ततोऽर्धरात्रसमये बालका मुनिदारकाः}
{वाल्मीकेः प्रियमाचख्युः सीतायाः प्रसवं शुभम्}
{तस्य रक्षां महातेजः कुरु भूतविनाशिनीम्} %||7-58-2||

\twolineshloka
{तेषां तद्वचनं श्रुत्वा मुनिर्हर्षमुपागमत्}
{भूतघ्नीं चाकरोत्ताभ्यां रक्षां रक्षोविनाशिनीम्} %||7-58-3||

\twolineshloka
{कुशमुष्टिमुपादाय लवं चैव तु स द्विजः}
{वाल्मीकिः प्रददौ ताभ्यां रक्षां भूतविनाशिनीम्} %||7-58-4||

\twolineshloka
{यस्तयोः पूर्वजो जातः स कुशैर्मन्त्रसंस्कृतैः}
{निर्मार्जनीयस्तु भवेत्कुश इत्यस्य नामतः} %||7-58-5||

\twolineshloka
{यश्चापरो भवेत्ताभ्यां लवेन सुसमाहितः}
{निर्मार्जनीयो वृद्धाभिर्लवश्चेति स नामतः} %||7-58-6||

\twolineshloka
{एवं कुशलवौ नाम्ना तावुभौ यमजातकौ}
{मत्कृतभ्यां च नामभ्यां ख्यातियुक्तौ भविष्यतः} %||7-58-7||

\twolineshloka
{ते रक्षां जगृहुस्तां च मुनिहस्तात्समाहिताः}
{अकुर्वंश्च ततो रक्षां तयोर्विगतकल्मषाः} %||7-58-8||

\twolineshloka
{तथा तां क्रियमाणां तु रक्षां गोत्रं च नाम च}
{सङ्कीर्तनं च रामस्य सीतायाः प्रसवौ शुभौ} %||7-58-9||

\twolineshloka
{अर्धरात्रे तु शत्रुघ्नः शुश्राव सुमहत्प्रियम्}
{पर्णशालां गतो रात्रौ दिष्ट्या दिष्ट्येति चाब्रवीत्} %||7-58-10||

\twolineshloka
{तथ तस्य प्रहृष्टस्य शत्रुघ्नस्य महात्मनः}
{व्यतीता वार्षिकी रात्रिः श्रावणी लघुविक्रमा} %||7-58-11||

\twolineshloka
{प्रभाते तु महावीर्यः कृत्वा पौर्वाह्णिकं क्रमम्}
{मुनिं प्राञ्जलिरामन्त्र्य प्रायात्पश्चान्मुखः पुनः} %||7-58-12||

\twolineshloka
{स गत्वा यमुनातीरं सप्तरात्रोषितः पथि}
{ऋषीणां पुण्यकीर्तीनामाश्रमे वासमभ्ययात्} %||7-58-13||

\twolineshloka
{स तत्र मुनिभिः सार्धं भार्गवप्रमुखैर्नृपः}
{कथाभिर्बहुरूपाभिर्वासं चक्रे महायशाः} %||7-59-14||


{॥इत्यार्षे श्रीमद्रामायणे वाल्मीकीये आदिकाव्ये उत्तरकाण्डे अष्टपञ्चाशत्तमः सर्गः॥}

\sect{एकोनषष्टितमः सर्गः}

\twolineshloka
{अथ रात्र्यां प्रवृत्तायां शत्रुघ्नो भृगुनन्दनम्}
{पप्रच्छ च्यवनं विप्रं लवणस्य बलाबलम्} %||7-59-1||

\twolineshloka
{शूलस्य च बलं ब्रह्मन्के च पूर्वं निपातिताः}
{अनेन शूलमुखेन द्वन्द्वयुद्धमुपागताः} %||7-59-2||

\twolineshloka
{तस्य तद्भाषितं श्रुत्वा शत्रुघ्नस्य महात्मनः}
{प्रत्युवाच महातेजाश्च्यवनो रघुनन्दनम्} %||7-59-3||

\twolineshloka
{असङ्ख्येयानि कर्माणि यान्यस्य पुरुषर्षभ}
{इक्ष्वाकुवंशप्रभवे यद्वृत्तं तच्छृणुष्व मे} %||7-59-4||

\twolineshloka
{अयोध्यायां पुरा राजा युवनाश्वसुतो बली}
{मान्धाता इति विख्यातस्त्रिषु लोकेषु वीर्यवान्} %||7-59-5||

\twolineshloka
{स कृत्वा पृथिवीं कृत्स्नां शासने पृथिवीपतिः}
{सुरलोकमथो जेतुमुद्योगमकरोन्नृपः} %||7-59-6||

\twolineshloka
{इन्द्रस्य तु भयं तीव्रं सुराणां च महात्मनाम्}
{मान्धातरि कृतोद्योगे देवलोकजिगीषया} %||7-59-7||

\twolineshloka
{अर्धासनेन शक्रस्य राज्यार्धेन च पार्थिवः}
{वन्द्यमानः सुरगणैः प्रतिज्ञामध्यरोहत} %||7-59-8||

\twolineshloka
{तस्य पापमभिप्रायं विदित्वा पाकशासनः}
{सान्त्वपूर्वमिदं वाक्यमुवाच युवनाश्वजम्} %||7-59-9||

\twolineshloka
{राजा त्वं मानुषे लोके न तावत्पुरुषर्षभ}
{अकृत्वा पृथिवीं वश्यां देवराज्यमिहेच्छसि} %||7-59-10||

\twolineshloka
{यदि वीर समग्रा ते मेदिनी निखिला वशे}
{देवराज्यं कुरुष्वेह सभृत्यबलवाहनः} %||7-59-11||

\twolineshloka
{इन्द्रमेवं ब्रुवाणं तु मान्धाता वाक्यमब्रवीत्}
{क्व मे शक्र प्रतिहतं शासनं पृथिवीतले} %||7-59-12||

\twolineshloka
{तमुवाच सहस्राक्षो लवणो नाम राक्षसः}
{मधुपुत्रो मधुवने नाज्ञां ते कुरुतेऽनघ} %||7-59-13||

\twolineshloka
{तच्छ्रुत्वा विप्रियं घोरं सहस्राक्षेण भाषितम्}
{व्रीडितोऽवाङ्मुखो राजा व्याहर्तुं न शशाक ह} %||7-59-14||

\twolineshloka
{आमन्त्र्य तु सहस्राक्षं ह्रिया किञ्चिदवाङ्मुखः}
{पुनरेवागमच्छ्रीमानिमं लोकं नरेश्वरः} %||7-59-15||

\twolineshloka
{स कृत्वा हृदयेऽमर्षं सभृत्यबलवाहनः}
{आजगाम मधोः पुत्रं वशे कर्तुमनिन्दितः} %||7-59-16||

\twolineshloka
{स काङ्क्षमाणो लवणं युद्धाय पुरुषर्षभः}
{दूतं सम्प्रेषयामास सकाशं लवणस्य सः} %||7-59-17||

\twolineshloka
{स गत्वा विप्रियाण्याह बहूनि मधुनः सुतम्}
{वदन्तमेवं तं दूतं भक्षयामास राक्षसः} %||7-59-18||

\twolineshloka
{चिरायमाणे दूते तु राजा क्रोधसमन्वितः}
{अर्दयामास तद्रक्षः शरवृष्ट्या समन्ततः} %||7-59-19||

\twolineshloka
{ततः प्रहस्य लवणः शूलं जग्राह पाणिना}
{वधाय सानुबन्धस्य मुमोचायुधमुत्तमम्} %||7-59-20||

\twolineshloka
{तच्छूलं दीप्यमानं तु सभृत्यबलवाहनम्}
{भस्मीकृत्य नृपं भूयो लवणस्यागमत्करम्} %||7-59-21||

\twolineshloka
{एवं स राजा सुमहान्हतः सबलवाहनः}
{शूलस्य च बलं वीर अप्रमेयमनुत्तमम्} %||7-59-22||

\twolineshloka
{श्वः प्रभाते तु लवणं वधिष्यसि न संशयः}
{अगृहीतायुधं क्षिप्रं ध्रुवो हि विजयस्तव} %||7-60-23||


{॥इत्यार्षे श्रीमद्रामायणे वाल्मीकीये आदिकाव्ये उत्तरकाण्डे एकोनषष्टितमः सर्गः॥}

\sect{षष्टितमः सर्गः}

\twolineshloka
{कथां कथयतां तेषां जयं चाकाङ्क्षतां शुभम्}
{व्यतीता रजनी शीघ्रं शत्रुघ्नस्य महात्मनः} %||7-60-1||

\twolineshloka
{ततः प्रभाते विमले तस्मिन्काले स राक्षसः}
{निर्गतस्तु पुराद्वीरो भक्षाहारप्रचोदितः} %||7-60-2||

\twolineshloka
{एतस्मिन्नन्तरे शूरः शत्रुघ्नो यमुनां नदीम्}
{तीर्त्वा मधुपुरद्वारि धनुष्पाणिरतिष्ठत} %||7-60-3||

\twolineshloka
{ततोऽर्धदिवसे प्राप्ते क्रूरकर्मा स राक्षसः}
{आगच्छद्बहुसहस्रं प्राणिनामुद्वहन्भरम्} %||7-60-4||

\twolineshloka
{ततो ददर्श शत्रुघ्नं स्थितं द्वारि धृतायुधम्}
{तमुवाच ततो रक्षः किमनेन करिष्यसि} %||7-60-5||

\twolineshloka
{ईदृशानां सहस्राणि सायुधानां नराधम}
{भक्षितानि मया रोषात्कालमाकाङ्क्षसे नु किम्} %||7-60-6||

\twolineshloka
{आहारश्चाप्यसम्पूर्णो ममायं पुरुषाधम}
{स्वयं प्रविष्टो नु मुखं कथमासाद्य दुर्मते} %||7-60-7||

\twolineshloka
{तस्यैवं भाषमाणस्य हसतश्च मुहुर्मुहुः}
{शत्रुघ्नो वीर्यसम्पन्नो रोषादश्रूण्यवर्तयत्} %||7-60-8||

\twolineshloka
{तस्य रोषाभिभूतस्य शत्रुघ्नस्य महात्मनः}
{तेजोमया मरीच्यस्तु सर्वगात्रैर्विनिष्पतन्} %||7-60-9||

\twolineshloka
{उवाच च सुसङ्क्रुद्धः शत्रुघ्नस्तं निशाचरम्}
{योद्धुमिच्छामि दुर्बुद्धे द्वन्द्वयुद्धं त्वया सह} %||7-60-10||

\twolineshloka
{पुत्रो दशरथस्याहं भ्राता रामस्य धीमतः}
{शत्रुघ्नो नाम शत्रुघ्नो वधाकाङ्क्षी तवागतः} %||7-60-11||

\twolineshloka
{तस्य मे युद्धकामस्य द्वन्द्वयुद्धं प्रदीयताम्}
{शत्रुस्त्वं सर्वजीवानां न मे जीवन्गमिष्यसि} %||7-60-12||

\twolineshloka
{तस्मिंस्तथा ब्रुवाणे तु राक्षसः प्रहसन्निव}
{प्रत्युवाच नरश्रेष्ठं दिष्ट्या प्राप्तोऽसि दुर्मते} %||7-60-13||

\twolineshloka
{मम मातृष्वसुर्भ्राता रावणो नाम राक्षसः}
{हतो रामेण दुर्बुद्धे स्त्रीहेतोः पुरुषाधम} %||7-60-14||

\twolineshloka
{तच्च सर्वं मया क्षान्तं रावणस्य कुलक्षयम्}
{अवज्ञां पुरतः कृत्वा मया यूयं विशेषतः} %||7-60-15||

\twolineshloka
{न हताश्च हि मे सर्वे परिभूतास्तृणं यथा}
{भूताश्चैव भविष्याश्च यूयं च पुरुषाधमाः} %||7-60-16||

\twolineshloka
{तस्य ते युद्धकामस्य युद्धं दास्यामि दुर्मते}
{ईप्सितं यादृशं तुभ्यं सज्जये यावदायुधम्} %||7-60-17||

\twolineshloka
{तमुवाचाथ शत्रुघ्नः क्व मे जीवन्गमिष्यसि}
{दुर्बलोऽप्यागतः शत्रुर्न मोक्तव्यः कृतात्मना} %||7-60-18||

\twolineshloka
{यो हि विक्लवया बुद्ध्या प्रसरं शत्रवे ददौ}
{स हतो मन्दबुद्धित्वाद्यथा कापुरुषस्तथा} %||7-61-19||


{॥इत्यार्षे श्रीमद्रामायणे वाल्मीकीये आदिकाव्ये उत्तरकाण्डे षष्टितमः सर्गः॥}

\sect{एकषष्टितमः सर्गः}

\twolineshloka
{तच्छ्रुत्वा भाषितं तस्य शत्रुघ्नस्य महात्मनः}
{क्रोधमाहारयत्तीव्रं तिष्ठ तिष्ठेति चाब्रवीत्} %||7-61-1||

\twolineshloka
{पाणौ पाणिं विनिष्पिष्य दन्तान्कटकटाय्य च}
{लवणो रघुशार्दूलमाह्वयामास चासकृत्} %||7-61-2||

\twolineshloka
{तं ब्रुवाणं तथा वाक्यं लवणं घोरविक्रमम्}
{शत्रुघ्नो देव शत्रुघ्न इदं वचनमब्रवीत्} %||7-61-3||

\twolineshloka
{शत्रुघ्नो न तदा जातो यदान्ये निर्जितास्त्वया}
{तदद्य बाणाभिहतो व्रज तं यमसादनम्} %||7-61-4||

\twolineshloka
{ऋषयोऽप्यद्य पापात्मन्मया त्वां निहतं रणे}
{पश्यन्तु विप्रा विद्वांसस्त्रिदशा इव रावणम्} %||7-61-5||

\twolineshloka
{त्वयि मद्बाणनिर्दग्धे पतितेऽद्य निशाचर}
{पुरं जनपदं चापि क्षेममेतद्भविष्यति} %||7-61-6||

\twolineshloka
{अद्य मद्बाहुनिष्क्रान्तः शरो वज्रनिभाननः}
{प्रवेक्ष्यते ते हृदयं पद्ममंशुरिवार्कजः} %||7-61-7||

\twolineshloka
{एवमुक्तो महावृक्षं लवणः क्रोधमूर्छितः}
{शत्रुघ्नोरसि चिक्षेप तं शूरः शतधाच्छिनत्} %||7-61-8||

\twolineshloka
{तद्दृष्ट्वा विफलं कर्म राक्षसः पुनरेव तु}
{पादपान्सुबहून्गृह्य शत्रुघ्ने व्यसृजद्बली} %||7-61-9||

\twolineshloka
{शत्रुघ्नश्चापि तेजस्वी वृक्षानापततो बहून्}
{त्रिभिश्चतुर्भिरेकैकं चिच्छेद नतपर्वभिः} %||7-61-10||

\twolineshloka
{ततो बाणमयं वर्षं व्यसृजद्राक्षसोरसि}
{शत्रुघ्नो वीर्यसम्पन्नो विव्यथे न च राक्षसः} %||7-61-11||

\twolineshloka
{ततः प्रहस्य लवणो वृक्षमुत्पाट्य लीलया}
{शिरस्यभ्यहनच्छूरं स्रस्ताङ्गः स मुमोह वै} %||7-61-12||

\twolineshloka
{तस्मिन्निपतिते वीरे हाहाकारो महानभूत्}
{ऋषीणां देव सङ्घानां गन्धर्वाप्सरसामपि} %||7-61-13||

\twolineshloka
{तमवज्ञाय तु हतं शत्रुघ्नं भुवि पातितम्}
{रक्षो लब्धान्तरमपि न विवेश स्वमालयम्} %||7-61-14||

\twolineshloka
{नापि शूलं प्रजग्राह तं दृष्ट्वा भुवि पातितम्}
{ततो हत इति ज्ञात्वा तान्भक्षान्समुदावहत्} %||7-61-15||

\twolineshloka
{मुहूर्ताल्लब्धसंज्ञस्तु पुनस्तस्थौ धृतायुधः}
{शत्रुघ्नो राक्षसद्वारि ऋषिभिः सम्प्रपूजितः} %||7-61-16||

\twolineshloka
{ततो दिव्यममोघं तं जग्राह शरमुत्तमम्}
{ज्वलन्तं तेजसा घोरं पूरयन्तं दिशो दश} %||7-61-17||

\twolineshloka
{वज्राननं वज्रवेगं मेरुमन्दर गौरवम्}
{नतं पर्वसु सर्वेषु संयुगेष्वपराजितम्} %||7-61-18||

\twolineshloka
{असृक्चन्दनदिग्धाङ्गं चारुपत्रं पतत्रिणम्}
{दानवेन्द्राचलेन्द्राणामसुराणां च दारुणम्} %||7-61-19||

\twolineshloka
{तं दीप्तमिव कालाग्निं युगान्ते समुपस्थिते}
{दृष्ट्वा सर्वाणि भूतानि परित्रासमुपागमन्} %||7-61-20||

\twolineshloka
{सदेवासुरगन्धर्वं समुनिं साप्सरोगणम्}
{जगद्धि सर्वमस्वस्थं पितामहमुपस्थितम्} %||7-61-21||

\twolineshloka
{ऊचुश्च देवदेवेशं वरदं प्रपितामहम्}
{कच्चिल्लोकक्षयो देव प्राप्तो वा युगसङ्कयः} %||7-61-22||

\twolineshloka
{नेदृशं दृष्टपूर्वं न श्रुतं वा प्रपितामह}
{देवानां भयसम्मोहो लोकानां सङ्क्षयः प्रभो} %||7-61-23||

\twolineshloka
{तेषां तद्वचनं श्रुत्वा ब्रह्मा लोकपितामनः}
{भयकारणमाचष्टे देवानामभयङ्करः} %||7-61-24||

\twolineshloka
{वधाय लवणस्याजौ शरः शत्रुघ्नधारितः}
{तेजसा यस्य सर्वे स्म सम्मूढाः सुरसत्तमाः} %||7-61-25||

\twolineshloka
{एषो हि पूर्वं देवस्य लोककर्तुः सनातनः}
{शरस्तेजोमयो वत्सा येन वै भयमागतम्} %||7-61-26||

\twolineshloka
{एष वै कैटभस्यार्थे मधुनश्च महाशरः}
{सृष्टो महात्मना तेन वधार्थं दैत्ययोस्तयोः} %||7-61-27||

\twolineshloka
{एवमेतं प्रजानीध्वं विष्णोस्तेजोमयं शरम्}
{एषा चैव तनुः पूर्वा विष्णोस्तस्य महात्मनः} %||7-61-28||

\twolineshloka
{इतो गच्छता पश्यध्वं वध्यमानं महात्मना}
{रामानुजेन वीरेण लवणं राक्षसोत्तमम्} %||7-61-29||

\twolineshloka
{तस्य ते देवदेवस्य निशम्य मधुरां गिरम्}
{आजग्मुर्यत्र युध्येते शत्रुघ्नलवणावुभौ} %||7-61-30||

\twolineshloka
{तं शरं दिव्यसङ्काशं शत्रुघ्नकरधारितम्}
{ददृशुः सर्वभूतानि युगान्ताग्निमिवोत्थितम्} %||7-61-31||

\twolineshloka
{आकाशमावृतं दृष्ट्वा देवैर्हि रघुनन्दनः}
{सिंहनादं मुहुः कृत्वा ददर्श लवणं पुनः} %||7-61-32||

\twolineshloka
{आहूतश्च ततस्तेन शत्रुघ्नेन महात्मना}
{लवणः क्रोधसंयुक्तो युद्धाय समुपस्थितः} %||7-61-33||

\threelineshloka
{आकर्णात्स विकृष्याथ तद्धनुर्धन्विनां वरः}
{स मुमोच महाबाणं लवणस्य महोरसि}
{उरस्तस्य विदार्याशु प्रविवेश रसातलम्} %||7-61-34||

\twolineshloka
{गत्वा रसातलं दिव्यं शरो विबुधपूजितः}
{पुनरेवागमत्तूर्णमिक्ष्वाकुकुलनन्दनम्} %||7-61-35||

\twolineshloka
{शत्रुघ्नशरनिर्भिन्नो लवणः स निशाचरः}
{पपात सहसा भूमौ वज्राहत इवाचलः} %||7-61-36||

\twolineshloka
{तच्च दिव्यं महच्छूलं हते लवणराक्षसे}
{पश्यतां सर्वभूतानां रुद्रस्य वशमन्वगात्} %||7-61-37||

\fourlineindentedshloka
{एकेषुपातेन भयं निहत्य}
{ लोकत्रयस्यास्य रघुप्रवीरः}
{विनिर्बभावुद्यतचापबाण}
{स्तमः प्रणुद्येव सहस्ररश्मिः} %||7-62-38||


{॥इत्यार्षे श्रीमद्रामायणे वाल्मीकीये आदिकाव्ये उत्तरकाण्डे एकषष्टितमः सर्गः॥}

\sect{द्विषष्टितमः सर्गः}

\twolineshloka
{हते तु लवणे देवाः सेन्द्राः साग्निपुरोगमाः}
{ऊचुः सुमधुरां वाणीं शत्रुघ्नां शत्रुतापनम्} %||7-62-1||

\twolineshloka
{दिष्ट्या ते विजयो वत्स दिष्ट्य लवणराक्षसः}
{हतः पुरुषशार्दूलवरं वरय राघव} %||7-62-2||

\twolineshloka
{वरदाः स्म महाबाहो सर्व एव समागताः}
{विजयाकाङ्क्षिणस्तुभ्यममोघं दर्शनं हि नः} %||7-62-3||

\twolineshloka
{देवानां भाषितं श्रुत्वा शूरो मूर्ध्नि कृताञ्जलिः}
{प्रत्युवाच महाबाहुः शत्रुघ्नः प्रयतात्मवान्} %||7-62-4||

\twolineshloka
{इमां मधुपुरीं रम्यां मधुरां देव निर्मिताम्}
{निवेशं प्रप्नुयां शीघ्रमेष मेऽस्तु वरो मतः} %||7-62-5||

\twolineshloka
{तं देवाः प्रीतमनसो बाढमित्येव राघवम्}
{भविष्यति पुरी रम्या शूरसेना न संशयः} %||7-62-6||

\twolineshloka
{ते तथोक्त्वा महात्मानो दिवमारुरुहुस्तदा}
{शत्रुघ्नोऽपि महातेजास्तां सेनां समुपानयत्} %||7-62-7||

\twolineshloka
{सा सेन शीघ्रमागच्छच्छ्रुत्वा शत्रुघ्नशासनम्}
{निवेशनं च शत्रुघ्नः शासनेन समारभत्} %||7-62-8||

\twolineshloka
{सा पुरी दिव्यसङ्काशा वर्षे द्वादशमे शुभा}
{निविष्टा शूरसेनानां विषयश्चाकुतोभयः} %||7-62-9||

\twolineshloka
{क्षेत्राणि सस्य युक्तानि काले वर्षति वासवः}
{अरोगा वीरपुरुषा शत्रुघ्नभुजपालिता} %||7-62-10||

\twolineshloka
{अर्धचन्द्रप्रतीकाशा यमुनातीरशोभिता}
{शोभिता गृहमुख्यैश्च शोभिता चत्वरापणैः} %||7-62-11||

\twolineshloka
{यच्च तेन महच्छून्यं लवणेन कृतं पुरा}
{शोभयामास तद्वीरो नानापण्यसमृद्धिभिः} %||7-62-12||

\twolineshloka
{तां समृद्धां समृद्धार्थः शत्रुघ्नो भरतानुजः}
{निरीक्ष्य परमप्रीतः परं हर्षमुपागमत्} %||7-62-13||

\twolineshloka
{तस्य बुद्धिः समुत्पन्ना निवेश्य मधुरां पुरीम्}
{रामपादौ निरीक्षेयं वर्षे द्वादशमे शुभे} %||7-63-14||


{॥इत्यार्षे श्रीमद्रामायणे वाल्मीकीये आदिकाव्ये उत्तरकाण्डे द्विषष्टितमः सर्गः॥}

\sect{त्रिषष्टितमः सर्गः}

\twolineshloka
{ततो द्वादशमे वर्षे शत्रुघ्नो रामपालिताम्}
{अयोध्यां चकमे गन्तुमल्पभृत्यबलानुगः} %||7-63-1||

\twolineshloka
{मन्त्रिणो बलमुख्यांश्च निवर्त्य च पुरोधसम्}
{जगाम रथमुख्येन हययुक्तेन भास्वता} %||7-63-2||

\twolineshloka
{स गत्वा गणितान्वासान्सप्ताष्टौ रघुनन्दनः}
{अयोध्यामगमत्तूर्णं राघवोत्सुकदर्शनः} %||7-63-3||

\twolineshloka
{स प्रविश्य पुरीं रम्यां श्रीमानिक्ष्वाकुनन्दनः}
{प्रविवेश महाबाहुर्यत्र रामो महाद्युतिः} %||7-63-4||

\twolineshloka
{सोऽभिवाद्य महात्मानं ज्वलन्तमिव तेजसा}
{उवाच प्राञ्जलिर्भूत्वा रामं सत्यपराक्रमम्} %||7-63-5||

\twolineshloka
{यदाज्ञप्तं महाराज सर्वं तत्कृतवानहम्}
{हतः स लवणः पापः पुरी सा च निवेशिता} %||7-63-6||

\twolineshloka
{द्वादशं च गतं वर्षं त्वां विना रघुनन्दन}
{नोत्सहेयमहं वस्तुं त्वया विरहितो नृप} %||7-63-7||

\twolineshloka
{स मे प्रसादं काकुत्स्थ कुरुष्वामितविक्रम}
{मातृहीनो यथा वत्सस्त्वां विना प्रवसाम्यहम्} %||7-63-8||

\twolineshloka
{एवं ब्रुवाणं शत्रुघ्नं परिष्वज्येदमब्रवीत्}
{मा विषादं कृथा वीर नैतत्क्षत्रिय चेष्टितम्} %||7-63-9||

\twolineshloka
{नावसीदन्ति राजानो विप्रवासेषु राघव}
{प्रजाश्च परिपाल्या हि क्षत्रधर्मेण राघव} %||7-63-10||

\twolineshloka
{काले काले च मां वीर अयोध्यामवलोकितुम्}
{आगच्छ त्वं नरश्रेष्ठ गन्तासि च पुरं तव} %||7-63-11||

\twolineshloka
{ममापि त्वं सुदयितः प्राणैरपि न संशयः}
{अवश्यं करणीयं च राज्यस्य परिपालनम्} %||7-63-12||

\twolineshloka
{तस्मात्त्वं वस काकुत्स्थ पञ्चरात्रं मया सह}
{ऊर्ध्वं गन्तासि मधुरां सभृत्यबलवाहनः} %||7-63-13||

\twolineshloka
{रामस्यैतद्वचः श्रुत्वा धर्मयुक्तं मनोऽनुगम्}
{शत्रुघ्नो दीनया वाचा बाढमित्येव चाब्रवीत्} %||7-63-14||

\twolineshloka
{स पञ्चरात्रं काकुत्स्थो राघवस्य यथाज्ञया}
{उष्य तत्र महेष्वासो गमनायोपचक्रमे} %||7-63-15||

\twolineshloka
{आमन्त्र्य तु महात्मानं रामं सत्यपराक्रमम्}
{भरतं लक्ष्मणं चैव महारथमुपारुहत्} %||7-63-16||

\twolineshloka
{दूरं ताभ्यामनुगतो लक्ष्मणेन महात्मना}
{भरतेन च शत्रुघ्नो जगामाशु पुरीं तदा} %||7-64-17||


{॥इत्यार्षे श्रीमद्रामायणे वाल्मीकीये आदिकाव्ये उत्तरकाण्डे त्रिषष्टितमः सर्गः॥}

\sect{चतुःषष्टितमः सर्गः}

\twolineshloka
{प्रस्थाप्य तु स शत्रुघ्नं भ्रातृभ्यां सह राघवः}
{प्रमुमोद सुखी राज्यं धर्मेण परिपालयन्} %||7-64-1||

\twolineshloka
{ततः कतिपयाहःसु वृद्धो जानपदो द्विजः}
{शवं बालमुपादाय राजद्वारमुपागमत्} %||7-64-2||

\twolineshloka
{रुदन्बहुविधा वाचः स्नेहाक्षरसमन्विताः}
{असकृत्पुत्रपुत्रेति वाक्यमेतदुवाच ह} %||7-64-3||

\twolineshloka
{किं नु मे दुष्कृतं कर्म पूर्वं देहान्तरे कृतम्}
{यदहं पुत्रमेकं त्वां पश्यामि निधनं गतम्} %||7-64-4||

\twolineshloka
{अप्राप्तयौवनं बालं पञ्चवर्षसमन्वितम्}
{अकाले कालमापन्नं दुःखाय मम पुत्रक} %||7-64-5||

\twolineshloka
{अल्पैरहोभिर्निधनं गमिष्यामि न संशयः}
{अहं च जननी चैव तव शोकेन पुत्रक} %||7-64-6||

\threelineshloka
{न स्मराम्यनृतं ह्युक्तं न च हिंसां स्मराम्यहम्}
{केन मे दुष्कृतेनाद्य बाल एव ममात्मजः}
{अकृत्वा पितृकार्याणि नीतो वैवस्वतक्षयम्} %||7-64-7||

\twolineshloka
{नेदृशं दृष्टपूर्वं मे श्रुतं वा घोरदर्शनम्}
{मृत्युरप्राप्तकालानां रामस्य विषये यथा} %||7-64-8||

\twolineshloka
{रामस्य दुष्कृतं किञ्चिन्महदस्ति न संशयः}
{त्वं राजञ्जीवयस्वैनं बालं मृत्युवशं गतम्} %||7-64-9||

\twolineshloka
{भ्रातृभिः सहितो राजन्दीर्घमायुरवाप्नुहि}
{उषिताः स्म सुखं राज्ये तवास्मिन्सुमहाबल} %||7-64-10||

\twolineshloka
{सम्प्रत्यनाथो विषय इक्ष्वाकूणां महात्मनाम्}
{रामं नाथमिहासाद्य बालान्तकरणं नृपम्} %||7-64-11||

\twolineshloka
{राजदोषैर्विपद्यन्ते प्रजा ह्यविधिपालिताः}
{असद्वृत्ते तु नृपतावकाले म्रियते जनः} %||7-64-12||

\twolineshloka
{यदा पुरेष्वयुक्तानि जना जनपदेशु च}
{कुर्वते न च रक्षास्ति तदाकालकृतं भयम्} %||7-64-13||

\twolineshloka
{सव्यक्तं राजदोषोऽयं भविष्यति न संशयः}
{पुरे जनपदे वापि तदा बालवधो ह्ययम्} %||7-64-14||

\twolineshloka
{एवं बहुविधैर्वाक्यैर्निन्दयानो मुहुर्मुहुः}
{राजानं दुःखसन्तप्तः सुतं तमुपगूहति} %||7-65-15||


{॥इत्यार्षे श्रीमद्रामायणे वाल्मीकीये आदिकाव्ये उत्तरकाण्डे चतुःषष्टितमः सर्गः॥}

\sect{पञ्चषष्टितमः सर्गः}

\twolineshloka
{तथा तु करुणं तस्य द्विजस्य परिदेवितम्}
{शुश्राव राघवः सर्वं दुःखशोकसमन्वितम्} %||7-65-1||

\twolineshloka
{स दुःखेन सुसन्तप्तो मन्त्रिणः समुपाह्वयत्}
{वसिष्ठं वामदेवं च भ्रातॄंश्च सहनैगमान्} %||7-65-2||

\twolineshloka
{ततो द्विजा वसिष्ठेन सार्धमष्टौ प्रवेशिताः}
{राजानं देवसङ्काशं वर्धस्वेति ततोऽब्रुवन्} %||7-65-3||

\twolineshloka
{मार्कण्डेयोऽथ मौद्गल्यो वामदेवश्च काश्यपः}
{कात्यायनोऽथ जाबालिर्गौतमो नारदस्तथा} %||7-65-4||

\twolineshloka
{एते द्विजर्षभाः सर्वे आगनेषूपवेशिताः}
{मन्त्रिणो नैगमाश्चैव यथार्हमनुकूलतः} %||7-65-5||

\twolineshloka
{तेषां समुपविष्टानां सर्वेषां दीप्ततेजसाम्}
{रघवः सर्वमाचष्टे द्विजो यस्मात्प्ररोदिति} %||7-65-6||

\twolineshloka
{तस्य तद्वचनं श्रुत्वा राज्ञो दीनस्य नारदः}
{प्रत्युवाच शुभं वाक्यमृषीणां संनिधौ नृपम्} %||7-65-7||

\twolineshloka
{शृणु राजन्यथाकाले प्राप्तोऽयं बालसङ्क्षयः}
{श्रुत्वा कर्तव्यतां वीर कुरुष्व रघुनन्दन} %||7-65-8||

\twolineshloka
{पुरा कृतयुगे राम ब्राह्मणा वै तपस्विनः}
{अब्राह्मणस्तदा राजन्न तपस्वी कथञ्चन} %||7-65-9||

\twolineshloka
{तस्मिन्युगे प्रज्वलिते ब्रह्मभूते अनावृते}
{अमृत्यवस्तदा सर्वे जज्ञिरे दीर्घदर्शिनः} %||7-65-10||

\twolineshloka
{ततस्त्रेतायुगं नाम मानवानां वपुष्मताम्}
{क्षत्रिया यत्र जायन्ते पूर्वेण तपसान्विताः} %||7-65-11||

\twolineshloka
{वीर्येण तपसा चैव तेऽधिकाः पूर्वजन्मनि}
{मानवा ये महात्मानस्तस्मिंस्त्रेतायुगे युगे} %||7-65-12||

\twolineshloka
{ब्रह्मक्षत्रं तु तत्सर्वं यत्पूर्वमपरं च यत्}
{युगयोरुभयोरासीत्समवीर्यसमन्वितम्} %||7-65-13||

\twolineshloka
{अपश्यन्तस्तु ते सर्वे विशेषमधिकं ततः}
{स्थापनं चक्रिरे तत्र चातुर्वर्ण्यस्य सर्वतः} %||7-65-14||

\twolineshloka
{अधर्मः पादमेकं तु पातयत्पृथिवीतले}
{अधर्मेण हि संयुक्तास्तेन मन्दाभवन्द्विजाः} %||7-65-15||

\twolineshloka
{ततः प्रादुष्कृतं पूर्वमायुषः परिनिष्ठितम्}
{शुभान्येवाचरँल्लोकाः सत्यधर्मपरायणाः} %||7-65-16||

\twolineshloka
{त्रेतायुगे त्ववर्तन्त ब्राह्मणाः क्षत्रियश्च ये}
{तपोऽतप्यन्त ते सर्वे शुश्रूषामपरे जनाः} %||7-65-17||

\twolineshloka
{स धर्मः परमस्तेषां वैश्यशूद्रमथागमत्}
{पूजां च सर्ववर्णानां शूद्राश्चक्रुर्विशेषतः} %||7-65-18||

\twolineshloka
{ततः पादमधर्मस्य द्वितीयमवतारयत्}
{ततो द्वापरसङ्ख्या सा युगस्य समजायत} %||7-65-19||

\twolineshloka
{तस्मिन्द्वापरसङ्ख्ये तु वर्तमाने युगक्षये}
{अधर्मश्चानृतं चैव ववृधे पुरुषर्षभ} %||7-65-20||

\twolineshloka
{तस्मिन्द्वापरसङ्ख्याते तपो वैश्यान्समाविशत्}
{न शूद्रो लभते धर्ममुग्रं तप्तं नरर्षभ} %||7-65-21||

\twolineshloka
{हीनवर्णो नरश्रेष्ठ तप्यते सुमहत्तपः}
{भविष्या शूद्रयोन्यां हि तपश्चर्या कलौ युगे} %||7-65-22||

\threelineshloka
{अधर्मः परमो राम द्वापरे शूद्रधारितः}
{स वै विषयपर्यन्ते तव राजन्महातपाः}
{शूद्रस्तप्यति दुर्बुद्धिस्तेन बालवधो ह्ययम्} %||7-65-23||

\threelineshloka
{यो ह्यधर्ममकार्यं वा विषये पार्थिवस्य हि}
{करोति राजशार्दूल पुरे वा दुर्मतिर्नरः}
{क्षिप्रं हि नरकं याति स च राजा न संशयः} %||7-65-24||

\twolineshloka
{स त्वं पुरुषशार्दूल मार्गस्व विषयं स्वकम्}
{दुष्कृतं यत्र पश्येथास्तत्र यत्नं समाचर} %||7-65-25||

\twolineshloka
{एवं ते धर्मवृद्धिश्च नृणां चायुर्विवर्धनम्}
{भविष्यति नरश्रेष्ठ बालस्यास्य च जीवितम्} %||7-66-26||


{॥इत्यार्षे श्रीमद्रामायणे वाल्मीकीये आदिकाव्ये उत्तरकाण्डे पञ्चषष्टितमः सर्गः॥}

\sect{षट्षष्टितमः सर्गः}

\twolineshloka
{नारदस्य तु तद्वाक्यं श्रुत्वामृतमयं यथा}
{प्रहर्षमतुलं लेभे लक्ष्मणं चेदमब्रवीत्} %||7-66-1||

\twolineshloka
{गच्छ सौम्य द्विजश्रेष्ठं समाश्वासय लक्ष्मण}
{बालस्य च शरीरं तत्तैलद्रोण्यां निधापय} %||7-66-2||

\twolineshloka
{गन्धैश्च परमोदारैस्तैलैश्च सुसुगन्धिभिः}
{यथा न क्षीयते बालस्तथा सौम्य विधीयताम्} %||7-66-3||

\twolineshloka
{यथा शरीरे बालस्य गुप्तस्याक्लिष्टकर्मणः}
{विपत्तिः परिभेदो वा भवेन्न च तथा कुरु} %||7-66-4||

\twolineshloka
{तथा सन्दिश्य काकुत्स्थो लक्ष्मणं शुभलक्षणम्}
{मनसा पुष्पकं दध्यावागच्छेति महायशाः} %||7-66-5||

\twolineshloka
{इङ्गितं स तु विज्ञाय पुष्पको हेमभूषितः}
{आजगाम मुहूर्तेन सम्पीपं राघवस्य वै} %||7-66-6||

\twolineshloka
{सोऽब्रवीत्प्रणतो भूत्वा अयमस्मि नराधिप}
{वश्यस्तव महाबाहो किङ्करः समुपस्थितः} %||7-66-7||

\twolineshloka
{भाषितं रुचिरं श्रुत्वा पुष्पकस्य नराधिपः}
{अभिवाद्य महर्षीस्तान्विमानं सोऽध्यरोहत} %||7-66-8||

\twolineshloka
{धनुर्गृहीत्वा तूणीं च खग्दं च रुचिरप्रभम्}
{निक्षिप्य नगरे वीरौ सौमित्रिभरतावुभौ} %||7-66-9||

\twolineshloka
{प्रायात्प्रतीचीं स मरून्विचिन्वंश्च समन्ततः}
{उत्तरामगमच्छ्रीमान्दिशं हिमवदावृतम्} %||7-66-10||

\twolineshloka
{अपश्यमानस्तत्रापि स्वल्पमप्यथ दुष्कृतम्}
{पूर्वामपि दिशं सर्वामथापश्यन्नराधिपः} %||7-66-11||

\twolineshloka
{दक्षिणां दिशमाक्रामत्ततो राजर्षिनन्दनः}
{शैवलस्योत्तरे पार्श्वे ददर्श सुमहत्सरः} %||7-66-12||

\twolineshloka
{तस्मिन्सरसि तप्यन्तं तापसं सुमहत्तपः}
{ददर्श राघवः श्रीमाँल्लम्बमानमधो मुखम्} %||7-66-13||

\twolineshloka
{अथैनं समुपागम्य तप्यन्तं तप उत्तमम्}
{उवाच राघवो वाक्यं धन्यस्त्वमसि सुव्रत} %||7-66-14||

\twolineshloka
{कस्यां योन्यां तपोवृद्धवर्तसे दृढविक्रम}
{कौतूहलात्त्वां पृच्छामि रामो दाशरथिर्ह्यहम्} %||7-66-15||

\twolineshloka
{मनीषितस्ते को न्वर्थः स्वर्गलाभो वराश्रयः}
{यमश्रित्य तपस्तप्तं श्रोतुमिच्छामि तापस} %||7-66-16||

\twolineshloka
{ब्राह्मणो वासि भद्रं ते क्षत्रियो वासि दुर्जयः}
{वैश्यो वा यदि वा शूद्रः सत्यमेतद्ब्रवीहि मे} %||7-67-17||


{॥इत्यार्षे श्रीमद्रामायणे वाल्मीकीये आदिकाव्ये उत्तरकाण्डे षट्षष्टितमः सर्गः॥}

\sect{सप्तषष्टितमः सर्गः}

\twolineshloka
{तस्य तद्वचनं श्रुत्वा रामस्याक्लिष्टकर्मणः}
{अवाक्शिरास्तथाभूतो वाक्यमेतदुवाच ह} %||7-67-1||

\twolineshloka
{शूद्रयोन्यां प्रसूतोऽस्मि तप उग्रं समास्थितः}
{देवत्वं प्रार्थये राम सशरीरो महायशः} %||7-67-2||

\twolineshloka
{न मिथ्याहं वदे राजन्देवलोकजिगीषया}
{शूद्रं मां विद्धि काकुत्स्थ शम्बूकं नाम नामतः} %||7-67-3||

\twolineshloka
{भाषतस्तस्य शूद्रस्य खड्गं सुरुचिरप्रभम्}
{निष्कृष्य कोशाद्विमलं शिरश्चिच्छेद राघवः} %||7-67-4||

{तस्मिन्मुहूर्ते बालोऽसौ जीवेन समयुज्यत॥\devanumber{5}॥} %||7-67-5|| (Check)

\addtocounter{shlokacount}{1}

\twolineshloka
{ततोऽगस्त्याश्रमपदं रामः कमललोचनः}
{स गत्वा विनयेनैव तं नत्वा मुमुदे सुखी} %||7-67-6||

\twolineshloka
{सोऽभिवाद्य महात्मानं ज्वलन्तमिव तेजसा}
{आतिथ्यं परमं प्राप्य निषसाद नराधिपः} %||7-67-7||

\twolineshloka
{तमुवाच महातेजाः कुम्भयोनिर्महातपाः}
{स्वागतं ते नरश्रेष्ठ दिष्ट्या प्राप्तोऽसि राघव} %||7-67-8||

\twolineshloka
{त्वं मे बहुमतो राम गुणैर्बहुभिरुत्तमैः}
{अतिथिः पूजनीयश्च माम राजन्हृदि स्थितः} %||7-67-9||

\twolineshloka
{सुरा हि कथयन्ति त्वामागतं शूद्रघातिनम्}
{ब्राह्मणस्य तु धर्मेण त्वया जीवापितः सुतः} %||7-67-10||

\twolineshloka
{उष्यतां चेह रजनीं सकाशे मम राघव}
{प्रभाते पुष्पकेण त्वं गन्ता स्वपुरमेव हि} %||7-67-11||

\threelineshloka
{इदं चाभरणं सौम्य निर्मितं विश्वकर्मणा}
{दिव्यं दिव्येन वपुषा दीप्यमानं स्वतेजसा}
{प्रतिगृह्णीष्व काकुत्स्थ मत्प्रियं कुरु राघव} %||7-67-12||

\twolineshloka
{दत्तस्य हि पुनर्दानं सुमहत्फलमुच्यते}
{तस्मात्प्रदास्ये विधिवत्तत्प्रतीच्छ नरर्षभ} %||7-67-13||

\twolineshloka
{तद्रामः प्रतिजग्राह मुनेस्तस्य महात्मनः}
{दिव्यमाभरणं चित्रं प्रदीप्तमिव भास्करम्} %||7-67-14||

\twolineshloka
{प्रतिगृह्य ततो रामस्तदाभरणमुत्तमम्}
{आगमं तस्य दिव्यस्य प्रष्टुमेवोपचक्रमे} %||7-67-15||

\twolineshloka
{अत्यद्भुतमिदं ब्रह्मन्वपुषा युक्तमुत्तमम्}
{कथं भगवता प्राप्तं कुतो वा केन वाहृतम्} %||7-67-16||

\twolineshloka
{कुतूहलतया ब्रह्मन्पृच्छामि त्वां महायशः}
{आश्चर्याणां बहूनां हि निधिः परमको भवान्} %||7-67-17||

\twolineshloka
{एवं ब्रुवति काकुत्स्थे मुनिर्वाक्यमथाब्रवीत्}
{शृणु राम यथावृत्तं पुरा त्रेतायुगे गते} %||7-68-18||


{॥इत्यार्षे श्रीमद्रामायणे वाल्मीकीये आदिकाव्ये उत्तरकाण्डे सप्तषष्टितमः सर्गः॥}

\sect{अष्टषष्टितमः सर्गः}

\twolineshloka
{पुरा त्रेतायुगे ह्यासीदरण्यं बहुविस्तरम्}
{समन्ताद्योजनशतं निर्मृगं पक्षिवर्जितम्} %||7-68-1||

\twolineshloka
{तस्मिन्निर्मानुषेऽरण्ये कुर्वाणस्तप उत्तमम्}
{अहमाक्रमितुं शौम्य तदरण्यमुपागमम्} %||7-68-2||

\twolineshloka
{तस्य रूपमरण्यस्य निर्देष्टुं न शशाक ह}
{फलमूलैः सुखास्वादैर्बहुरूपैश्च पादपैः} %||7-68-3||

\twolineshloka
{तस्यारण्यस्य मध्ये तु सरो योजनमायतम्}
{पद्मोत्पलसमाकीर्णं समतिक्रान्तशैवलम्} %||7-68-4||

\twolineshloka
{तदाश्चर्यमिवात्यर्थं सुखास्वादमनुत्तमम्}
{अरजस्कं तथाक्षोभ्यं श्रीमत्पक्षिगणायुतम्} %||7-68-5||

\twolineshloka
{तस्मिन्सरःसमीपे तु महदद्भुतमाश्रमम्}
{पुराणं पुण्यमत्यर्थं तपस्विजनवर्जितम्} %||7-68-6||

\twolineshloka
{तत्राहमवसं रात्रिं नैदाघीं पुरुषर्षभ}
{प्रभाते काल्यमुत्थाय सरस्तदुपचक्रमे} %||7-68-7||

\twolineshloka
{अथापश्यं शवं तत्र सुपुष्टमजरं क्वचित्}
{तिष्ठन्तं परया लक्ष्म्या तस्मिंस्तोयाशये नृप} %||7-68-8||

\twolineshloka
{तमर्थं चिन्तयानोऽहं मुहूर्तं तत्र राघव}
{विष्ठितोऽस्मि सरस्तीरे किं न्विदं स्यादिति प्रभो} %||7-68-9||

\twolineshloka
{अथापश्यं मुहूर्तात्तु दिव्यमद्भुतदर्शनम्}
{विमानं परमोदारं हंसयुक्तं मनोजवम्} %||7-68-10||

\threelineshloka
{अत्यर्थं स्वर्गिणं तत्र विमाने रघुनन्दन}
{उपास्तेऽप्सरसां वीर सहस्रं दिव्यभूषणम्}
{गान्ति गेयानि रम्याणि वादयन्ति तथापराः} %||7-68-11||

\twolineshloka
{पश्यतो मे तदा राम विमानादवरुह्य च}
{तं शवं भक्षयामास स स्वर्गी रघुनन्दन} %||7-68-12||

\twolineshloka
{ततो भुक्त्वा यथाकामं मांसं बहु च सुष्ठु च}
{अवतीर्य सरः स्वर्गी संस्प्रष्टुमुपचक्रमे} %||7-68-13||

\twolineshloka
{उपस्पृश्य यथान्यायं स स्वर्गी पुरुषर्षभ}
{आरोढुमुपचक्राम विमानवरमुत्तमम्} %||7-68-14||

\twolineshloka
{तमहं देवसङ्काशमारोहन्तमुदीक्ष्य वै}
{अथाहमब्रुवं वाक्यं तमेव पुरुषर्षभ} %||7-68-15||

\twolineshloka
{को भवान्देवसङ्काश आहारश्च विगर्हितः}
{त्वयायं भुज्यते सौम्य किं कर्थं वक्तुमर्हसि} %||7-68-16||

\twolineshloka
{आश्चर्यमीदृशो भावो भास्वरो देवसम्मतः}
{आहारो गर्हितः सौम्य श्रोतुमिच्छामि तत्त्वतः} %||7-69-17||


{॥इत्यार्षे श्रीमद्रामायणे वाल्मीकीये आदिकाव्ये उत्तरकाण्डे अष्टषष्टितमः सर्गः॥}

\sect{एकोनसप्ततितमः सर्गः}

\twolineshloka
{भुक्त्वा तु भाषितं वाक्यं मम राम शुभाक्षरम्}
{प्राञ्जलिः प्रत्युवाचेदं स स्वर्गी रघुनन्दन} %||7-69-1||

\twolineshloka
{शृणु ब्रह्मन्यथावृत्तं ममैतत्सुखदुःखयोः}
{दुरतिक्रमणीयं हि यथा पृच्छसि मां द्विज} %||7-69-2||

\twolineshloka
{पुरा वैदर्भको राजा पिता मम महायशाः}
{सुदेव इति विख्यातस्त्रिषु लोकेषु वीर्यवान्} %||7-69-3||

\twolineshloka
{तस्य पुत्रद्वयं ब्रह्मन्द्वाभ्यां स्त्रीभ्यामजायत}
{अहं श्वेत इति ख्यातो यवीयान्सुरथोऽभवत्} %||7-69-4||

\twolineshloka
{ततः पितरि स्वर्याते पौरा मामभ्यषेचयन्}
{तत्राहं कृतवान्राज्यं धर्मेण सुसमाहितः} %||7-69-5||

\twolineshloka
{एवं वर्षसहस्राणि समतीतानि सुव्रत}
{राज्यं कारयतो ब्रह्मन्प्रजा धर्मेण रक्षतः} %||7-69-6||

\twolineshloka
{सोऽहं निमित्ते कस्मिंश्चिद्विज्ञातायुर्द्विजोत्तम}
{कालधर्मं हृदि न्यस्य ततो वनमुपागमम्} %||7-69-7||

\twolineshloka
{सोऽहं वनमिदं दुर्गं मृगपक्षिविवर्जितम्}
{तपश्चर्तुं प्रविष्टोऽस्मि समीपे सरसः शुभे} %||7-69-8||

\twolineshloka
{भ्रातरं सुरथं राज्ये अभिषिच्य नराधिपम्}
{इदं सरः समासाद्य तपस्तप्तं मया चिरम्} %||7-69-9||

\twolineshloka
{सोऽहं वर्षसहस्राणि तपस्त्रीणि महामुने}
{तप्त्वा सुदुष्करं प्राप्तो ब्रह्मलोकमनुत्तमम्} %||7-69-10||

\twolineshloka
{ततो मां स्वर्गसंस्थं वै क्षुत्पिपासे द्विजोत्तम}
{बाधेते परमोदार ततोऽहं व्यथितेन्द्रियः} %||7-69-11||

\twolineshloka
{गत्वा त्रिभुवणश्रेष्ठं पितामहमुवाच ह}
{भगवन्ब्रह्मलोकोऽयं क्षुत्पिपासाविवर्जितः} %||7-69-12||

\twolineshloka
{कस्येयं कर्मणः प्राप्तिः क्षुत्पिपासावशोऽस्मि यत्}
{आहारः कश्च मे देव तन्मे ब्रूहि पितामह} %||7-69-13||

\twolineshloka
{पितामहस्तु मामाह तवाहारः सुदेवज}
{स्वादूनि स्वानि मांसानि तानि भक्षय नित्यशः} %||7-69-14||

\twolineshloka
{स्वशरीरं त्वया पुष्टं कुर्वता तप उत्तमम्}
{अनुप्तं रोहते श्वेत न कदाचिन्महामते} %||7-69-15||

\twolineshloka
{दत्तं न तेऽस्ति सूक्ष्मोऽपि वने सत्त्वनिषेविते}
{तेन स्वर्गगतो वत्स बाध्यसे क्षुत्पिपासया} %||7-69-16||

\twolineshloka
{स त्वं सुपुष्टमाहारैः स्वशरीरमनुत्तमम्}
{भक्षयस्वामृतरसं सा ते तृप्तिर्भविष्यति} %||7-69-17||

\twolineshloka
{यदा तु तद्वनं श्वेत अगस्त्यः सुमहानृषिः}
{आक्रमिष्यति दुर्धर्षस्तदा कृच्छाद्विमोक्ष्यसे} %||7-69-18||

\twolineshloka
{स हि तारयितुं सौम्य शक्तः सुरगणानपि}
{किं पुनस्त्वां महाबाहो क्षुत्पिपासावशं गतम्} %||7-69-19||

\twolineshloka
{सोऽहं भगवतः श्रुत्वा देवदेवस्य निश्चयम्}
{आहारं गर्हितं कुर्मि स्वशरीरं द्विजोत्तम} %||7-69-20||

\twolineshloka
{बहून्वर्षगणान्ब्रह्मन्भुज्यमानमिदं मया}
{क्षयं नाभ्येति ब्रह्मर्षे तृप्तिश्चापि ममोत्तमा} %||7-69-21||

\twolineshloka
{तस्य मे कृच्छ्रभूतस्य कृच्छ्रादस्माद्विमोक्षय}
{अन्येषामगतिर्ह्यत्र कुम्भयोनिमृते द्विजम्} %||7-69-22||

\twolineshloka
{इदमाभरणं सौम्य तारणार्थं द्विजोत्तम}
{प्रतिगृह्णीष्व ब्रह्मर्षे प्रसादं कर्तुमर्हसि} %||7-69-23||

\twolineshloka
{तस्याहं स्वर्गिणो वाक्यं श्रुत्वा दुःखसमन्वितम्}
{तारणायोपजग्राह तदाभरणमुत्तमम्} %||7-69-24||

\twolineshloka
{मया प्रतिगृहीते तु तस्मिन्नाभरणे शुभे}
{मानुषः पूर्वको देहो राजर्षेः स ननाश ह} %||7-69-25||

\twolineshloka
{प्रनष्टे तु शरीरेऽसौ राजर्षिः परया मुदा}
{तृप्तः प्रमुदितो राजा जगाम त्रिदिवं पुनः} %||7-69-26||

\twolineshloka
{तेनेदं शक्रतुल्येन दिव्यमाभरणं मम}
{तस्मिन्निमित्ते काकुत्स्थ दत्तमद्भुतदर्शनम्} %||7-70-27||


{॥इत्यार्षे श्रीमद्रामायणे वाल्मीकीये आदिकाव्ये उत्तरकाण्डे एकोनसप्ततितमः सर्गः॥}

\sect{सप्ततितमः सर्गः}

\twolineshloka
{तदद्भुततमं वाक्यं श्रुत्वागस्त्यस्य राघवः}
{गौरवाद्विस्मयाच्चैव भूयः प्रष्टुं प्रचक्रमे} %||7-70-1||

\twolineshloka
{भगवंस्तद्वनं घोरं तपस्तप्यति यत्र सः}
{श्वेतो वैदर्भको राजा कथं तदमृगद्विजम्} %||7-70-2||

\twolineshloka
{निःसत्त्वं च वनं जातं शून्यं मनुजवर्जितम्}
{तपश्चर्तुं प्रविष्टः स श्रोतुमिच्छामि तत्त्वतः} %||7-70-3||

\twolineshloka
{रामस्य भाषितं श्रुत्वा कौतूहलसमन्वितम्}
{वाक्यं परमतेजस्वी वक्तुमेवोपचक्रमे} %||7-70-4||

\twolineshloka
{पुरा कृतयुगे राम मनुर्दण्डधरः प्रभुः}
{तस्य पुत्रो महानासीदिक्ष्वाकुः कुलवर्धनः} %||7-70-5||

\twolineshloka
{तं पुत्रं पूर्वके राज्ये निक्षिप्य भुवि दुर्जयम्}
{पृथिव्यां राजवंशानां भव कर्तेत्युवाच ह} %||7-70-6||

\twolineshloka
{तथेति च प्रतिज्ञातं पितुः पुत्रेण राघव}
{ततः परमसंहृष्टो मनुः पुनरुवाच ह} %||7-70-7||

\twolineshloka
{प्रीतोऽस्मि परमोदारकर्ता चासि न संशयः}
{दण्डेन च प्रजा रक्ष मा च दण्डमकारणे} %||7-70-8||

\twolineshloka
{अपराधिषु यो दण्डः पात्यते मानवेषु वै}
{स दण्डो विधिवन्मुक्तः स्वर्गं नयति पार्थिवम्} %||7-70-9||

\twolineshloka
{तस्माद्दण्डे महाबाहो यत्नवान्भव पुत्रक}
{धर्मो हि परमो लोके कुर्वतस्ते भविष्यति} %||7-70-10||

\twolineshloka
{इति तं बहु सन्दिश्य मनुः पुत्रं समाधिना}
{जगाम त्रिदिवं हृष्टो ब्रह्मलोकमनुत्तमम्} %||7-70-11||

\twolineshloka
{प्रयाते त्रिदिवे तस्मिन्निक्ष्वाकुरमितप्रभः}
{जनयिष्ये कथं पुत्रानिति चिन्तापरोऽभवत्} %||7-70-12||

\twolineshloka
{कर्मभिर्बहुरूपैश्च तैस्तैर्मनुसुतः सुतान्}
{जनयामास धर्मात्मा शतं देवसुतोपमान्} %||7-70-13||

\twolineshloka
{तेषामवरजस्तात सर्वेषां रघुनन्दन}
{मूढश्चाकृतिविद्यश्च न शुश्रूषति पूर्वजान्} %||7-70-14||

\twolineshloka
{नाम तस्य च दण्डेति पिता चक्रेऽल्पतेजसः}
{अवश्यं दण्डपतनं शरीरेऽस्य भविष्यति} %||7-70-15||

\twolineshloka
{स पश्यमानस्तं दोषं घोरं पुत्रस्य राघव}
{विन्ध्यशैवलयोर्मध्ये राज्यं प्रादादरिन्दम} %||7-70-16||

\twolineshloka
{स दण्डस्तत्र राजाभूद्रम्ये पर्वतरोधसि}
{पुरं चाप्रतिमं राम न्यवेशयदनुत्तमम्} %||7-70-17||

\twolineshloka
{पुरस्य चाकरोन्नाम मधुमन्तमिति प्रभो}
{पुरोहितं चोशनसं वरयामास सुव्रतम्} %||7-70-18||

\twolineshloka
{एवं स राजा तद्राज्यं कारयत्सपुरोहितः}
{प्रहृष्टमनुजाकीर्णं देवराज्यं यथा दिवि} %||7-71-19||


{॥इत्यार्षे श्रीमद्रामायणे वाल्मीकीये आदिकाव्ये उत्तरकाण्डे सप्ततितमः सर्गः॥}

\sect{एकसप्ततितमः सर्गः}

\twolineshloka
{एतदाख्याय रामाय महर्षिः कुम्भसम्भवः}
{अस्यामेवापरं वाक्यं कथायामुपचक्रमे} %||7-71-1||

\twolineshloka
{ततः स दण्डः काकुत्स्थ बहुवर्षगणायुतम्}
{अकरोत्तत्र मन्दात्मा राज्यं निहतकण्टकम्} %||7-71-2||

\twolineshloka
{अथ काले तु कस्मिंश्चिद्राजा भार्गवमाश्रमम्}
{रमणीयमुपाक्रामच्चैत्रे मासि मनोरमे} %||7-71-3||

\twolineshloka
{तत्र भार्गवकन्यां स रूपेणाप्रतिमां भुवि}
{विचरन्तीं वनोद्देशे दण्डोऽपश्यदनुत्तमाम्} %||7-71-4||

\twolineshloka
{स दृष्ट्वा तां सुदुर्मेधा अनङ्गशरपीडितः}
{अभिगम्य सुसंविग्नः कन्यां वचनमब्रवीत्} %||7-71-5||

\twolineshloka
{कुतस्त्वमसि सुश्रोणि कस्य वासि सुता शुभे}
{पीडितोऽहमनङ्गेन पृच्छामि त्वां सुमध्यमे} %||7-71-6||

\twolineshloka
{तस्य त्वेवं ब्रुवाणस्य मोहोन्मत्तस्य कामिनः}
{भार्गवी प्रत्युवाचेदं वचः सानुनयं नृपम्} %||7-71-7||

\twolineshloka
{भार्गवस्य सुतां विद्धि देवस्याक्लिष्टकर्मणः}
{अरजां नाम राजेन्द्र ज्येष्ठामाश्रमवासिनीम्} %||7-71-8||

\twolineshloka
{गुरुः पिता मे राजेन्द्र त्वं च शिष्यो महात्मनः}
{व्यसनं सुमहत्क्रुद्धः स ते दद्यान्महातपाः} %||7-71-9||

\twolineshloka
{यदि वात्र मया कार्यं धर्मदृष्टेन सत्पथा}
{वरयस्व नृप श्रेष्ठ पितरं मे महाद्युतिम्} %||7-71-10||

\twolineshloka
{अन्यथा तु फलं तुभ्यं भवेद्घोराभिसंहितम्}
{क्रोधेन हि पिता मेऽसौ त्रैलोक्यमपि निर्दहेत्} %||7-71-11||

\twolineshloka
{एवं ब्रुवाणामरजां दण्डः कामशरार्दितः}
{प्रत्युवाच मदोन्मत्तः शिरस्याधाय सोऽञ्जलिम्} %||7-71-12||

\twolineshloka
{प्रसादं कुरु सुश्रोणि न कालं क्षेप्तुमर्हसि}
{त्वत्कृते हि मम प्राणा विदीर्यन्ते शुभानने} %||7-71-13||

\twolineshloka
{त्वां प्राप्य हि वधो वापि पापं वापि सुदारुणम्}
{भक्तं भजस्व मां भीरु भजमानं सुविह्वलम्} %||7-71-14||

\twolineshloka
{एवमुक्त्वा तु तां कन्यां दोर्भ्यां गृह्य बलाद्बली}
{विस्फुरन्तीं यथाकामं मैथुनायोपचक्रमे} %||7-71-15||

\twolineshloka
{तमनर्थं महाघोरं दण्डः कृत्वा सुदारुणम्}
{नगरं प्रययौ चाशु मधुमन्तमनुत्तमम्} %||7-71-16||

\twolineshloka
{अरजापि रुदन्ती सा आश्रमस्याविदूरतः}
{प्रतीक्षते सुसन्त्रस्ता पितरं देवसंनिभम्} %||7-72-17||


{॥इत्यार्षे श्रीमद्रामायणे वाल्मीकीये आदिकाव्ये उत्तरकाण्डे एकसप्ततितमः सर्गः॥}

\sect{द्विसप्ततितमः सर्गः}

\twolineshloka
{स मुहूर्तादुपश्रुत्य देवर्षिरमितप्रभः}
{स्वमाश्रमं शिष्य वृतः क्षुधार्तः संन्यवर्तत} %||7-72-1||

\twolineshloka
{सोऽपश्यदरजां दीनां रजसा समभिप्लुताम्}
{ज्योत्स्नामिवारुणग्रस्तां प्रत्यूषे न विराजतीम्} %||7-72-2||

\twolineshloka
{तस्य रोषः समभवत्क्षुधार्तस्य विशेषतः}
{निर्दहन्निव लोकांस्त्रीञ्शिष्यांश्चेदमुवाच ह} %||7-72-3||

\twolineshloka
{पश्यध्वं विपरीतस्य दण्डस्याविदितात्मनः}
{विपत्तिं घोरसङ्काशां क्रुद्धामग्निशिखामिव} %||7-72-4||

\twolineshloka
{क्षयोऽस्य दुर्मतेः प्राप्तः सानुगस्य दुरात्मनः}
{यः प्रदीप्तां हुताशस्य शिखां वै स्प्रष्टुमिच्छति} %||7-72-5||

\twolineshloka
{यस्मात्स कृतवान्पापमीदृशं घोरदर्शनम्}
{तस्मात्प्राप्स्यति दुर्मेधाः फलं पापस्य कर्मणः} %||7-72-6||

\twolineshloka
{सप्तरात्रेण राजासौ सभृत्यबलवाहनः}
{पापकर्मसमाचारो वधं प्राप्स्यति दुर्मतिः} %||7-72-7||

\twolineshloka
{समन्ताद्योजनशतं विषयं चास्य दुर्मतेः}
{धक्ष्यते पांसुवर्षेण महता पाकशासनः} %||7-72-8||

\twolineshloka
{सर्वसत्त्वानि यानीह स्थावराणि चराणि च}
{महता पांसुवर्षेण नाशं यास्यन्ति सर्वशः} %||7-72-9||

\twolineshloka
{दण्डस्य विषयो यावत्तावत्सर्वसमुच्छ्रयः}
{पांसुभुत इवालक्ष्यः सप्तरात्राद्भविष्यति} %||7-72-10||

\twolineshloka
{इत्युक्त्वा क्रोधसन्तपस्तमाश्रमनिवासिनम्}
{जनं जनपदान्तेषु स्थीयतामिति चाब्रवीत्} %||7-72-11||

\twolineshloka
{श्रुत्वा तूशसनो वाक्यं स आश्रमावसथो जनः}
{निष्क्रान्तो विषयात्तस्य स्थानं चक्रेऽथ बाह्यतः} %||7-72-12||

\twolineshloka
{स तथोक्त्वा मुनिजनमरजामिदमब्रवीत्}
{इहैव वस दुर्मेधे आश्रमे सुसमाहिता} %||7-72-13||

\twolineshloka
{इदं योजनपर्यन्तं सरः सुरुचिरप्रभम्}
{अरजे विज्वरा भुङ्क्ष्व कालश्चात्र प्रतीक्ष्यताम्} %||7-72-14||

\twolineshloka
{त्वत्समीपे तु ये सत्त्वा वासमेष्यन्ति तां निशाम्}
{अवध्याः पांसुवर्षेण ते भविष्यन्ति नित्यदा} %||7-72-15||

\twolineshloka
{इत्युक्त्वा भार्गवो वासमन्यत्र समुपाक्रमत्}
{सप्ताहाद्भस्मसाद्भूतं यथोक्तं ब्रह्मवादिना} %||7-72-16||

\twolineshloka
{तस्यासौ दण्डविषयो विन्ध्यशैवलसानुषु}
{शप्तो ब्रह्मर्षिणा तेन पुरा वैधर्मके कृते} %||7-72-17||

\twolineshloka
{ततः प्रभृति काकुत्स्थ दण्डकारण्यमुच्यते}
{तपस्विनः स्थिता यत्र जनस्थानमथोऽभवत्} %||7-72-18||

\twolineshloka
{एतत्ते सर्वमाख्यातं यन्मां पृच्छसि राघव}
{सन्ध्यामुपासितुं वीर समयो ह्यतिवर्तते} %||7-72-19||

\twolineshloka
{एते महर्षयः सर्वे पूर्णकुम्भाः समन्ततः}
{कृतोदको नरव्याघ्र आदित्यं पर्युपासते} %||7-72-20||

\twolineshloka
{स तैरृषिभिरभ्यस्तः सहितैर्ब्रह्मसत्तमैः}
{रविरस्तं गतो राम गच्छोदकमुपस्पृश} %||7-73-21||


{॥इत्यार्षे श्रीमद्रामायणे वाल्मीकीये आदिकाव्ये उत्तरकाण्डे द्विसप्ततितमः सर्गः॥}

\sect{त्रिसप्ततितमः सर्गः}

\twolineshloka
{ऋषेर्वचनमाज्ञाय रामः सन्ध्यामुपासितुम्}
{उपाक्रामत्सरः पुण्यमप्सरोभिर्निषेवितम्} %||7-73-1||

\twolineshloka
{तत्रोदकमुपस्पृष्श्य सन्ध्यामन्वास्य पश्चिमाम्}
{आश्रमं प्राविशद्रामः कुम्भयोनेर्महात्मनः} %||7-73-2||

\twolineshloka
{अस्यागस्त्यो बहुगुणं फलमूलं तथौषधीः}
{शाकानि च पवित्राणि भोजनार्थमकल्पयत्} %||7-73-3||

\twolineshloka
{स भुक्तवान्नरश्रेष्ठस्तदन्नममृतोपमम्}
{प्रीतश्च परितुष्टश्च तां रात्रिं समुपावसत्} %||7-73-4||

\twolineshloka
{प्रभाते काल्यमुत्थाय कृत्वाह्निकमरिन्दमः}
{ऋषिं समभिचक्राम गमनाय रघूत्तमः} %||7-73-5||

\twolineshloka
{अभिवाद्याब्रवीद्रामो महर्षिं कुम्भसम्भवम्}
{आपृच्छे त्वां गमिष्यामि मामनुज्ञातुमर्हसि} %||7-73-6||

\twolineshloka
{धन्योऽस्म्यनुगृहीतोऽस्मि दर्शनेन महात्मनः}
{द्रष्टुं चैवागमिष्यामि पावनार्थमिहात्मनः} %||7-73-7||

\twolineshloka
{तथा वदति काकुत्स्थे वाक्यमद्भुतदर्शनम्}
{उवाच परमप्रीतो धर्मनेत्रस्तपोधनः} %||7-73-8||

\twolineshloka
{अत्यद्भुतमिदं वाक्यं तव राम शुभाक्षरम्}
{पावनः सर्वलोकानां त्वमेव रघुनन्दन} %||7-73-9||

\twolineshloka
{मुहूर्तमपि राम त्वां ये नु पश्यन्ति केचन}
{पाविताः स्वर्गभूतास्ते पूज्यन्ते दिवि दैवतैः} %||7-73-10||

\twolineshloka
{ये च त्वां घोरचक्षुर्भिरीक्षन्ते प्राणिनो भुवि}
{हतास्ते यमदण्डेन सद्यो निरयगामिनः} %||7-73-11||

\twolineshloka
{गच्छ चारिष्टमव्यग्रः पन्थानमकुतोभयम्}
{प्रशाधि राज्यं धर्मेण गतिर्हि जगतो भवान्} %||7-73-12||

\twolineshloka
{एवमुक्तस्तु मुनिना प्राञ्जलिः प्र्पग्रहो नृपः}
{अभ्यवादयत प्राज्ञस्तमृषिं पुण्यशीलिनम्} %||7-73-13||

\twolineshloka
{अभिवाद्य मुनिश्रेष्ठं तांश्च सर्वांस्तपोधनान्}
{अध्यारोहत्तदव्यग्रः पुष्पकं हेमभूषितम्} %||7-73-14||

\twolineshloka
{तं प्रयान्तं मुनिगणा आशीर्वादैः समन्ततः}
{अपूजयन्महेन्द्राभं सहस्राक्षमिवामराः} %||7-73-15||

\twolineshloka
{स्वस्थः स ददृशे रामः पुष्पके हेमभूषिते}
{शशी मेघसमीपस्थो यथा जलधरागमे} %||7-73-16||

\twolineshloka
{ततोऽर्धदिवसे प्राप्ते पूज्यमानस्ततस्ततः}
{अयोध्यां प्राप्य काकुत्स्थो विमानादवरोहत} %||7-73-17||

\twolineshloka
{ततो विसृज्य रुचिरं पुष्पकं कामगामिनम्}
{कक्ष्यान्तरविनिक्षिप्तं द्वाःस्थं रामोऽब्रवीद्वचः} %||7-73-18||

\twolineshloka
{लक्ष्मणं भरतं चैव गत्वा तौ लघुविक्रमौ}
{ममागमनमाख्याय शब्दापय च मां चिरम्} %||7-74-19||


{॥इत्यार्षे श्रीमद्रामायणे वाल्मीकीये आदिकाव्ये उत्तरकाण्डे त्रिसप्ततितमः सर्गः॥}

\sect{चतुःसप्ततितमः सर्गः}

\twolineshloka
{तच्छ्रुत्वा भाषितं तस्य रामस्याक्लिष्टकर्मणः}
{द्वाःस्थः कुमारावाहूय राघवाय न्यवेदयत्} %||7-74-1||

\twolineshloka
{दृष्ट्वा तु राघवः प्राप्तौ प्रियौ भरतलक्ष्मणौ}
{परिष्वज्य ततो रामो वाक्यमेतदुवाच ह} %||7-74-2||

\twolineshloka
{कृतं मया यथातथ्यं द्विजकार्यमनुत्तमम्}
{धर्मसेतुमतो भूयः कर्तुमिच्छामि राघवौ} %||7-74-3||

\twolineshloka
{युवाभ्यामात्मभूताभ्यां राजसूयमनुत्तमम्}
{सहितो यष्टुमिच्छामि तत्र धर्मो हि शाश्वतः} %||7-74-4||

\twolineshloka
{इष्ट्वा तु राजसूयेन मित्रः शत्रुनिबर्हणः}
{सुहुतेन सुयज्ञेन वरुणत्वमुपागमत्} %||7-74-5||

\twolineshloka
{सोमश्च राजसूयेन इष्ट्वा धर्मेण धर्मवित्}
{प्राप्तश्च सर्वलोकानां कीर्तिं स्थानं च शाश्वतम्} %||7-74-6||

\twolineshloka
{अस्मिन्नहनि यच्छ्रेयश्चिन्त्यतां तन्मया सह}
{हितं चायति युक्तं च प्रयतौ वक्तुमर्हथ} %||7-74-7||

\twolineshloka
{श्रुता तु राघवस्यैतद्वाक्यं वाक्यविशारदः}
{भरतः प्राञ्जलिर्भूत्वा वाक्यमेतदुवाच ह} %||7-74-8||

\twolineshloka
{त्वयि धर्मः परः साधो त्वयि सर्वा वसुन्धरा}
{प्रतिष्ठिता महाबाहो यशश्चामितविक्रम} %||7-74-9||

\twolineshloka
{महीपालाश्च सर्वे त्वां प्रजापतिमिवामराः}
{निरीक्षन्ते महात्मानो लोकनाथं यथा वयम्} %||7-74-10||

\twolineshloka
{प्रजाश्च पितृवद्राजन्पश्यन्ति त्वां महाबल}
{पृथिव्यां गतिभूतोऽसि प्राणिनामपि राघव} %||7-74-11||

\twolineshloka
{स त्वमेवंविधं यज्ञमाहर्तासि कथं नृप}
{पृथिव्यां राजवंशानां विनाशो यत्र दृश्यते} %||7-74-12||

\twolineshloka
{पृथिव्यां ये च पुरुषा राजन्पौरुषमागताः}
{सर्वेषां भविता तत्र क्षयः सर्वान्तकोपमः} %||7-74-13||

\twolineshloka
{स त्वं पुरुषशार्दूल गुणैरतुलविक्रम}
{पृथिवीं नार्हसे हन्तुं वशे हि तव वर्तते} %||7-74-14||

\twolineshloka
{भरतस्य तु तद्वाक्यं श्रुत्वामृतमयं यथा}
{प्रहर्षमतुलं लेभे रामः सत्यपराक्रमः} %||7-74-15||

\twolineshloka
{उवाच च शुभां वाणीं कैकेय्या नन्दिवर्धनम्}
{प्रीतोऽस्मि परितुष्टोऽस्मि तवाद्य वचनेन हि} %||7-74-16||

\twolineshloka
{इदं वचनमक्लीबं त्वया धर्मसमाहितम्}
{व्याहृतं पुरुषव्याघ्र पृथिव्याः परिपालनम्} %||7-74-17||

\twolineshloka
{एष तस्मादभिप्रायाद्राजसूयात्क्रतूत्तमान्}
{निवर्तयामि धर्मज्ञ तव सुव्याहृतेन वै} %||7-74-18||

\twolineshloka
{प्रजानां पालनं धर्मो राज्ञां यज्ञेन सम्मितः}
{तस्माच्छृणोमि ते वाक्यं साधूक्तं सुसमाहितम्} %||7-75-19||


{॥इत्यार्षे श्रीमद्रामायणे वाल्मीकीये आदिकाव्ये उत्तरकाण्डे चतुःसप्ततितमः सर्गः॥}

\sect{पञ्चसप्ततितमः सर्गः}

\twolineshloka
{तथोक्तवति रामे तु भरते च महात्मनि}
{लक्ष्मणोऽपि शुभं वाक्यमुवाच रघुनन्दनम्} %||7-75-1||

\twolineshloka
{अश्वमेधो महायज्ञः पावनः सर्वपाप्मनाम्}
{पावनस्तव दुर्धर्षो रोचतां क्रतुपुङ्गवः} %||7-75-2||

\twolineshloka
{श्रूयते हि पुरावृत्तं वासवे सुमहात्मनि}
{ब्रह्महत्यावृतः शक्रो हयमेधेन पावितः} %||7-75-3||

\twolineshloka
{पुरा किल महाबाहो देवासुरसमागमे}
{वृत्रो नाम महानासीद्दैतेयो लोकसम्मतः} %||7-75-4||

\twolineshloka
{विस्तीर्णा योजनशतमुच्छ्रितस्त्रिगुणं ततः}
{अनुरागेण लोकांस्त्रीन्स्नेहात्पश्यति सर्वतः} %||7-75-5||

\twolineshloka
{धर्मज्ञश्च कृतज्ञश्च बुद्ध्या च परिनिष्ठितः}
{शशास पृथिवीं सर्वां धर्मेण सुसमाहितः} %||7-75-6||

\twolineshloka
{तस्मिन्प्रशासति तदा सर्वकामदुघा मही}
{रसवन्ति प्रसूतानि मूलानि च फलानि च} %||7-75-7||

\twolineshloka
{अकृष्टपच्या पृथिवी सुसम्पन्ना महात्मनः}
{स राज्यं तादृशं भुङ्क्ते स्फीतमद्भुतदर्शनम्} %||7-75-8||

\twolineshloka
{तस्य बुद्धिः समुत्पन्ना तपः कुर्यामनुत्तमम्}
{तपो हि परमं श्रेयस्तपो हि परमं सुखम्} %||7-75-9||

\twolineshloka
{स निक्षिप्य सुतं ज्येष्ठं पौरेषु परमेश्वरम्}
{तप उग्रमुपातिष्ठत्तापयन्सर्वदेवताः} %||7-75-10||

\twolineshloka
{तपस्तप्यति वृत्रे तु वासवः परमार्तवत्}
{विष्णुं समुपसङ्क्रम्य वाक्यमेतदुवाच ह} %||7-75-11||

\twolineshloka
{तपस्यता महाबाहो लोका वृत्रेण निर्जिताः}
{बलवान्स हि धर्मात्मा नैनं शक्ष्यामि बाधितुम्} %||7-75-12||

\twolineshloka
{यद्यसौ तप आतिष्ठेद्भूय एव सुरेश्वर}
{यावल्लोका धरिष्यन्ति तावदस्य वशानुगाः} %||7-75-13||

\twolineshloka
{त्वं चैनं परमोदारमुपेक्षसि महाबल}
{क्षणं हि न भवेद्वृत्रः क्रुद्धे त्वयि सुरेश्वर} %||7-75-14||

\twolineshloka
{यदा हि प्रीतिसंयोगं त्वया विष्णो समागतः}
{तदा प्रभृति लोकानां नाथत्वमुपलब्धवान्} %||7-75-15||

\twolineshloka
{स त्वं प्रसादं लोकानां कुरुष्व सुमहायशः}
{त्वत्कृतेन हि सर्वं स्यात्प्रशान्तमजरं जगत्} %||7-75-16||

\twolineshloka
{इमे हि सर्वे विष्णो त्वां निरीक्षन्ते दिवौकसः}
{वृत्रघतेन महता एषां साह्यं कुरुष्व ह} %||7-75-17||

\twolineshloka
{त्वया हि नित्यशः साह्यं कृतमेषां महात्मनाम्}
{असह्यमिदमन्येषामगतीनां गतिर्भवान्} %||7-76-18||


{॥इत्यार्षे श्रीमद्रामायणे वाल्मीकीये आदिकाव्ये उत्तरकाण्डे पञ्चसप्ततितमः सर्गः॥}

\sect{षट्सप्ततितमः सर्गः}

\twolineshloka
{लक्ष्मणस्य तु तद्वाक्यं श्रुत्वा शत्रुनिबर्हणः}
{वृत्रघातमशेषेण कथयेत्याह लक्ष्मणम्} %||7-76-1||

\twolineshloka
{राघवेणैवमुक्तस्तु सुमित्रानन्दवर्धनः}
{भूय एव कथां दिव्यां कथयामास लक्ष्मणः} %||7-76-2||

\twolineshloka
{सहस्राक्षवचः श्रुत्वा सर्वेषां च दिवौकसाम्}
{विष्णुर्देवानुवाचेदं सर्वानिन्द्रपुरोगमान्} %||7-76-3||

\twolineshloka
{पूर्वं सौहृदबद्धोऽस्मि वृत्रस्य सुमहात्मनः}
{तेन युष्मत्प्रियार्थं वै नाहं हन्मि महासुरम्} %||7-76-4||

\twolineshloka
{अवश्यं करणीयं च भवतां सुखमुत्तमम्}
{तस्मादुपायमाख्यास्ये येन वृत्रं हनिष्यथ} %||7-76-5||

\twolineshloka
{त्रिधा भूतं करिष्येऽहमात्मानं सुरसत्तमाः}
{तेन वृत्रं सहस्राक्षो हनिष्यति न संशयः} %||7-76-6||

\twolineshloka
{एकोंऽशो वासवं यातु द्वितीयो वज्रमेव तु}
{तृतीयो भूतलं शक्रस्ततो वृत्रं हनिष्यति} %||7-76-7||

\twolineshloka
{तथा ब्रुवति देवेशे देवा वाक्यमथाब्रुवन्}
{एवमेतन्न सन्देहो यथा वदसि दैत्यहन्} %||7-76-8||

\twolineshloka
{भद्रं तेऽस्तु गमिष्यामो वृत्रासुरवधैषिणः}
{भजस्व परमोदारवासवं स्वेन तेजसा} %||7-76-9||

\twolineshloka
{ततः सर्वे महात्मानः सहस्राक्षपुरोगमाः}
{तदरण्यमुपाक्रामन्यत्र वृत्रो महासुरः} %||7-76-10||

\twolineshloka
{तेऽपश्यंस्तेजसा भूतं तपन्तमसुरोत्तमम्}
{पिबन्तमिव लोकांस्त्रीन्निर्दहन्तमिवाम्बरम्} %||7-76-11||

\twolineshloka
{दृष्ट्वैव चासुरश्रेष्ठं देवास्त्रासमुपागमन्}
{कथमेनं वधिष्यामः कथं न स्यात्पराजयः} %||7-76-12||

\twolineshloka
{तेषां चिन्तयतां तत्र सहस्राक्षः पुरन्दरः}
{वज्रं प्रगृह्य बाहुभ्यां प्रहिणोद्वृत्रमूर्धनि} %||7-76-13||

\twolineshloka
{कालाग्निनेव घोरेण दीप्तेनेव महार्चिषा}
{प्रतप्तं वृत्रशिरसि जगत्त्रासमुपागमत्} %||7-76-14||

\twolineshloka
{असम्भाव्यं वधं तस्य वृत्रस्य विबुधाधिपः}
{चिन्तयानो जगामाशु लोकस्यान्तं महायशाः} %||7-76-15||

\twolineshloka
{तमिन्द्रं ब्रह्महत्याशु गच्छन्तमनुगच्छति}
{अपतच्चास्य गात्रेषु तमिन्द्रं दुःखमाविशत्} %||7-76-16||

\twolineshloka
{हतारयः प्रनष्टेन्द्रा देवाः साग्निपुरोगमाः}
{विष्णुं त्रिभुवणश्रेष्ठं मुहुर्मुहुरपूजयन्} %||7-76-17||

\twolineshloka
{त्वं गतिः परमा देव पूर्वजो जगतः प्रभुः}
{रथार्थं सर्वभूतानां विष्णुत्वमुपजग्मिवान्} %||7-76-18||

\twolineshloka
{हतश्चायं त्वया वृत्रो ब्रह्महत्या च वासवम्}
{बाधते सुरशार्दूल मोक्षं तस्य विनिर्दिश} %||7-76-19||

\twolineshloka
{तेषां तद्वचनं श्रुत्वा देवानां विष्णुरब्रवीत्}
{मामेव यजतां शक्रः पावयिष्यामि वज्रिणम्} %||7-76-20||

\twolineshloka
{पुण्येन हयमेधेन मामिष्ट्वा पाकशासनः}
{पुनरेष्यति देवानामिन्द्रत्वमकुतोभयः} %||7-76-21||

\twolineshloka
{एवं सन्दिश्य देवानां तां वाणीममृतोपमा}
{जगाम विष्णुर्देवेशः स्तूयमानस्त्रिविष्टपम्} %||7-77-22||


{॥इत्यार्षे श्रीमद्रामायणे वाल्मीकीये आदिकाव्ये उत्तरकाण्डे षट्सप्ततितमः सर्गः॥}

\sect{सप्तसप्ततितमः सर्गः}

\twolineshloka
{तथा वृत्रवधं सर्वमखिलेन स लक्ष्मणः}
{कथयित्वा नरश्रेष्ठः कथाशेषमुपाक्रमत्} %||7-77-1||

\twolineshloka
{ततो हते महावीर्ये वृत्रे देवभयङ्करे}
{ब्रह्महत्यावृतः शक्रः संज्ञां लेभे न वृत्रहा} %||7-77-2||

\twolineshloka
{सोऽन्तमाश्रित्य लोकानां नष्टसंज्ञो विचेतनः}
{कालं तत्रावसत्कञ्चिद्वेष्टमानो यथोरगः} %||7-77-3||

\twolineshloka
{अथ नष्टे सहस्राक्षे उद्विग्नमभवज्जगत्}
{भूमिश्च ध्वस्तसङ्काशा निःस्नेहा शुष्ककानना} %||7-77-4||

\twolineshloka
{निःस्रोतसश्चाम्बुवाहा ह्रदाश्च सरितस्तथा}
{सङ्क्षोभश्चैव सत्त्वानामनावृष्टिकृतोऽभवत्} %||7-77-5||

\twolineshloka
{क्षीयमाणे तु लोकेऽस्मिन्सम्भ्रान्तमनसः सुराः}
{यदुक्तं विष्णुना पूर्वं तं यज्ञं समुपानयन्} %||7-77-6||

\twolineshloka
{ततः सर्वे सुरगणाः सोपाध्यायाः सहर्षिभिः}
{तं देशं सहिता जग्मुर्यत्रेन्द्रो भयमोहितः} %||7-77-7||

\twolineshloka
{ते तु दृष्ट्वा सहस्राक्षं मोहितं ब्रह्महत्यया}
{तं पुरस्कृत्य देवेशमश्वमेधं प्रचक्रिरे} %||7-77-8||

\twolineshloka
{ततोऽश्वमेधः सुमहान्महेन्द्रस्य महात्मनः}
{ववृधे ब्रह्महत्यायाः पावनार्थं नरेश्वर} %||7-77-9||

\twolineshloka
{ततो यज्ञसमाप्तौ तु ब्रह्महत्या महात्मनः}
{अभिगम्याब्रवीद्वाक्यं क्व मे स्थानं विधास्यथ} %||7-77-10||

\twolineshloka
{ते तामूचुस्ततो देवास्तुष्टाः प्रीतिसमन्विताः}
{चतुर्धा विभजात्मानमात्मनैव दुरासदे} %||7-77-11||

\twolineshloka
{देवानां भाषितं श्रुत्वा ब्रह्महत्या महात्मनाम्}
{संनिधौ स्थानमन्यत्र वरयामास दुर्वसा} %||7-77-12||

\twolineshloka
{एकेनांशेन वत्स्यामि पूर्णोदासु नदीषु वै}
{द्वितीयेन तु वृक्षेषु सत्यमेतद्ब्रवीमि वः} %||7-77-13||

\twolineshloka
{योऽयमंशस्तृतीयो मे स्त्रीषु यौवनशालिषु}
{त्रिरात्रं दर्पपर्णासु वसिष्ये दर्पघातिनी} %||7-77-14||

\twolineshloka
{हन्तारो ब्राह्मणान्ये तु प्रेक्षापूर्वमदूषकान्}
{तांश्चतुर्थेन भागेन संश्रयिष्ये सुरर्षभाः} %||7-77-15||

\twolineshloka
{प्रत्यूचुस्तां ततो देवा यथा वदसि दुर्वसे}
{तथा भवतु तत्सर्वं साधयस्व यथेप्सितम्} %||7-77-16||

\twolineshloka
{ततः प्रीत्यान्विता देवाः सहस्राक्षं ववन्दिरे}
{विज्वरः पूतपाप्मा च वासवः समपद्यत} %||7-77-17||

\twolineshloka
{प्रशान्तं च जगत्सर्वं सहस्राक्षे प्रतिष्ठते}
{यज्ञं चाद्भुतसङ्काशं तदा शक्रोऽभ्यपूजयत्} %||7-77-18||

\twolineshloka
{ईदृशो ह्यश्वमेधस्य प्रभावो रघुनन्दन}
{यजस्व सुमहाभाग हयमेधेन पार्थिव} %||7-78-19||


{॥इत्यार्षे श्रीमद्रामायणे वाल्मीकीये आदिकाव्ये उत्तरकाण्डे सप्तसप्ततितमः सर्गः॥}

\sect{अष्टसप्ततितमः सर्गः}

\twolineshloka
{तच्छ्रुत्वा लक्ष्मणेनोक्तं वाक्यं वाक्यविशारदः}
{प्रत्युवाच महातेजाः प्रहसन्राघवो वचः} %||7-78-1||

\twolineshloka
{एवमेतन्नरश्रेष्ठ यथा वदसि लक्ष्मण}
{वृत्रघातमशेषेण वाजिमेधफलं च यत्} %||7-78-2||

\twolineshloka
{श्रूयते हि पुरा सौम्य कर्दमस्य प्रजापतेः}
{पुत्रो बाह्लीश्वरः श्रीमानिलो नाम सुधार्मिकः} %||7-78-3||

\twolineshloka
{स राजा पृथिवीं सर्वां वशे कृत्वा महायशाः}
{राज्यं चैव नरव्याघ्र पुत्रवत्पर्यपालयत्} %||7-78-4||

\twolineshloka
{सुरैश्च परमोदारैर्दैतेयैश्च महासुरैः}
{नागराक्षसगन्धर्वैर्यक्षैश्च सुमहात्मभिः} %||7-78-5||

\twolineshloka
{पूज्यते नित्यशः सौम्य भयार्तै रघुनन्दन}
{अबिभ्यंश्च त्रयो लोकाः सरोषस्य महात्मनः} %||7-78-6||

\twolineshloka
{स राजा तादृशो ह्यासीद्धर्मे वीर्ये च निष्ठितः}
{बुद्ध्या च परमोदारो बाह्लीकानां महायशाः} %||7-78-7||

\twolineshloka
{स प्रचक्रे महाबाहुर्मृगयां रुचिरे वने}
{चैत्रे मनोरमे मासि सभृत्यबलवाहनः} %||7-78-8||

\twolineshloka
{प्रजघ्ने स नृपोऽरण्ये मृगाञ्शतसहस्रशः}
{हत्वैव तृप्तिर्नाभूच्च राज्ञस्तस्य महात्मनः} %||7-78-9||

\twolineshloka
{नानामृगाणामयुतं वध्यमानं महात्मना}
{यत्र जातो महासेनस्तं देशमुपचक्रमे} %||7-78-10||

\twolineshloka
{तस्मिंस्तु देवदेवेशः शैलराजसुतां हरः}
{रमयामास दुर्धर्षैः सर्वैरनुचरैः सह} %||7-78-11||

\twolineshloka
{कृत्वा स्त्रीभूतमात्मानमुमेशो गोपतिध्वजः}
{देव्याः प्रियचिकीर्षुः स तस्मिन्पर्वतनिर्झरे} %||7-78-12||

\twolineshloka
{ये च तत्र वनोद्देशे सत्त्वाः पुरुषवादिनः}
{यच्च किञ्चन तत्सर्वं नारीसंज्ञं बभूव ह} %||7-78-13||

\twolineshloka
{एतस्मिन्नन्तरे राजा स इलः कर्दमात्मजः}
{निघ्नन्मृगसहस्राणि तं देशमुपचक्रमे} %||7-78-14||

\twolineshloka
{स दृष्ट्वा स्त्रीकृतं सर्वं सव्यालमृगपक्षिणम्}
{आत्मानं सानुगं चैव स्त्रीभूतं रघुनन्दन} %||7-78-15||

\twolineshloka
{तस्य दुःखं महत्त्वासीद्दृष्ट्वात्मानं तथा गतम्}
{उमापतेश्च तत्कर्म ज्ञात्वा त्रासमुपागमत्} %||7-78-16||

\twolineshloka
{ततो देवं महात्मानं शितिकण्ठं कपर्दिनम्}
{जगाम शरणं राजा सभृत्यबलवाहनः} %||7-78-17||

\twolineshloka
{ततः प्रहस्य वरदः सह देव्या महायशाः}
{प्रजापतिसुतं वाक्यमुवाच वरदः स्वयम्} %||7-78-18||

\twolineshloka
{उत्तिष्ठोत्तिष्ठ राजर्षे कार्दमेय महाबल}
{पुरुषत्वमृते सौम्य वरं वरय सुव्रत} %||7-78-19||

\twolineshloka
{ततः स राजा शोकार्ताः प्रत्याख्यातो महात्मना}
{न स जग्राह स्त्रीभूतो वरमन्यं सुरोत्तमात्} %||7-78-20||

\twolineshloka
{ततः शोकेन महता शैलराजसुतां नृपः}
{प्रणिपत्य महादेवीं सर्वेणैवान्तरात्मना} %||7-78-21||

\twolineshloka
{ईशे वराणां वरदे लोकानामसि भामिनि}
{अमोघदर्शने देवि भजे सौम्ये नमोऽस्तु ते} %||7-78-22||

\twolineshloka
{हृद्गतं तस्य राजर्षेर्विज्ञाय हरसंनिधौ}
{प्रत्युवाच शुभं वाक्यं देवी रुद्रस्य सम्मता} %||7-78-23||

\twolineshloka
{अर्धस्य देवो वरदो वरार्धस्य तथा ह्यहम्}
{तस्मादर्धं गृहाण त्वं स्त्रीपुंसोर्यावदिच्छसि} %||7-78-24||

\twolineshloka
{तदद्भुततमं श्रुत्वा देव्या वरमनुत्तमम्}
{सम्प्रहृष्टमना भूत्वा राजा वाक्यमथाब्रवीत्} %||7-78-25||

\twolineshloka
{यदि देवि प्रसन्ना मे रूपेणाप्रतिमा भुवि}
{मासं स्त्रीत्वमुपासित्वा मासं स्यां पुरुषः पुनः} %||7-78-26||

\twolineshloka
{ईप्सितं तस्य विज्ञाय देवी सुरुचिरानना}
{प्रत्युवाच शुभं वाक्यमेवमेतद्भविष्यति} %||7-78-27||

\twolineshloka
{राजन्पुरुषभूतस्त्वं स्त्रीभावं न स्मरिष्यसि}
{स्त्रीभूतश्चापरं मासं न स्मरिष्यसि पौरुषम्} %||7-78-28||

\twolineshloka
{एवं स राजा पुरुषो मामं भूत्वाथ कार्दमिः}
{त्रैलोक्यसुन्दरी नारी मासमेकमिलाभवत्} %||7-79-29||


{॥इत्यार्षे श्रीमद्रामायणे वाल्मीकीये आदिकाव्ये उत्तरकाण्डे अष्टसप्ततितमः सर्गः॥}

\sect{एकोनाशीतितमः सर्गः}

\twolineshloka
{तां कथामिलसम्बद्धां रामेण समुदीरिताम्}
{लक्ष्मणो भरतश्चैव श्रुत्वा परमविस्मितौ} %||7-79-1||

\twolineshloka
{तौ रामं प्राञ्जलीभूत्वा तस्य राज्ञो महात्मनः}
{विस्तरं तस्य भावस्य तदा पप्रच्छतुः पुनः} %||7-79-2||

\twolineshloka
{कथं स राजा स्त्रीभूतो वर्तयामास दुर्गतिम्}
{पुरुषो वा यदा भूतः कां वृत्तिं वर्तयत्यसौ} %||7-79-3||

\twolineshloka
{तयोस्तद्भाषितं श्रुत्वा कौतूहलसमन्वितम्}
{कथयामास काकुत्ष्ठस्तस्य राज्ञो यथा गतम्} %||7-79-4||

\twolineshloka
{तमेव प्रथमं मासं स्त्रीभूत्वा लोकसुन्दरी}
{ताभिः परिवृता स्त्रीभिर्येऽस्य पूर्वं पदानुगाः} %||7-79-5||

\twolineshloka
{तत्काननं विगाह्याशु विजह्रे लोकसुन्दरी}
{द्रुमगुल्मलताकीर्णं पद्भ्यां पद्मदलेक्षणा} %||7-79-6||

\twolineshloka
{वाहनानि च सर्वाणि सन्त्यक्त्वा वै समन्ततः}
{पर्वताभोगविवरे तस्मिन्रेमे इला तदा} %||7-79-7||

\twolineshloka
{अथ तस्मिन्वनोद्देशे पर्वतस्याविदूरतः}
{सरः सुरुचिरप्रख्यं नानापक्षिगणायुतम्} %||7-79-8||

\twolineshloka
{ददर्श सा इला तस्मिन्बुधं सोमसुतं तदा}
{ज्वलन्तं स्वेन वपुषा पूर्णं सोममिवोदितम्} %||7-79-9||

\twolineshloka
{तपन्तं च तपस्तीव्रमम्भोमध्ये दुरासदम्}
{यशक्सरं कामगमं तारुण्ये पर्यवस्थितम्} %||7-79-10||

\twolineshloka
{सा तं जलाशयं सर्वं क्षोभयामास विस्मिता}
{सह तैः पूर पुरुषैः स्त्रीभूतै रघुनन्दन} %||7-79-11||

\twolineshloka
{बुधस्तु तां निरीक्ष्यैव कामबाणाभिपीडितः}
{नोपलेभे तदात्मानं चचाल च तदाम्भसि} %||7-79-12||

\twolineshloka
{इलां निरीक्षमाणः स त्रैलोक्याभ्यधिकां शुभाम्}
{चिन्तां समभ्यतिक्रामत्का न्वियं देवताधिका} %||7-79-13||

\twolineshloka
{न देवीषु न नागीषु नासुरीष्वप्सरःसु च}
{दृष्टपूर्वा मया काचिद्रूपेणैतेन शोभिता} %||7-79-14||

\twolineshloka
{सदृशीयं मम भवेद्यदि नान्यपरिग्रहा}
{इति बुद्धिं समास्थाय जलात्स्थलमुपागमत्} %||7-79-15||

\twolineshloka
{स आश्रमं समुपागम्य चतस्रः प्रमदास्ततः}
{शब्दापयत धर्मात्मा ताश्चैनं च ववन्दिरे} %||7-79-16||

\twolineshloka
{स ताः पप्रच्छ धर्मात्म कस्यैषा लोकसुन्दरी}
{किमर्थमागता चेह सत्यमाख्यात माचिरम्} %||7-79-17||

\twolineshloka
{शुभं तु तस्य तद्वाक्यं मधुरं मधुराक्षरम्}
{श्रुत्वा तु ताः स्त्रियः सर्वा ऊचुर्मधुरया गिरा} %||7-79-18||

\twolineshloka
{अस्माकमेषा सुश्रोणी प्रभुत्वे वर्तते सदा}
{अपतिः काननान्तेषु सहास्माभिरटत्यसौ} %||7-79-19||

\twolineshloka
{तद्वाक्यमव्यक्तपदं तासां स्त्रीणां निशम्य तु}
{विद्यामावर्तनीं पुण्यामावर्तयत स द्विजः} %||7-79-20||

\twolineshloka
{सोऽर्थं विदित्वा निखिलं तस्य राज्ञो यथागतम्}
{सर्वा एव स्त्रियस्ताश्च बभाषे मुनिपुङ्गवः} %||7-79-21||

\twolineshloka
{अत्र किं पुरुषा भद्रा अवसञ्शैलरोधसि}
{वत्स्यथास्मिन्गिरौ यूयमवकाशो विधीयताम्} %||7-79-22||

\twolineshloka
{मूलपुत्रफलैः सर्वा वर्तयिष्यथ नित्यदा}
{स्त्रियः किम्पुरुषान्नाम भर्तॄन्समुपलप्स्यथ} %||7-79-23||

\twolineshloka
{ताः श्रुत्वा सोमपुत्रस्य वाचं किम्पुरुषीकृताः}
{उपासां चक्रिरे शैलं बह्व्यस्ता बहुधा तदा} %||7-80-24||


{॥इत्यार्षे श्रीमद्रामायणे वाल्मीकीये आदिकाव्ये उत्तरकाण्डे एकोनाशीतितमः सर्गः॥}

\sect{अशीतितमः सर्गः}

\twolineshloka
{श्रुत्वा किम्पुरुषोत्पत्तिं लक्ष्मणो भरतस्तदा}
{आश्चर्यमिति चाब्रूतामुभौ रामं जनेश्वरम्} %||7-80-1||

\twolineshloka
{अथ रामः कथामेतां भूय एव महायशाः}
{कथयामास धर्मात्मा प्रजापतिसुतस्य वै} %||7-80-2||

\twolineshloka
{सर्वास्ता विद्रुता दृष्ट्वा किंनरीरृषिसत्तमः}
{उवाच रूपसम्पन्नां तां स्त्रियं प्रहसन्निव} %||7-80-3||

\twolineshloka
{सोमस्याहं सुदयितः सुतः सुरुचिरानने}
{भजस्व मां वरारोहे भक्त्या स्निग्धेन चक्षुषा} %||7-80-4||

\twolineshloka
{तस्य तद्वचनं श्रुत्वा शून्ये स्वजनवर्जिता}
{इला सुरुचिरप्रख्यं प्रत्युवाच महाग्रहम्} %||7-80-5||

\twolineshloka
{अहं कामकरी सौम्य तवास्मि वशवर्तिनी}
{प्रशाधि मां सोमसुत यथेच्छसि तथा कुरु} %||7-80-6||

\twolineshloka
{तस्यास्तदद्भुतप्रख्यं श्रुत्वा हर्षसमन्वितः}
{स वै कामी सह तया रेमे चन्द्रमसः सुतः} %||7-80-7||

\twolineshloka
{बुधस्य माधवो मासस्तामिलां रुचिराननाम्}
{गतो रमयतोऽत्यर्थं क्षणवत्तस्य कामिनः} %||7-80-8||

\twolineshloka
{अथ मासे तु सम्पूर्णे पूर्णेन्दुसदृशाननः}
{प्रजापतिसुतः श्रीमाञ्शयने प्रत्यबुध्यत} %||7-80-9||

\twolineshloka
{सोऽपश्यत्सोमजं तत्र तप्यन्तं सलिलाशये}
{ऊर्ध्वबाहुं निरालम्बं तं राजा प्रत्यभाषत} %||7-80-10||

\twolineshloka
{भगवन्पर्वतं दुर्गं प्रविष्टोऽस्मि सहानुगः}
{न च पश्यामि तत्सैन्यं क्व नु ते मामका गताः} %||7-80-11||

\twolineshloka
{तच्छ्रुत्वा तस्य राजर्षेर्नष्टसंज्ञस्य भाषितम्}
{प्रत्युवाच शुभं वाक्यं सान्त्वयन्परया गिरा} %||7-80-12||

\twolineshloka
{अश्मवर्षेण महता भृत्यास्ते विनिपातिताः}
{त्वं चाश्रमपदे सुप्तो वातवर्षभयार्दितः} %||7-80-13||

\twolineshloka
{समाश्वसिहि भद्रं ते निर्भयो विगतज्वरः}
{फलमूलाशनो वीर वस चेह यथासुखम्} %||7-80-14||

\twolineshloka
{स राजा तेन वाक्येन प्रत्याश्वस्तो महायशाः}
{प्रत्युवाच शुभं वाक्यं दीनो भृत्यजनक्षयात्} %||7-80-15||

\twolineshloka
{त्यक्ष्याम्यहं स्वकं राज्यं नाहं भृत्यैर्विना कृतः}
{वर्तयेयं क्षणं ब्रह्मन्समनुज्ञातुमर्हसि} %||7-80-16||

\twolineshloka
{सुतो धर्मपरो ब्रह्मञ्ज्येष्ठो मम महायशाः}
{शशबिन्दुरिति ख्यातः स मे राज्यं प्रपत्स्यते} %||7-80-17||

\twolineshloka
{न हि शक्ष्याम्यहं गत्वा भृत्यदारान्सुखान्वितान्}
{प्रतिवक्तुं महातेजः किञ्चिदप्यशुभं वचः} %||7-80-18||

\twolineshloka
{तथा ब्रुवति राजेन्द्रे बुधः परममद्भुतम्}
{सान्त्वपूर्वमथोवाच वासस्त इह रोचताम्} %||7-80-19||

\twolineshloka
{न सन्तापस्त्वया कार्यः कार्दमेय महाबल}
{संवत्सरोषितस्येह कारयिष्यामि ते हितम्} %||7-80-20||

\twolineshloka
{तस्य तद्वचनं श्रुत्वा बुधस्याक्लिष्टकर्मणः}
{वासाय विदधे बुद्धिं यदुक्तं ब्रह्मवादिना} %||7-80-21||

\twolineshloka
{मासं स स्त्री तदा भूत्वा रमयत्यनिशं शुभा}
{मासं पुरुषभावेन धर्मबुद्धिं चकार सः} %||7-80-22||

\twolineshloka
{ततः स नवमे मासि इला सोमसुतात्मजम्}
{जनयामास सुश्रोणी पुरूरवसमात्मजम्} %||7-80-23||

\twolineshloka
{जातमात्रं तु सुश्रोणी पितुर्हस्ते न्यवेशयत्}
{बुधस्य समवर्णाभमिलापुत्रं महाबलम्} %||7-80-24||

\twolineshloka
{बुधोऽपि पुरुषीभूतं समाश्वास्य नराधिपम्}
{कथाभी रमयामास धर्मयुक्ताभिरात्मवान्} %||7-81-25||


{॥इत्यार्षे श्रीमद्रामायणे वाल्मीकीये आदिकाव्ये उत्तरकाण्डे अशीतितमः सर्गः॥}

\sect{एकाशीतितमः सर्गः}

\twolineshloka
{तथोक्तवति रामे तु तस्य जन्म तदद्भुतम्}
{उवाच लक्ष्मणो भूयो भरतश्च महायशाः} %||7-81-1||

\twolineshloka
{सा प्रिया सोमपुत्रस्य संवत्सरमथोषिता}
{अकरोत्किं नरश्रेष्ठ तत्त्वं शंसितुमर्हसि} %||7-81-2||

\twolineshloka
{तयोस्तद्वाक्यमाधुर्यं निशम्य परिपृच्छतोः}
{रामः पुनरुवाचेमां प्रजापतिसुते कथाम्} %||7-81-3||

\twolineshloka
{पुरुषत्वं गते शूरे बुधः परमबुद्धिमान्}
{संवर्तं परमोदारमाजुहाव महायशाः} %||7-81-4||

\twolineshloka
{च्यवनं भृगुपुत्रं च मुनिं चारिष्टनेमिनम्}
{प्रमोदनं मोदकरं ततो दुर्वाससं मुनिम्} %||7-81-5||

\twolineshloka
{एतान्सर्वान्समानीय वाक्यज्ञस्तत्त्वदर्शिनः}
{उवाच सर्वान्सुहृदो धैर्येण सुसमाहितः} %||7-81-6||

\twolineshloka
{अयं राजा महाबाहुः कर्दमस्य इलः सुतः}
{जानीतैनं यथा भूतं श्रेयो ह्यस्य विधीयताम्} %||7-81-7||

\twolineshloka
{तेषां संवदतामेव तमाश्रममुपागमत्}
{कर्दमः सुमहातेजा द्विजैः सह महात्मभिः} %||7-81-8||

\twolineshloka
{पुलस्त्यश्च क्रतुश्चैव वषट्कारस्तथैव च}
{ओङ्कारश्च महातेजास्तमाश्रममुपागमन्} %||7-81-9||

\twolineshloka
{ते सर्वे हृष्टमनसः परस्परसमागमे}
{हितैषिणो बाह्लि पतेः पृथग्वाक्यमथाब्रुवन्} %||7-81-10||

\twolineshloka
{कर्दमस्त्वब्रवीद्वाक्यं सुतार्थं परमं हितम्}
{द्विजाः शृणुत मद्वाक्यं यच्छ्रेयः पार्थिवस्य हि} %||7-81-11||

\twolineshloka
{नान्यं पश्यामि भैषज्यमन्तरेण वृषध्वजम्}
{नाश्वमेधात्परो यज्ञः प्रियश्चैव महात्मनः} %||7-81-12||

\threelineshloka
{तस्माद्यजामहे सर्वे पार्थिवार्थे दुरासदम्}
{कर्दमेनैवमुक्तास्तु सर्व एव द्विजर्षभाः}
{रोचयन्ति स्म तं यज्ञं रुद्रस्याराधनं प्रति} %||7-81-13||

\twolineshloka
{संवर्तस्य तु राजर्षिः शिष्यः परपुरंजयः}
{मरुत्त इति विख्यातस्तं यज्ञं समुपाहरत्} %||7-81-14||

\twolineshloka
{ततो यज्ञो महानासीद्बुधाश्रमसमीपतः}
{रुद्रश्च परमं तोषमाजगाम महायशाः} %||7-81-15||

\twolineshloka
{अथ यज्ञसमाप्तौ तु प्रीतः परमया मुदा}
{उमापतिर्द्विजान्सर्वानुवाचेदमिलां प्रति} %||7-81-16||

\twolineshloka
{प्रीतोऽस्मि हयमेधेन भक्त्या च द्विजसत्तमाः}
{अस्य बाह्लिपतेश्चैव किं करोमि प्रियं शुभम्} %||7-81-17||

\twolineshloka
{तथा वदति देवेशे द्विजास्ते सुसमाहिताः}
{प्रसादयन्ति देवेशं यथा स्यात्पुरुषस्त्विला} %||7-81-18||

\twolineshloka
{ततः प्रीतमना रुद्रः पुरुषत्वं ददौ पुनः}
{इलायै सुमहातेजा दत्त्वा चान्तरधीयत} %||7-81-19||

\twolineshloka
{निवृत्ते हयमेधे तु गते चादर्शनं हरे}
{यथागतं द्विजाः सर्वे अगच्छन्दीर्घदर्शिनः} %||7-81-20||

\twolineshloka
{राजा तु बाह्लिमुत्सृज्य मध्यदेशे ह्यनुत्तमम्}
{निवेशयामास पुरं प्रतिष्ठानं यशस्करम्} %||7-81-21||

\twolineshloka
{शशबिन्दुस्तु राजासीद्बाह्ल्यां परपुरंजयः}
{प्रतिष्ठान इलो राजा प्रजापतिसुतो बली} %||7-81-22||

\twolineshloka
{स काले प्राप्तवाँल्लोकमिलो ब्राह्ममनुत्तमम्}
{ऐलः पुरूरवा राजा प्रतिष्ठानमवाप्तवान्} %||7-81-23||

\twolineshloka
{ईदृशो ह्यश्वमेधस्य प्रभावः पुरुषर्षभौ}
{स्त्रीभूतः पौरुषं लेभे यच्चान्यदपि दुर्लभम्} %||7-82-24||


{॥इत्यार्षे श्रीमद्रामायणे वाल्मीकीये आदिकाव्ये उत्तरकाण्डे एकाशीतितमः सर्गः॥}

\sect{द्व्यशीतितमः सर्गः}

\twolineshloka
{एतदाख्याय काकुत्स्थो भ्रातृह्याममितप्रभः}
{लक्ष्मणं पुनारेवाह धर्मयुक्तमिदं वचः} %||7-82-1||

\twolineshloka
{वसिष्ठं वामदेवं च जाबालिमथ कश्यपम्}
{द्विजांश्च सर्वप्रवरानश्वमेधपुरस्कृतान्} %||7-82-2||

\twolineshloka
{एतान्सर्वान्समाहूय मन्त्रयित्वा च लक्ष्मण}
{हयं लक्ष्मणसम्पन्नं विमोक्ष्यामि समाधिना} %||7-82-3||

\twolineshloka
{तद्वाक्यं राघवेणोक्तं श्रुत्वा त्वरितविक्रमः}
{द्विजान्सर्वान्समाहूय दर्शयामास राघवम्} %||7-82-4||

\twolineshloka
{ते दृष्ट्वा देवसङ्काशं कृतपादाभिवन्दनम्}
{राघवं सुदुराधर्षमाशीर्भिः समपूजयन्} %||7-82-5||

\twolineshloka
{प्राञ्जलिस्तु ततो भूत्वा राघवो द्विजसात्तमान्}
{उवाच धर्मसंयुक्तमश्वमेधाश्रितं वचः} %||7-82-6||

\twolineshloka
{स तेषां द्विजमुख्यानां वाक्यमद्भुतदर्शनम्}
{अश्वमेधाश्रितं श्रुत्वा भृशं प्रीतोऽभवत्तदा} %||7-82-7||

\twolineshloka
{विज्ञाय तु मतं तेषां रामो लक्ष्मणमब्रवीत्}
{प्रेषयस्व महाबाहो सुग्रीवाय महात्मने} %||7-82-8||

\twolineshloka
{शीघ्रं महद्भिर्हरिभिर्बहिभिश्च तदाश्रयैः}
{सार्धमागच्छ भद्रं ते अनुभोक्तुं मखोत्तमम्} %||7-82-9||

\twolineshloka
{विभीषणश्च रक्षोभिः कामगैर्बहुभिर्वृतः}
{अश्वमेधं महाबाहुः प्राप्नोतु लघुविक्रमः} %||7-82-10||

\twolineshloka
{राजानश्च नरव्याघ्र ये मे प्रियचिकीर्षवः}
{सानुगाः क्षिप्रमायान्तु यज्ञभूमिमनुत्तमाम्} %||7-82-11||

\twolineshloka
{देशान्तरगता ये च द्विजा धर्मपरायणाः}
{निमन्त्रयस्व तान्सर्वानश्वमेधाय लक्ष्मण} %||7-82-12||

\twolineshloka
{ऋषयश्चा महाबाहो आहूयन्तां तपोधनाः}
{देशान्तरगता ये च सदाराश्च महर्षयः} %||7-82-13||

\twolineshloka
{यज्ञवाटश्च सुमहान्गोमत्या नैमिषे वने}
{आज्ञाप्यतां महाबाहो तद्धि पुण्यमनुत्तमम्} %||7-82-14||

\twolineshloka
{शतं वाहसहस्राणां तण्डुलानां वपुष्मताम्}
{अयुतं तिलमुद्गस्य प्रयात्वग्रे महाबल} %||7-82-15||

\twolineshloka
{सुवर्णकोट्यो बहुला हिरण्यस्य शतोत्तराः}
{अग्रतो भरतः कृत्वा गच्छत्वग्रे महामतिः} %||7-82-16||

\twolineshloka
{अन्तरापणवीथ्यश्च सर्वांश्च नटनर्तकान्}
{नैगमान्बालवृद्धांश्च द्विजांश्च सुसमाहितान्} %||7-82-17||

\twolineshloka
{कर्मान्तिकांश्च कुशलाञ्शिल्पिनश्च सुपण्डितान्}
{मातरश्चैव मे सर्वाः कुमारान्तःपुराणि च} %||7-82-18||

\twolineshloka
{काञ्चनीं मम पत्नीं च दीक्षार्हां यज्ञकर्मणि}
{अग्रतो भरतः कृत्वा गच्छत्वग्रे महामतिः} %||7-83-19||


{॥इत्यार्षे श्रीमद्रामायणे वाल्मीकीये आदिकाव्ये उत्तरकाण्डे द्व्यशीतितमः सर्गः॥}

\sect{त्र्यशीतितमः सर्गः}

\twolineshloka
{तत्सर्वमखिलेनाशु प्रस्थाप्य भरताग्रजः}
{हयं लक्ष्मणसम्पन्नं कृष्णसारं मुमोच ह} %||7-83-1||

\twolineshloka
{ऋत्विग्भिर्लक्ष्मणं सार्धमश्वे च विनियुज्य सः}
{ततोऽभ्यगच्छत्काकुत्स्थः सह सैन्येन नैमिषम्} %||7-83-2||

\twolineshloka
{यज्ञवाटं महाबाहुर्दृष्ट्वा परममद्भुतम्}
{प्रहर्षमतुलं लेभे श्रीमानिति च सोऽब्रवीत्} %||7-83-3||

\twolineshloka
{नैमिषे वसतस्तस्य सर्व एव नराधिपाः}
{आजग्मुः सर्वराष्ट्रेभ्यस्तान्रामः प्रत्यपूजयत्} %||7-83-4||

\twolineshloka
{उपकार्यान्महार्हांश्च पार्थिवानां महात्मनाम्}
{सानुगानां नरश्रेष्ठो व्यादिदेश महाद्युतिः} %||7-83-5||

\twolineshloka
{अन्नपानानि वस्त्राणि सानुगानां महात्मनाम्}
{भरतः सन्ददावाशु शत्रुघ्नसहितस्तदा} %||7-83-6||

\twolineshloka
{वानराश्च महात्मानः सुग्रीवसहितास्तदा}
{विप्राणां प्रणताः सर्वे चक्रिरे परिवेषणम्} %||7-83-7||

\twolineshloka
{विभीषणश्च रक्षोभिः स्रग्विभिर्बहुभिर्वृतः}
{ऋषीणामुग्रतपसां किङ्करः पर्युपस्थितः} %||7-83-8||

\twolineshloka
{एवं सुविहितो यज्ञो हयमेधोऽभ्यवर्तत}
{लक्ष्मणेनाभिगुप्ता च हयचर्या प्रवर्तिता} %||7-83-9||

\threelineshloka
{नान्यः शब्दोऽभवत्तत्र हयमेधे महात्मनः}
{छन्दतो देहि विस्रब्धो यावत्तुष्यन्ति याचकाः}
{तावद्वानररक्षोभिर्दत्तमेवाभ्यदृश्यत} %||7-83-10||

\twolineshloka
{न कश्चिन्मलिनस्तत्र दीनो वाप्यथ वा कृशः}
{तस्मिन्यज्ञवरे राज्ञो हृष्टपुष्टजनावृते} %||7-83-11||

\twolineshloka
{ये च तत्र महात्मानो मुनयश्चिरजीविनः}
{नास्मरंस्तादृशं यज्ञं दानौघसमलङ्कृतम्} %||7-83-12||

\twolineshloka
{रजतानां सुवर्णानां रत्नानामथ वाससाम्}
{अनिशं दीयमानानां नान्तः समुपदृश्यते} %||7-83-13||

\twolineshloka
{न शक्रस्य न सोमस्य यमस्य वरुणस्य वा}
{ईदृशो दृष्टपूर्वो न एवमूचुस्तपोधनाः} %||7-83-14||

\twolineshloka
{सर्वत्र वानरास्तस्थुः सर्वत्रैव च राक्षसाः}
{वासो धनानि कामिभ्यः पूर्णहस्ता ददुर्भृशम्} %||7-83-15||

\twolineshloka
{ईदृशो राजसिंहस्य यज्ञः सर्वगुणान्वितः}
{संवत्सरमथो साग्रं वर्तते न च हीयते} %||7-84-16||


{॥इत्यार्षे श्रीमद्रामायणे वाल्मीकीये आदिकाव्ये उत्तरकाण्डे त्र्यशीतितमः सर्गः॥}

\sect{चतुरशीतितमः सर्गः}

\twolineshloka
{वर्तमाने तथाभूते यज्ञे परमकेऽद्भुते}
{सशिष्य आजगामाशु वाल्मीकिर्मुनिपुङ्गवः} %||7-84-1||

\twolineshloka
{स दृष्ट्वा दिव्यसङ्काशं यज्ञमद्भुतदर्शनम्}
{एकान्ते ऋषिवाटानां चकार उटजाञ्शुभान्} %||7-84-2||

\twolineshloka
{स शिष्यावब्रवीद्धृष्टो युवां गत्वा समाहितौ}
{कृत्स्नं रामायणं काव्यं गायतां परया मुदा} %||7-84-3||

\twolineshloka
{ऋषिवाटेषु पुण्येषु ब्राह्मणावसथेषु च}
{रथ्यासु राजमार्गेषु पार्थिवानां गृहेषु च} %||7-84-4||

\twolineshloka
{रामस्य भवनद्वारि यत्र कर्म च वर्तते}
{ऋत्विजामग्रतश्चैव तत्र गेयं विशेषतः} %||7-84-5||

\twolineshloka
{इमानि च फलान्यत्र स्वादूनि विविधानि च}
{जातानि पर्वताग्रेषु आस्वाद्यास्वाद्य गीयताम्} %||7-84-6||

\twolineshloka
{न यास्यथः श्रमं वत्सौ भक्षयित्वा फलानि वै}
{मूलानि च सुमृष्टानि नगरात्परिहास्यथ} %||7-84-7||

\twolineshloka
{यदि शब्दापयेद्रामः श्रवणाय महीपतिः}
{ऋषीणामुपविष्टानां ततो गेयं प्रवर्तताम्} %||7-84-8||

\twolineshloka
{दिवसे विंशतिः सर्गा गेया वै परया मुदा}
{प्रमाणैर्बहुभिस्तत्र यथोद्दिष्टं मया पुरा} %||7-84-9||

\twolineshloka
{लोभश्चापि न कर्तव्यः स्वल्पोऽपि धनकाङ्क्षया}
{किं धनेनाश्रमस्थानां फलमूलोपभोगिनाम्} %||7-84-10||

\twolineshloka
{यदि पृच्छेत्स काकुत्स्थो युवां कस्येति दारकौ}
{वाल्मीकेरथ शिष्यौ हि ब्रूतामेवं नराधिपम्} %||7-84-11||

\twolineshloka
{इमास्तन्त्रीः सुमधुराः स्थानं वा पूर्वदर्शितम्}
{मूर्छयित्वा सुमधुरं गायेतां विगतज्वरौ} %||7-84-12||

\twolineshloka
{आदिप्रभृति गेयं स्यान्न चावज्ञाय पार्थिवम्}
{पिता हि सर्वभूतानां राजा भवति धर्मतः} %||7-84-13||

\twolineshloka
{तद्युवां हृष्टमनसौ श्वः प्रभाते समाधिना}
{गायेतां मधुरं गेयं तन्त्रीलयसमन्वितम्} %||7-84-14||

\twolineshloka
{इति सन्दिश्य बहुशो मुनिः प्राचेतसस्तदा}
{वाल्मीकिः परमोदारस्तूष्णीमासीन्महायशाः} %||7-84-15||

\fourlineindentedshloka
{तामद्भुतां तौ हृदये कुमारौ}
{ निवेश्य वाणीमृषिभाषितां शुभाम्}
{समुत्सुकौ तौ सुखमूषतुर्निशाम्}
{ यथाश्विनौ भार्गवनीतिसंस्कृतौ} %||7-85-16||


{॥इत्यार्षे श्रीमद्रामायणे वाल्मीकीये आदिकाव्ये उत्तरकाण्डे चतुरशीतितमः सर्गः॥}

\sect{पञ्चाशीतितमः सर्गः}

\twolineshloka
{तौ रजन्यां प्रभातायां स्नातौ हुतहुताशनौ}
{यथोक्तमृषिणा पूर्वं तत्र तत्राभ्यगायताम्} %||7-85-1||

\twolineshloka
{तां स शुश्राव काकुत्स्थः पूर्वचर्यां ततस्ततः}
{अपूर्वां पाठ्य जातिं च गेयेन समलङ्कृताम्} %||7-85-2||

\twolineshloka
{प्रमाणैर्बहुभिर्बद्धां तन्त्रीलयसमन्विताम्}
{बालाभ्यां राघवः श्रुत्वा कौतूहलपरोऽभवत्} %||7-85-3||

\twolineshloka
{अथ कर्मान्तरे राजा समानीय महामुनीन्}
{पार्थिवांश्च नरव्याघ्रः पण्डितान्नैगमांस्तथा} %||7-85-4||

\twolineshloka
{पौराणिकाञ्शब्दवितो ये च वृद्धा द्विजातयः}
{एतान्सर्वान्समानीय गातारौ समवेशयत्} %||7-85-5||

\twolineshloka
{हृष्टा ऋषिगणास्तत्र पार्थिवाश्च महौजसः}
{पिबन्त इव चक्षुर्भ्यां राजानं गायकौ च तौ} %||7-85-6||

\twolineshloka
{परस्परमथोचुस्ते सर्व एव समं ततः}
{उभौ रामस्य सदृशौ बिम्बाद्बिम्बमिवोद्धृतौ} %||7-85-7||

\twolineshloka
{जटिलौ यदि न स्यातां न वल्कलधरौ यदि}
{विशेषं नाधिगच्छामो गायतो राघवस्य च} %||7-85-8||

\twolineshloka
{तेषां संवदतामेवं श्रोतॄणां हर्षवर्धनम्}
{गेयं प्रचक्रतुस्तत्र तावुभौ मुनिदारकौ} %||7-85-9||

\twolineshloka
{ततः प्रवृत्तं मधुरं गान्धर्वमतिमानुषम्}
{न च तृप्तिं ययुः सर्वे श्रोतारो गेय सम्पदा} %||7-85-10||

\twolineshloka
{प्रवृत्तमादितः पूर्वं सर्गान्नारददर्शनात्}
{ततः प्रभृति सर्गांश्च यावद्विंशत्यगायताम्} %||7-85-11||

\twolineshloka
{ततोऽपराह्णसमये राघवः समभाषत}
{श्रुत्वा विंशतिसर्गांस्तान्भरतं भ्रातृवत्सलः} %||7-85-12||

\twolineshloka
{अष्टादश सहस्राणि सुवर्णस्य महात्मनोः}
{ददस्व शीघ्रं काकुत्स्थ बालयोर्मा वृथा श्रमः} %||7-85-13||

\twolineshloka
{दीयमानं सुवर्णं तन्नागृह्णीतां कुशीलवौ}
{ऊचतुश्च महात्मानौ किमनेनेति विस्मितौ} %||7-85-14||

\twolineshloka
{वन्येन फलमूलेन निरतु स्वो वनौकसौ}
{सुवर्णेन हिरण्येन किं करिष्यावहे वने} %||7-85-15||

\twolineshloka
{तथा तयोः प्रब्रुवतोः कौतूहलसमन्विताः}
{श्रोतारश्चैव रामश्च सर्व एव सुविस्मिताः} %||7-85-16||

\twolineshloka
{तस्य चैवागमं रामः काव्यस्य श्रोतुमुत्सुकः}
{पप्रच्छ तौ महातेजास्तावुभौ मुनिदारकौ} %||7-85-17||

\twolineshloka
{किम्प्रमाणमिदं काव्यं का प्रतिष्ठा महात्मनः}
{कर्ता काव्यस्य महतः को वासौ मुनिपुङ्गवः} %||7-85-18||

\threelineshloka
{पृच्छन्तं राघवं वाक्यमूचतुर्मुनिदारकौ}
{वाल्मीकिर्भगवान्कर्ता सम्प्राप्तो यज्ञसंनिधिम्}
{येनेदं चरितं तुभ्यमशेषं सम्प्रदर्शितम्} %||7-85-19||

\twolineshloka
{आदिप्रभृति राजेन्द्र पञ्चसर्ग शतानि च}
{प्रतिष्ठा जीवितं यावत्तावद्राजञ्शुभाशुभम्} %||7-85-20||

\twolineshloka
{यदि बुद्धिः कृता राजञ्श्रवणाय महारथ}
{कर्मान्तरे क्षणी हूतस्तच्छृणुष्व सहानुजः} %||7-85-21||

\twolineshloka
{बाढमित्यब्रवीद्रामस्तौ चानुज्ञाप्य राघवम्}
{प्रहृष्टौ जग्मतुर्वासं यत्रासौ मुनिपुङ्गवः} %||7-85-22||

\twolineshloka
{रामोऽपि मुनिभिः सार्धं पार्थिवैश्च महात्मभिः}
{श्रुत्वा तद्गीतमाधुर्यं कर्मशालामुपागमत्} %||7-86-23||


{॥इत्यार्षे श्रीमद्रामायणे वाल्मीकीये आदिकाव्ये उत्तरकाण्डे पञ्चाशीतितमः सर्गः॥}

\sect{षडशीतितमः सर्गः}

\twolineshloka
{रामो बहून्यहान्येव तद्गीतं परमाद्भुतम्}
{शुश्राव मुनिभिः सार्धं राजभिः सह वानरैः} %||7-86-1||

\twolineshloka
{तस्मिन्गीते तु विज्ञाय सीतापुत्रौ कुशीलवौ}
{तस्याः परिषदो मध्ये रामो वचनमब्रवीत्} %||7-86-2||

{मद्वचो ब्रूत गच्छध्वमिति भगवतोऽन्तिकम्॥\devanumber{3}॥} %||7-86-3|| (Check)

\addtocounter{shlokacount}{1}

\twolineshloka
{यदि शुद्धसमाचारा यदि वा वीतकल्मषा}
{करोत्विहात्मनः शुद्धिमनुमान्य महामुनिम्} %||7-86-4||

\twolineshloka
{छन्दं मुनेस्तु विज्ञाय सीतायाश्च मनोगतम्}
{प्रत्ययं दातुकामायास्ततः शंसत मे लघु} %||7-86-5||

\twolineshloka
{श्वः प्रभाते तु शपथं मैथिली जनकात्मजा}
{करोतु परिषन्मध्ये शोधनार्थं ममेह च} %||7-86-6||

\twolineshloka
{श्रुत्वा तु राघवस्यैतद्वचः परममद्भुतम्}
{दूताः सम्प्रययुर्वाटं यत्रास्ते मुनिपुङ्गवः} %||7-86-7||

\twolineshloka
{ते प्रणम्य महात्मानं ज्वलन्तममितप्रभम्}
{ऊचुस्ते राम वाक्यानि मृदूनि मधुराणि च} %||7-86-8||

\twolineshloka
{तेषां तद्भाषितं श्रुत्वा रामस्य च मनोगतम्}
{विज्ञाय सुमहातेजा मुनिर्वाक्यमथाब्रवीत्} %||7-86-9||

\twolineshloka
{एवं भवतु भद्रं वो यथा तुष्यति राघवः}
{तथा करिष्यते सीता दैवतं हि पतिः स्त्रियाः} %||7-86-10||

\twolineshloka
{तथोक्ता मुनिना सर्वे रामदूता महौजसः}
{प्रत्येत्य राघवं सर्वे मुनिवाक्यं बभाषिरे} %||7-86-11||

\twolineshloka
{ततः प्रहृष्टः काकुत्स्थः श्रुत्वा वाक्यं महात्मनः}
{ऋषींस्तत्र समेतांश्च राज्ञश्चैवाभ्यभाषत} %||7-86-12||

\twolineshloka
{भगवन्तः सशिष्या वै सानुगश्च नराधिपाः}
{पश्यन्तु सीताशपथं यश्चैवान्योऽभिकाङ्क्षते} %||7-86-13||

\twolineshloka
{तस्य तद्वचनं श्रुत्वा राघवस्य महात्मनः}
{सर्वेषमृषिमुख्यानां साधुवादो महानभूत्} %||7-86-14||

\twolineshloka
{राजानश्च महात्मानः प्रशंसन्ति स्म राघवम्}
{उपपन्नं नरश्रेष्ठ त्वय्येव भुवि नान्यतः} %||7-86-15||

\twolineshloka
{एवं विनिश्चयं कृत्वा श्वोभूत इति राघवः}
{विसर्जयामास तदा सर्वांस्ताञ्शत्रुसूदनः} %||7-87-16||


{॥इत्यार्षे श्रीमद्रामायणे वाल्मीकीये आदिकाव्ये उत्तरकाण्डे षडशीतितमः सर्गः॥}

\sect{सप्ताशीतितमः सर्गः}

\twolineshloka
{तस्यां रजन्यां व्युष्टायां यज्ञवाटगतो नृपः}
{ऋषीन्सर्वान्महातेजाः शब्दापयति राघवः} %||7-87-1||

\twolineshloka
{वसिष्ठो वामदेवश्च जाबालिरथ काश्यपः}
{विश्वामित्रो दीर्घतपा दुर्वासाश्च महातपाः} %||7-87-2||

\twolineshloka
{अगस्त्योऽथ तथाशक्तिर्भार्गवश्चैव वामनः}
{मार्कण्डेयश्च दीर्घायुर्मौद्गल्यश्च महातपाः} %||7-87-3||

\twolineshloka
{भार्गवश्च्यवनश्चैव शतानन्दश्च धर्मवित्}
{भरद्वाजश्च तेजस्वी अग्निपुत्रश्च सुप्रभः} %||7-87-4||

\twolineshloka
{एते चान्ये च मुनयो बहवः संशितव्रताः}
{राजानश्च नरव्याघ्राः सर्व एव समागताः} %||7-87-5||

\twolineshloka
{राक्षसाश्च महावीर्या वानराश्च महाबलाः}
{समाजग्मुर्महात्मानः सर्व एव कुतूहलात्} %||7-87-6||

\twolineshloka
{क्षत्रियाश्चैव वैश्याश्च शूद्राश्चैव सहस्रशः}
{सीताशपथवीक्षार्थं सर्व एव समागताः} %||7-87-7||

\twolineshloka
{तथा समागतं सर्वमश्वभूतमिवाचलम्}
{श्रुत्वा मुनिवरस्तूर्णं ससीतः समुपागमत्} %||7-87-8||

\twolineshloka
{तमृषिं पृष्ठतः सीता सान्वगच्छदवाङ्मुखी}
{कृताञ्जलिर्बाष्पगला कृत्वा रामं मनोगतम्} %||7-87-9||

\twolineshloka
{तां दृष्ट्वा श्रीमिवायान्तीं ब्रह्माणमनुगामिनीम्}
{वाल्मीकेः पृष्ठतः सीतां साधुकारो महानभूत्} %||7-87-10||

\twolineshloka
{ततो हलहला शब्दः सर्वेषामेवमाबभौ}
{दुःखजेन विशालेन शोकेनाकुलितात्मनाम्} %||7-87-11||

\twolineshloka
{साधु सीतेति केचित्तु साधु रामेति चापरे}
{उभावेव तु तत्रान्ये साधु साध्विति चाब्रुवन्} %||7-87-12||

\twolineshloka
{ततो मध्यं जनौघानां प्रविश्य मुनिपुङ्गवः}
{सीतासहायो वाल्मीकिरिति होवाच राघवम्} %||7-87-13||

\twolineshloka
{इयं दाशरथे सीता सुव्रता धर्मचारिणी}
{अपापा ते परित्यक्ता ममाश्रमसमीपतः} %||7-87-14||

\twolineshloka
{लोकापवादभीतस्य तव राम महाव्रत}
{प्रत्ययं दास्यते सीता तामनुज्ञातुमर्हसि} %||7-87-15||

\twolineshloka
{इमौ च जानकी पुत्रावुभौ च यमजातकौ}
{सुतौ तवैव दुर्धर्षो सत्यमेतद्ब्रवीमि ते} %||7-87-16||

\twolineshloka
{प्रचेतसोऽहं दशमः पुत्रो राघवनन्दन}
{न स्मराम्यनृतं वाक्यं तथेमौ तव पुत्रकौ} %||7-87-17||

\twolineshloka
{बहुवर्षसहस्राणि तपश्चर्या मया कृता}
{तस्याः फलमुपाश्नीयामपापा मैथिली यथा} %||7-87-18||

\twolineshloka
{अहं पञ्चसु भूतेषु मनःषष्ठेषु राघव}
{विचिन्त्य सीतां शुद्धेति न्यगृह्णां वननिर्झरे} %||7-87-19||

\twolineshloka
{इयं शुद्धसमाचारा अपापा पतिदेवता}
{लोकापवादभीतस्य दास्यति प्रत्ययं तव} %||7-88-20||


{॥इत्यार्षे श्रीमद्रामायणे वाल्मीकीये आदिकाव्ये उत्तरकाण्डे सप्ताशीतितमः सर्गः॥}

\sect{अष्टाशीतितमः सर्गः}

\twolineshloka
{वाल्मीकिनैवमुक्तस्तु राघवः प्रत्यभाषत}
{प्राञ्जलिर्जगतो मध्ये दृष्ट्वा तां देववर्णिनीम्} %||7-88-1||

\twolineshloka
{एवमेतन्महाभाग यथा वदसि धर्मवित्}
{प्रत्ययो हि मम ब्रह्मंस्तव वाक्यैरकल्मषैः} %||7-88-2||

\threelineshloka
{प्रत्ययो हि पुरा दत्तो वैदेह्या सुरसंनिधौ}
{सेयं लोकभयाद्ब्रह्मन्नपापेत्यभिजानता}
{परित्यक्ता मया सीता तद्भवान्क्षन्तुमर्हति} %||7-88-3||

\twolineshloka
{जानामि चेमौ पुत्रौ मे यमजातौ कुशीलवौ}
{शुद्धायां जगतो मध्ये मैथिल्यां प्रीतिरस्तु मे} %||7-88-4||

\twolineshloka
{अभिप्रायं तु विज्ञाय रामस्य सुरसत्तमाः}
{पितामहं पुरस्कृत्य सर्व एव समागताः} %||7-88-5||

\threelineshloka
{आदित्या वसवो रुद्रा विश्वे देशा मरुद्गणाः}
{अश्विनावृषिगन्धर्वा अप्सराणां गणास्तथा}
{साध्याश्च देवाः सर्वे ते सर्वे च परमर्षयः} %||7-88-6||

\twolineshloka
{ततो वायुः शुभः पुण्यो दिव्यगन्धो मनोरमः}
{तं जनौघं सुरश्रेष्ठो ह्लादयामास सर्वतः} %||7-88-7||

\twolineshloka
{तदद्भुतमिवाचिन्त्यं निरीक्षन्ते समाहिताः}
{मानवाः सर्वराष्ट्रेभ्यः पूर्वं कृतयुगे यथा} %||7-88-8||

\twolineshloka
{सर्वान्समागतान्दृष्ट्वा सीता काषायवासिनी}
{अब्रवीत्प्राञ्जलिर्वाक्यमधोदृष्टिरवान्मुखी} %||7-88-9||

\twolineshloka
{यथाहं राघवादन्यं मनसापि न चिन्तये}
{तथा मे माधवी देवी विवरं दातुमर्हति} %||7-88-10||

\twolineshloka
{तथा शपन्त्यां वैदेह्यां प्रादुरासीत्तदद्भुतम्}
{भूतलादुत्थितं दिव्यं सिंहासनमनुत्तमम्} %||7-88-11||

\twolineshloka
{ध्रियमाणं शिरोभिस्तन्नागैरमितविक्रमैः}
{दिव्यं दिव्येन वपुषा सर्वरत्नविभूषितम्} %||7-88-12||

\twolineshloka
{तस्मिंस्तु धरणी देवी बाहुभ्यां गृह्य मैथिलीम्}
{स्वागतेनाभिनन्द्यैनामासने चोपवेषयत्} %||7-88-13||

\twolineshloka
{तामासनगतां दृष्ट्वा प्रविशन्तीं रसातलम्}
{पुण्यवृष्टिरविच्छिन्ना दिव्या सीतामवाकिरत्} %||7-88-14||

\twolineshloka
{साधुकारश्च सुमहान्देवानां सहसोत्थितः}
{साधु साध्विति वै सीते यस्यास्ते शीलमीदृशम्} %||7-88-15||

\twolineshloka
{एवं बहुविधा वाचो ह्यन्तरिक्षगताः सुराः}
{व्याजह्रुर्हृष्टमनसो दृष्ट्वा सीताप्रवेशनम्} %||7-88-16||

\twolineshloka
{यज्ञवाटगताश्चापि मुनयः सर्व एव ते}
{राजानश्च नरव्याघ्रा विस्मयान्नोपरेमिरे} %||7-88-17||

\twolineshloka
{अन्तरिक्षे च भूमौ च सर्वे स्थावरजङ्गमाः}
{दानवाश्च महाकायाः पाताले पन्नगाधिपाः} %||7-88-18||

\twolineshloka
{केचिद्विनेदुः संहृष्टाः केचिद्ध्यानपरायणाः}
{केचिद्रामं निरीक्षन्ते केचित्सीतामचेतनाः} %||7-88-19||

\twolineshloka
{सीताप्रवेशनं दृष्ट्वा तेषामासीत्समागमः}
{तं मुहूर्तमिवात्यर्थं सर्वं सम्मोहितं जगत्} %||7-89-20||


{॥इत्यार्षे श्रीमद्रामायणे वाल्मीकीये आदिकाव्ये उत्तरकाण्डे अष्टाशीतितमः सर्गः॥}

\sect{एकोननवतितमः सर्गः}

\threelineshloka
{तदावसाने यज्ञस्य रामः परमदुर्मनाः}
{अपश्यमानो वैदेहीं मेने शून्यमिदं जगत्}
{शोकेन परमायत्तो न शान्तिं मनसागमत्} %||7-89-1||

\twolineshloka
{विसृज्य पार्थिवान्सर्वानृक्षवानरराक्षसान्}
{जनौघं ब्रह्ममुख्यानां वित्तपूर्णं व्यसर्जयत्} %||7-89-2||

\twolineshloka
{ततो विसृज्य तान्सर्वान्रामो राजीवलोचनः}
{हृदि कृत्वा तदा सीतामयोध्यां प्रविवेश सः} %||7-89-3||

\twolineshloka
{न सीतायाः परां भार्यां वव्रे स रघुनन्दनः}
{यज्ञे यज्ञे च पत्न्यर्थं जानकी काञ्चनी भवत्} %||7-89-4||

\twolineshloka
{दशवर्षसहस्राणि वाजिमेधमुपाकरोत्}
{वाजपेयान्दशगुणांस्तथा बहुसुवर्णकान्} %||7-89-5||

\twolineshloka
{अग्निष्टोमातिरात्राभ्यां गोसवैश्च महाधनैः}
{ईजे क्रतुभिरन्यैश्च स श्रीमानाप्तदक्षिणैः} %||7-89-6||

\twolineshloka
{एवं स कालः सुमहान्राज्यस्थस्य महात्मनः}
{धर्मे प्रयतमानस्य व्यतीयाद्राघवस्य तु} %||7-89-7||

\twolineshloka
{ऋक्षवानररक्षांसि स्थिता रामस्य शासने}
{अनुरज्यन्ति राजानो अहन्यहनि राघवम्} %||7-89-8||

\twolineshloka
{काले वर्षति पर्जन्यः सुभिक्षं विमला दिशः}
{हृष्टपुष्टजनाकीर्णं पुरं जनपदस्तथा} %||7-89-9||

\twolineshloka
{नाकाले म्रियते कश्चिन्न व्याधिः प्राणिनां तदा}
{नाधर्मश्चाभवत्कश्चिद्रामे राज्यं प्रशासति} %||7-89-10||

\twolineshloka
{अथ दीर्घस्य कालस्य राममाता यशस्विनी}
{पुत्रपौत्रैः परिवृता कालधर्ममुपागमत्} %||7-89-11||

\twolineshloka
{अन्वियाय सुमित्रापि कैकेयी च यशस्विनी}
{धर्मं कृत्वा बहुविधं त्रिदिवे पर्यवस्थिता} %||7-89-12||

\twolineshloka
{सर्वाः प्रतिष्ठिताः स्वर्गे राज्ञा दशरथेन च}
{समागता महाभागाः सह धर्मं च लेभिरे} %||7-89-13||

\twolineshloka
{तासां रामो महादानं काले काले प्रयच्छति}
{मातॄणामविशेषेण ब्राह्मणेषु तपस्विषु} %||7-89-14||

\twolineshloka
{पित्र्याणि बहुरत्नानि यज्ञान्परमदुस्तरान्}
{चकार रामो धर्मात्मा पितॄन्देवान्विवर्धयन्} %||7-90-15||


{॥इत्यार्षे श्रीमद्रामायणे वाल्मीकीये आदिकाव्ये उत्तरकाण्डे एकोननवतितमः सर्गः॥}

\sect{नवतितमः सर्गः}

\twolineshloka
{कस्यचित्त्वथ कालस्य युधाजित्केकयो नृपः}
{स्वगुरुं प्रेषयामास राघवाय महात्मने} %||7-90-1||

\twolineshloka
{गार्ग्यमङ्गिरसः पुत्रं ब्रह्मर्षिममितप्रभम्}
{दश चाश्वसहस्राणि प्रीतिदानमनुत्तमम्} %||7-90-2||

\twolineshloka
{कम्बलानि च रत्नानि चित्रवस्त्रमथोत्तमम्}
{रामाय प्रददौ राजा बहून्याभरणानि च} %||7-90-3||

\twolineshloka
{श्रुत्वा तु राघवो गार्ग्यं महर्षिं समुपागतम्}
{मातुलस्याश्वपतिनः प्रियं दूतमुपागतम्} %||7-90-4||

\twolineshloka
{प्रत्युद्गम्य च काकुत्स्थः क्रोशमात्रं सहानुगः}
{गार्ग्यं सम्पूजयामास धनं तत्प्रतिगृह्य च} %||7-90-5||

\twolineshloka
{पृष्ट्वा च प्रीतिदं सर्वं कुशलं मातुलस्य च}
{उपविष्टं महाभागं रामः प्रष्टुं प्रचक्रमे} %||7-90-6||

\twolineshloka
{किमाह मतुलो वाक्यं यदर्थं भगवानिह}
{प्राप्तो वाक्यविदां श्रेष्ठ साक्षादिव बृहस्पतिः} %||7-90-7||

\twolineshloka
{रामस्य भाषितं श्रुत्वा ब्रह्मर्षिः कार्यविस्तरम्}
{वक्तुमद्भुतसङ्काशं राघवायोपचक्रमे} %||7-90-8||

\twolineshloka
{मातुलस्ते महाबाहो वाक्यमाह नरर्षभ}
{युधाजित्प्रीतिसंयुक्तं श्रूयतां यदि रोचते} %||7-90-9||

\twolineshloka
{अयं गन्धर्वविषयः फलमूलोपशोभितः}
{सिन्धोरुभयतः पार्श्वे देशः परमशोभनः} %||7-90-10||

\twolineshloka
{तं च रक्षन्ति गन्धर्वाः सायुधा युद्धकोविदाः}
{शैलूषस्य सुता वीरास्तिस्रः कोट्यो महाबलाः} %||7-90-11||

\twolineshloka
{तान्विनिर्जित्य काकुत्स्थ गन्धर्वविषयं शुभम्}
{निवेशय महाबाहो द्वे पुरे सुसमाहितः} %||7-90-12||

\twolineshloka
{अन्यस्य न गतिस्तत्र देशश्चायं सुशोभनः}
{रोचतां ते महाबाहो नाहं त्वामनृतं वदे} %||7-90-13||

\twolineshloka
{तच्छ्रुत्वा राघवः प्रीतो महर्षेर्मातुलस्य च}
{उवाच बाढमित्येवं भरतं चान्ववैक्षत} %||7-90-14||

\twolineshloka
{सोऽब्रवीद्राघवः प्रीतः प्राञ्जलिप्रग्रहो द्विजम्}
{इमौ कुमारौ तं देशं ब्रह्मर्षे विजयिष्यतः} %||7-90-15||

\twolineshloka
{भरतस्यात्मजौ वीरौ तक्षः पुष्कल एव च}
{मातुलेन सुगुप्तौ तौ धर्मेण च समाहितौ} %||7-90-16||

\twolineshloka
{भरतं चाग्रतः कृत्वा कुमारौ सबलानुगौ}
{निहत्य गन्धर्वसुतान्द्वे पुरे विभजिष्यतः} %||7-90-17||

\twolineshloka
{निवेश्य ते पुरवरे आत्माजौ संनिवेश्य च}
{आगमिष्यति मे भूयः सकाशमतिधार्मिकः} %||7-90-18||

\twolineshloka
{ब्रह्मर्षिमेवमुक्त्वा तु भरतं सबलानुगम्}
{आज्ञापयामास तदा कुमारौ चाभ्यषेचयत्} %||7-90-19||

\twolineshloka
{नक्षत्रेण च सौम्येन पुरस्कृत्याङ्गिरः सुतम्}
{भरतः सह सैन्येन कुमाराभ्यां च निर्ययौ} %||7-90-20||

\twolineshloka
{सा सेना शक्रयुक्तेव नरगान्निर्ययावथ}
{राघवानुगता दूरं दुराधर्षा सुरासुरैः} %||7-90-21||

\twolineshloka
{मांसाशीनि च सत्त्वानि रक्षांसि सुमहान्ति च}
{अनुजग्मुश्च भरतं रुधिरस्य पिपासया} %||7-90-22||

\twolineshloka
{भूतग्रामाश्च बहवो मांसभक्षाः सुदारुणाः}
{गन्धर्वपुत्रमांसानि भोक्तुकामाः सहस्रशः} %||7-90-23||

\twolineshloka
{सिंहव्याघ्रसृगालानां खेचराणां च पक्षिणाम्}
{बहूनि वै सहस्राणि सेनाया ययुरग्रतः} %||7-90-24||

\twolineshloka
{अध्यर्धमासमुषिता पथि सेना निरामया}
{हृष्टपुष्टजनाकीर्णा केकयं समुपागमत्} %||7-91-25||


{॥इत्यार्षे श्रीमद्रामायणे वाल्मीकीये आदिकाव्ये उत्तरकाण्डे नवतितमः सर्गः॥}

\sect{एकनवतितमः सर्गः}

\twolineshloka
{श्रुत्वा सेनापतिं प्राप्तं भरतं केकयाधिपः}
{युधाजिद्गार्ग्यसहितं परां प्रीतिमुपागमत्} %||7-91-1||

\twolineshloka
{स निर्ययौ जनौघेन महता केकयाधिपः}
{त्वरमाणोऽभिचक्राम गन्धर्वान्देवरूपिणः} %||7-91-2||

\twolineshloka
{भरतश्च युधाजिच्च समेतौ लघुविक्रमौ}
{गन्धर्वनगरं प्राप्तौ सबलौ सपदानुगौ} %||7-91-3||

\twolineshloka
{श्रुत्वा तु भरतं प्राप्तं गन्धर्वास्ते समागताः}
{योद्धुकामा महावीर्या विनदन्तः समन्ततः} %||7-91-4||

\twolineshloka
{ततः समभवद्युद्धं तुमुलं लोमहर्षणम्}
{सप्तरात्रं महाभीमं न चान्यतरयोर्जयः} %||7-91-5||

\twolineshloka
{ततो रामानुजः क्रुद्धः कालस्यास्त्रं सुदारुणम्}
{संवर्तं नाम भरतो गन्धर्वेष्वभ्ययोजयत्} %||7-91-6||

\twolineshloka
{ते बद्धाः कालपाशेन संवर्तेन विदारिताः}
{क्षणेनाभिहतास्तिस्रस्तत्र कोट्यो महात्मना} %||7-91-7||

\twolineshloka
{तं घातं घोरसङ्काशं न स्मरन्ति दिवौकसः}
{निमेषान्तरमात्रेण तादृशानां महात्मनाम्} %||7-91-8||

\threelineshloka
{हतेषु तेषु वीरेषु भरतः कैकयीसुतः}
{निवेशयामास तदा समृद्धे द्वे पुरोत्तमे}
{तक्षं तक्षशिलायां तु पुष्करं पुष्करावतौ} %||7-91-9||

\twolineshloka
{गन्धर्वदेशो रुचिरो गान्धारविषयश्च सः}
{वर्षैः पञ्चभिराकीर्णो विषयैर्नागरैस्तथा} %||7-91-10||

\twolineshloka
{धनरत्नौघसम्पूर्णो काननैरुपशोभिते}
{अन्योन्यसङ्घर्षकृते स्पर्धया गुणविस्तरे} %||7-91-11||

\twolineshloka
{उभे सुरुचिरप्रख्ये व्यवहारैरकल्मषैः}
{उद्यानयानौघवृते सुविभक्तान्तरापणे} %||7-91-12||

\twolineshloka
{उभे पुरवरे रम्ये विस्तरैरुपशोभिते}
{गृहमुख्यैः सुरुचिरैर्विमानैः समवर्णिभिः} %||7-91-13||

\threelineshloka
{शोभिते शोभनीयैश्च देवायतनविस्तरैः}
{निवेश्य पञ्चभिर्वर्षैर्भरतो राघवानुजः}
{पुनरायान्महाबाहुरयोध्यां कैकयीसुतः} %||7-91-14||

\twolineshloka
{सोऽभिवाद्य महात्मानं साक्षाद्धर्ममिवापरम्}
{राघवं भरतः श्रीमान्ब्रह्माणमिव वासवः} %||7-91-15||

\twolineshloka
{शशंस च यथावृत्तं गन्धर्ववधमुत्तमम्}
{निवेशनं च देशस्य श्रुत्वा प्रीतोऽस्य राघवः} %||7-92-16||


{॥इत्यार्षे श्रीमद्रामायणे वाल्मीकीये आदिकाव्ये उत्तरकाण्डे एकनवतितमः सर्गः॥}

\sect{द्विनवतितमः सर्गः}

\twolineshloka
{तच्छ्रुत्वा हर्षमापेदे राघवो भ्रातृभिः सह}
{वाक्यं चाद्भुतसङ्काशं भ्रातॄन्प्रोवाच राघवः} %||7-92-1||

\twolineshloka
{इमौ कुमारौ सौमित्रे तव धर्मविशारदौ}
{अङ्गदश्चन्द्रकेतुश्च राज्यार्हौ दृढधन्विनौ} %||7-92-2||

\twolineshloka
{इमौ राज्येऽभिषेक्ष्यामि देशः साधु विधीयताम्}
{रमणीयो ह्यसम्बाधो रमेतां यत्र धन्विनौ} %||7-92-3||

\twolineshloka
{न राज्ञां यत्र पीदा स्यान्नाश्रमाणां विनाशनम्}
{स देशो दृश्यतां सौम्य नापराध्यामहे यथा} %||7-92-4||

\twolineshloka
{तथोक्तवति रामे तु भरतः प्रत्युवाच ह}
{अयं कारापथो देशः सुरमण्यो निरामयः} %||7-92-5||

\twolineshloka
{निवेश्यतां तत्र पुरमङ्गदस्य महात्मनः}
{चन्द्रकेतोश्च रुचिरं चन्द्रकान्तं निरामयम्} %||7-92-6||

\twolineshloka
{तद्वाक्यं भरतेनोक्तं प्रतिजग्राह राघवः}
{तं च कृता वशे देशमङ्गदस्य न्यवेशयत्} %||7-92-7||

\twolineshloka
{अङ्गदीया पुरी रम्या अङ्गदस्य निवेशिता}
{रमणीया सुगुप्ता च रामेणाक्लिष्टकर्मणा} %||7-92-8||

\twolineshloka
{चन्द्रकेतुस्तु मल्लस्य मल्लभूम्यां निवेशिता}
{चन्द्रकान्तेति विख्याता दिव्या स्वर्गपुरी यथा} %||7-92-9||

\twolineshloka
{ततो रामः परां प्रीतिं भरतो लक्ष्मणस्तथा}
{ययुर्युधि दुराधर्षा अभिषेकं च चक्रिरे} %||7-92-10||

\twolineshloka
{अभिषिच्य कुमारौ द्वौ प्रस्थाप्य सबलानुगौ}
{अङ्गदं पश्चिमा भूमिं चन्द्रकेतुमुदङ्मुखम्} %||7-92-11||

\twolineshloka
{अङ्गदं चापि सौमित्रिर्लक्ष्मणोऽनुजगाम ह}
{चन्द्रकेतोस्तु भरतः पार्ष्णिग्राहो बभूव ह} %||7-92-12||

\twolineshloka
{लक्ष्मणस्त्वङ्गदीयायां संवत्सरमथोषितः}
{पुत्रे स्थिते दुराधर्षे अयोध्यां पुनरागमत्} %||7-92-13||

\twolineshloka
{भरतोऽपि तथैवोष्य संवत्सरमथाधिकम्}
{अयोध्यां पुनरगम्य रामपादावुपागमत्} %||7-92-14||

\twolineshloka
{उभौ सौमित्रिभरतौ रामपादावनुव्रतौ}
{कालं गतमपि स्नेहान्न जज्ञातेऽतिधार्मिकौ} %||7-92-15||

\twolineshloka
{एवं वर्षसहस्राणि दशतेषां ययुस्तदा}
{धर्मे प्रयतमानानां पौरकार्येषु नित्यदा} %||7-92-16||

\fourlineindentedshloka
{विहृत्य लाकं परिपूर्णमानसाः}
{ श्रिया वृता धर्मपथे परे स्थिताः}
{त्रयः समिद्धा इव दीप्ततेजसा}
{ हुताग्नयः साधु महाध्वरे त्रयः} %||7-93-17||


{॥इत्यार्षे श्रीमद्रामायणे वाल्मीकीये आदिकाव्ये उत्तरकाण्डे द्विनवतितमः सर्गः॥}

\sect{त्रिनवतितमः सर्गः}

\twolineshloka
{कस्यचित्त्वथ कालस्य रामे धर्मपथे स्थिते}
{कालस्तापसरूपेण राजद्वारमुपागमत्} %||7-93-1||

\twolineshloka
{सोऽब्रवील्लक्ष्मणं वाक्यं धृतिमन्तं यशस्विनम्}
{मां निवेदय रामाय सम्प्राप्तं कार्यगौरवात्} %||7-93-2||

\twolineshloka
{दूतो ह्यतिबलस्याहं महर्षेरमितौजसः}
{रामं दिदृक्षुरायातः कार्येण हि महाबल} %||7-93-3||

\twolineshloka
{तस्य तद्वचनं श्रुत्वा सौमित्रिस्त्वरयान्वितः}
{न्यवेदयत रामाय तापसस्य विवक्षितम्} %||7-93-4||

\twolineshloka
{जयस्व राजन्धर्मेण उभौ लोकौ महाद्युते}
{दूतस्त्वां द्रष्टुमायातस्तपस्वी भास्करप्रभः} %||7-93-5||

\twolineshloka
{तद्वाक्यं लक्ष्मणेनोक्तं श्रुत्वा राम उवाच ह}
{प्रवेश्यतां मुनिस्तात महौजास्तस्य वाक्यधृक्} %||7-93-6||

\twolineshloka
{सौमित्रिस्तु तथेत्युक्त्वा प्रावेशयत तं मुनिम्}
{ज्वलन्तमिव तेजोभिः प्रदहन्तमिवांशुभिः} %||7-93-7||

\twolineshloka
{सोऽभिगम्य रघुश्रेष्ठं दीप्यमानं स्वतेजसा}
{ऋषिर्मधुरया वाचा वर्धस्वेत्याह राघवम्} %||7-93-8||

\twolineshloka
{तस्मै रामो महातेजाः पूजामर्घ्य पुरोगमाम्}
{ददौ कुशलमव्यग्रं प्रष्टुं चैवोपचक्रमे} %||7-93-9||

\twolineshloka
{पृष्ठश्च कुशलं तेन रामेण वदतां वरः}
{आसने काञ्चने दिव्ये निषसाद महायशाः} %||7-93-10||

\twolineshloka
{तमुवाच ततो रामः स्वागतं ते महामुने}
{प्रापयस्व च वाक्यानि यतो दूतस्त्वमागतः} %||7-93-11||

\twolineshloka
{चोदितो राजसिंहेन मुनिर्वाक्यमुदीरयत्}
{द्वन्द्वमेतत्प्रवक्तव्यं न च चक्षुर्हतं वचः} %||7-93-12||

\twolineshloka
{यः शृणोति निरीक्षेद्वा स वध्यस्तव राघव}
{भवेद्वै मुनिमुख्यस्य वचनं यद्यवेक्षसे} %||7-93-13||

\twolineshloka
{तथेति च प्रतिज्ञाय रामो लक्ष्मणमब्रवीत्}
{द्वारि तिष्ठ महाबाहो प्रतिहारं विसर्जय} %||7-93-14||

\twolineshloka
{स मे वध्यः खलु भवेत्कथां द्वन्द्वसमीरिताम्}
{ऋषेर्मम च सौमित्रे पश्येद्वा शृणुया च यः} %||7-93-15||

\twolineshloka
{ततो निक्षिप्य काकुत्स्थो लक्ष्मणं द्वारसङ्ग्रहे}
{तमुवाच मुनिं वाक्यं कथयस्वेति राघवः} %||7-93-16||

\twolineshloka
{यत्ते मनीषितं वाक्यं येन वासि समाहितः}
{कथयस्व विशङ्कस्त्वं ममापि हृदि वर्तते} %||7-94-17||


{॥इत्यार्षे श्रीमद्रामायणे वाल्मीकीये आदिकाव्ये उत्तरकाण्डे त्रिनवतितमः सर्गः॥}

\sect{चतुर्नवतितमः सर्गः}

\twolineshloka
{शृणु राम महाबाहो यदर्थमहमाहतः}
{पितामहेन देवेन प्रेषितोऽस्मि महाबल} %||7-94-1||

\twolineshloka
{तवाहं पूर्वके भावे पुत्रः परपुरंजय}
{मायासम्भावितो वीर कालः सर्वसमाहरः} %||7-94-2||

\twolineshloka
{पितामहश्च भगवानाह लोकपतिः प्रभुः}
{समयस्ते महाबाहो स्वर्लोकान्परिरक्षितुम्} %||7-94-3||

\twolineshloka
{सङ्क्षिप्य च पुरा लोकान्मायया स्वयमेव हि}
{महार्णवे शयानोऽप्सु मां त्वं पूर्वमजीजनः} %||7-94-4||

\twolineshloka
{भोगवन्तं ततो नागमनन्तमुदके शयम्}
{मायया जनयित्वा त्वं द्वौ च सत्त्वौ महाबलौ} %||7-94-5||

\twolineshloka
{मधुं च कैटभं चैव ययोरस्थिचयैर्वृता}
{इयं पर्वतसम्बाधा मेदिनी चाभवन्मही} %||7-94-6||

\twolineshloka
{पद्मे दिव्यार्कसङ्काशे नाभ्यामुत्पाद्य मामपि}
{प्राजापत्यं त्वया कर्म सर्वं मयि निवेशितम्} %||7-94-7||

\twolineshloka
{सोऽहं संन्यस्तभारो हि त्वामुपासे जगत्पतिम्}
{रक्षां विधत्स्व भूतेषु मम तेजः करो भवान्} %||7-94-8||

\twolineshloka
{ततस्त्वमपि दुर्धर्षस्तस्माद्भावात्सनातनात्}
{रक्षार्थं सर्वभूतानां विष्णुत्वमुपजग्मिवान्} %||7-94-9||

\twolineshloka
{अदित्यां वीर्यवान्पुत्रो भ्रातॄणां हर्षवर्धनः}
{समुत्पन्नेषु कृत्येषु लोकसाह्याय कल्पसे} %||7-94-10||

\twolineshloka
{स त्वं वित्रास्यमानासु प्रजासु जगतां वर}
{रावणस्य वधाकाङ्क्षी मानुषेषु मनोऽदधाः} %||7-94-11||

\twolineshloka
{दशवर्षसहस्राणि दशवर्षशतानि च}
{कृत्वा वासस्य नियतिं स्वयमेवात्मनः पुरा} %||7-94-12||

\twolineshloka
{स त्वं मनोमयः पुत्रः पूर्णायुर्मानुषेष्विह}
{कालो नरवरश्रेष्ठ समीपमुपवर्तितुम्} %||7-94-13||

\twolineshloka
{यदि भूयो महाराज प्रजा इच्छस्युपासितुम्}
{वस वा वीर भद्रं ते एवमाह पितामहः} %||7-94-14||

\twolineshloka
{अथ वा विजिगीषा ते सुरलोकाय राघव}
{सनाथा विष्णुना देवा भवन्तु विगतज्वराः} %||7-94-15||

\twolineshloka
{श्रुत्वा पितामहेनोक्तं वाक्यं कालसमीरितम्}
{राघवः प्रहसन्वाक्यं सर्वसंहारमब्रवीत्} %||7-94-16||

\twolineshloka
{श्रुतं मे देवदेवस्य वाक्यं परममद्भुतम्}
{प्रीतिर्हि महती जाता तवागमनसम्भवा} %||7-94-17||

\twolineshloka
{भद्रं तेऽस्तु गमिष्यामि यत एवाहमागतः}
{हृद्गतो ह्यसि सम्प्राप्तो न मेऽस्त्यत्र विचारणा} %||7-94-18||

\twolineshloka
{मया हि सर्वकृत्येषु देवानां वशवर्तिनाम्}
{स्थातव्यं सर्वसंहारे यथा ह्याह पितामहः} %||7-95-19||


{॥इत्यार्षे श्रीमद्रामायणे वाल्मीकीये आदिकाव्ये उत्तरकाण्डे चतुर्नवतितमः सर्गः॥}

\sect{पञ्चनवतितमः सर्गः}

\twolineshloka
{तथा तयोः कथयतोर्दुर्वासा भगवानृषिः}
{रामस्य दर्शनाकाङ्क्षी राजद्वारमुपागमत्} %||7-95-1||

\twolineshloka
{सोऽभिगम्य च सौमित्रिमुवाच ऋषिसत्तमः}
{रामं दर्शय मे शीघ्रं पुरा मेऽर्थोऽतिवर्तते} %||7-95-2||

\twolineshloka
{मुनेस्तु भाषितं श्रुत्वा लक्ष्मणः परवीरहा}
{अभिवाद्य महात्मानं वाक्यमेतदुवाच ह} %||7-95-3||

\twolineshloka
{किं कार्यं ब्रूहि भगवन्को वार्थः किं करोम्यहम्}
{व्यग्रो हि राघवो ब्रह्मन्मुहूर्तं वा प्रतीक्षताम्} %||7-95-4||

\twolineshloka
{तच्छ्रुत्वा ऋषिशार्दूलः क्रोधेन कलुषीकृतः}
{उवाच लक्ष्मणं वाक्यं निर्दहन्निव चक्षुषा} %||7-95-5||

\twolineshloka
{अस्मिन्क्षणे मां सौमित्रे रामाय प्रतिवेदय}
{विषयं त्वां पुरं चैव शपिष्ये राघवं तथा} %||7-95-6||

\twolineshloka
{भरतं चैव सौमित्रे युष्माकं या च सन्ततिः}
{न हि शक्ष्याम्यहं भूयो मन्युं धारयितुं हृदि} %||7-95-7||

\twolineshloka
{तच्छ्रुत्वा घोरसङ्काशं वाक्यं तस्य महात्मनः}
{चिन्तयामास मनसा तस्य वाक्यस्य निश्चयम्} %||7-95-8||

\twolineshloka
{एकस्य मरणं मेऽस्तु मा भूत्सर्वविनाशनम्}
{इति बुद्ध्या विनिश्चित्य राघवाय न्यवेदयत्} %||7-95-9||

\twolineshloka
{लक्ष्मणस्य वचः श्रुत्वा रामः कालं विसृज्य च}
{निष्पत्य त्वरितं राजा अत्रेः पुत्रं ददर्श ह} %||7-95-10||

\twolineshloka
{सोऽभिवाद्य महात्मानं ज्वलन्तमिव तेजसा}
{किं कार्यमिति काकुत्स्थः कृताञ्जलिरभाषत} %||7-95-11||

\twolineshloka
{तद्वाक्यं राघवेण्णोक्तं श्रुत्वा मुनिवरः प्रभुः}
{प्रत्याह रामं दुर्वासाः श्रूयतां धर्मवत्सल} %||7-95-12||

\twolineshloka
{अद्य वर्षसहस्रस्य समाप्तिर्मम राघव}
{सोऽहं भोजनमिच्छामि यथासिद्धं तवानघ} %||7-95-13||

\twolineshloka
{तच्छ्रुत्वा वचनं रामो हर्षेण महतान्वितः}
{भोजनं मुनिमुख्याय यथासिद्धमुपाहरत्} %||7-95-14||

\twolineshloka
{स तु भुक्त्वा मुनिश्रेष्ठस्तदन्नममृतोपमम्}
{साधु रामेति सम्भाष्य स्वमाश्रममुपागमत्} %||7-95-15||

\twolineshloka
{तस्मिन्गते महातेजा राघवः प्रीतमानसः}
{संस्मृत्य कालवाक्यानि ततो दुःखमुपेयिवान्} %||7-95-16||

\twolineshloka
{दुःखेन च सुसन्तप्तः स्मृत्वा तद्घोरदर्शनम्}
{अवान्मुखो दीनमना व्याहर्तुं न शशाक ह} %||7-95-17||

\twolineshloka
{ततो बुद्ध्या विनिश्चित्य कालवाक्यानि राघवः}
{नैतदस्तीति चोक्त्वा स तूष्णीमासीन्महायशाः} %||7-96-18||


{॥इत्यार्षे श्रीमद्रामायणे वाल्मीकीये आदिकाव्ये उत्तरकाण्डे पञ्चनवतितमः सर्गः॥}

\sect{षण्णवतितमः सर्गः}

\twolineshloka
{अवाङ्मुखमथो दीनं दृष्ट्वा सोममिवाप्लुतम्}
{राघवं लक्ष्मणो वाक्यं हृष्टो मधुरमब्रवीत्} %||7-96-1||

\twolineshloka
{न सन्तापं महाबाहो मदर्थं कर्तुमर्हसि}
{पूर्वनिर्माणबद्धा हि कालस्य गतिरीदृशी} %||7-96-2||

\twolineshloka
{जहि मां सौम्य विस्रब्दः प्रतिज्ञां परिपालय}
{हीनप्रतिज्ञाः काकुत्स्थ प्रयान्ति नरकं नराः} %||7-96-3||

\twolineshloka
{यदि प्रीतिर्महाराज यद्यनुग्राह्यता मयि}
{जहि मां निर्विशङ्कस्त्वं धर्मं वर्धय राघव} %||7-96-4||

\twolineshloka
{लक्ष्मणेन तथोक्तस्तु रामः प्रचलितेन्द्रियः}
{मन्त्रिणः समुपानीय तथैव च पुरोधसम्} %||7-96-5||

\twolineshloka
{अब्रवीच्च यथावृत्तं तेषां मध्ये नराधिपः}
{दुर्वासोऽभिगमं चैव प्रतिज्ञां तापसस्य च} %||7-96-6||

\twolineshloka
{तच्छ्रुत्वा मन्त्रिणः सर्वे सोपाध्यायाः समासत}
{वसिष्ठस्तु महातेजा वाक्यमेतदुवाच ह} %||7-96-7||

\twolineshloka
{दृष्टमेतन्महाबाहो क्षयं ते लोमहर्षणम्}
{लक्ष्मणेन वियोगश्च तव राम महायशः} %||7-96-8||

\twolineshloka
{त्यजैनं बलवान्कालो मा प्रतिज्ञां वृथा कृथाः}
{विनष्टायां प्रतिज्ञायां धर्मो हि विलयं व्रजेत्} %||7-96-9||

\twolineshloka
{ततो धर्मे विनष्टे तु त्रैलोक्ये सचराचरम्}
{सदेवर्षिगणं सर्वं विनश्येत न संशयः} %||7-96-10||

\twolineshloka
{स त्वं पुरुषशार्दूल त्रैलोक्यस्याभिपालनम्}
{लक्ष्मणस्य वधेनाद्य जगत्स्वस्थं कुरुष्व ह} %||7-96-11||

\twolineshloka
{तेषां तत्समवेतानां वाक्यं धर्मार्थसंहितम्}
{श्रुत्वा परिषदो मध्ये रामो लक्ष्मणमब्रवीत्} %||7-96-12||

\twolineshloka
{विसर्जये त्वां सौमित्रे मा भूद्धर्मविपर्ययः}
{त्यागो वधो वा विहितः साधूनामुभयं समम्} %||7-96-13||

\twolineshloka
{रामेण भाषिते वाक्ये बाष्पव्याकुलितेक्षणः}
{लक्ष्मणस्त्वरितः प्रायात्स्वगृहं न विवेश ह} %||7-96-14||

\twolineshloka
{स गत्वा सरयूतीरमुपस्पृश्य कृताञ्जलिः}
{निगृह्य सर्वस्रोतांसि निःश्वासं न मुमोच ह} %||7-96-15||

\twolineshloka
{अनुच्छ्वसन्तं युक्तं तं सशक्राः साप्सरोगणाः}
{देवाः सर्षिगणाः सर्वे पुष्पैरवकिरंस्तदा} %||7-96-16||

\twolineshloka
{अदृश्यं सर्वमनुजैः सशरीरं महाबलम्}
{प्रगृह्य लक्ष्मणं शक्रो दिवं सम्प्रविवेश ह} %||7-96-17||

\twolineshloka
{ततो विष्णोश्चतुर्भागमागतं सुरसत्तमाः}
{हृष्टाः प्रमुदिताः सर्वेऽपूजयनृषिभिः सह} %||7-97-18||


{॥इत्यार्षे श्रीमद्रामायणे वाल्मीकीये आदिकाव्ये उत्तरकाण्डे षण्णवतितमः सर्गः॥}

\sect{सप्तनवतितमः सर्गः}

\twolineshloka
{विसृज्य लक्ष्मणं रामो दुःखशोकसमन्वितः}
{पुरोधसं मन्त्रिणश्च नैगमांश्चेदमब्रवीत्} %||7-97-1||

\twolineshloka
{अद्य राज्येऽभिषेक्ष्यामि भरतं धर्मवत्सलम्}
{अयोध्यायां पतिं वीरं ततो यास्याम्यहं वनम्} %||7-97-2||

\twolineshloka
{प्रवेशयत सम्भारान्मा भूत्कालात्ययो यथा}
{अद्यैवाहं गमिष्यामि लक्ष्मणेन गतां गतिम्} %||7-97-3||

\twolineshloka
{तच्छ्रुत्वा राघवेणोक्तं सर्वाः प्रकृतयो भृशम्}
{मूर्धभिः प्रणता भूमौ गतसत्त्वा इवाभवन्} %||7-97-4||

\twolineshloka
{भरतश्च विसंज्ञोऽभूच्छ्रुत्वा रामस्य भाषितम्}
{राज्यं विगर्हयामास राघवं चेदमब्रवीत्} %||7-97-5||

\twolineshloka
{सत्येन हि शपे राजन्स्वर्गलोके न चैव हि}
{न कामये यथा राज्यं त्वां विना रघुनन्दन} %||7-97-6||

\twolineshloka
{इमौ कुशीलवौ राजन्नभिषिञ्च नराधिप}
{कोसलेषु कुशं वीरमुत्तरेषु तथा लवम्} %||7-97-7||

\twolineshloka
{शत्रुघ्नस्य तु गच्छन्तु दूतास्त्वरितविक्रमाः}
{इदं गमनमस्माकं स्वर्गायाख्यान्तु माचिरम्} %||7-97-8||

\twolineshloka
{तच्छ्रुत्वा भरतेनोक्तं दृष्ट्वा चापि ह्यधो मुखान्}
{पौरान्दुःखेन सन्तप्तान्वसिष्ठो वाक्यमब्रवीत्} %||7-97-9||

\twolineshloka
{वत्स राम इमाः पश्य धरणीं प्रकृतीर्गताः}
{ज्ञात्वैषामीप्सितं कार्यं मा चैषां विप्रियं कृथाः} %||7-97-10||

\twolineshloka
{वसिष्ठस्य तु वाक्येन उत्थाप्य प्रकृतीजनम्}
{किं करोमीति काकुत्स्थः सर्वान्वचनमब्रवीत्} %||7-97-11||

\twolineshloka
{ततः सर्वाः प्रकृतयो रामं वचनमब्रुवन्}
{गच्छन्तमनुगच्छामो यतो राम गमिष्यसि} %||7-97-12||

\twolineshloka
{एषा नः परमा प्रीतिरेष धर्मः परो मतः}
{हृद्गता नः सदा तुष्टिस्तवानुगमने दृढा} %||7-97-13||

\twolineshloka
{पौरेषु यदि ते प्रीतिर्यदि स्नेहो ह्यनुत्तमः}
{सपुत्रदाराः काकुत्स्थ समं गच्छाम सत्पथम्} %||7-97-14||

\twolineshloka
{तपोवनं वा दुर्गं वा नदीमम्भोनिधिं तथा}
{वयं ते यदि न त्याज्याः सर्वान्नो नय ईश्वर} %||7-97-15||

\twolineshloka
{स तेषां निश्चयं ज्ञात्वा कृतान्तं च निरीक्ष्यच}
{पौराणां दृढभक्तिं च बाढमित्येव सोऽब्रवीत्} %||7-97-16||

\twolineshloka
{एवं विनिश्चयं कृत्वा तस्मिन्नहनि राघवः}
{कोसलेषु कुशं वीरमुत्तरेषु तथा लवम्} %||7-97-17||

\twolineshloka
{अभिषिञ्चन्महात्मानावुभावेव कुशीलवौ}
{रथानां तु सहस्राणि त्रीणि नागायुतानि च} %||7-97-18||

\twolineshloka
{दश चाश्वसहस्राणि एकैकस्य धनं ददौ}
{बहुरत्नौ बहुधनौ हृष्टपुष्टजनावृतौ} %||7-97-19||

\twolineshloka
{अभिषिच्य तु तौ वीरौ प्रस्थाप्य स्वपुरे तथा}
{दूतान्सम्प्रेषयामास शत्रुघ्नाय महात्मने} %||7-98-20||


{॥इत्यार्षे श्रीमद्रामायणे वाल्मीकीये आदिकाव्ये उत्तरकाण्डे सप्तनवतितमः सर्गः॥}

\sect{अष्टनवतितमः सर्गः}

\twolineshloka
{ते दूता रामवाक्येन चोदिता लघुविक्रमाः}
{प्रजग्मुर्मधुरां शीघ्रं चक्रुर्वासं न चाध्वनि} %||7-98-1||

\twolineshloka
{ततस्त्रिभिरहो रात्रैः सम्प्राप्य मधुरामथ}
{शत्रुघ्नाय यथावृत्तमाचख्युः सर्वमेव तत्} %||7-98-2||

\twolineshloka
{लक्ष्मणस्य परित्यागं प्रतिज्ञां राघवस्य च}
{पुत्रयोरभिषेकं च पौरानुगमनं तथा} %||7-98-3||

\twolineshloka
{कुशस्य नगरी रम्या विन्ध्यपर्वतरोधसि}
{कुशावतीति नाम्ना सा कृता रामेण धीमता} %||7-98-4||

\twolineshloka
{श्राविता च पुरी रम्या श्रावतीति लवस्य च}
{अयोध्यां विजनां चैव भरतं राघवानुगम्} %||7-98-5||

\twolineshloka
{एवं सर्वं निवेद्याशु शत्रुघ्नाय महात्मने}
{विरेमुस्ते ततो दूतास्त्वर राजन्निति ब्रुवन्} %||7-98-6||

\twolineshloka
{श्रुत्वा तं घोरसङ्काशं कुलक्षयमुपस्थितम्}
{प्रकृतीस्तु समानीय काञ्चनं च पुरोहितम्} %||7-98-7||

\twolineshloka
{तेषां सर्वं यथावृत्तमाख्याय रघुनन्दनः}
{आत्मनश्च विपर्यासं भविष्यं भ्रातृभिः सह} %||7-98-8||

\twolineshloka
{ततः पुत्रद्वयं वीरः सोऽभ्यषिञ्चन्नराधिपः}
{सुबाहुर्मधुरां लेभे शत्रुघाती च वैदिशम्} %||7-98-9||

\twolineshloka
{द्विधाकृत्वा तु तां सेनां माधुरीं पुत्रयोर्द्वयोः}
{धनधान्यसमायुक्तौ स्थापयामास पार्थिवौ} %||7-98-10||

\twolineshloka
{ततो विसृज्य राजानं वैदिशे शत्रुघातिनम्}
{जगाम त्वरितोऽयोध्यां रथेनैकेन राघवः} %||7-98-11||

\twolineshloka
{स ददर्श महात्मानं ज्वलन्तमिव पावकम्}
{क्षौमसूक्ष्माम्बरधरं मुनिभिः सार्धमक्षयैः} %||7-98-12||

\twolineshloka
{सोऽभिवाद्य ततो रामं प्राञ्जलिः प्रयतेन्द्रियः}
{उवाच वाक्यं धर्मज्ञो धर्ममेवानुचिन्तयन्} %||7-98-13||

\twolineshloka
{कृत्वाभिषेकं सुतयोर्युक्तं राघवयोर्धनैः}
{तवानुगमने राजन्विद्धि मां कृतनिश्चयम्} %||7-98-14||

\twolineshloka
{न चान्यदत्र वक्तव्यं दुस्तरं तव शासनम्}
{त्यक्तुं नार्हसि मां वीर भक्तिमन्तं विशेषतः} %||7-98-15||

\twolineshloka
{तस्य तां बुद्धिमक्लीबां विज्ञाय रघुनन्दनः}
{बाढमित्येव शत्रुघ्नं रामो वचनमब्रवीत्} %||7-98-16||

\twolineshloka
{तस्य वाक्यस्य वाक्यान्ते वानराः कामरूपिणः}
{ऋक्षराक्षससङ्घाश्च समापेतुरनेकशः} %||7-98-17||

\twolineshloka
{देवपुत्रा ऋषिसुता गन्धर्वाणां सुतास्तथा}
{राम क्षयं विदित्वा ते सर्व एव समागताः} %||7-98-18||

\twolineshloka
{ते राममभिवाद्याहुः सर्व एव समागताः}
{तवानुगमने राजन्सम्प्राप्ताः स्म महायशः} %||7-98-19||

\twolineshloka
{यदि राम विनास्माभिर्गच्छेस्त्वं पुरुषर्षभ}
{यमदण्डमिवोद्यम्य त्वया स्म विनिपातिताः} %||7-98-20||

\twolineshloka
{एवं तेषां वचः श्रुत्वा ऋष्कवानररक्षसाम्}
{विभीषणमथोवाच मधुरं श्लक्ष्णया गिरा} %||7-98-21||

\twolineshloka
{यावत्प्रजा धरिष्यन्ति तावत्त्वं वै विभीषण}
{राक्षसेन्द्र महावीर्य लङ्कास्थः स्वं धरिष्यसि} %||7-98-22||

{प्रजाः संरक्ष धर्मेण नोत्तरं वक्तुमर्हसि॥\devanumber{23}॥} %||7-98-23|| (Check)

\addtocounter{shlokacount}{1}

\twolineshloka
{तमेवमुक्त्वा काकुत्स्थो हनूमन्तमथाब्रवीत्}
{जीविते कृतबुद्धिस्त्वं मा प्रतिज्ञां विलोपय} %||7-98-24||

\twolineshloka
{मत्कथाः प्रचरिष्यन्ति यावल्लोके हरीश्वर}
{तावत्त्वं धारयन्प्राणान्प्रतिज्ञामनुपालय} %||7-98-25||

\twolineshloka
{तथैवमुक्त्वा काकुत्स्थः सर्वांस्तानृक्षवानरान्}
{मया सार्धं प्रयातेति तदा तान्राघवोऽब्रवीत्} %||7-99-26||


{॥इत्यार्षे श्रीमद्रामायणे वाल्मीकीये आदिकाव्ये उत्तरकाण्डे अष्टनवतितमः सर्गः॥}

\sect{एकोनशततमः सर्गः}

\twolineshloka
{प्रभातायां तु शर्वर्यां पृथुवक्षा महायशाः}
{रामः कमलपत्राक्षः पुरोधसमथाब्रवीत्} %||7-99-1||

\twolineshloka
{अग्निहोत्रं व्रजत्वग्रे सर्पिर्ज्वलितपावकम्}
{वाजपेयातपत्रं च शोभयानं महापथम्} %||7-99-2||

\twolineshloka
{ततो वसिष्ठस्तेजस्वी सर्वं निरवशेषतः}
{चकार विधिवद्धर्म्यं महाप्रास्थानिकं विधिम्} %||7-99-3||

\twolineshloka
{ततः क्षौमाम्बरधरो ब्रह्म चावर्तयन्परम्}
{कुशान्गृहीत्वा पाणिभ्यां प्रसज्य प्रययावथ} %||7-99-4||

\twolineshloka
{अव्याहरन्क्वचित्किञ्चिन्निश्चेष्टो निःसुखः पथि}
{निर्जगाम गृहात्तस्माद्दीप्यमानो यथांशुमान्} %||7-99-5||

\twolineshloka
{रामस्य पार्श्वे सव्ये तु पद्मा श्रीः सुसमाहिता}
{दक्षिणे ह्रीर्विशालाक्षी व्यवसायस्तथाग्रतः} %||7-99-6||

\twolineshloka
{शरा नानाविधाश्चापि धनुरायतविग्रहम्}
{अनुव्रजन्ति काकुत्स्थं सर्वे पुरुषविग्रहाः} %||7-99-7||

\twolineshloka
{वेदा ब्राह्मणरूपेण सावित्री सर्वरक्षिणी}
{ओङ्कारोऽथ वषट्कारः सर्वे राममनुव्रताः} %||7-99-8||

\twolineshloka
{ऋषयश्च महात्मानः सर्व एव महीसुराः}
{अन्वगच्छन्त काकुत्स्थं स्वर्गद्वारमुपागतम्} %||7-99-9||

\twolineshloka
{तं यान्तमनुयान्ति स्म अन्तःपुरचराः स्त्रियः}
{सवृद्धबालदासीकाः सवर्षवरकिङ्कराः} %||7-99-10||

\twolineshloka
{सान्तःपुरश्च भरतः शत्रुघ्नसहितो ययौ}
{रामव्रतमुपागम्य राघवं समनुव्रताः} %||7-99-11||

\twolineshloka
{ततो विप्रा महात्मानः साग्निहोत्राः समाहिताः}
{सपुत्रदाराः काकुत्स्थमन्वगच्छन्महामतिम्} %||7-99-12||

\twolineshloka
{मन्त्रिणो भृत्यवर्गाश्च सपुत्राः सहबान्धवाः}
{सानुगा राघवं सर्वे अन्वगच्छन्प्रहृष्टवत्} %||7-99-13||

\twolineshloka
{ततः सर्वाः प्रकृतयो हृष्टपुष्टजनावृताः}
{अनुजग्मुः प्रगच्छन्तं राघवं गुणरञ्जिताः} %||7-99-14||

\twolineshloka
{स्नातं प्रमुदितं सर्वं हृष्टपुष्पमनुत्तमम्}
{दृप्तं किलिकिलाशब्दैः सर्वं राममनुव्रतम्} %||7-99-15||

\twolineshloka
{न तत्र कश्चिद्दीनोऽभूद्व्रीडितो वापि दुःखितः}
{हृष्टं प्रमुदितं सर्वं बभूव परमाद्भुतम्} %||7-99-16||

\twolineshloka
{द्रष्टुकामोऽथ निर्याणं राज्ञो जानपदो जनः}
{सम्प्राप्तः सोऽपि दृष्ट्वैव सह सर्वैरनुव्रतः} %||7-99-17||

\twolineshloka
{ऋक्षवानररक्षांसि जनाश्च पुरवासिनः}
{अगछन्परया भक्त्या पृष्ठतः सुसमाहिताः} %||7-100-18||


{॥इत्यार्षे श्रीमद्रामायणे वाल्मीकीये आदिकाव्ये उत्तरकाण्डे एकोनशततमः सर्गः॥}

\sect{शततमः सर्गः}

\twolineshloka
{अध्यर्धयोजनं गत्वा नदीं पश्चान्मुखाश्रिताम्}
{सरयूं पुण्यसलिलां ददर्श रघुनन्दनः} %||7-100-1||

\twolineshloka
{अथ तस्मिन्मुहूर्ते तु ब्रह्मा लोकपितामहः}
{सर्वैः परिवृतो देवैरृषिभिश्च महात्मभिः} %||7-100-2||

\twolineshloka
{आययौ यत्र काकुत्स्थः स्वर्गाय समुपस्थितः}
{विमानशतकोटीभिर्दिव्याभिरभिसंवृतः} %||7-100-3||

{पपात पुष्पवृष्टिश्च वायुमुक्ता महौघवत्॥\devanumber{4}॥} %||7-100-4|| (Check)

\addtocounter{shlokacount}{1}

\twolineshloka
{तस्मिंस्तूर्यशताकीर्णे गन्धर्वाप्सरसङ्कुले}
{सरयूसलिलं रामः पद्भ्यां समुपचक्रमे} %||7-100-5||

\twolineshloka
{ततः पितामहो वाणीमन्तरिक्षादभाषत}
{आगच्छ विष्णो भद्रं ते दिष्ट्या प्राप्तोऽसि राघव} %||7-100-6||

\twolineshloka
{भ्रातृभिः सह देवाभैः प्रविशस्व स्वकां तनुम्}
{वैष्णवीं तां महातेजस्तदाकाशं सनातनम्} %||7-100-7||

\twolineshloka
{त्वं हि लोकगतिर्देव न त्वां केचित्प्रजानते}
{ऋते मायां विशालाक्ष तव पूर्वपरिग्रहाम्} %||7-100-8||

\twolineshloka
{त्वमचिन्त्यं महद्भूतमक्षयं सर्वसङ्ग्रहम्}
{यामिच्छसि महातेजस्तां तनुं प्रविश स्वयम्} %||7-100-9||

\twolineshloka
{पितामहवचः श्रुत्वा विनिश्चित्य महामतिः}
{विवेश वैष्णवं तेजः सशरीरः सहानुजः} %||7-100-10||

\twolineshloka
{ततो विष्णुगतं देवं पूजयन्ति स्म देवताः}
{साध्या मरुद्गणाश्चैव सेन्द्राः साग्निपुरोगमाः} %||7-100-11||

\twolineshloka
{ये च दिव्या ऋषिगणा गन्धर्वाप्सरसश्च याः}
{सुपर्णनागयक्षाश्च दैत्यदानवराक्षसाः} %||7-100-12||

\twolineshloka
{सर्वं हृष्टं प्रमुदितं सर्वं पूर्णमनोरथम्}
{साधु साध्विति तत्सर्वं त्रिदिवं गतकल्मषम्} %||7-100-13||

\twolineshloka
{अथ विष्णुर्महातेजाः पितामहमुवाच ह}
{एषां लोकाञ्जनौघानां दातुमर्हसि सुव्रत} %||7-100-14||

\twolineshloka
{इमे हि सर्वे स्नेहान्मामनुयाता मनस्विनः}
{भक्ता भाजयितव्याश्च त्यक्तात्मानश्च मत्कृते} %||7-100-15||

\twolineshloka
{तच्छ्रुत्वा विष्णुवचनं ब्रह्मा लोकगुरुः प्रभुः}
{लोकान्सान्तानिकान्नाम यास्यन्तीमे समागताः} %||7-100-16||

\threelineshloka
{यच्च तिर्यग्गतं किञ्चिद्राममेवानुचिन्तयत्}
{प्राणांस्त्यक्ष्यति भक्त्या वै सन्ताने तु निवत्स्यति}
{सर्वैरेव गुणैर्युक्ते ब्रह्मलोकादनन्तरे} %||7-100-17||

\twolineshloka
{वानराश्च स्वकां योनिमृक्षाश्चैव तथा ययुः}
{येभ्यो विनिःसृता ये ये सुरादिभ्यः सुसम्भवाः} %||7-100-18||

\twolineshloka
{ऋषिभ्यो नागयक्षेभ्यस्तांस्तानेव प्रपेदिरे}
{तथोक्तवति देवेशे गोप्रतारमुपागताः} %||7-100-19||

\twolineshloka
{भेजिरे सरयूं सर्वे हर्षपूर्णाश्रुविक्लवाः}
{अवगाह्य जलं यो यः प्राणी ह्यासीत्प्रहृष्टवत्} %||7-100-20||

\twolineshloka
{मानुषं देहमुत्सृज्य विमानं सोऽध्यरोहत}
{तिर्यग्योनिगताश्चापि सम्प्राप्ताः सरयूजलम्} %||7-100-21||

\twolineshloka
{दिव्या दिव्येन वपुषा देवा दीप्ता इवाभवन्}
{गत्वा तु सरयूतोयं स्थावराणि चराणि च} %||7-100-22||

\twolineshloka
{प्राप्य तत्तोयविक्लेदं देवलोकमुपागमन्}
{देवानां यस्य या योनिर्वानरा ऋष्क राक्षसाः} %||7-100-23||

\twolineshloka
{तामेव विविशुः सर्वे देवान्निक्षिप्य चाम्भसि}
{तथा स्वर्गगतं सर्वं कृत्वा लोकगुरुर्दिवम्} %||7-100-24||

{जगाम त्रिदशैः सार्धं हृष्टैर्हृष्टो महामतिः॥\devanumber{25}॥} %||7-101-25|| (Check)

\addtocounter{shlokacount}{1}


{॥इत्यार्षे श्रीमद्रामायणे वाल्मीकीये आदिकाव्ये उत्तरकाण्डे शततमः सर्गः॥}

